% Options for packages loaded elsewhere
% Options for packages loaded elsewhere
\PassOptionsToPackage{unicode}{hyperref}
\PassOptionsToPackage{hyphens}{url}
\PassOptionsToPackage{dvipsnames,svgnames,x11names}{xcolor}
%
\documentclass[
  australian,
  12pt,
  letterpaper,
  DIV=11,
  numbers=noendperiod,
  twoside]{scrreprt}
\usepackage{xcolor}
\usepackage[top=22mm,bottom=25mm,left=25mm,right=25mm]{geometry}
\usepackage{amsmath,amssymb}
\setcounter{secnumdepth}{3}
\usepackage{iftex}
\ifPDFTeX
  \usepackage[OT1]{fontenc}
  \usepackage[utf8]{inputenc}
  \usepackage{textcomp} % provide euro and other symbols
\else % if luatex or xetex
  \usepackage{unicode-math} % this also loads fontspec
  \defaultfontfeatures{Scale=MatchLowercase}
  \defaultfontfeatures[\rmfamily]{Ligatures=TeX,Scale=1}
\fi
\usepackage[]{fix-cm}
\ifPDFTeX\else
  % xetex/luatex font selection
\fi
% Use upquote if available, for straight quotes in verbatim environments
\IfFileExists{upquote.sty}{\usepackage{upquote}}{}
\IfFileExists{microtype.sty}{% use microtype if available
  \usepackage[]{microtype}
  \UseMicrotypeSet[protrusion]{basicmath} % disable protrusion for tt fonts
}{}
\makeatletter
\@ifundefined{KOMAClassName}{% if non-KOMA class
  \IfFileExists{parskip.sty}{%
    \usepackage{parskip}
  }{% else
    \setlength{\parindent}{0pt}
    \setlength{\parskip}{6pt plus 2pt minus 1pt}}
}{% if KOMA class
  \KOMAoptions{parskip=half}}
\makeatother
% Make \paragraph and \subparagraph free-standing
\makeatletter
\ifx\paragraph\undefined\else
  \let\oldparagraph\paragraph
  \renewcommand{\paragraph}{
    \@ifstar
      \xxxParagraphStar
      \xxxParagraphNoStar
  }
  \newcommand{\xxxParagraphStar}[1]{\oldparagraph*{#1}\mbox{}}
  \newcommand{\xxxParagraphNoStar}[1]{\oldparagraph{#1}\mbox{}}
\fi
\ifx\subparagraph\undefined\else
  \let\oldsubparagraph\subparagraph
  \renewcommand{\subparagraph}{
    \@ifstar
      \xxxSubParagraphStar
      \xxxSubParagraphNoStar
  }
  \newcommand{\xxxSubParagraphStar}[1]{\oldsubparagraph*{#1}\mbox{}}
  \newcommand{\xxxSubParagraphNoStar}[1]{\oldsubparagraph{#1}\mbox{}}
\fi
\makeatother

\usepackage{color}
\usepackage{fancyvrb}
\newcommand{\VerbBar}{|}
\newcommand{\VERB}{\Verb[commandchars=\\\{\}]}
\DefineVerbatimEnvironment{Highlighting}{Verbatim}{commandchars=\\\{\}}
% Add ',fontsize=\small' for more characters per line
\usepackage{framed}
\definecolor{shadecolor}{RGB}{241,243,245}
\newenvironment{Shaded}{\begin{snugshade}}{\end{snugshade}}
\newcommand{\AlertTok}[1]{\textcolor[rgb]{0.68,0.00,0.00}{#1}}
\newcommand{\AnnotationTok}[1]{\textcolor[rgb]{0.37,0.37,0.37}{#1}}
\newcommand{\AttributeTok}[1]{\textcolor[rgb]{0.40,0.45,0.13}{#1}}
\newcommand{\BaseNTok}[1]{\textcolor[rgb]{0.68,0.00,0.00}{#1}}
\newcommand{\BuiltInTok}[1]{\textcolor[rgb]{0.00,0.23,0.31}{#1}}
\newcommand{\CharTok}[1]{\textcolor[rgb]{0.13,0.47,0.30}{#1}}
\newcommand{\CommentTok}[1]{\textcolor[rgb]{0.37,0.37,0.37}{#1}}
\newcommand{\CommentVarTok}[1]{\textcolor[rgb]{0.37,0.37,0.37}{\textit{#1}}}
\newcommand{\ConstantTok}[1]{\textcolor[rgb]{0.56,0.35,0.01}{#1}}
\newcommand{\ControlFlowTok}[1]{\textcolor[rgb]{0.00,0.23,0.31}{\textbf{#1}}}
\newcommand{\DataTypeTok}[1]{\textcolor[rgb]{0.68,0.00,0.00}{#1}}
\newcommand{\DecValTok}[1]{\textcolor[rgb]{0.68,0.00,0.00}{#1}}
\newcommand{\DocumentationTok}[1]{\textcolor[rgb]{0.37,0.37,0.37}{\textit{#1}}}
\newcommand{\ErrorTok}[1]{\textcolor[rgb]{0.68,0.00,0.00}{#1}}
\newcommand{\ExtensionTok}[1]{\textcolor[rgb]{0.00,0.23,0.31}{#1}}
\newcommand{\FloatTok}[1]{\textcolor[rgb]{0.68,0.00,0.00}{#1}}
\newcommand{\FunctionTok}[1]{\textcolor[rgb]{0.28,0.35,0.67}{#1}}
\newcommand{\ImportTok}[1]{\textcolor[rgb]{0.00,0.46,0.62}{#1}}
\newcommand{\InformationTok}[1]{\textcolor[rgb]{0.37,0.37,0.37}{#1}}
\newcommand{\KeywordTok}[1]{\textcolor[rgb]{0.00,0.23,0.31}{\textbf{#1}}}
\newcommand{\NormalTok}[1]{\textcolor[rgb]{0.00,0.23,0.31}{#1}}
\newcommand{\OperatorTok}[1]{\textcolor[rgb]{0.37,0.37,0.37}{#1}}
\newcommand{\OtherTok}[1]{\textcolor[rgb]{0.00,0.23,0.31}{#1}}
\newcommand{\PreprocessorTok}[1]{\textcolor[rgb]{0.68,0.00,0.00}{#1}}
\newcommand{\RegionMarkerTok}[1]{\textcolor[rgb]{0.00,0.23,0.31}{#1}}
\newcommand{\SpecialCharTok}[1]{\textcolor[rgb]{0.37,0.37,0.37}{#1}}
\newcommand{\SpecialStringTok}[1]{\textcolor[rgb]{0.13,0.47,0.30}{#1}}
\newcommand{\StringTok}[1]{\textcolor[rgb]{0.13,0.47,0.30}{#1}}
\newcommand{\VariableTok}[1]{\textcolor[rgb]{0.07,0.07,0.07}{#1}}
\newcommand{\VerbatimStringTok}[1]{\textcolor[rgb]{0.13,0.47,0.30}{#1}}
\newcommand{\WarningTok}[1]{\textcolor[rgb]{0.37,0.37,0.37}{\textit{#1}}}

\usepackage{longtable,booktabs,array}
\newcounter{none} % for unnumbered tables
\usepackage{calc} % for calculating minipage widths
% Correct order of tables after \paragraph or \subparagraph
\usepackage{etoolbox}
\makeatletter
\patchcmd\longtable{\par}{\if@noskipsec\mbox{}\fi\par}{}{}
\makeatother
% Allow footnotes in longtable head/foot
\IfFileExists{footnotehyper.sty}{\usepackage{footnotehyper}}{\usepackage{footnote}}
\makesavenoteenv{longtable}
\usepackage{graphicx}
\makeatletter
\newsavebox\pandoc@box
\newcommand*\pandocbounded[1]{% scales image to fit in text height/width
  \sbox\pandoc@box{#1}%
  \Gscale@div\@tempa{\textheight}{\dimexpr\ht\pandoc@box+\dp\pandoc@box\relax}%
  \Gscale@div\@tempb{\linewidth}{\wd\pandoc@box}%
  \ifdim\@tempb\p@<\@tempa\p@\let\@tempa\@tempb\fi% select the smaller of both
  \ifdim\@tempa\p@<\p@\scalebox{\@tempa}{\usebox\pandoc@box}%
  \else\usebox{\pandoc@box}%
  \fi%
}
% Set default figure placement to htbp
\def\fps@figure{htbp}
\makeatother



\ifLuaTeX
\usepackage[bidi=basic,shorthands=off]{babel}
\else
\usepackage[bidi=default,shorthands=off]{babel}
\fi
\ifLuaTeX
  \usepackage{selnolig} % disable illegal ligatures
\fi


\setlength{\emergencystretch}{3em} % prevent overfull lines

\providecommand{\tightlist}{%
  \setlength{\itemsep}{0pt}\setlength{\parskip}{0pt}}



 


\usepackage{pmboxdraw}
\usepackage[scaled=0.90]{roboto-mono}

% Roman-numbered page numbering
\AtBeginDocument{\pagenumbering{roman}}

% Headers
\usepackage[automark,headsepline]{scrlayer-scrpage}
\clearpairofpagestyles

\automark[section]{chapter}

% Left/even page: chapter
\lehead{\leftmark}

% Right/odd page: section
\rohead{\rightmark}

% Page number on every page, including chapter-opening pages
\cfoot*{\pagemark}

\setkomafont{pageheadfoot}{\small\sffamily}
\KOMAoption{captions}{tableheading}
\makeatletter
\@ifpackageloaded{bookmark}{}{\usepackage{bookmark}}
\makeatother
\makeatletter
\@ifpackageloaded{caption}{}{\usepackage{caption}}
\AtBeginDocument{%
\ifdefined\contentsname
  \renewcommand*\contentsname{Table of contents}
\else
  \newcommand\contentsname{Table of contents}
\fi
\ifdefined\listfigurename
  \renewcommand*\listfigurename{List of Figures}
\else
  \newcommand\listfigurename{List of Figures}
\fi
\ifdefined\listtablename
  \renewcommand*\listtablename{List of Tables}
\else
  \newcommand\listtablename{List of Tables}
\fi
\ifdefined\figurename
  \renewcommand*\figurename{Figure}
\else
  \newcommand\figurename{Figure}
\fi
\ifdefined\tablename
  \renewcommand*\tablename{Table}
\else
  \newcommand\tablename{Table}
\fi
}
\@ifpackageloaded{float}{}{\usepackage{float}}
\floatstyle{ruled}
\@ifundefined{c@chapter}{\newfloat{codelisting}{h}{lop}}{\newfloat{codelisting}{h}{lop}[chapter]}
\floatname{codelisting}{Listing}
\newcommand*\listoflistings{\listof{codelisting}{List of Listings}}
\makeatother
\makeatletter
\makeatother
\makeatletter
\@ifpackageloaded{caption}{}{\usepackage{caption}}
\@ifpackageloaded{subcaption}{}{\usepackage{subcaption}}
\makeatother
\usepackage{bookmark}
\IfFileExists{xurl.sty}{\usepackage{xurl}}{} % add URL line breaks if available
\urlstyle{same}
\hypersetup{
  pdftitle={Foundations of Autonomous Robotics},
  pdfauthor={Chien-Hsiang Yeh},
  pdflang={en-AU},
  colorlinks=true,
  linkcolor={blue},
  filecolor={Maroon},
  citecolor={Blue},
  urlcolor={Blue},
  pdfcreator={LaTeX via pandoc}}


\title{Foundations of Autonomous Robotics}
\usepackage{etoolbox}
\makeatletter
\providecommand{\subtitle}[1]{% add subtitle to \maketitle
  \apptocmd{\@title}{\par {\large #1 \par}}{}{}
}
\makeatother
\subtitle{Mathematics, Physics, Algorithms, and Engineering}
\author{Chien-Hsiang Yeh}
\date{31 July 2026}
\begin{document}
\maketitle

\renewcommand*\contentsname{Table of contents}
{
\setcounter{tocdepth}{2}
\tableofcontents
}

\bookmarksetup{startatroot}

\chapter*{Preface}\label{preface}
\addcontentsline{toc}{chapter}{Preface}

\markboth{Preface}{Preface}

This book is my collection of learning notes for understanding
intelligence, especially intelligence embodied in agents that perceive
their environment, predict what may happen, choose actions, and act in
the physical world. Writing the notes is part of the learning process:
the aim is to make each idea precise enough that I can explain it,
derive it, implement it, and test it.

My main algorithmic interests are \textbf{reinforcement learning} and
\textbf{world models}. Reinforcement learning studies how an agent can
improve its decisions through interaction and feedback. A world model is
an internal model that predicts how the environment or system may
evolve, allowing an agent to evaluate possible actions before or while
taking them. These algorithms depend on more fundamental ideas from
mathematics, dynamics, control, estimation, planning, perception, and
software engineering, so the notes build those foundations first.

\bookmarksetup{startatroot}

\chapter*{Unified Notation and
Conventions}\label{unified-notation-and-conventions}
\addcontentsline{toc}{chapter}{Unified Notation and Conventions}

\markboth{Unified Notation and Conventions}{Unified Notation and
Conventions}

This note is the canonical notation registry for the project. Chapters
may introduce local symbols, but they must not redefine symbols listed
here. When a textbook uses different notation, translate it into these
conventions.

\section*{Core conventions}\label{core-conventions}

\markright{Core conventions}

\begin{enumerate}
\def\labelenumi{\arabic{enumi}.}
\tightlist
\item
  Scalars use italic lowercase letters, vectors use bold lowercase
  letters, and matrices use bold uppercase letters.
\item
  Vectors are column vectors unless stated otherwise.
\item
  SI units are used in calculations. Angles are expressed in radians.
\item
  Coordinate frames are right-handed. A leading left superscript
  identifies the frame in which coordinates are expressed.
\item
  Continuous time uses \(t\); sampled time uses the integer index \(k\).
\item
  A wide hat over a twist, as in \(\widehat{\boldsymbol\xi}\), denotes
  the declared Lie-algebra hat operator and is not an estimate.
\item
  Bold symbols and scalar symbols with the same letter are distinct; for
  example, \(\mathbf R\) is a rotation matrix whereas \(R\) may be a
  signed turning radius.
\item
  A symbol may have a different local meaning only in a clearly
  separated mathematical domain where its type and meaning are
  explicitly redeclared; for example, \(m\) may denote a matrix
  dimension in a generic linear-algebra statement and robot mass in a
  dynamics model.
\end{enumerate}

\section*{Part I: Motion, Mechanics, and
Control}\label{part-i-motion-mechanics-and-control}

\markright{Part I: Motion, Mechanics, and Control}

\subsection*{\texorpdfstring{Notation introduced in
Chapter~\ref{sec-linear-algebra-foundations}}{Notation introduced in Chapter~}}\label{notation-introduced-in-sec-linear-algebra-foundations}

{\def\LTcaptype{none} % do not increment counter
\begin{longtable}[]{@{}
  >{\raggedright\arraybackslash}p{(\linewidth - 4\tabcolsep) * \real{0.3333}}
  >{\raggedright\arraybackslash}p{(\linewidth - 4\tabcolsep) * \real{0.3333}}
  >{\raggedright\arraybackslash}p{(\linewidth - 4\tabcolsep) * \real{0.3333}}@{}}
\toprule\noalign{}
\begin{minipage}[b]{\linewidth}\raggedright
Symbol
\end{minipage} & \begin{minipage}[b]{\linewidth}\raggedright
Meaning
\end{minipage} & \begin{minipage}[b]{\linewidth}\raggedright
Units or convention
\end{minipage} \\
\midrule\noalign{}
\endhead
\bottomrule\noalign{}
\endlastfoot
\(a,\alpha\) & Scalar quantity & Depends on the quantity \\
\(m,n,i,j\) & Integer dimensions or indices & Declared locally \\
\(\mathbb R\) & Set of real numbers & --- \\
\(\mathbb R^n,\mathbb R^{m\times n}\) & Real coordinate columns and real
matrices & Dimensions shown by the superscripts \\
\(\mathbb R_{\geq0},\mathbb R_{>0}\) & Nonnegative and strictly positive
real numbers & --- \\
\(\mathbf{v}\in\mathbb{R}^n\) & Column vector with \(n\) components &
Depends on the quantity \\
\(\mathbf{A}\in\mathbb{R}^{m\times n}\) & Matrix with \(m\) rows and
\(n\) columns & Depends on the map \\
\(\mathbf{I}_n\) & \(n\times n\) identity matrix & Dimensionless \\
\(\mathbf{0}_n\) & Zero vector with \(n\) components & Depends on the
quantity \\
\((\cdot)^\mathsf{T}\) & Transpose & --- \\
\((\cdot)^{-1}\) & Inverse, when it exists & --- \\
\(\lVert\mathbf{v}\rVert_2\) & Euclidean norm of \(\mathbf{v}\) & Same
units as \(\mathbf{v}\) \\
\(\det(\mathbf A)\) & Determinant of a square matrix & Depends on the
matrix units \\
\(\operatorname{rank}(\mathbf A)\) & Dimension of the column space of
\(\mathbf A\) & Nonnegative integer \\
\(\operatorname{nullity}(\mathbf A)\) & Dimension of the null space of
\(\mathbf A\) & Nonnegative integer \\
\(\operatorname{span}(\mathbf v_1,\ldots,\mathbf v_k)\) & Set of all
linear combinations of the listed vectors & Same vector space as the
vectors \\
\(\operatorname{Col}(\mathbf A),\mathcal N(\mathbf A)\) & Column space
and null space of \(\mathbf A\) & Subspaces of the output and input
spaces \\
\(\mathbf u\times\mathbf v\) & Three-dimensional cross product of
vectors expressed in the same right-handed orthonormal frame & Product
of the input units \\
\(O(n),SO(n)\) & Orthogonal group and special orthogonal group of real
\(n\times n\) matrices & \(\mathbf Q^\mathsf T\mathbf Q=\mathbf I_n\);
\(SO(n)\) additionally requires \(\det\mathbf Q=1\) \\
\end{longtable}
}

\subsection*{\texorpdfstring{Notation introduced in
Chapter~\ref{sec-geometry-coordinate-frames}}{Notation introduced in Chapter~}}\label{notation-introduced-in-sec-geometry-coordinate-frames}

{\def\LTcaptype{none} % do not increment counter
\begin{longtable}[]{@{}
  >{\raggedright\arraybackslash}p{(\linewidth - 4\tabcolsep) * \real{0.3333}}
  >{\raggedright\arraybackslash}p{(\linewidth - 4\tabcolsep) * \real{0.3333}}
  >{\raggedright\arraybackslash}p{(\linewidth - 4\tabcolsep) * \real{0.3333}}@{}}
\toprule\noalign{}
\begin{minipage}[b]{\linewidth}\raggedright
Symbol
\end{minipage} & \begin{minipage}[b]{\linewidth}\raggedright
Meaning
\end{minipage} & \begin{minipage}[b]{\linewidth}\raggedright
Units or convention
\end{minipage} \\
\midrule\noalign{}
\endhead
\bottomrule\noalign{}
\endlastfoot
\(\{a\}\) & Coordinate frame named \(a\) & Right-handed \\
\(O_a\) & Origin of frame \(\{a\}\) & --- \\
\(\mathbf{e}_{x,a},\mathbf{e}_{y,a},\mathbf{e}_{z,a}\) & Unit basis
vectors of frame \(\{a\}\) & Dimensionless \\
\({}^{a}\mathbf{v}\) & Coordinates of geometric vector \(\mathbf{v}\)
expressed in frame \(\{a\}\) & Depends on \(\mathbf{v}\) \\
\({}^{a}\mathbf{p}_P\) & Position of point \(P\) relative to \(O_a\),
expressed in frame \(\{a\}\) & \(\mathrm{m}\) \\
\({}^{a}\mathbf r_{OP}\) & Relative-position vector from point \(O\) to
point \(P\), expressed in frame \(\{a\}\) & \(\mathrm{m}\) \\
\(\mathbf{R}_{ab}\) & Rotation mapping coordinates from frame \(\{b\}\)
to frame \(\{a\}\) & \(\mathbf{R}_{ab}\in SO(2)\) or \(SO(3)\) \\
\({}^{a}\mathbf{p}_b\) & Position of origin \(O_b\) relative to \(O_a\),
expressed in frame \(\{a\}\) & \(\mathrm{m}\) \\
\(\boldsymbol\omega\) & Three-dimensional angular-velocity vector &
\(\mathrm{rad\,s^{-1}}\) \\
\end{longtable}
}

\subsection*{\texorpdfstring{Notation introduced in
Chapter~\ref{sec-differential-drive-model}}{Notation introduced in Chapter~}}\label{notation-introduced-in-sec-differential-drive-model}

{\def\LTcaptype{none} % do not increment counter
\begin{longtable}[]{@{}
  >{\raggedright\arraybackslash}p{(\linewidth - 4\tabcolsep) * \real{0.2331}}
  >{\raggedright\arraybackslash}p{(\linewidth - 4\tabcolsep) * \real{0.4294}}
  >{\raggedright\arraybackslash}p{(\linewidth - 4\tabcolsep) * \real{0.3374}}@{}}
\toprule\noalign{}
\begin{minipage}[b]{\linewidth}\raggedright
Symbol
\end{minipage} & \begin{minipage}[b]{\linewidth}\raggedright
Meaning
\end{minipage} & \begin{minipage}[b]{\linewidth}\raggedright
Units or convention
\end{minipage} \\
\midrule\noalign{}
\endhead
\bottomrule\noalign{}
\endlastfoot
\(t\) & Continuous time & \(\mathrm{s}\) \\
\(k\in\mathbb N_0\) & Sample index, where
\(\mathbb N_0=\{0,1,2,\ldots\}\) & Dimensionless integer \\
\(t_k\) & Time at sample \(k\) & \(\mathrm{s}\) \\
\(\Delta t=t_{k+1}-t_k\) & Sampling or numerical-integration interval &
\(\mathrm{s}\) \\
\(\dot{\mathbf{x}},\ddot{\mathbf{x}}\) & First and second derivatives
with respect to time & Quantity per \(\mathrm{s}\) and per
\(\mathrm{s}^2\) \\
\(\{w\},\{b\}\) & World and robot body frames & --- \\
\(O_w,O_b\) & Origins of the world and body frames & --- \\
\(\mathbf{q}=[x\;y\;\theta]^\mathsf{T}\) & Planar robot pose coordinates
in the world frame & \(\mathrm{m},\mathrm{m},\mathrm{rad}\) \\
\(\theta\) & Counter-clockwise heading of \(\{b\}\) relative to
\(\{w\}\) & \(\mathrm{rad}\) \\
\(\mathbf u=[v\;\omega]^\mathsf T\) & Planar body-velocity input &
\(\mathrm{m\,s^{-1}},\mathrm{rad\,s^{-1}}\) \\
\(v,\omega\) & Signed body-forward speed and counter-clockwise yaw rate
& \(\mathrm{m\,s^{-1}},\mathrm{rad\,s^{-1}}\) \\
\(\phi_L,\phi_R\) & Left and right wheel angles & \(\mathrm{rad}\) \\
\(\dot{\phi}_L,\dot{\phi}_R\) & Left and right wheel angular speeds &
\(\mathrm{rad\,s^{-1}}\) \\
\(v_L,v_R\) & Signed longitudinal ground speeds of the left and right
wheels & \(\mathrm{m\,s^{-1}}\) \\
\(r,L\) & Effective wheel radius and track width between wheel contact
lines & \(\mathrm{m}\) \\
\(C\) & Instantaneous centre of curvature for planar motion & Geometric
point \\
\(\rho\) & Nonnegative radial distance declared in the local geometric
context & \(\mathrm{m}\) \\
\(R\) & Signed turning radius; positive for an ICC on the robot's left &
\(\mathrm{m}\); distinct from rotation matrix \(\mathbf R\) \\
\(\Delta\ell,\Delta\theta\) & Signed arc-length and heading changes over
an interval & \(\mathrm{m},\mathrm{rad}\) \\
\(\mathbf q_k^{\mathrm E}\) & Forward Euler approximation of the pose at
sample \(k\) & \(\mathrm{m},\mathrm{m},\mathrm{rad}\) \\
\(E(\Delta t)\) & Declared numerical-integration error magnitude at step
size \(\Delta t\) & Depends on the error definition \\
\(e_p\) & One-step position-error magnitude in the declared integration
comparison & \(\mathrm m\) \\
\(O(h^p)\) & Quantity bounded in magnitude by a constant times \(h^p\)
as \(h\to0\) & Units follow the bounded quantity \\
\(\operatorname{sinc}(z)\) & Unnormalised sinc function, \(\sin z/z\)
for \(z\neq0\) and \(1\) for \(z=0\) & Dimensionless argument and
value \\
\end{longtable}
}

\subsection*{\texorpdfstring{Notation introduced in
Chapter~\ref{sec-planar-rigid-body-motion}}{Notation introduced in Chapter~}}\label{notation-introduced-in-sec-planar-rigid-body-motion}

{\def\LTcaptype{none} % do not increment counter
\begin{longtable}[]{@{}
  >{\raggedright\arraybackslash}p{(\linewidth - 4\tabcolsep) * \real{0.3333}}
  >{\raggedright\arraybackslash}p{(\linewidth - 4\tabcolsep) * \real{0.3333}}
  >{\raggedright\arraybackslash}p{(\linewidth - 4\tabcolsep) * \real{0.3333}}@{}}
\toprule\noalign{}
\begin{minipage}[b]{\linewidth}\raggedright
Symbol
\end{minipage} & \begin{minipage}[b]{\linewidth}\raggedright
Meaning
\end{minipage} & \begin{minipage}[b]{\linewidth}\raggedright
Units or convention
\end{minipage} \\
\midrule\noalign{}
\endhead
\bottomrule\noalign{}
\endlastfoot
\({}^{a}\overline{\mathbf p}_P\) & Homogeneous point coordinate obtained
by appending \(1\) to \({}^{a}\mathbf p_P\) & Position entries in
\(\mathrm m\); final entry dimensionless \\
\(SO(2)\) & Special orthogonal group of planar rotation matrices &
\(\mathbf R^\mathsf T\mathbf R=\mathbf I_2\), \(\det\mathbf R=1\) \\
\(SE(2)\) & Special Euclidean group of planar rigid transformations &
Planar rotation and translation represented together \\
\(\mathbf{T}_{ab}\) & Homogeneous transform mapping coordinates from
frame \(\{b\}\) to frame \(\{a\}\) & \(\mathbf{T}_{ab}\in SE(2)\) \\
\(T_{\mathbf I_2}SO(2)\) & Tangent space of \(SO(2)\) at the identity
rotation \(\mathbf I_2\) & Linear space of rotation-curve derivatives
through \(\mathbf I_2\) \\
\(T_{\mathbf I_3}SE(2)\) & Tangent space of \(SE(2)\) at the identity
transformation \(\mathbf I_3\) & Linear space of transformation-curve
derivatives through \(\mathbf I_3\) \\
\(\mathbf S_2\) & Planar skew generator
\(\begin{bmatrix}0&-1\\1&0\end{bmatrix}\) & Dimensionless \\
\(\mathfrak{so}(2)\) & Lie algebra and identity tangent space of
\(SO(2)\) & \(2\times2\) real skew-symmetric matrices \\
\(\mathfrak{se}(2)\) & Lie algebra and identity tangent space of
\(SE(2)\) & Structured \(3\times3\) infinitesimal planar rigid-motion
matrices \\
\(\boldsymbol\xi=[\boldsymbol\nu^\mathsf T\;\omega]^\mathsf T\) & Planar
twist coordinate column with linear part
\(\boldsymbol\nu\in\mathbb R^2\) and angular part \(\omega\in\mathbb R\)
& \(\mathrm{m\,s^{-1}}\) and \(\mathrm{rad\,s^{-1}}\) \\
\(\widehat{\boldsymbol\xi}\) & Hat-operator matrix representing
\(\boldsymbol\xi\) in \(\mathfrak{se}(2)\) & Upper-left block in
\(\mathrm{s^{-1}}\); translation column in \(\mathrm{m\,s^{-1}}\) \\
\((\cdot)^\vee\) & Vee operator, inverse of the declared hat operator &
Returns the twist coordinate column \\
\(\boldsymbol\xi_b,\boldsymbol\xi_w\) & Body and spatial twist
coordinates for the same instantaneous motion & Body axes/body origin
and world axes/world origin conventions \\
\(\boldsymbol\nu_b,\boldsymbol\nu_w\) & Linear parts of the body and
spatial twists & \(\mathrm{m\,s^{-1}}\) \\
\(\operatorname{Ad}_{\mathbf T}\) & Adjoint mapping that changes twist
coordinates according to rigid transform \(\mathbf T\) & Matrix
dimensions and mixed units follow the twist convention \\
\(\ell\) & Declared characteristic length used to compare linear and
angular twist components & \(\mathrm m\) \\
\end{longtable}
}

\subsection*{\texorpdfstring{Notation introduced in
Chapter~\ref{sec-spatial-rigid-body-motion}}{Notation introduced in Chapter~}}\label{notation-introduced-in-sec-spatial-rigid-body-motion}

{\def\LTcaptype{none} % do not increment counter
\begin{longtable}[]{@{}
  >{\raggedright\arraybackslash}p{(\linewidth - 4\tabcolsep) * \real{0.3333}}
  >{\raggedright\arraybackslash}p{(\linewidth - 4\tabcolsep) * \real{0.3333}}
  >{\raggedright\arraybackslash}p{(\linewidth - 4\tabcolsep) * \real{0.3333}}@{}}
\toprule\noalign{}
\begin{minipage}[b]{\linewidth}\raggedright
Symbol
\end{minipage} & \begin{minipage}[b]{\linewidth}\raggedright
Meaning
\end{minipage} & \begin{minipage}[b]{\linewidth}\raggedright
Units or convention
\end{minipage} \\
\midrule\noalign{}
\endhead
\bottomrule\noalign{}
\endlastfoot
\(\mathcal C\) & Configuration space: the set of all configurations
allowed by the model & May be curved and need not be a vector space \\
\(\mathbf q\in\mathcal C\) & Configuration or a chosen column of
generalised coordinates & Components and units depend on the system \\
\(n_{\mathrm{dof}}\) & Number of configuration degrees of freedom &
Nonnegative integer \\
\(n_v,n_c\) & Numbers of scalar variables and independent regular
equality constraints in a local count & Nonnegative integers \\
\(\mathcal M\) & Smooth manifold: a space that locally resembles a
Euclidean coordinate space & Dimension declared locally \\
\(T_{\mathbf G}\mathcal M\) & Tangent space of manifold \(\mathcal M\)
at \(\mathbf G\) & Linear space of curve derivatives through
\(\mathbf G\) \\
\(SO(3)\) & Special orthogonal group of spatial rotation matrices &
\(\mathbf R^\mathsf T\mathbf R=\mathbf I_3\), \(\det\mathbf R=1\) \\
\(SE(3)\) & Special Euclidean group of spatial rigid transformations &
Spatial rotation and translation represented together \\
\([\mathbf{v}]_\times\) & Skew-symmetric matrix satisfying
\([\mathbf{v}]_\times\mathbf{w}=\mathbf{v}\times\mathbf{w}\) & Same
units as \(\mathbf{v}\) \\
\end{longtable}
}

\subsection*{\texorpdfstring{Notation introduced in
\textbf{?@sec-dynamics-actuator-limits}}{Notation introduced in ?@sec-dynamics-actuator-limits}}\label{notation-introduced-in-sec-dynamics-actuator-limits}

{\def\LTcaptype{none} % do not increment counter
\begin{longtable}[]{@{}
  >{\raggedright\arraybackslash}p{(\linewidth - 4\tabcolsep) * \real{0.3333}}
  >{\raggedright\arraybackslash}p{(\linewidth - 4\tabcolsep) * \real{0.3333}}
  >{\raggedright\arraybackslash}p{(\linewidth - 4\tabcolsep) * \real{0.3333}}@{}}
\toprule\noalign{}
\begin{minipage}[b]{\linewidth}\raggedright
Symbol
\end{minipage} & \begin{minipage}[b]{\linewidth}\raggedright
Meaning
\end{minipage} & \begin{minipage}[b]{\linewidth}\raggedright
Units or convention
\end{minipage} \\
\midrule\noalign{}
\endhead
\bottomrule\noalign{}
\endlastfoot
\(s(t),v(t),a(t)\) & Signed one-dimensional position, velocity, and
acceleration & \(\mathrm m,\mathrm{m\,s^{-1}},\mathrm{m\,s^{-2}}\) \\
\(m\) & Robot mass & \(\mathrm{kg}\) \\
\(\mathbf F,F\) & Force vector and force magnitude & \(\mathrm N\) \\
\(F_j,F_{\mathrm{net}}\) & Signed component of external force \(j\) and
their signed net sum & \(\mathrm N\) \\
\(F_{\mathrm{drive}},F_{\mathrm{resist}}\) & Nonnegative drive and
resistance force magnitudes in the declared direction & \(\mathrm N\) \\
\(\boldsymbol\rho\) & Position vector from a chosen moment point to a
force application point & \(\mathrm m\) \\
\(\boldsymbol\tau_O,\tau\) & Torque vector about point \(O\) and its
magnitude & \(\mathrm{N\,m}\) \\
\(\tau_i,F_i\) & Signed axle torque and longitudinal ground force for
wheel \(i\in\{L,R\}\) & \(\mathrm{N\,m},\mathrm N\) \\
\(I_i\) & Wheel mass moment of inertia about its axle &
\(\mathrm{kg\,m^2}\) \\
\(\ddot\phi_i\) & Wheel angular acceleration &
\(\mathrm{rad\,s^{-2}}\) \\
\(G\) & Positive gearbox reduction ratio & Dimensionless \\
\(\eta\) & Gearbox efficiency for the declared power-flow direction &
Dimensionless, \(0<\eta\leq1\) \\
\(\dot\phi_m,\dot\phi_w\) & Motor-shaft and wheel-side angular speeds &
\(\mathrm{rad\,s^{-1}}\) \\
\(\tau_m,\tau_w\) & Motor-shaft and wheel-side torques &
\(\mathrm{N\,m}\) \\
\(P=\tau\dot\phi\) & Rotational mechanical power &
\(\mathrm W=\mathrm{N\,m\,s^{-1}}\) \\
\end{longtable}
}

\cleardoublepage
\pagenumbering{arabic}

\part{Motion, Mechanics, and Control}

\chapter{Linear Algebra
Foundations}\label{sec-linear-algebra-foundations}

\section*{Why linear algebra matters}\label{why-linear-algebra-matters}
\addcontentsline{toc}{section}{Why linear algebra matters}

\markright{Why linear algebra matters}

Autonomous robots represent directions, velocities, forces, sensor
measurements, model parameters, and control commands with vectors.
Matrices transform those vectors, combine coordinate descriptions, solve
systems of equations, and expose whether information has been preserved
or lost. Linear algebra provides the language needed to reason about
these operations before they are used in geometry, kinematics,
estimation, optimisation, and control.

This chapter develops that language from first principles. It
distinguishes geometric vectors from coordinate columns, explains
matrices as linear maps, connects null spaces and rank to invertibility,
defines determinants as signed volume scaling, and derives the
properties of orthogonal matrices used later for rotations.

\section*{Prerequisites and
assumptions}\label{prerequisites-and-assumptions}
\addcontentsline{toc}{section}{Prerequisites and assumptions}

\markright{Prerequisites and assumptions}

\subsection*{Prerequisites}\label{prerequisites}
\addcontentsline{toc}{subsection}{Prerequisites}

The reader should be comfortable with:

\begin{itemize}
\tightlist
\item
  arithmetic and elementary algebra;
\item
  Cartesian coordinates;
\item
  solving simple equations;
\item
  sine and cosine; and
\item
  the project notation conventions.
\end{itemize}

\subsection*{Assumptions}\label{assumptions}
\addcontentsline{toc}{subsection}{Assumptions}

\begin{itemize}
\tightlist
\item
  Scalars are real numbers unless stated otherwise.
\item
  Vectors are written as columns.
\item
  Matrix and vector dimensions are declared when symbols are introduced.
\item
  Coordinate bases used for Euclidean lengths and angles are
  orthonormal.
\item
  Physical units are carried by the represented quantity; abstract basis
  and transformation matrices are dimensionless unless stated otherwise.
\end{itemize}

\section*{Dependency map}\label{dependency-map}
\addcontentsline{toc}{section}{Dependency map}

\markright{Dependency map}

\begin{Shaded}
\begin{Highlighting}[]
\NormalTok{Vectors and linear combinations}
\NormalTok{        |}
\NormalTok{        v}
\NormalTok{Bases and coordinate columns}
\NormalTok{        |}
\NormalTok{        v}
\NormalTok{Matrices as linear maps}
\NormalTok{        |}
\NormalTok{        v}
\NormalTok{Null spaces, rank, and inverses}
\NormalTok{        |}
\NormalTok{        v}
\NormalTok{Determinants and orthogonal matrices}
\NormalTok{        |}
\NormalTok{        v}
\NormalTok{Coordinate frames, rotations, kinematics, estimation, and control}
\end{Highlighting}
\end{Shaded}

\section{Vectors: quantities with magnitude and
direction}\label{vectors-quantities-with-magnitude-and-direction}

\subsection{Geometric vector versus coordinate
column}\label{geometric-vector-versus-coordinate-column}

A geometric vector is independent of any coordinate frame. A coordinate
column is the numerical description of that vector in a chosen basis.

Let \(\{a\}\) denote a right-handed coordinate frame with orthonormal
basis vectors \(\mathbf{e}_{x,a}\), \(\mathbf{e}_{y,a}\), and
\(\mathbf{e}_{z,a}\). For a geometric vector \(\mathbf{v}\), the symbol
\({}^{a}\mathbf{v}\) denotes its coordinate column expressed in frame
\(\{a\}\). In three-dimensional space,

\[
{}^{a}\mathbf{v}
=
\begin{bmatrix}
v_x \\
v_y \\
v_z
\end{bmatrix}.
\]

The entries mean that

\[
\mathbf{v}
=
v_x\mathbf{e}_{x,a}
+v_y\mathbf{e}_{y,a}
+v_z\mathbf{e}_{z,a}.
\]

The vector \(\mathbf{v}\) is the physical or geometric object, whereas
the column \({}^{a}\mathbf{v}\) depends on frame \(\{a\}\). Omitting the
frame label is safe only when the surrounding context makes the frame
unambiguous.

\subsection{Addition and scalar
multiplication}\label{addition-and-scalar-multiplication}

For vectors expressed in the same frame,

\[
{}^{a}\mathbf{u}
+{}^{a}\mathbf{v}
=
\begin{bmatrix}
u_x+v_x \\
u_y+v_y \\
u_z+v_z
\end{bmatrix}.
\]

For a scalar \(\alpha\),

\[
\alpha\,{}^{a}\mathbf{v}
=
\begin{bmatrix}
\alpha v_x \\
\alpha v_y \\
\alpha v_z
\end{bmatrix}.
\]

The algebra is valid only when the quantities are physically compatible.
For example, adding two positions expressed in different frames is not
meaningful until one has been transformed. Adding a velocity in metres
per second to a position in metres is dimensionally invalid even if both
are three-element columns.

\subsection{Length and unit
direction}\label{sec-linear-algebra-vector-norm}

The Euclidean norm is

\[
\left\|{}^{a}\mathbf{v}\right\|_2
=
\sqrt{v_x^2+v_y^2+v_z^2}.
\]

If \(\mathbf{v}\neq\mathbf{0}\), let \({}^{a}\mathbf{e}_v\) denote its
unit direction expressed in frame \(\{a\}\):

\[
{}^{a}\mathbf{e}_v
=
\frac{{}^{a}\mathbf{v}}
{\left\|{}^{a}\mathbf{v}\right\|_2}.
\]

The norm has the same physical units as the vector, while the unit
vector is dimensionless because the units cancel.

\subsection{Dot product and angle}\label{dot-product-and-angle}

Let two vectors be expressed as coordinate columns in the same
orthonormal frame:

\[
{}^{a}\mathbf{u}
=
\begin{bmatrix}
u_x\\
u_y\\
u_z
\end{bmatrix},
\qquad
{}^{a}\mathbf{v}
=
\begin{bmatrix}
v_x\\
v_y\\
v_z
\end{bmatrix}.
\]

The superscript \(\mathsf{T}\) denotes the \textbf{transpose} operation.
It turns the column \({}^{a}\mathbf{u}\in\mathbb{R}^{3\times 1}\) into
the row

\[
\left({}^{a}\mathbf{u}\right)^\mathsf{T}
=
{}^{a}\mathbf{u}^\mathsf{T}
=
\begin{bmatrix}
u_x & u_y & u_z
\end{bmatrix}
\in\mathbb{R}^{1\times 3}.
\]

Here, \(\mathsf{T}\) is an operation on the coordinate column; it is not
a power. Multiplying this row by the column \({}^{a}\mathbf{v}\)
produces a scalar:

\[
{}^{a}\mathbf{u}^\mathsf{T}\,{}^{a}\mathbf{v}
=
\begin{bmatrix}
u_x & u_y & u_z
\end{bmatrix}
\begin{bmatrix}
v_x\\
v_y\\
v_z
\end{bmatrix}
=
u_xv_x+u_yv_y+u_zv_z.
\]

This scalar is the \textbf{dot product}, or Euclidean inner product, of
the two vectors:

\[
\mathbf{u}\boldsymbol{\cdot}\mathbf{v}
\mathrel{:=}
{}^{a}\mathbf{u}^\mathsf{T}\,{}^{a}\mathbf{v}.
\]

The norm introduced in Section~\ref{sec-linear-algebra-vector-norm} can
equivalently be defined from the dot product of a vector with itself:

\[
\begin{aligned}
\left\|{}^{a}\mathbf{u}\right\|_2
&\mathrel{:=}
\sqrt{{}^{a}\mathbf{u}^\mathsf{T}\,{}^{a}\mathbf{u}}\\
&=
\sqrt{u_x^2+u_y^2+u_z^2}.
\end{aligned}
\]

Equivalently, the squared norm is

\[
\left\|{}^{a}\mathbf{u}\right\|_2^2
=
{}^{a}\mathbf{u}^\mathsf{T}\,{}^{a}\mathbf{u}.
\]

Thus, dotting a vector with itself gives its squared length, and taking
the nonnegative square root gives its length. Because frame \(\{a\}\) is
orthonormal, this coordinate norm equals the geometric length
\(\|\mathbf{u}\|_2\).

The coordinate formula above has a geometric interpretation. To
establish it, we first derive the law of cosines.

\subsubsection{Law of cosines}\label{law-of-cosines}

Consider a triangle with two sides of lengths \(r\) and \(s\) and
included angle \(\theta\). Place its vertices at

\[
\mathbf{p}_0=(0,0),
\qquad
\mathbf{p}_1=(r,0),
\qquad
\mathbf{p}_2=(s\cos\theta,s\sin\theta).
\]

Let \(c\) be the length of the side from \(\mathbf{p}_1\) to
\(\mathbf{p}_2\). The horizontal and vertical displacements along that
side are \(s\cos\theta-r\) and \(s\sin\theta\). By the Pythagorean
theorem,

\[
\begin{aligned}
c^2
&=(s\cos\theta-r)^2+(s\sin\theta)^2\\
&=s^2\cos^2\theta-2rs\cos\theta+r^2
  +s^2\sin^2\theta\\
&=r^2+s^2\left(\cos^2\theta+\sin^2\theta\right)
  -2rs\cos\theta\\
&=r^2+s^2-2rs\cos\theta.
\end{aligned}
\]

Therefore, the law of cosines is

\[
\boxed{c^2=r^2+s^2-2rs\cos\theta}.
\]

\subsubsection{Deriving the dot-product angle
formula}\label{deriving-the-dot-product-angle-formula}

For nonzero vectors \(\mathbf{u}\) and \(\mathbf{v}\), place their tails
at the same point and let \(\theta\) be the angle between them. The
vectors form a triangle whose side lengths are

\[
r=\|\mathbf{u}\|_2,
\qquad
s=\|\mathbf{v}\|_2,
\qquad
c=\|\mathbf{u}-\mathbf{v}\|_2.
\]

The law of cosines therefore gives

\[
\|\mathbf{u}-\mathbf{v}\|_2^2
=
\|\mathbf{u}\|_2^2+\|\mathbf{v}\|_2^2
-2\|\mathbf{u}\|_2\|\mathbf{v}\|_2\cos\theta.
\]

Because frame \(\{a\}\) is orthonormal, Euclidean lengths are preserved
by its coordinate representation. Expanding the same squared length
using coordinate columns gives

\[
\begin{aligned}
\|\mathbf{u}-\mathbf{v}\|_2^2
&=
\left({}^{a}\mathbf{u}-{}^{a}\mathbf{v}\right)^\mathsf{T}
\left({}^{a}\mathbf{u}-{}^{a}\mathbf{v}\right)\\
&=
\|\mathbf{u}\|_2^2+\|\mathbf{v}\|_2^2
-2\,{}^{a}\mathbf{u}^\mathsf{T}\,{}^{a}\mathbf{v}.
\end{aligned}
\]

Equating the two expressions for \(\|\mathbf{u}-\mathbf{v}\|_2^2\) and
cancelling their common terms gives

\[
{}^{a}\mathbf{u}^\mathsf{T}\,{}^{a}\mathbf{v}
=
\|\mathbf{u}\|_2\|\mathbf{v}\|_2\cos\theta,
\]

which is the dot-product angle formula. Solving it for the angle gives

\[
\theta
=
\cos^{-1}
\left(
\frac{{}^{a}\mathbf{u}^\mathsf{T}\,{}^{a}\mathbf{v}}
{\|\mathbf{u}\|_2\|\mathbf{v}\|_2}
\right),
\]

provided neither vector is zero.

The dot product is useful for projection, alignment, and checking
orthogonality. If

\[
{}^{a}\mathbf{u}^\mathsf{T}\,{}^{a}\mathbf{v}=0,
\]

then nonzero \(\mathbf{u}\) and \(\mathbf{v}\) are perpendicular.

\subsection{Cross product and useful
identities}\label{sec-linear-algebra-cross-product}

The dot product returns a scalar that measures alignment.
Three-dimensional geometry also needs an operation that constructs a
vector perpendicular to two given directions. Let
\(\mathbf u=[u_1\;u_2\;u_3]^\mathsf T\in\mathbb R^3\) and
\(\mathbf v=[v_1\;v_2\;v_3]^\mathsf T\in\mathbb R^3\) be coordinate
columns expressed in the same right-handed orthonormal frame. Their
\textbf{cross product} is the coordinate column

\[
\boxed{
\mathbf u\times\mathbf v
\mathrel{:=}
\begin{bmatrix}
u_2v_3-u_3v_2\\
u_3v_1-u_1v_3\\
u_1v_2-u_2v_1
\end{bmatrix}
\in\mathbb R^3
}.
\]

The output is perpendicular to both inputs because direct substitution
gives

\[
\mathbf u^\mathsf T(\mathbf u\times\mathbf v)
=
0,
\qquad
\mathbf v^\mathsf T(\mathbf u\times\mathbf v)
=
0.
\]

Its direction follows the right-hand rule. Expanding its squared norm
gives

\[
\begin{aligned}
\|\mathbf u\times\mathbf v\|_2^2
&=
(u_2v_3-u_3v_2)^2
+
(u_3v_1-u_1v_3)^2
+
(u_1v_2-u_2v_1)^2\\
&=
\|\mathbf u\|_2^2\|\mathbf v\|_2^2
-
(\mathbf u^\mathsf T\mathbf v)^2\\
&=
\|\mathbf u\|_2^2\|\mathbf v\|_2^2\sin^2\theta,
\end{aligned}
\]

where \(\theta\in[0,\pi]\) is the angle between the nonzero inputs and
the last equality uses
\(\mathbf u^\mathsf T\mathbf v=\|\mathbf u\|_2\|\mathbf v\|_2\cos\theta\).
Taking the nonnegative square root gives

\[
\boxed{
\|\mathbf u\times\mathbf v\|_2
=
\|\mathbf u\|_2\|\mathbf v\|_2\sin\theta
}.
\]

If the input components have physical units, the cross-product
components have the product of those units. The coordinate definition
also gives the useful identities

\[
\boxed{
\mathbf u\times\mathbf v
=
-\mathbf v\times\mathbf u,
\qquad
\mathbf u\times\mathbf u
=
\mathbf 0_3
}.
\]

Cross products are not associative:
\(\mathbf a\times(\mathbf b\times\mathbf c)\) and
\((\mathbf a\times\mathbf b)\times\mathbf c\) are generally different.
For \(\mathbf a,\mathbf b,\mathbf c\in\mathbb R^3\) expressed in the
same right-handed orthonormal frame, the \textbf{vector triple-product
identity} is

\[
\boxed{
\mathbf a\times(\mathbf b\times\mathbf c)
=
\mathbf b(\mathbf a^\mathsf T\mathbf c)
-
\mathbf c(\mathbf a^\mathsf T\mathbf b)
}.
\]

To verify the identity from the cross-product definition, set
\(\mathbf d\mathrel{:=}\mathbf b\times\mathbf c\). Substituting the
components of \(\mathbf d\) into \(\mathbf a\times\mathbf d\) and
regrouping by the components of \(\mathbf b\) and \(\mathbf c\) gives

\[
\begin{aligned}
\mathbf a\times\mathbf d
&=
\begin{bmatrix}
a_2(b_1c_2-b_2c_1)-a_3(b_3c_1-b_1c_3)\\
a_3(b_2c_3-b_3c_2)-a_1(b_1c_2-b_2c_1)\\
a_1(b_3c_1-b_1c_3)-a_2(b_2c_3-b_3c_2)
\end{bmatrix}\\
&=
\begin{bmatrix}
b_1(\mathbf a^\mathsf T\mathbf c)-c_1(\mathbf a^\mathsf T\mathbf b)\\
b_2(\mathbf a^\mathsf T\mathbf c)-c_2(\mathbf a^\mathsf T\mathbf b)\\
b_3(\mathbf a^\mathsf T\mathbf c)-c_3(\mathbf a^\mathsf T\mathbf b)
\end{bmatrix}\\
&=
\mathbf b(\mathbf a^\mathsf T\mathbf c)
-
\mathbf c(\mathbf a^\mathsf T\mathbf b).
\end{aligned}
\]

A case used later for rotations sets \(\mathbf b=\mathbf a\) and
\(\mathbf c=\mathbf y\):

\[
\boxed{
\mathbf a\times(\mathbf a\times\mathbf y)
=
\mathbf a(\mathbf a^\mathsf T\mathbf y)
-
(\mathbf a^\mathsf T\mathbf a)\mathbf y
}.
\]

If \(\mathbf a\) is a unit vector, then
\(\mathbf a^\mathsf T\mathbf a=1\).

\subsection{Linear combinations, span, and
basis}\label{linear-combinations-span-and-basis}

Let \(\mathbf v_1,\ldots,\mathbf v_k\in\mathbb R^n\) and let
\(\alpha_1,\ldots,\alpha_k\in\mathbb R\). A \textbf{linear combination}
of these vectors is any vector of the form

\[
\alpha_1\mathbf v_1+\cdots+\alpha_k\mathbf v_k.
\]

The \textbf{span} of \(\mathbf v_1,\ldots,\mathbf v_k\) is the set of
every linear combination that can be formed from them:

\[
\operatorname{span}\{\mathbf v_1,\ldots,\mathbf v_k\}
\mathrel{:=}
\left\{
\sum_{j=1}^{k}\alpha_j\mathbf v_j:
\alpha_1,\ldots,\alpha_k\in\mathbb R
\right\}.
\]

The vectors are \textbf{linearly independent} if

\[
\alpha_1\mathbf v_1+\cdots+\alpha_k\mathbf v_k=\mathbf 0_n
\]

implies \(\alpha_1=\cdots=\alpha_k=0\), where
\(\mathbf 0_n\in\mathbb R^n\) is the zero vector. If a nonzero set of
coefficients produces \(\mathbf 0_n\), the vectors are linearly
dependent.

A \textbf{basis} of \(\mathbb R^n\) is a set of \(n\) linearly
independent vectors that spans \(\mathbb R^n\). If
\(\mathbf e_1,\ldots,\mathbf e_n\) is a basis, every
\(\mathbf x\in\mathbb R^n\) has a unique expansion

\[
\mathbf x=x_1\mathbf e_1+\cdots+x_n\mathbf e_n.
\]

The scalar coefficients \(x_1,\ldots,x_n\in\mathbb R\) form the
coordinate column of \(\mathbf x\) in that basis. Uniqueness is
essential: if the basis vectors were dependent, the same vector could
have more than one coefficient description.

The \textbf{standard basis} of \(\mathbb R^n\) is
\(\mathbf e_1,\ldots,\mathbf e_n\), where \(\mathbf e_i\) has entry
\(1\) in position \(i\) and entry \(0\) in every other position. In the
standard basis,

\[
\mathbf x=x_1\mathbf e_1+\cdots+x_n\mathbf e_n
\quad\Longleftrightarrow\quad
\mathbf x=
\begin{bmatrix}
x_1&\cdots&x_n
\end{bmatrix}^{\mathsf T}.
\]

\section*{Quick test A: vectors and
bases}\label{quick-test-a-vectors-and-bases}
\addcontentsline{toc}{section}{Quick test A: vectors and bases}

\markright{Quick test A: vectors and bases}

Complete this test without looking back at the section.

\begin{enumerate}
\def\labelenumi{\arabic{enumi}.}
\tightlist
\item
  \textbf{Recall:} What is the difference between a geometric vector
  \(\mathbf v\) and a coordinate column \({}^{a}\mathbf v\)?
\item
  \textbf{Explanation:} Why must coordinate columns be associated with a
  basis?
\item
  \textbf{Derivation:} Use the law of cosines and an algebraic expansion
  of \(\|\mathbf u-\mathbf v\|_2^2\) to derive
  \({}^{a}\mathbf u^{\mathsf T}{}^{a}\mathbf v=\|\mathbf u\|_2\|\mathbf v\|_2\cos\theta\).
\item
  \textbf{Application:} Find the norm and unit direction of
  \(\mathbf x=[3\;\;4]^\mathsf T\).
\item
  \textbf{Limitations:} Why is the angle formula undefined when either
  vector is zero?
\item
  \textbf{Basis:} Explain why two parallel nonzero vectors cannot form a
  basis of \(\mathbb R^2\).
\item
  \textbf{Cross product:} For a unit vector \(\mathbf a\in\mathbb R^3\),
  use the vector triple-product identity to simplify
  \(\mathbf a\times(\mathbf a\times\mathbf y)\).
\end{enumerate}

\subsection*{Answers and hints}\label{answers-and-hints}
\addcontentsline{toc}{subsection}{Answers and hints}

\begin{enumerate}
\def\labelenumi{\arabic{enumi}.}
\tightlist
\item
  \(\mathbf v\) is the coordinate-independent geometric object;
  \({}^{a}\mathbf v\) is its numerical representation in basis or frame
  \(\{a\}\).
\item
  The entries are coefficients multiplying particular basis vectors.
  Without the basis, the numbers do not determine physical or geometric
  directions.
\item
  Equate the law-of-cosines expression for
  \(\|\mathbf u-\mathbf v\|_2^2\) with its dot-product expansion and
  cancel the common norm terms.
\item
  \(\|\mathbf x\|_2=5\) and the unit direction is
  \([0.6\;\;0.8]^\mathsf T\).
\item
  The denominator \(\|\mathbf u\|_2\|\mathbf v\|_2\) becomes zero, and a
  zero vector has no direction.
\item
  One vector is a scalar multiple of the other, so their span is only a
  line and their coefficient representation is not unique.
\item
  \(\mathbf a\times(\mathbf a\times\mathbf y)=\mathbf a(\mathbf a^\mathsf T\mathbf y)-\mathbf y\)
  because \(\mathbf a^\mathsf T\mathbf a=1\).
\end{enumerate}

\section{Matrices: linear maps and basis
changes}\label{matrices-linear-maps-and-basis-changes}

\subsection{Matrix dimensions}\label{matrix-dimensions}

A matrix \(\mathbf{A}\in\mathbb{R}^{m\times n}\) maps an \(n\)-element
coordinate vector to an \(m\)-element coordinate vector:

\[
\mathbf{y}=\mathbf{A}\mathbf{x}
\]

where \(\mathbf{y} \in \mathbb{R}^m\) and
\(\mathbf{x} \in \mathbb{R}^n\).

If

\[
\mathbf{A}=
\begin{bmatrix}
a_{11} & a_{12} \\
a_{21} & a_{22}
\end{bmatrix},
\qquad
\mathbf{x}=
\begin{bmatrix}
x_1\\
x_2
\end{bmatrix},
\]

then

\[
\mathbf{A}\mathbf{x}
=
\begin{bmatrix}
a_{11}x_1+a_{12}x_2\\
a_{21}x_1+a_{22}x_2
\end{bmatrix}.
\]

Equivalently, if \(\mathbf{a}_1\) and \(\mathbf{a}_2\) are the columns
of \(\mathbf{A}\),

\[
\mathbf{A}\mathbf{x}=x_1\mathbf{a}_1+x_2\mathbf{a}_2.
\]

This column interpretation is central to rotation matrices: their
columns are the axes of one frame expressed in another frame.

\subsection{Identity, transpose, null space, rank, and
inverse}\label{identity-transpose-null-space-rank-and-inverse}

Let \(\mathbf{I}_n\) denote the \(n\times n\) identity matrix. It leaves
compatible coordinates unchanged:

\[
\mathbf{I}_n\mathbf{x}=\mathbf{x}.
\]

The transpose swaps rows and columns:

\[
\left(\mathbf{A}^\mathsf{T}\right)_{ij}=a_{ji}.
\]

Let \(\mathbf A\in\mathbb R^{m\times n}\). Its \textbf{column space},
denoted \(\operatorname{Col}(\mathbf A)\), is the set of all linear
combinations of its columns. The \textbf{rank}, denoted
\(\operatorname{rank}(\mathbf A)\), is the number of linearly
independent columns, equivalently the dimension of
\(\operatorname{Col}(\mathbf A)\).

The \textbf{null space} of \(\mathbf A\) is

\[
\mathcal N(\mathbf A)
\mathrel{:=}
\left\{\mathbf x\in\mathbb R^n:\mathbf A\mathbf x=\mathbf 0_m\right\},
\]

where \(\mathbf 0_m\in\mathbb R^m\) is the zero vector. The null space
contains all input directions that the mapping sends to zero.

\subsubsection{Rank-nullity
theorem}\label{sec-linear-algebra-rank-nullity}

The \textbf{nullity} of \(\mathbf A\) is the dimension of its null
space:

\[
\operatorname{nullity}(\mathbf A)
\mathrel{:=}
\dim\mathcal N(\mathbf A).
\]

For every matrix \(\mathbf A\in\mathbb R^{m\times n}\), the
\textbf{rank-nullity theorem} states

\[
\boxed{
\operatorname{rank}(\mathbf A)
+
\operatorname{nullity}(\mathbf A)
=
n
}.
\]

The number \(n\) is the dimension of the input space \(\mathbb R^n\).
The theorem says that its independent input directions divide into two
counts: \(\operatorname{nullity}(\mathbf A)\) directions that the
mapping sends to zero and \(\operatorname{rank}(\mathbf A)\) directions
that produce independent output directions.

To prove the theorem, let

\[
k
\mathrel{:=}
\operatorname{nullity}(\mathbf A),
\]

and choose a basis \(\mathbf n_1,\ldots,\mathbf n_k\in\mathbb R^n\) for
\(\mathcal N(\mathbf A)\). Because these vectors are linearly
independent, they can be extended to a basis of the whole input space by
repeatedly adding vectors outside their current span. Write the
resulting basis as

\[
\mathbf n_1,\ldots,\mathbf n_k,
\mathbf v_1,\ldots,\mathbf v_{n-k}.
\]

The output vectors

\[
\mathbf A\mathbf v_1,\ldots,\mathbf A\mathbf v_{n-k}
\]

span \(\operatorname{Col}(\mathbf A)\). To see this, express any input
\(\mathbf x\in\mathbb R^n\) in the chosen input basis:

\[
\mathbf x
=
\sum_{i=1}^{k}a_i\mathbf n_i
+
\sum_{j=1}^{n-k}b_j\mathbf v_j.
\]

Because every \(\mathbf n_i\) belongs to the null space,
\(\mathbf A\mathbf n_i=\mathbf0_m\). Applying \(\mathbf A\) therefore
gives

\[
\mathbf A\mathbf x
=
\sum_{j=1}^{n-k}b_j\mathbf A\mathbf v_j,
\]

so every possible output is a linear combination of the listed output
vectors.

These output vectors are also linearly independent. If

\[
\sum_{j=1}^{n-k}c_j\mathbf A\mathbf v_j
=
\mathbf0_m,
\]

then

\[
\mathbf A
\left(
\sum_{j=1}^{n-k}c_j\mathbf v_j
\right)
=
\mathbf0_m.
\]

Thus, \(\sum_j c_j\mathbf v_j\) belongs to \(\mathcal N(\mathbf A)\) and
can also be written as a linear combination of
\(\mathbf n_1,\ldots,\mathbf n_k\). The complete list
\(\mathbf n_1,\ldots,\mathbf n_k,\mathbf v_1,\ldots,\mathbf v_{n-k}\) is
a basis and is therefore linearly independent, so this is possible only
when every \(c_j=0\). Hence, the \(n-k\) vectors
\(\mathbf A\mathbf v_j\) form a basis of
\(\operatorname{Col}(\mathbf A)\), and

\[
\operatorname{rank}(\mathbf A)
=
n-k.
\]

Since \(k=\operatorname{nullity}(\mathbf A)\),

\[
\operatorname{rank}(\mathbf A)
+
\operatorname{nullity}(\mathbf A)
=
(n-k)+k
=
n,
\]

which proves the theorem.

For example, let

\[
\mathbf A
=
\begin{bmatrix}
1&0&1\\
0&1&1
\end{bmatrix}
\in\mathbb R^{2\times3}.
\]

The first two columns are independent and the third is their sum, so
\(\operatorname{rank}(\mathbf A)=2\). Solving
\(\mathbf A\mathbf x=\mathbf0_2\) gives

\[
\mathcal N(\mathbf A)
=
\left\{
\alpha
\begin{bmatrix}
-1\\-1\\1
\end{bmatrix}
\;\middle|\;
\alpha\in\mathbb R
\right\},
\]

so \(\operatorname{nullity}(\mathbf A)=1\). The input dimension is
therefore partitioned as

\[
\underbrace{3}_{\text{input dimension}}
=
\underbrace{2}_{\text{rank}}
+
\underbrace{1}_{\text{nullity}}.
\]

\subsubsection{Injective mappings and unique
recovery}\label{injective-mappings-and-unique-recovery}

Let \(\mathbf A\in\mathbb R^{m\times n}\) define the linear map

\[
T_{\mathbf A}:\mathbb R^n\rightarrow\mathbb R^m,
\qquad
T_{\mathbf A}(\mathbf x)=\mathbf A\mathbf x.
\]

The map \(T_{\mathbf A}\) is \textbf{injective}, or \textbf{one-to-one},
if equal outputs can occur only when the inputs are equal. Formally, for
every \(\mathbf x^{(1)},\mathbf x^{(2)}\in\mathbb R^n\), where the
parenthesised superscripts \((1)\) and \((2)\) label two inputs and are
not exponents,

\[
\mathbf A\mathbf x^{(1)}=\mathbf A\mathbf x^{(2)}
\quad\Longrightarrow\quad
\mathbf x^{(1)}=\mathbf x^{(2)}.
\]

Equivalently, distinct inputs must produce distinct outputs:

\[
\mathbf x^{(1)}\neq\mathbf x^{(2)}
\quad\Longrightarrow\quad
\mathbf A\mathbf x^{(1)}\neq\mathbf A\mathbf x^{(2)}.
\]

The map is \textbf{not injective} if at least one pair of distinct
inputs produces the same output:

\[
\exists\,\mathbf x^{(1)},\mathbf x^{(2)}\in\mathbb R^n:
\quad
\mathbf x^{(1)}\neq\mathbf x^{(2)},
\qquad
\mathbf A\mathbf x^{(1)}=\mathbf A\mathbf x^{(2)}.
\]

To connect injectivity to the null space, suppose two inputs have the
same output. Subtracting the outputs gives

\[
\begin{aligned}
\mathbf A\mathbf x^{(1)}&=\mathbf A\mathbf x^{(2)}\\
\Longrightarrow\quad
\mathbf A(\mathbf x^{(1)}-\mathbf x^{(2)})&=\mathbf0_m.
\end{aligned}
\]

If \(\mathcal N(\mathbf A)=\{\mathbf0_n\}\), then
\(\mathbf x^{(1)}-\mathbf x^{(2)}=\mathbf0_n\), so the inputs must be
equal. Conversely, if a nonzero vector
\(\mathbf z\in\mathcal N(\mathbf A)\) exists, then for any
\(\mathbf x\in\mathbb R^n\),

\[
\mathbf A(\mathbf x+\mathbf z)
=\mathbf A\mathbf x+\mathbf A\mathbf z
=\mathbf A\mathbf x.
\]

Because \(\mathbf z\neq\mathbf0_n\), the inputs \(\mathbf x\) and
\(\mathbf x+\mathbf z\) are different but produce the same output.
Therefore,

\[
\boxed{
T_{\mathbf A}\text{ is injective}
\quad\Longleftrightarrow\quad
\mathcal N(\mathbf A)=\{\mathbf0_n\}
\quad\Longleftrightarrow\quad
\mathbf A\text{ has full column rank}
}.
\]

For a square matrix, these conditions are also equivalent to
nonsingularity and invertibility. If the map is injective, its output
uniquely identifies its input; if it is not injective, some input
information cannot be recovered from the output.

Let \(\mathbf A\in\mathbb R^{n\times n}\) be a square matrix with
columns \(\mathbf a_1,\ldots,\mathbf a_n\in\mathbb R^n\), and let
\(\mathbf x=[x_1\;\cdots\;x_n]^\mathsf T\in\mathbb R^n\). Matrix-vector
multiplication forms the linear combination

\[
\mathbf A\mathbf x
=
x_1\mathbf a_1+\cdots+x_n\mathbf a_n
=
\sum_{j=1}^{n}x_j\mathbf a_j.
\]

For every matrix, \(\mathbf x=\mathbf 0\) implies
\(\mathbf A\mathbf x=\mathbf 0\). The square matrix \(\mathbf A\) is
\textbf{nonsingular} when the converse is also true:

\[
\mathbf A\mathbf x=\mathbf 0
\quad\Longleftrightarrow\quad
\mathbf x=\mathbf 0.
\]

Equivalently,

\[
\mathbf x\neq\mathbf 0
\quad\Longrightarrow\quad
\mathbf A\mathbf x\neq\mathbf 0.
\]

Because \(\mathbf A\mathbf x=\sum_{j=1}^{n}x_j\mathbf a_j\),
nonsingularity means that

\[
x_1\mathbf a_1+\cdots+x_n\mathbf a_n=\mathbf 0
\]

only when \(x_1=\cdots=x_n=0\). This is precisely the definition that
the columns \(\mathbf a_1,\ldots,\mathbf a_n\) are linearly independent.

A square matrix that is not nonsingular is called \textbf{singular}. For
a singular matrix, at least one nonzero \(\mathbf x\) satisfies
\(\mathbf A\mathbf x=\mathbf 0\), so the mapping loses information about
that input direction.

More generally, if \(\mathbf A\in\mathbb R^{m\times n}\), then its
columns are \(\mathbf a_1,\ldots,\mathbf a_n\in\mathbb R^m\) and
\(\mathbf A\mathbf x=\sum_{j=1}^{n}x_j\mathbf a_j\). The condition
\(\mathbf A\mathbf x=\mathbf 0_m\Rightarrow\mathbf x=\mathbf 0_n\) means
that \(\mathcal N(\mathbf A)=\{\mathbf 0_n\}\) and that all \(n\)
columns are linearly independent. In that case
\(\operatorname{rank}(\mathbf A)=n\), which is called \textbf{full
column rank}. The term \emph{nonsingular} is normally reserved for the
square case \(m=n\).

Only a nonsingular square matrix has an inverse
\(\mathbf A^{-1}\in\mathbb R^{n\times n}\). The inverse reverses the
action of \(\mathbf A\) and satisfies

\[
\mathbf{A}^{-1}\mathbf{A}=\mathbf{A}\mathbf{A}^{-1}=\mathbf{I}_n.
\]

\subsection{Multiplication order}\label{multiplication-order}

Let \(n\) be a positive integer. Let
\(\mathbf A,\mathbf B,\mathbf C\in\mathbb R^{n\times n}\) be three
matrices, and let \(\mathbf x,\mathbf y,\mathbf z\in\mathbb R^n\) be
coordinate vectors. Define the intermediate vector \(\mathbf z\) and
output vector \(\mathbf y\) by

\[
\mathbf z=\mathbf B\mathbf x
\qquad\text{and}\qquad
\mathbf y=\mathbf A\mathbf z.
\]

Substituting the expression for \(\mathbf z\) into the expression for
\(\mathbf y\) gives

\[
\mathbf y
=\mathbf A(\mathbf B\mathbf x)
=(\mathbf A\mathbf B)\mathbf x.
\]

Thus, the rightmost matrix acts first: \(\mathbf B\) maps \(\mathbf x\)
to \(\mathbf z\), and then \(\mathbf A\) maps \(\mathbf z\) to
\(\mathbf y\).

The third matrix \(\mathbf C\) allows three transformations to be
composed. Matrix multiplication is \textbf{associative}, meaning that
changing the grouping does not change the product:

\[
(\mathbf A\mathbf B)\mathbf C
=
\mathbf A(\mathbf B\mathbf C).
\]

Both sides are matrices in \(\mathbb R^{n\times n}\). However, matrix
multiplication is generally \textbf{not commutative}, meaning that
changing the order can change the result:

\[
\mathbf A\mathbf B\neq\mathbf B\mathbf A.
\]

This order becomes critical when several coordinate transformations are
composed.

\section*{Quick test B: matrices, rank, and
inverses}\label{quick-test-b-matrices-rank-and-inverses}
\addcontentsline{toc}{section}{Quick test B: matrices, rank, and
inverses}

\markright{Quick test B: matrices, rank, and inverses}

\begin{enumerate}
\def\labelenumi{\arabic{enumi}.}
\tightlist
\item
  \textbf{Recall:} For \(\mathbf A\in\mathbb R^{m\times n}\), what are
  the dimensions of an input \(\mathbf x\) and output
  \(\mathbf y=\mathbf A\mathbf x\)?
\item
  \textbf{Explanation:} What does column \(j\) of \(\mathbf A\) tell us
  about the mapping?
\item
  \textbf{Derivation:} Explain why \(\mathbf A\mathbf x=\mathbf 0\) has
  only the zero solution exactly when the columns of \(\mathbf A\) are
  linearly independent.
\item
  \textbf{Application:} Determine the null space and rank of
  \(\mathbf A=\begin{bmatrix}1&2\\2&4\end{bmatrix}\).
\item
  \textbf{Limitations:} Why can a nonsquare matrix not have a two-sided
  inverse of the same kind as a nonsingular square matrix?
\item
  \textbf{Composition:} If \(\mathbf z=\mathbf B\mathbf x\) and
  \(\mathbf y=\mathbf A\mathbf z\), which matrix acts first in
  \(\mathbf y=\mathbf A\mathbf B\mathbf x\)?
\item
  \textbf{Injectivity:} Define an injective matrix mapping and explain
  why it is equivalent to \(\mathcal N(\mathbf A)=\{\mathbf0\}\).
\item
  \textbf{Rank-nullity:} State the rank-nullity theorem for
  \(\mathbf A\in\mathbb R^{m\times n}\) and verify it for the matrix in
  Question 4.
\end{enumerate}

\subsection*{Answers and hints}\label{answers-and-hints-1}
\addcontentsline{toc}{subsection}{Answers and hints}

\begin{enumerate}
\def\labelenumi{\arabic{enumi}.}
\tightlist
\item
  \(\mathbf x\in\mathbb R^n\) and \(\mathbf y\in\mathbb R^m\).
\item
  It is the output produced by applying \(\mathbf A\) to the
  corresponding standard basis input; general outputs are linear
  combinations of the columns.
\item
  The equation is the linear combination
  \(x_1\mathbf a_1+\cdots+x_n\mathbf a_n=\mathbf 0\). Only zero
  coefficients solve it precisely when the columns are independent.
\item
  The second column is twice the first, so the rank is \(1\). Solving
  \(x_1+2x_2=0\) gives
  \(\mathcal N(\mathbf A)=\{t[-2\;\;1]^\mathsf T:t\in\mathbb R\}\).
\item
  A two-sided inverse would have to undo mappings in both domain and
  codomain, which requires equal dimensions and no lost directions.
\item
  \(\mathbf B\) acts first, followed by \(\mathbf A\).
\item
  A map is injective if
  \(\mathbf A\mathbf x^{(1)}=\mathbf A\mathbf x^{(2)}\) implies
  \(\mathbf x^{(1)}=\mathbf x^{(2)}\). If
  \(\mathcal N(\mathbf A)=\{\mathbf0\}\), subtracting equal outputs
  gives \(\mathbf A(\mathbf x^{(1)}-\mathbf x^{(2)})=\mathbf0\), so
  \(\mathbf x^{(1)}-\mathbf x^{(2)}=\mathbf0\). Conversely, a nonzero
  null vector \(\mathbf z\) makes \(\mathbf x\) and
  \(\mathbf x+\mathbf z\) distinct inputs with the same output.
\item
  The theorem is
  \(\operatorname{rank}(\mathbf A)+\operatorname{nullity}(\mathbf A)=n\).
  The matrix in Question 4 has two input columns, rank one, and a
  one-dimensional null space, so \(1+1=2\).
\end{enumerate}

\section{Determinants and signed
volume}\label{determinants-and-signed-volume}

\subsection{What the determinant
measures}\label{what-the-determinant-measures}

The determinant is defined for a square matrix
\(\mathbf A\in\mathbb R^{n\times n}\) and returns a scalar:

\[
\det:\mathbb R^{n\times n}\rightarrow\mathbb R.
\]

Geometrically, \(|\det(\mathbf A)|\) is the factor by which the linear
map \(\mathbf x\mapsto\mathbf A\mathbf x\) scales \(n\)-dimensional
volume. The sign records whether the ordered basis orientation is
preserved or reversed.

For
\(\mathbf A=\begin{bmatrix}a&b\\c&d\end{bmatrix}\in\mathbb R^{2\times2}\),

\[
\det(\mathbf A)=ad-bc.
\]

Its magnitude is the area of the parallelogram formed by the columns of
\(\mathbf A\). For
\(\mathbf A=[\mathbf a_1\;\mathbf a_2\;\mathbf a_3]\in\mathbb R^{3\times3}\),
the determinant is the scalar triple product

\[
\det(\mathbf A)
=
\mathbf a_1^{\mathsf T}(\mathbf a_2\times\mathbf a_3).
\]

Its magnitude is the volume of the parallelepiped formed by the columns.

\subsubsection{Determinant sign, the right-hand rule, and
reflection}\label{determinant-sign-the-right-hand-rule-and-reflection}

The columns of \(\mathbf A=[\mathbf a_1\;\cdots\;\mathbf a_n]\) are the
images of the standard basis vectors because
\(\mathbf A\mathbf e_i=\mathbf a_i\). The determinant sign therefore
compares the orientation of the ordered transformed columns with the
orientation of the ordered standard basis.

In three dimensions, the standard basis is right-handed:

\[
\mathbf e_1\times\mathbf e_2=\mathbf e_3.
\]

For
\(\mathbf A=[\mathbf a_1\;\mathbf a_2\;\mathbf a_3]\in\mathbb R^{3\times3}\),
the determinant can be written as

\[
\det(\mathbf A)
=
(\mathbf a_1\times\mathbf a_2)^{\mathsf T}\mathbf a_3.
\]

The cross product \(\mathbf a_1\times\mathbf a_2\) points in the
right-hand-rule normal direction determined by the first two columns.
Consequently:

\begin{itemize}
\tightlist
\item
  If \(\det(\mathbf A)>0\), then \(\mathbf a_3\) has a positive
  component in that right-hand-rule direction. The ordered columns
  \((\mathbf a_1,\mathbf a_2,\mathbf a_3)\) are right-handed, so the
  transformation preserves orientation.
\item
  If \(\det(\mathbf A)<0\), then \(\mathbf a_3\) points to the opposite
  side of the plane spanned by \(\mathbf a_1\) and \(\mathbf a_2\). The
  ordered columns are left-handed, so the transformation reverses
  orientation.
\item
  If \(\det(\mathbf A)=0\), then the three columns are coplanar or
  otherwise linearly dependent. The three-dimensional volume collapses
  to zero, so a full three-dimensional orientation is no longer defined.
\end{itemize}

A reflection reverses orientation. For example, reflection across the
\(y\)-\(z\) plane is represented by

\[
\mathbf S
=
\begin{bmatrix}
-1&0&0\\
0&1&0\\
0&0&1
\end{bmatrix},
\qquad
\det(\mathbf S)=-1.
\]

This transformation sends \(\mathbf e_1\) to \(-\mathbf e_1\) while
leaving \(\mathbf e_2\) and \(\mathbf e_3\) unchanged. The transformed
basis is left-handed because

\[
(-\mathbf e_1)\times\mathbf e_2=-\mathbf e_3,
\]

whereas its third basis vector is \(+\mathbf e_3\). Thus, the negative
determinant detects the orientation reversal introduced by the
reflection.

A matrix with negative determinant need not be a pure reflection: it may
also stretch, scale, or shear. The negative sign means that its total
action contains an orientation-reversing, or reflection-like, component.
By contrast, a proper rotation preserves right-handedness and has
determinant \(+1\).

\subsection{Algebraic definition}\label{algebraic-definition}

The determinant is the unique scalar-valued function of the columns of a
square matrix with three defining properties:

\begin{enumerate}
\def\labelenumi{\arabic{enumi}.}
\tightlist
\item
  \textbf{Multilinearity:} it is linear in each column while the other
  columns are held fixed.
\item
  \textbf{Alternation:} swapping two columns changes the sign;
  consequently, a matrix with repeated columns has determinant zero.
\item
  \textbf{Normalisation:} \(\det(\mathbf I_n)=1\).
\end{enumerate}

\subsubsection{Multilinearity in one
column}\label{multilinearity-in-one-column}

Write the columns of \(\mathbf A\in\mathbb R^{n\times n}\) as
\(\mathbf A=[\mathbf a_1\;\cdots\;\mathbf a_n]\), where each
\(\mathbf a_i\in\mathbb R^n\). Choose a column index
\(j\in\{1,\ldots,n\}\) and hold every column except column \(j\) fixed.
For any vectors \(\mathbf u,\mathbf v\in\mathbb R^n\) and scalars
\(\alpha,\beta\in\mathbb R\), linearity in column \(j\) means

\[
\begin{aligned}
&\det(\mathbf a_1,\ldots,\mathbf a_{j-1},\alpha\mathbf u+\beta\mathbf v,\mathbf a_{j+1},\ldots,\mathbf a_n)\\
&\quad=
\alpha\det(\mathbf a_1,\ldots,\mathbf a_{j-1},\mathbf u,\mathbf a_{j+1},\ldots,\mathbf a_n)\\
&\qquad+
\beta\det(\mathbf a_1,\ldots,\mathbf a_{j-1},\mathbf v,\mathbf a_{j+1},\ldots,\mathbf a_n).
\end{aligned}
\]

Thus, a sum placed in one column can be separated into a sum of
determinants, and a scalar multiplying one column can be moved outside
the determinant. The word \textbf{multilinear} means that this linearity
property holds separately for every column. It does not mean that the
determinant is linear in the whole matrix at once; in general,

\[
\det(\mathbf A+\mathbf B)\neq\det(\mathbf A)+\det(\mathbf B).
\]

For a two-dimensional check, let \(\mathbf u=[u_1\;u_2]^\mathsf T\),
\(\mathbf v=[v_1\;v_2]^\mathsf T\), and
\(\mathbf w=[w_1\;w_2]^\mathsf T\). Linearity in the first column
follows directly from

\[
\begin{aligned}
\det\!\begin{bmatrix}\alpha u_1+\beta v_1&w_1\\\alpha u_2+\beta v_2&w_2\end{bmatrix}
&=(\alpha u_1+\beta v_1)w_2-w_1(\alpha u_2+\beta v_2)\\
&=\alpha(u_1w_2-w_1u_2)+\beta(v_1w_2-w_1v_2)\\
&=\alpha\det[\mathbf u\;\mathbf w]+\beta\det[\mathbf v\;\mathbf w].
\end{aligned}
\]

\subsubsection{Alternation and repeated
columns}\label{alternation-and-repeated-columns}

Let \(\mathbf A'\) be obtained from \(\mathbf A\) by swapping columns
\(i\) and \(j\), where \(i\neq j\). Alternation means

\[
\det(\mathbf A')=-\det(\mathbf A).
\]

A single swap reverses the orientation of the ordered columns, so it
reverses the sign of the signed volume without changing its magnitude.

Now suppose columns \(i\) and \(j\) are equal to the same vector
\(\mathbf c\in\mathbb R^n\). Let the determinant of this matrix be
\(D\). Swapping the two equal columns leaves the matrix numerically
unchanged, so its determinant is still \(D\). Alternation simultaneously
requires the determinant to become \(-D\). Therefore,

\[
D=-D
\quad\Longrightarrow\quad
2D=0
\quad\Longrightarrow\quad
D=0.
\]

Hence, every matrix with two identical columns has determinant zero.
Geometrically, two identical column directions cannot span an
\(n\)-dimensional parallelepiped, so its \(n\)-dimensional volume
collapses to zero.

Equivalently, if \(\mathbf A=[a_{ij}]\in\mathbb R^{n\times n}\), then

\[
\det(\mathbf A)
=
\sum_{\sigma\in S_n}
\operatorname{sgn}(\sigma)
\prod_{j=1}^{n}a_{\sigma(j),j}.
\]

Here, \(S_n\) is the set of all permutations of the indices
\(1,\ldots,n\). The sign \(\operatorname{sgn}(\sigma)\) equals \(+1\)
when the permutation can be produced by an even number of pairwise swaps
and \(-1\) when it requires an odd number.

If the columns of \(\mathbf A\) are linearly dependent, one column can
be written as a linear combination of the others. Multilinearity and
alternation then force \(\det(\mathbf A)=0\). Conversely,
\(\det(\mathbf A)=0\) means that the signed volume has collapsed, so the
columns do not span \(\mathbb R^n\). Therefore,

\[
\boxed{
\det(\mathbf A)\neq0
\quad\Longleftrightarrow\quad
\mathbf A\text{ is nonsingular}
}.
\]

\subsubsection{Example: a singular matrix loses one input
component}\label{example-a-singular-matrix-loses-one-input-component}

Consider the linear map \(\mathbf x\mapsto\mathbf A\mathbf x\) from
\(\mathbb R^2\) to \(\mathbb R^2\), where

\[
\mathbf A
=
\begin{bmatrix}
1&0\\
0&0
\end{bmatrix}.
\]

Its determinant is

\[
\det(\mathbf A)=0.
\]

For an input coordinate column

\[
\mathbf x
=
\begin{bmatrix}
x_1\\
x_2
\end{bmatrix}
\in\mathbb R^2,
\]

the output is

\[
\mathbf A\mathbf x
=
\begin{bmatrix}
x_1\\
0
\end{bmatrix}
\in\mathbb R^2.
\]

The output retains \(x_1\) but completely removes \(x_2\). For example,

\[
\mathbf A
\begin{bmatrix}
3\\
1
\end{bmatrix}
=
\begin{bmatrix}
3\\
0
\end{bmatrix},
\qquad
\mathbf A
\begin{bmatrix}
3\\
7
\end{bmatrix}
=
\begin{bmatrix}
3\\
0
\end{bmatrix}.
\]

The two inputs are different, but their outputs are identical. From the
output \([3\;0]^{\mathsf T}\), we can recover \(x_1=3\), but we cannot
determine whether \(x_2\) was \(1\), \(7\), or any other real number. In
fact,

\[
\mathcal N(\mathbf A)
=
\left\{
 t
 \begin{bmatrix}
 0\\
 1
 \end{bmatrix}
 :t\in\mathbb R
\right\},
\]

so the entire input direction \([0\;1]^{\mathsf T}\) is lost.

Geometrically, \(\mathbf A\) maps the entire two-dimensional plane onto
the one-dimensional \(x\)-axis:

\[
\mathbb R^2
\longrightarrow
\left\{
\begin{bmatrix}
y_1\\
0
\end{bmatrix}
:y_1\in\mathbb R
\right\}.
\]

Every two-dimensional region is flattened into a line segment, whose
two-dimensional area is zero. This agrees with \(\det(\mathbf A)=0\).

Recall that \textbf{not injective} means that there exist at least two
distinct inputs \(\mathbf x^{(1)},\mathbf x^{(2)}\in\mathbb R^n\) with
the same output:

\[
\mathbf x^{(1)}\neq\mathbf x^{(2)},
\qquad
\mathbf A\mathbf x^{(1)}=\mathbf A\mathbf x^{(2)}.
\]

The statement does not mean that every pair of inputs has the same
output; it means that at least one distinct pair cannot be distinguished
from the output.

If \(\mathbf A\) is singular, then there exists
\(\mathbf z \in \mathcal N(\mathbf A)\) such that
\(\mathbf A \mathbf z = \mathbf 0\). This implies
\(\mathbf A(\mathbf x+\mathbf z) = \mathbf A\mathbf x\) for all
\(\mathbf x \in \mathbb R^n\). Hence, a singular matrix is not injective
and different inputs can produce the same output.

For a square matrix \(\mathbf A\in\mathbb R^{n\times n}\), the general
equivalences are

\[
\boxed{
\begin{aligned}
\det(\mathbf A)=0
&\Longleftrightarrow \mathbf A\text{ is singular}\\
&\Longleftrightarrow \mathcal N(\mathbf A)\neq\{\mathbf0\}\\
&\Longleftrightarrow \mathbf A\text{ is not injective}\\
&\Longleftrightarrow \text{different inputs can produce the same output}\\
&\Longleftrightarrow \text{the input cannot be uniquely recovered}.
\end{aligned}
}
\]

\subsection{Why determinants multiply}\label{why-determinants-multiply}

Let \(\mathbf A,\mathbf B\in\mathbb R^{n\times n}\). The product
\(\mathbf A\mathbf B\) represents the composition

\[
\mathbf x \mapsto \mathbf B\mathbf x \mapsto \mathbf A\mathbf B\mathbf x.
\]

Geometrically, \(\mathbf B\) first scales signed volume by
\(\det(\mathbf B)\), and \(\mathbf A\) then scales the result by
\(\det(\mathbf A)\). This suggests

\[
\det(\mathbf A\mathbf B)=\det(\mathbf A)\det(\mathbf B).
\]

A formal proof follows from the defining properties of the determinant.
Hold \(\mathbf A\) fixed and define

\[
F:\mathbb R^{n\times n}\rightarrow\mathbb R,
\qquad
F(\mathbf B)\mathrel{:=}\det(\mathbf A\mathbf B).
\]

Write \(\mathbf B=[\mathbf b_1\;\cdots\;\mathbf b_n]\), where each
\(\mathbf b_j\in\mathbb R^n\). Then

\[
\mathbf A\mathbf B
=
\begin{bmatrix}
\mathbf A\mathbf b_1&\cdots&\mathbf A\mathbf b_n
\end{bmatrix}.
\]

To check linearity in column \(j\), let
\(\mathbf u,\mathbf v\in\mathbb R^n\) and \(\alpha,\beta\in\mathbb R\).
Because multiplication by \(\mathbf A\) is linear,

\[
\mathbf A(\alpha\mathbf u+\beta\mathbf v)
=
\alpha\mathbf A\mathbf u+\beta\mathbf A\mathbf v.
\]

The determinant is linear in column \(j\), so

\[
\begin{aligned}
&F(\mathbf b_1,\ldots,\alpha\mathbf u+\beta\mathbf v,\ldots,\mathbf b_n)\\
&=\alpha F(\mathbf b_1,\ldots,\mathbf u,\ldots,\mathbf b_n)
+\beta F(\mathbf b_1,\ldots,\mathbf v,\ldots,\mathbf b_n).
\end{aligned}
\]

Thus, \(F\) is multilinear in the columns of \(\mathbf B\). If two
columns of \(\mathbf B\) are equal, their images under \(\mathbf A\) are
equal, so \(\mathbf A\mathbf B\) has repeated columns and
\(F(\mathbf B)=0\). Swapping two columns of \(\mathbf B\) swaps the
corresponding columns of \(\mathbf A\mathbf B\) and changes the
determinant sign.

Hence, \(F\) is an alternating multilinear scalar-valued function of the
columns of \(\mathbf B\). We claim that

\[
F(\mathbf B)=c\det(\mathbf B),
\qquad
c\mathrel{:=}F(\mathbf e_1,\ldots,\mathbf e_n).
\]

The following expansion proves the claim.

Let \(\mathbf e_1,\ldots,\mathbf e_n\) be the standard basis of
\(\mathbb R^n\). Because it is a basis, every column \(\mathbf b_j\) of
\(\mathbf B\) has the unique expansion

\[
\mathbf b_j=\sum_{i=1}^{n}b_{ij}\mathbf e_i.
\]

Using multilinearity once for each of the \(n\) columns gives

\[
F(\mathbf B)
=
\sum_{i_1=1}^{n}\cdots\sum_{i_n=1}^{n}
\left(\prod_{j=1}^{n}b_{i_j,j}\right)
F(\mathbf e_{i_1},\ldots,\mathbf e_{i_n}).
\]

This sum initially contains \(n^n\) terms. If two indices are equal,
such as \(i_p=i_q\), then the corresponding argument list contains the
same basis vector twice. Alternation makes that term zero:

\[
F(\ldots,\mathbf e_{i_p},\ldots,\mathbf e_{i_p},\ldots)=0.
\]

A term can therefore survive only when \(i_1,\ldots,i_n\) are all
different. Because every \(i_j\) belongs to \(\{1,\ldots,n\}\), the
surviving index list must contain every index exactly once. Thus, it has
the form

\[
(i_1,\ldots,i_n)=(\sigma(1),\ldots,\sigma(n))
\]

for some permutation \(\sigma\in S_n\).

Alternation determines the value of every surviving basis term relative
to the ordered basis term. Reordering
\((\mathbf e_{\sigma(1)},\ldots,\mathbf e_{\sigma(n)})\) into
\((\mathbf e_1,\ldots,\mathbf e_n)\) requires an even or odd number of
swaps, so

\[
F(\mathbf e_{\sigma(1)},\ldots,\mathbf e_{\sigma(n)})
=
\operatorname{sgn}(\sigma)
F(\mathbf e_1,\ldots,\mathbf e_n).
\]

Substituting this result into the expansion for \(F(\mathbf B)\) leaves
only the permutation terms:

\[
\begin{aligned}
F(\mathbf B)
&=\sum_{\sigma\in S_n}
\operatorname{sgn}(\sigma)
\prod_{j=1}^{n}b_{\sigma(j),j}  F(\mathbf e_1,\ldots,\mathbf e_n)\\
\end{aligned}
\]

The scalar appearing in the claim is

\[
c\mathrel{:=}F(\mathbf e_1,\ldots,\mathbf e_n).
\]

The columns of \(\mathbf I_n\) are the ordered standard basis vectors,
so the same reference value can be written as

\[
c=F(\mathbf e_1,\ldots,\mathbf e_n)=F(\mathbf I_n).
\]

Substituting the surviving terms into the expansion gives

\[
\begin{aligned}
F(\mathbf B)
&=c\sum_{\sigma\in S_n}
\operatorname{sgn}(\sigma)
\prod_{j=1}^{n}b_{\sigma(j),j}\\
&=c\det(\mathbf B).
\end{aligned}
\]

A two-dimensional example makes the one-free-value idea explicit. In
this example, \(F(\mathbf u,\mathbf v)\) is shorthand for
\(F([\mathbf u\;\mathbf v])\), where \(\mathbf u\) and \(\mathbf v\) are
the two columns of the input matrix. Let
\(\mathbf u=u_1\mathbf e_1+u_2\mathbf e_2\) and
\(\mathbf v=v_1\mathbf e_1+v_2\mathbf e_2\). Using multilinearity in
both columns gives

\[
\begin{aligned}
F(\mathbf u,\mathbf v)
&=u_1v_1F(\mathbf e_1,\mathbf e_1)
+u_1v_2F(\mathbf e_1,\mathbf e_2)\\
&\quad+u_2v_1F(\mathbf e_2,\mathbf e_1)
+u_2v_2F(\mathbf e_2,\mathbf e_2).
\end{aligned}
\]

The repeated-basis terms are zero, \(F(\mathbf e_1,\mathbf e_2)=c\), and
\(F(\mathbf e_2,\mathbf e_1)=-c\). Therefore,

\[
\begin{aligned}
F(\mathbf u,\mathbf v)
&=u_1v_2c-u_2v_1c\\
&=c(u_1v_2-u_2v_1)\\
&=c\det[\mathbf u\;\mathbf v].
\end{aligned}
\]

After \(c\) is known, this equation fixes \(F(\mathbf u,\mathbf v)\) for
every pair \(\mathbf u,\mathbf v\in\mathbb R^2\).

For the particular function \(F(\mathbf B)=\det(\mathbf A\mathbf B)\),

\[
\begin{aligned}
c
&=F(\mathbf I_n)\\
&=\det(\mathbf A\mathbf I_n)\\
&=\det(\mathbf A).
\end{aligned}
\]

Consequently,

\[
\boxed{
\det(\mathbf A\mathbf B)
=
\det(\mathbf A)\det(\mathbf B)
}.
\]

\subsubsection{Why transposition does not change the
determinant}\label{why-transposition-does-not-change-the-determinant}

The related identity

\[
\boxed{\det(\mathbf A^{\mathsf T})=\det(\mathbf A)}
\]

does not follow from multiplicativity alone. It follows directly from
the permutation definition of the determinant.

Let \(\mathbf A=[a_{ij}]\in\mathbb R^{n\times n}\). Because
\((\mathbf A^{\mathsf T})_{ij}=a_{ji}\),

\[
\begin{aligned}
\det(\mathbf A^{\mathsf T})
&=\sum_{\sigma\in S_n}\operatorname{sgn}(\sigma)
\prod_{j=1}^{n}(\mathbf A^{\mathsf T})_{\sigma(j),j}\\
&=\sum_{\sigma\in S_n}\operatorname{sgn}(\sigma)
\prod_{j=1}^{n}a_{j,\sigma(j)}.
\end{aligned}
\]

A permutation \(\sigma\in S_n\) is a bijection from the index set
\(\{1,\ldots,n\}\) to itself. Because it is bijective, it has a unique
\textbf{inverse permutation}, denoted by \(\sigma^{-1}\), that reverses
the mapping performed by \(\sigma\). Define

\[
\tau\mathrel{:=}\sigma^{-1}.
\]

The defining identities are

\[
\tau(\sigma(j))=j,
\qquad
\sigma(\tau(k))=k.
\]

Equivalently,

\[
k=\sigma(j)
\quad\Longleftrightarrow\quad
j=\tau(k).
\]

For example, suppose \(n=3\) and

\[
\sigma(1)=2,
\qquad
\sigma(2)=3,
\qquad
\sigma(3)=1.
\]

The inverse permutation reverses these assignments:

\[
\tau(1)=3,
\qquad
\tau(2)=1,
\qquad
\tau(3)=2.
\]

For instance, \(\sigma(1)=2\) is reversed by \(\tau(2)=1\).

As \(\sigma\) runs through all permutations in \(S_n\),
\(\tau=\sigma^{-1}\) also runs through all permutations in \(S_n\). This
is because inversion pairs every \(\sigma\) with exactly one \(\tau\),
and the pairing is reversible: \(\sigma=\tau^{-1}\). Therefore,
replacing the summation index \(\sigma\) by \(\tau\) neither loses nor
duplicates a permutation.

The inverse permutation also has the same sign as the original
permutation. Suppose \(\sigma\) can be produced by \(r\) pairwise swaps.
To undo \(\sigma\), perform those same \(r\) swaps in reverse order.
Thus, \(\sigma^{-1}\) uses the same number of swaps as \(\sigma\). They
therefore have the same even-or-odd parity, so

\[
\operatorname{sgn}(\tau)
=
\operatorname{sgn}(\sigma^{-1})
=
\operatorname{sgn}(\sigma).
\]

Now consider the product

\[
\prod_{j=1}^{n}a_{j,\sigma(j)}.
\]

Introduce the new index

\[
k\mathrel{:=}\sigma(j).
\]

Because \(\sigma\) is a permutation, as \(j\) visits \(1,\ldots,n\), the
index \(k=\sigma(j)\) also visits every value \(1,\ldots,n\) exactly
once. From \(\tau=\sigma^{-1}\), the relation \(k=\sigma(j)\) implies
\(j=\tau(k)\). Each factor can therefore be rewritten as

\[
a_{j,\sigma(j)}=a_{\tau(k),k}.
\]

Consequently,

\[
\prod_{j=1}^{n}a_{j,\sigma(j)}
=
\prod_{k=1}^{n}a_{\tau(k),k}.
\]

The two products contain the same scalar factors, possibly in a
different order. Scalar multiplication is commutative, so their order
does not affect the product. In the three-dimensional example above,

\[
\begin{aligned}
\prod_{j=1}^{3}a_{j,\sigma(j)}
&=a_{1,2}a_{2,3}a_{3,1},\\
\prod_{k=1}^{3}a_{\tau(k),k}
&=a_{3,1}a_{1,2}a_{2,3}.
\end{aligned}
\]

These products are equal because they contain exactly the same three
factors.

Therefore,

\[
\begin{aligned}
\det(\mathbf A^{\mathsf T})
&=\sum_{\tau\in S_n}\operatorname{sgn}(\tau)
\prod_{k=1}^{n}a_{\tau(k),k}\\
&=\det(\mathbf A).
\end{aligned}
\]

Thus, exchanging the rows and columns changes how the same determinant
terms are indexed but does not change their values or signs.

\subsubsection{Determinant of an
inverse}\label{determinant-of-an-inverse}

For nonsingular \(\mathbf A\), multiplicativity does give the inverse
identity. Taking determinants of \(\mathbf A^{-1}\mathbf A=\mathbf I_n\)
yields

\[
\begin{aligned}
1
&=\det(\mathbf I_n)\\
&=\det(\mathbf A^{-1}\mathbf A)\\
&=\det(\mathbf A^{-1})\det(\mathbf A).
\end{aligned}
\]

Since \(\det(\mathbf A)\neq0\),

\[
\boxed{
\det(\mathbf A^{-1})
=
\frac{1}{\det(\mathbf A)}
}.
\]

\section*{Quick test C: determinants}\label{quick-test-c-determinants}
\addcontentsline{toc}{section}{Quick test C: determinants}

\markright{Quick test C: determinants}

\begin{enumerate}
\def\labelenumi{\arabic{enumi}.}
\tightlist
\item
  \textbf{Recall:} What does \(|\det(\mathbf A)|\) measure
  geometrically?
\item
  \textbf{Explanation:} What does the sign of a determinant record?
\item
  \textbf{Derivation:} Explain why repeated columns force the
  determinant to be zero.
\item
  \textbf{Application:} Compute
  \(\det\begin{bmatrix}2&1\\3&4\end{bmatrix}\) and interpret its
  magnitude.
\item
  \textbf{Limitations:} Why is the determinant not defined in the same
  way for a rectangular matrix?
\item
  \textbf{Composition:} Explain geometrically why
  \(\det(\mathbf A\mathbf B)=\det(\mathbf A)\det(\mathbf B)\).
\item
  \textbf{Reindexing:} If \(\sigma(1)=2\), \(\sigma(2)=3\), and
  \(\sigma(3)=1\), find \(\tau=\sigma^{-1}\) and verify
  \(\prod_{j=1}^{3}a_{j,\sigma(j)}=\prod_{k=1}^{3}a_{\tau(k),k}\).
\end{enumerate}

\subsection*{Answers and hints}\label{answers-and-hints-2}
\addcontentsline{toc}{subsection}{Answers and hints}

\begin{enumerate}
\def\labelenumi{\arabic{enumi}.}
\tightlist
\item
  The factor by which the associated linear map scales \(n\)-dimensional
  volume.
\item
  A positive determinant preserves the ordered right-handed orientation;
  a negative determinant reverses it to a left-handed orientation and
  therefore indicates a reflection component; a zero determinant
  collapses the volume, so no full-dimensional orientation remains.
\item
  Swapping identical columns leaves the matrix unchanged but must negate
  the determinant, so the determinant equals its own negative and is
  zero.
\item
  The determinant is \(2(4)-1(3)=5\), so areas are scaled by a factor of
  \(5\) and orientation is preserved.
\item
  A rectangular map changes dimension, so it does not map an
  \(n\)-dimensional volume to another volume in the same
  \(n\)-dimensional space.
\item
  \(\mathbf B\) scales signed volume first and \(\mathbf A\) scales the
  resulting signed volume second, so the factors multiply.
\item
  The inverse is \(\tau(1)=3\), \(\tau(2)=1\), and \(\tau(3)=2\). The
  left product is \(a_{1,2}a_{2,3}a_{3,1}\), while the reindexed product
  is \(a_{3,1}a_{1,2}a_{2,3}\); these are equal because scalar factors
  commute.
\end{enumerate}

\section{Orthogonal matrices}\label{orthogonal-matrices}

Let \(\mathbf Q\in\mathbb R^{n\times n}\). The matrix \(\mathbf Q\) is
\textbf{orthogonal} if

\[
\mathbf Q^{\mathsf T}\mathbf Q=\mathbf I_n.
\]

The entries of \(\mathbf Q^{\mathsf T}\mathbf Q\) are pairwise dot
products between columns of \(\mathbf Q\). Therefore, this equation is
equivalent to saying that the columns form an orthonormal basis of
\(\mathbb R^n\).

Because \(\mathbf Q^{\mathsf T}\mathbf Q=\mathbf I_n\), the transpose
satisfies the defining equation for the inverse:

\[
\boxed{\mathbf Q^{-1}=\mathbf Q^{\mathsf T}}.
\]

For \(\mathbf x,\mathbf y\in\mathbb R^n\),

\[
(\mathbf Q\mathbf x)^{\mathsf T}(\mathbf Q\mathbf y)
=
\mathbf x^{\mathsf T}\mathbf Q^{\mathsf T}\mathbf Q\mathbf y
=
\mathbf x^{\mathsf T}\mathbf y.
\]

Thus, orthogonal matrices preserve dot products, lengths, and angles.

Taking determinants of \(\mathbf Q^{\mathsf T}\mathbf Q=\mathbf I_n\)
gives

\[
\det(\mathbf Q)^2=1,
\]

so

\[
\det(\mathbf Q)=+1
\qquad\text{or}\qquad
\det(\mathbf Q)=-1.
\]

The \textbf{orthogonal group} is

\[
O(n)
\mathrel{:=}
\left\{\mathbf Q\in\mathbb R^{n\times n}:\mathbf Q^{\mathsf T}\mathbf Q=\mathbf I_n\right\}.
\]

The \textbf{special orthogonal group} is

\[
SO(n)
\mathrel{:=}
\left\{\mathbf Q\in O(n):\det(\mathbf Q)=1\right\}.
\]

Matrices in \(SO(n)\) preserve lengths, angles, and orientation.
Orthogonal matrices with determinant \(-1\) preserve lengths and angles
but reverse orientation through a reflection. Rotation matrices used for
right-handed coordinate frames belong to \(SO(2)\) or \(SO(3)\).

\section*{Quick test D: orthogonal
matrices}\label{quick-test-d-orthogonal-matrices}
\addcontentsline{toc}{section}{Quick test D: orthogonal matrices}

\markright{Quick test D: orthogonal matrices}

\begin{enumerate}
\def\labelenumi{\arabic{enumi}.}
\tightlist
\item
  \textbf{Recall:} What equation defines an orthogonal matrix?
\item
  \textbf{Explanation:} Why is the inverse of an orthogonal matrix its
  transpose?
\item
  \textbf{Derivation:} Show that
  \(\|\mathbf Q\mathbf x\|_2=\|\mathbf x\|_2\).
\item
  \textbf{Application:} Determine whether
  \(\begin{bmatrix}0&1\\1&0\end{bmatrix}\) belongs to \(SO(2)\).
\item
  \textbf{Limitations:} Why is
  \(\mathbf Q^{\mathsf T}\mathbf Q=\mathbf I_n\) insufficient by itself
  to identify a proper rotation?
\end{enumerate}

\subsection*{Answers and hints}\label{answers-and-hints-3}
\addcontentsline{toc}{subsection}{Answers and hints}

\begin{enumerate}
\def\labelenumi{\arabic{enumi}.}
\tightlist
\item
  \(\mathbf Q^{\mathsf T}\mathbf Q=\mathbf I_n\).
\item
  The transpose already satisfies the defining inverse equation.
\item
  Square both norms and use
  \(\mathbf Q^{\mathsf T}\mathbf Q=\mathbf I_n\).
\item
  It is orthogonal but has determinant \(-1\), so it is a reflection in
  \(O(2)\) and does not belong to \(SO(2)\).
\item
  Orthogonality permits determinant \(-1\) reflections as well as
  determinant \(+1\) rotations.
\end{enumerate}

\section*{Cumulative chapter test}\label{cumulative-chapter-test}
\addcontentsline{toc}{section}{Cumulative chapter test}

\markright{Cumulative chapter test}

\begin{enumerate}
\def\labelenumi{\arabic{enumi}.}
\tightlist
\item
  Define a basis and explain why coordinate columns are unique in a
  basis.
\item
  Given \(\mathbf A\in\mathbb R^{3\times2}\) and
  \(\mathbf x\in\mathbb R^2\), state the dimension of
  \(\mathbf A\mathbf x\) and interpret it using the columns of
  \(\mathbf A\).
\item
  Explain the relationship between \(\mathcal N(\mathbf A)\),
  independent columns, and full column rank.
\item
  Prove that a nonsingular square matrix cannot map a nonzero vector to
  zero.
\item
  Define the determinant using multilinearity, alternation, and
  normalisation.
\item
  Explain why \(\det(\mathbf A)=0\) indicates lost information.
\item
  Derive \(\det(\mathbf A^{-1})=1/\det(\mathbf A)\).
\item
  Distinguish \(O(3)\) from \(SO(3)\) and explain which set contains
  rigid-body rotation matrices.
\item
  State the rank-nullity theorem and explain what its two terms count.
\item
  State the vector triple-product identity and explain why the order and
  parentheses in a nested cross product matter.
\end{enumerate}

\subsection*{Cumulative answers and
hints}\label{cumulative-answers-and-hints}
\addcontentsline{toc}{subsection}{Cumulative answers and hints}

\begin{enumerate}
\def\labelenumi{\arabic{enumi}.}
\item
  A basis is a linearly independent spanning set. Independence prevents
  multiple coefficient descriptions, while spanning guarantees that a
  description exists.
\item
  \(\mathbf A\mathbf x\in\mathbb R^3\) and equals
  \(x_1\mathbf a_1+x_2\mathbf a_2\).
\item
  \(\mathcal N(\mathbf A)=\{\mathbf0\}\) exactly when the columns are
  independent; for \(n\) columns this gives
  \(\operatorname{rank}(\mathbf A)=n\).
\item
  Nonsingularity is equivalent to
  \(\mathbf A\mathbf x=\mathbf0\Rightarrow\mathbf x=\mathbf0\); the
  contrapositive is
  \(\mathbf x\neq\mathbf0\Rightarrow\mathbf A\mathbf x\neq\mathbf0\).
\item
  The determinant is the unique scalar function that is linear in each
  column, changes sign under a column swap, and equals \(1\) on
  \(\mathbf I_n\).
\item
  For \(\mathbf A\in\mathbb R^{n\times n}\), \(\det(\mathbf A)=0\) means
  that \(\mathbf A\) is singular, so its null space contains a nonzero
  vector \(\mathbf z\in\mathbb R^n\) satisfying
  \(\mathbf A\mathbf z=\mathbf0\). For any \(\mathbf x\in\mathbb R^n\),

  \[
  \mathbf A(\mathbf x+\mathbf z)
  =\mathbf A\mathbf x+\mathbf A\mathbf z
  =\mathbf A\mathbf x.
  \]

  Therefore, the distinct inputs \(\mathbf x\) and
  \(\mathbf x+\mathbf z\) produce the same output. The output cannot
  reveal the input component in the null-space direction \(\mathbf z\),
  so the input cannot be uniquely recovered. This is the meaning of
  \textbf{lost information} here: an independent input direction or
  degree of freedom has disappeared, not that Shannon information has
  been calculated. Geometrically, the map collapses \(n\)-dimensional
  volume into a lower-dimensional set, which is why its volume-scaling
  factor is zero.
\item
  Take determinants of \(\mathbf A^{-1}\mathbf A=\mathbf I_n\) and use
  multiplicativity.
\item
  \(O(3)\) contains all orthogonal \(3\times3\) matrices with
  determinant \(\pm1\); \(SO(3)\) contains the determinant-\(1\) members
  and therefore the proper rigid-body rotations.
\item
  For \(\mathbf A\in\mathbb R^{m\times n}\),
  \(\operatorname{rank}(\mathbf A)+\operatorname{nullity}(\mathbf A)=n\).
  Rank counts independent output directions produced by the map, while
  nullity counts independent input directions sent to zero.
\item
  \(\mathbf a\times(\mathbf b\times\mathbf c)=\mathbf b(\mathbf a^\mathsf T\mathbf c)-\mathbf c(\mathbf a^\mathsf T\mathbf b)\).
  The cross product is not associative, so changing the parentheses
  generally changes the result.
\end{enumerate}

\section*{References}\label{references}
\addcontentsline{toc}{section}{References}

\markright{References}

\begin{itemize}
\tightlist
\item
  Richard Szeliski, \emph{Computer Vision: Algorithms and Applications},
  2nd ed., Appendix A, especially the discussion of orthonormal
  matrices, rank, singular values, and matrix factorisations.
\item
  Kevin M. Lynch and Frank C. Park, \emph{Modern Robotics}, Chapter 3,
  especially Section 3.2.1 on rotation matrices and the distinction
  between orthogonality and right-handed orientation.
\end{itemize}

\chapter{Geometry and Coordinate
Frames}\label{sec-geometry-coordinate-frames}

\section*{Why coordinate frames
matter}\label{why-coordinate-frames-matter}
\addcontentsline{toc}{section}{Why coordinate frames matter}

\markright{Why coordinate frames matter}

A robot represents positions, directions, velocities, measurements, and
commands with numbers. Those numbers acquire physical meaning only when
their coordinate frame and units are known. A controller that combines
quantities expressed in incompatible frames can perform a mathematically
valid calculation and produce a physically incorrect command.

Coordinate-frame methods allow us to:

\begin{itemize}
\tightlist
\item
  represent a robot's pose relative to a world frame;
\item
  distinguish world-frame and body-frame quantities;
\item
  transform sensor measurements into a shared frame;
\item
  express target and tracking errors consistently;
\item
  convert body-relative motion into world-relative motion; and
\item
  detect invalid transformations before they reach a controller.
\end{itemize}

The central rule is simple: interpret every coordinate column together
with the frame in which it is expressed.

\section*{Prerequisites and
assumptions}\label{prerequisites-and-assumptions-1}
\addcontentsline{toc}{section}{Prerequisites and assumptions}

\markright{Prerequisites and assumptions}

\subsection*{Prerequisites}\label{prerequisites-1}
\addcontentsline{toc}{subsection}{Prerequisites}

The reader should understand the material in
Chapter~\ref{sec-linear-algebra-foundations}, especially:

\begin{itemize}
\tightlist
\item
  geometric vectors and coordinate columns;
\item
  dot products, norms, and orthogonality;
\item
  matrices as linear maps and compositions;
\item
  linear independence, rank, and inverses;
\item
  determinants and signed orientation; and
\item
  orthogonal matrices.
\end{itemize}

The reader should also be comfortable with sine, cosine, and angles
measured in radians.

\subsection*{Assumptions}\label{assumptions-1}
\addcontentsline{toc}{subsection}{Assumptions}

\begin{itemize}
\tightlist
\item
  Physical space is modelled as Euclidean space.
\item
  Robot links and sensor mounts are treated as rigid bodies.
\item
  Coordinate frames are right-handed and use orthonormal basis vectors.
\item
  Length is measured in metres, time in seconds, and angles in radians.
\item
  Coordinate vectors are written as columns.
\item
  Numerical examples ignore measurement noise unless stated otherwise.
\end{itemize}

These idealisations isolate the geometry. A physical robot also
introduces measurement noise, calibration error, structural compliance,
and model mismatch.

\section*{Dependency map}\label{dependency-map-1}
\addcontentsline{toc}{section}{Dependency map}

\markright{Dependency map}

\begin{Shaded}
\begin{Highlighting}[]
\NormalTok{Linear algebra foundations}
\NormalTok{        |}
\NormalTok{        v}
\NormalTok{Cross products and rigid{-}body velocity}
\NormalTok{        |}
\NormalTok{        v}
\NormalTok{Coordinate frames and units}
\NormalTok{        |}
\NormalTok{        v}
\NormalTok{Rotation matrices}
\NormalTok{        |}
\NormalTok{        v}
\NormalTok{Point transformations}
\NormalTok{        |}
\NormalTok{        v}
\NormalTok{Differential{-}drive motion and feedback control}
\end{Highlighting}
\end{Shaded}

\section{Cross product and rigid-body
velocity}\label{cross-product-and-rigid-body-velocity}

Let \(\mathbf u\) and \(\mathbf v\) be geometric vectors with coordinate
columns \({}^{a}\mathbf u,{}^{a}\mathbf v\in\mathbb R^3\) expressed in
the same right-handed orthonormal frame \(\{a\}\). Their cross product,
expressed in frame \(\{a\}\), is defined by

\[
{}^{a}\mathbf{u}\times{}^{a}\mathbf{v}
=
\begin{bmatrix}
u_yv_z-u_zv_y \\
u_zv_x-u_xv_z \\
u_xv_y-u_yv_x
\end{bmatrix}
\in\mathbb R^3.
\]

Its direction follows the right-hand rule. If both vectors are nonzero,
let \(\theta\in[0,\pi]\) be the smaller angle between them. The
cross-product magnitude is

\[
\left\|{}^{a}\mathbf{u}\times{}^{a}\mathbf{v}\right\|_2
=
\|\mathbf{u}\|_2\|\mathbf{v}\|_2\sin\theta.
\]

The cross product is perpendicular to both input vectors because

\[
{}^{a}\mathbf{u}^\mathsf{T}
\left({}^{a}\mathbf{u}\times{}^{a}\mathbf{v}\right)=0,
\qquad
{}^{a}\mathbf{v}^\mathsf{T}
\left({}^{a}\mathbf{u}\times{}^{a}\mathbf{v}\right)=0.
\]

\subsection{Example: velocity of a point on a rotating rigid
body}\label{example-velocity-of-a-point-on-a-rotating-rigid-body}

Let \(O\) and \(P\) be two points fixed to the same rigid body. Point
\(O\) is the chosen reference point; it need not be stationary, and it
need not be the body's centre. The vector \({}^{a}\mathbf{r}_{OP}\) runs
from \(O\) to \(P\). Although this vector is fixed relative to the body,
its coordinates generally change as the body rotates. Express every
quantity in the same orthonormal frame \(\{a\}\):

\begin{itemize}
\tightlist
\item
  \({}^{a}\boldsymbol{\omega}\) as the body's angular-velocity vector,
  measured in \(\mathrm{rad\,s^{-1}}\);
\item
  \({}^{a}\mathbf{r}_{OP}\) as the position of \(P\) relative to \(O\),
  measured in \(\mathrm{m}\);
\item
  \({}^{a}\mathbf{v}_O\) as the linear velocity of \(O\), measured in
  \(\mathrm{m\,s^{-1}}\); and
\item
  \({}^{a}\mathbf{v}_P\) as the linear velocity of \(P\), measured in
  \(\mathrm{m\,s^{-1}}\).
\end{itemize}

At any instant, the motion of \(P\) can be understood as the sum of two
simultaneous contributions:

\begin{enumerate}
\def\labelenumi{\arabic{enumi}.}
\tightlist
\item
  \textbf{Translation with \(O\).} If the body were not rotating, \(P\)
  would share the reference-point velocity \({}^{a}\mathbf v_O\).
\item
  \textbf{Motion relative to \(O\) caused by rotation.} Rotation adds a
  velocity whose direction and magnitude depend on
  \({}^{a}\boldsymbol{\omega}\) and \({}^{a}\mathbf r_{OP}\).
\end{enumerate}

The angular-velocity vector points along the instantaneous rotation axis
according to the right-hand rule, and its magnitude is the angular
speed. Define the rotational contribution to the velocity of \(P\) by

\[
{}^{a}\mathbf{v}_{P/O}
\mathrel{:=}
{}^{a}\boldsymbol{\omega}\times{}^{a}\mathbf{r}_{OP}.
\]

The cross product determines both the direction and the magnitude of
this contribution. In particular,

\[
{}^{a}\boldsymbol{\omega}^{\mathsf T}{}^{a}\mathbf{v}_{P/O}=0,
\qquad
{}^{a}\mathbf{r}_{OP}^{\mathsf T}{}^{a}\mathbf{v}_{P/O}=0.
\]

Therefore, \({}^{a}\mathbf{v}_{P/O}\) is perpendicular to the rotation
axis and to \({}^{a}\mathbf r_{OP}\). It points tangent to the
instantaneous circular path generated by the rotation. Adding the
translational and rotational contributions gives the rigid-body velocity
relation

\[
\boxed{
{}^{a}\mathbf{v}_P
=
{}^{a}\mathbf{v}_O
+{}^{a}\boldsymbol{\omega}\times{}^{a}\mathbf{r}_{OP}
}.
\]

\subsubsection{Velocity-vector
construction}\label{velocity-vector-construction}

The construction in Figure~\ref{fig-rigid-body-point-velocity} shows
this addition for one instant:

\begin{itemize}
\tightlist
\item
  The green arrow \({}^{a}\mathbf r_{OP}\) locates \(P\) relative to
  \(O\).
\item
  The magenta arrow is the reference-point velocity
  \({}^{a}\mathbf v_O\). It is deliberately not parallel to
  \({}^{a}\mathbf r_{OP}\).
\item
  The orange arrow is the rotational contribution
  \({}^{a}\boldsymbol{\omega}\times{}^{a}\mathbf r_{OP}={}^{a}\mathbf v_{P/O}\).
  The right-angle marker shows that it is perpendicular to
  \({}^{a}\mathbf r_{OP}\).
\item
  The dashed arrows are translated copies of the two velocity
  contributions. Moving a vector without rotating or resizing it does
  not change the vector, so the copies can be arranged head to tail.
\item
  The blue arrow joins the beginning of the first contribution to the
  end of the second. It is therefore the resultant velocity
  \({}^{a}\mathbf v_P\).
\end{itemize}

In this planar view, \({}^{a}\boldsymbol{\omega}\) points along
\(+z_a\), out of the page. The right-hand rule then gives the displayed
tangential direction for
\({}^{a}\boldsymbol{\omega}\times{}^{a}\mathbf r_{OP}\). Notice that
only the rotational contribution must be perpendicular to
\({}^{a}\mathbf r_{OP}\). The total velocity \({}^{a}\mathbf v_P\)
generally is not, because \({}^{a}\mathbf v_O\) may point in any
direction.

\begin{figure}[H]

\centering{

\includegraphics[width=1\linewidth,height=\textheight,keepaspectratio]{index_files/mediabag/figures/rigid_body_point_velocity.pdf}

}

\caption{\label{fig-rigid-body-point-velocity}Velocity construction for
two points on a rigid body. The rotational contribution is perpendicular
to the position vector, while the reference-point velocity is
independent of that perpendicular relationship.}

\end{figure}%

For pure rotation about a fixed axis through \(O\), the reference point
is stationary, so \({}^{a}\mathbf{v}_O=\mathbf{0}\). Therefore,

\[
\boxed{{}^{a}\mathbf{v}_P
=
{}^{a}\boldsymbol{\omega}\times{}^{a}\mathbf{r}_{OP}}.
\]

If \(\phi\) is the angle between \({}^{a}\boldsymbol{\omega}\) and
\({}^{a}\mathbf{r}_{OP}\), then the relative speed caused by rotation is

\[
\left\|{}^{a}\mathbf{v}_{P/O}\right\|_2
=
\left\|{}^{a}\boldsymbol{\omega}\right\|_2
\left\|{}^{a}\mathbf{r}_{OP}\right\|_2
\sin\phi
=
\omega r_\perp,
\]

where \(r_\perp\) is the perpendicular distance from \(P\) to the
rotation axis. A point on the axis has \(r_\perp=0\) and therefore has
no relative linear velocity due to rotation. Under the pure-rotation
assumption, \({}^{a}\mathbf{v}_P={}^{a}\mathbf{v}_{P/O}\).

For planar counter-clockwise rotation about the positive \(z\)-axis,

\[
{}^{a}\boldsymbol{\omega}
=
\begin{bmatrix}
0\\0\\\omega
\end{bmatrix},
\qquad
{}^{a}\mathbf{r}_{OP}
=
\begin{bmatrix}
x\\y\\0
\end{bmatrix}.
\]

The point velocity is

\[
{}^{a}\mathbf{v}_P
=
{}^{a}\boldsymbol{\omega}\times{}^{a}\mathbf{r}_{OP}
=
\begin{bmatrix}
-\omega y\\
\omega x\\
0
\end{bmatrix}.
\]

For example, if \(\omega=2\,\mathrm{rad\,s^{-1}}\) and
\({}^{a}\mathbf{r}_{OP}=[0.5\;\;0\;\;0]^\mathsf{T}\,\mathrm{m}\), then

\[
{}^{a}\mathbf{v}_P
=
\begin{bmatrix}
0\\1\\0
\end{bmatrix}
\mathrm{m\,s^{-1}}.
\]

The point lies on the positive \(x\)-axis, while its instantaneous
velocity points along the positive \(y\)-axis, tangent to its circular
path.

\section*{Quick test A: cross products and rigid-body
velocity}\label{quick-test-a-cross-products-and-rigid-body-velocity}
\addcontentsline{toc}{section}{Quick test A: cross products and
rigid-body velocity}

\markright{Quick test A: cross products and rigid-body velocity}

Complete this test without looking back at the section.

\begin{enumerate}
\def\labelenumi{\arabic{enumi}.}
\tightlist
\item
  \textbf{Recall:} What rule determines the direction of
  \(\mathbf u\times\mathbf v\)?
\item
  \textbf{Explanation:} Why is \(\mathbf u\times\mathbf v\)
  perpendicular to both input vectors?
\item
  \textbf{Derivation:} Substitute the component formula for
  \({}^{a}\mathbf u\times{}^{a}\mathbf v\) into
  \({}^{a}\mathbf u^{\mathsf T}({}^{a}\mathbf u\times{}^{a}\mathbf v)\)
  and show that the terms cancel.
\item
  \textbf{Application:} Compute
  \([1\;\;2\;\;0]^\mathsf T\times[0\;\;1\;\;0]^\mathsf T\).
\item
  \textbf{Limitations:} What happens to the cross product when the input
  vectors are parallel, and why does that prevent it from defining a
  unique normal direction?
\item
  \textbf{Transfer:} At one instant, a rigid body has
  \({}^{a}\boldsymbol{\omega}=[0\;\;0\;\;2]^\mathsf{T}\,\mathrm{rad\,s^{-1}}\),
  \({}^{a}\mathbf r_{OP}=[1\;\;0\;\;0]^\mathsf{T}\,\mathrm{m}\), and
  \({}^{a}\mathbf v_O=[0.5\;\;0.5\;\;0]^\mathsf{T}\,\mathrm{m\,s^{-1}}\).
  Find \({}^{a}\mathbf v_{P/O}\) and \({}^{a}\mathbf v_P\). Which
  velocity is perpendicular to \({}^{a}\mathbf r_{OP}\)?
\end{enumerate}

\subsection*{Answers and hints}\label{answers-and-hints-4}
\addcontentsline{toc}{subsection}{Answers and hints}

\begin{enumerate}
\def\labelenumi{\arabic{enumi}.}
\tightlist
\item
  The right-hand rule, applied in the order from \(\mathbf u\) to
  \(\mathbf v\).
\item
  Its dot product with either input is zero.
\item
  Expanding gives
  \(u_x(u_yv_z-u_zv_y)+u_y(u_zv_x-u_xv_z)+u_z(u_xv_y-u_yv_x)=0\) after
  pairwise cancellation.
\item
  The result is \([0\;\;0\;\;1]^\mathsf T\).
\item
  Parallel vectors have \(\sin\theta=0\), so their cross product is the
  zero vector. A zero vector has no direction and therefore cannot
  specify a unique normal.
\item
  \({}^{a}\mathbf v_{P/O}=[0\;\;2\;\;0]^\mathsf T\,\mathrm{m\,s^{-1}}\)
  and
  \({}^{a}\mathbf v_P=[0.5\;\;2.5\;\;0]^\mathsf T\,\mathrm{m\,s^{-1}}\).
  The rotational contribution is perpendicular to
  \({}^{a}\mathbf r_{OP}\); the total velocity is not.
\end{enumerate}

\section{Coordinate frames}\label{coordinate-frames}

\subsection{What defines a frame?}\label{what-defines-a-frame}

A three-dimensional coordinate frame \(\{a\}\) consists of:

\begin{enumerate}
\def\labelenumi{\arabic{enumi}.}
\tightlist
\item
  an origin \(O_a\), which is the physical point from which positions in
  frame \(\{a\}\) are measured; and
\item
  three right-handed orthonormal basis vectors \(\mathbf e_{x,a}\),
  \(\mathbf e_{y,a}\), and \(\mathbf e_{z,a}\), which define the
  positive \(x_a\)-, \(y_a\)-, and \(z_a\)-axis directions.
\end{enumerate}

The subscript \(a\) associates the origin and basis vectors with frame
\(\{a\}\). Orthonormal means that each basis vector has unit length and
that distinct basis vectors are perpendicular. Right-handed means that
\(\mathbf e_{x,a}\times\mathbf e_{y,a}=\mathbf e_{z,a}\).

A free vector depends only on the basis orientation used to describe it.
A point also depends on the frame origin. This distinction explains why
changing the coordinate frame of a free vector requires only a rotation,
whereas changing the coordinates of a point generally requires a
rotation and a translation.

\subsection{Orientation convention}\label{orientation-convention}

Let \(\{a\}\) and \(\{b\}\) be three-dimensional right-handed
orthonormal frames. Define \(\mathbf R_{ab}\in\mathbb R^{3\times3}\) as
the rotation matrix that maps coordinate columns from frame \(\{b\}\) to
frame \(\{a\}\). Its columns are the basis vectors of \(\{b\}\)
expressed in frame \(\{a\}\):

\[
\mathbf{R}_{ab}
=
\begin{bmatrix}
{}^{a}\mathbf{e}_{x,b} &
{}^{a}\mathbf{e}_{y,b} &
{}^{a}\mathbf{e}_{z,b}
\end{bmatrix},
\]

where each column \({}^{a}\mathbf e_{i,b}\in\mathbb R^3\) is an axis of
frame \(\{b\}\) written using the basis of frame \(\{a\}\), for
\(i\in\{x,y,z\}\).

The following notation distinguishes a geometric vector from its
coordinate representations:

{\def\LTcaptype{none} % do not increment counter
\begin{longtable}[]{@{}
  >{\raggedright\arraybackslash}p{(\linewidth - 2\tabcolsep) * \real{0.5000}}
  >{\raggedright\arraybackslash}p{(\linewidth - 2\tabcolsep) * \real{0.5000}}@{}}
\toprule\noalign{}
\begin{minipage}[b]{\linewidth}\raggedright
Symbol
\end{minipage} & \begin{minipage}[b]{\linewidth}\raggedright
Meaning
\end{minipage} \\
\midrule\noalign{}
\endhead
\bottomrule\noalign{}
\endlastfoot
\(\mathbf v\) & A geometric vector, independent of the coordinate frame
used to describe it. \\
\(\mathbf e_{x,b}\) & The geometric unit vector along the positive
\(x_b\)-axis of frame \(\{b\}\). \\
\([\mathbf v]_b={}^{b}\mathbf v\in\mathbb R^3\) & The numerical
coordinate column of \(\mathbf v\) in frame \(\{b\}\). \\
\([\mathbf v]_a={}^{a}\mathbf v\in\mathbb R^3\) & The numerical
coordinate column of \(\mathbf v\) in frame \(\{a\}\). \\
\([\mathbf e_{x,b}]_a={}^{a}\mathbf e_{x,b}\in\mathbb R^3\) & The
coordinate column of frame \(\{b\}\)'s geometric \(x\)-axis, expressed
in frame \(\{a\}\). \\
\end{longtable}
}

The bracket operator \([\,\cdot\,]_a\) means ``take the coordinate
column of the enclosed geometric vector in frame \(\{a\}\).'' The
geometric vector \(\mathbf v\) can be expanded using either frame's
basis:

\[
\boxed{
\mathbf v
=
v_x^{(b)}\mathbf e_{x,b}
+
v_y^{(b)}\mathbf e_{y,b}
+
v_z^{(b)}\mathbf e_{z,b}
=
v_x^{(a)}\mathbf e_{x,a}
+
v_y^{(a)}\mathbf e_{y,a}
+
v_z^{(a)}\mathbf e_{z,a}
}.
\]

Here, \(v_i^{(b)}\in\mathbb R\) and \(v_i^{(a)}\in\mathbb R\) are the
scalar components of \(\mathbf v\) along axis \(i\in\{x,y,z\}\) of
frames \(\{b\}\) and \(\{a\}\), respectively. The corresponding
coordinate columns are

\[
{}^{b}\mathbf v
=
[\mathbf v]_b
=
\begin{bmatrix}
v_x^{(b)}\\
v_y^{(b)}\\
v_z^{(b)}
\end{bmatrix},
\qquad
{}^{a}\mathbf v
=
[\mathbf v]_a
=
\begin{bmatrix}
v_x^{(a)}\\
v_y^{(a)}\\
v_z^{(a)}
\end{bmatrix}.
\]

Apply the coordinate operator \([\,\cdot\,]_a\) to the expansion of
\(\mathbf v\) in the frame-\(\{b\}\) basis. Because taking coordinates
is a linear operation,

\[
[\mathbf v]_a
=
v_x^{(b)}[\mathbf e_{x,b}]_a
+
v_y^{(b)}[\mathbf e_{y,b}]_a
+
v_z^{(b)}[\mathbf e_{z,b}]_a.
\]

Using the equivalent superscript notation gives

\[
{}^{a}\mathbf v
=
v_x^{(b)}\,{}^{a}\mathbf e_{x,b}
+
v_y^{(b)}\,{}^{a}\mathbf e_{y,b}
+
v_z^{(b)}\,{}^{a}\mathbf e_{z,b}.
\]

Writing this linear combination as matrix-vector multiplication gives

\[
{}^{a}\mathbf v
=
\underbrace{
\begin{bmatrix}
{}^{a}\mathbf e_{x,b} &
{}^{a}\mathbf e_{y,b} &
{}^{a}\mathbf e_{z,b}
\end{bmatrix}
}_{\mathbf R_{ab}}
\underbrace{
\begin{bmatrix}
v_x^{(b)}\\
v_y^{(b)}\\
v_z^{(b)}
\end{bmatrix}
}_{{}^{b}\mathbf v}.
\]

Therefore,

\[
\boxed{{}^{a}\mathbf v=\mathbf R_{ab}\,{}^{b}\mathbf v}.
\]

The subscript order can be read as:

\begin{quote}
\(\mathbf{R}_{ab}\) maps coordinates \textbf{from frame \(\{b\}\) to
frame \(\{a\}\)}.
\end{quote}

The \hyperref[sec-linear-algebra-foundations]{linear algebra foundations
chapter} develops orthogonal matrices, determinants, and inverses in
general. Applying those results to the orthonormal, right-handed columns
of \(\mathbf R_{ab}\) gives:

\[
\boxed{
\mathbf{R}_{ab}^\mathsf{T}\mathbf{R}_{ab}=\mathbf{I}_3,
\qquad
\det(\mathbf{R}_{ab})=1,
\qquad
\mathbf{R}_{ab}^{-1}=\mathbf{R}_{ab}^\mathsf{T}=\mathbf{R}_{ba}
}.
\]

The following derivations explain each property.

\subsubsection{Orthonormal columns}\label{orthonormal-columns}

Let \(\mathbf r_1,\mathbf r_2,\mathbf r_3\in\mathbb R^3\) denote the
columns of \(\mathbf R_{ab}\):

\[
\mathbf R_{ab}
=
\begin{bmatrix}
\mathbf r_1 & \mathbf r_2 & \mathbf r_3
\end{bmatrix},
\]

where \(\mathbf r_1={}^{a}\mathbf e_{x,b}\),
\(\mathbf r_2={}^{a}\mathbf e_{y,b}\), and
\(\mathbf r_3={}^{a}\mathbf e_{z,b}\). Multiplying the transpose by the
original matrix forms every pairwise dot product between these columns:

\[
\mathbf R_{ab}^{\mathsf T}\mathbf R_{ab}
=
\begin{bmatrix}
\mathbf r_1^{\mathsf T}\mathbf r_1 & \mathbf r_1^{\mathsf T}\mathbf r_2 & \mathbf r_1^{\mathsf T}\mathbf r_3\\
\mathbf r_2^{\mathsf T}\mathbf r_1 & \mathbf r_2^{\mathsf T}\mathbf r_2 & \mathbf r_2^{\mathsf T}\mathbf r_3\\
\mathbf r_3^{\mathsf T}\mathbf r_1 & \mathbf r_3^{\mathsf T}\mathbf r_2 & \mathbf r_3^{\mathsf T}\mathbf r_3
\end{bmatrix}.
\]

Because the frame axes are orthonormal, every column has unit length,

\[
\mathbf r_i^{\mathsf T}\mathbf r_i=1,
\]

and distinct columns are perpendicular,

\[
\mathbf r_i^{\mathsf T}\mathbf r_j=0
\qquad\text{for }i\neq j,
\]

where \(i,j\in\{1,2,3\}\) are column indices. Therefore,

\[
\boxed{
\mathbf R_{ab}^{\mathsf T}\mathbf R_{ab}
=
\begin{bmatrix}
1&0&0\\
0&1&0\\
0&0&1
\end{bmatrix}
=
\mathbf I_3
}.
\]

This identity means that \(\mathbf R_{ab}\) preserves dot products and
therefore preserves vector lengths and angles.

\subsubsection{Inverse and reverse coordinate
mapping}\label{inverse-and-reverse-coordinate-mapping}

The inverse \(\mathbf R_{ab}^{-1}\in\mathbb R^{3\times3}\) is defined by

\[
\mathbf R_{ab}^{-1}\mathbf R_{ab}=\mathbf I_3.
\]

Because \(\mathbf R_{ab}^{\mathsf T}\mathbf R_{ab}=\mathbf I_3\), the
transpose satisfies the defining property of the inverse. Hence,

\[
\boxed{\mathbf R_{ab}^{-1}=\mathbf R_{ab}^{\mathsf T}}.
\]

The coordinate transformation from frame \(\{b\}\) to frame \(\{a\}\) is

\[
{}^{a}\mathbf v=\mathbf R_{ab}\,{}^{b}\mathbf v.
\]

Multiplying by the inverse gives

\[
{}^{b}\mathbf v=\mathbf R_{ab}^{-1}\,{}^{a}\mathbf v.
\]

Define \(\mathbf R_{ba}\in\mathbb R^{3\times3}\) as the matrix that maps
coordinates in the reverse direction, from frame \(\{a\}\) to frame
\(\{b\}\). By definition,

\[
{}^{b}\mathbf v=\mathbf R_{ba}\,{}^{a}\mathbf v.
\]

The two reverse mappings must be the same, so

\[
\boxed{
\mathbf R_{ba}
=
\mathbf R_{ab}^{-1}
=
\mathbf R_{ab}^{\mathsf T}
}.
\]

\subsubsection{Determinant and
orientation}\label{determinant-and-orientation}

The notation \(\det(\mathbf R)\) denotes the \textbf{determinant} of a
square matrix \(\mathbf R\). Geometrically, \(|\det(\mathbf R)|\) is the
volume-scaling factor of the corresponding linear transformation. For an
orthogonal matrix \(\mathbf R\), taking determinants of
\(\mathbf R^{\mathsf T}\mathbf R=\mathbf I_3\) gives

\[
\det(\mathbf R^{\mathsf T}\mathbf R)
=
\det(\mathbf I_3),
\]

and therefore by
\(\det(\mathbf R^{\mathsf T}\mathbf R) = \det(\mathbf R^{\mathsf T})\det(\mathbf R)\)
and \(\det(R^{\mathsf T})=\det(\mathbf R)\) gives

\[
\det(\mathbf R)^2=1.
\]

Thus, an orthogonal matrix can have determinant \(+1\) or \(-1\). Here,
\textbf{orientation} means the handedness of the ordered coordinate
axes. A right-handed frame satisfies

\[
\mathbf e_x\times\mathbf e_y=\mathbf e_z.
\]

The determinant sign has the following interpretation:

\begin{itemize}
\tightlist
\item
  \(\det(\mathbf R)=+1\) preserves handedness, so a right-handed ordered
  basis remains right-handed.
\item
  \(\det(\mathbf R)=-1\) reverses handedness, so the transformation
  contains a reflection and produces a left-handed ordered basis from a
  right-handed one.
\end{itemize}

For example, consider

\[
\mathbf M
=
\begin{bmatrix}
-1&0&0\\
0&1&0\\
0&0&1
\end{bmatrix}.
\]

The matrix \(\mathbf M\in\mathbb R^{3\times3}\) reflects the
\(x\)-coordinate. Its columns remain orthonormal, so
\(\mathbf M^{\mathsf T}\mathbf M=\mathbf I_3\), but
\(\det(\mathbf M)=-1\). It reverses handedness and is therefore a
reflection rather than a rigid-body rotation.

Because frames \(\{a\}\) and \(\{b\}\) are both right-handed, their
rotation matrix must satisfy

\[
\boxed{
\mathbf R_{ab}^{\mathsf T}\mathbf R_{ab}=\mathbf I_3
\quad\text{and}\quad
\det(\mathbf R_{ab})=1
}.
\]

The determinant sign does not indicate whether a rotation is clockwise
or counter-clockwise. It indicates whether the transformation preserves
or reverses the handedness of space.

\subsection{Points require
translation}\label{points-require-translation}

Let \(O_a\) and \(O_b\) be the origins of frames \(\{a\}\) and
\(\{b\}\). Define \({}^{a}\mathbf p_b\in\mathbb R^3\) as the position of
\(O_b\) relative to \(O_a\), expressed in frame \(\{a\}\) and measured
in metres.

For a physical point \(P\), let \({}^{b}\mathbf p_P\in\mathbb R^3\) be
the position of \(P\) relative to \(O_b\), expressed in frame \(\{b\}\),
and let \({}^{a}\mathbf p_P\in\mathbb R^3\) be the position of \(P\)
relative to \(O_a\), expressed in frame \(\{a\}\). Both position columns
are measured in metres. Their relationship is

\[
{}^{a}\mathbf{p}_P
=
\mathbf{R}_{ab}\,{}^{b}\mathbf{p}_P
+{}^{a}\mathbf{p}_b.
\]

The rotation \(\mathbf R_{ab}\,{}^{b}\mathbf p_P\) changes the basis
used to express the displacement from \(O_b\) to \(P\). The translation
\({}^{a}\mathbf p_b\) then accounts for the displacement from \(O_a\) to
\(O_b\). A homogeneous transformation combines these two operations into
one matrix equation.

\section{Worked frame-change example}\label{worked-frame-change-example}

Consider two planar frames \(\{a\}\) and \(\{b\}\) whose axes differ by
an angle \(\theta\in\mathbb R\), measured in radians. A positive
\(\theta\) means that frame \(\{b\}\) is rotated counter-clockwise
relative to frame \(\{a\}\). Define
\(\mathbf R_{ab}(\theta)\in\mathbb R^{2\times2}\) as the rotation matrix
that maps two-dimensional coordinate columns from frame \(\{b\}\) to
frame \(\{a\}\):

\[
{}^{a}\mathbf v
=
\mathbf R_{ab}(\theta)\,{}^{b}\mathbf v,
\]

where \(\mathbf v\) is a geometric free vector and
\({}^{a}\mathbf v,{}^{b}\mathbf v\in\mathbb R^2\) are its coordinate
columns in frames \(\{a\}\) and \(\{b\}\).

Let \(\mathbf e_{x,a}\) and \(\mathbf e_{y,a}\) be the geometric unit
axes of frame \(\{a\}\), and let \(\mathbf e_{x,b}\) and
\(\mathbf e_{y,b}\) be the geometric unit axes of frame \(\{b\}\). By
definition, the columns of \(\mathbf R_{ab}\) are the axes of frame
\(\{b\}\) expressed in frame \(\{a\}\):

\[
\mathbf R_{ab}
=
\begin{bmatrix}
{}^{a}\mathbf e_{x,b}&{}^{a}\mathbf e_{y,b}
\end{bmatrix}.
\]

The \(x_b\)-axis makes angle \(\theta\) with the \(x_a\)-axis. Its
projection onto \(\mathbf e_{x,a}\) is \(\cos\theta\), and its
projection onto \(\mathbf e_{y,a}\) is \(\sin\theta\). Therefore,

\[
{}^{a}\mathbf e_{x,b}
=
\begin{bmatrix}
\cos\theta\\
\sin\theta
\end{bmatrix}.
\]

Because the planar frames are right-handed and orthonormal, the
\(y_b\)-axis is obtained by rotating the \(x_b\)-axis counter-clockwise
by \(\pi/2\). Its angle measured from the \(x_a\)-axis is therefore
\(\theta+\pi/2\), so

\[
{}^{a}\mathbf e_{y,b}
=
\begin{bmatrix}
\cos(\theta+\pi/2)\\
\sin(\theta+\pi/2)
\end{bmatrix}.
\]

Using

\[
\cos(\theta+\pi/2)=-\sin\theta,
\qquad
\sin(\theta+\pi/2)=\cos\theta,
\]

gives

\[
{}^{a}\mathbf e_{y,b}
=
\begin{bmatrix}
-\sin\theta\\
\cos\theta
\end{bmatrix}.
\]

Placing these two coordinate columns side by side produces

\[
\boxed{
\mathbf R_{ab}(\theta)
=
\begin{bmatrix}
\cos\theta&-\sin\theta\\
\sin\theta&\cos\theta
\end{bmatrix}
}.
\]

The column construction also explains why this matrix maps coordinates
from frame \(\{b\}\) to frame \(\{a\}\). If

\[
{}^{b}\mathbf v
=
\begin{bmatrix}
v_x^{(b)}\\
v_y^{(b)}
\end{bmatrix},
\]

where \(v_x^{(b)},v_y^{(b)}\in\mathbb R\) are the components of
\(\mathbf v\) along the \(x_b\)- and \(y_b\)-axes, then

\[
\begin{aligned}
{}^{a}\mathbf v
&=v_x^{(b)}\,{}^{a}\mathbf e_{x,b}
+v_y^{(b)}\,{}^{a}\mathbf e_{y,b}\\
&=
\begin{bmatrix}
{}^{a}\mathbf e_{x,b}&{}^{a}\mathbf e_{y,b}
\end{bmatrix}
\begin{bmatrix}
v_x^{(b)}\\
v_y^{(b)}
\end{bmatrix}\\
&=\mathbf R_{ab}(\theta)\,{}^{b}\mathbf v.
\end{aligned}
\]

For the worked example, \(\theta=\pi/2\,\mathrm{rad}\). Substitution
gives

\[
\mathbf R_{ab}(\pi/2)
=
\begin{bmatrix}
0&-1\\
1&0
\end{bmatrix}.
\]

The first column \([0\;1]^{\mathsf T}\) says that the positive
\(x_b\)-axis points along the positive \(y_a\)-axis. The second column
\([-1\;0]^{\mathsf T}\) says that the positive \(y_b\)-axis points along
the negative \(x_a\)-axis.

Let \({}^{b}\mathbf v\in\mathbb R^2\) be the coordinate column of a free
vector \(\mathbf v\), expressed in frame \(\{b\}\) and measured in
metres:

\[
{}^{b}\mathbf{v}
=
\begin{bmatrix}
2\\
1
\end{bmatrix}
\mathrm{m}.
\]

The coordinate column of the same geometric vector in frame \(\{a\}\) is
\({}^{a}\mathbf v\in\mathbb R^2\):

\[
{}^{a}\mathbf{v}
=
\mathbf{R}_{ab}(\theta)\,{}^{b}\mathbf{v}
=
\begin{bmatrix}
0 & -1\\
1 & 0
\end{bmatrix}
\begin{bmatrix}
2\\
1
\end{bmatrix}
=
\begin{bmatrix}
-1\\
2
\end{bmatrix}
\mathrm{m}.
\]

The coordinates changed, but the physical length did not:

\[
\left\|{}^{b}\mathbf{v}\right\|_2
=
\sqrt{2^2+1^2}
=
\sqrt{5}
=
\sqrt{(-1)^2+2^2}
=
\left\|{}^{a}\mathbf{v}\right\|_2.
\]

\subsection{Interactive frame-rotation
graph}\label{interactive-frame-rotation-graph}

Fix the frame-\(\{b\}\) coordinate column

\[
{}^{b}\mathbf v
=
\begin{bmatrix}
2\\
1
\end{bmatrix}.
\]

For each angle \(\theta\in[-\pi,\pi]\), the corresponding
frame-\(\{a\}\) coordinate column is

\[
{}^{a}\mathbf v(\theta)
=
\mathbf R_{ab}(\theta)\,{}^{b}\mathbf v
=
\begin{bmatrix}
2\cos\theta-\sin\theta\\
2\sin\theta+\cos\theta
\end{bmatrix}.
\]

In the HTML book, move the angle slider to rotate frame \(\{b\}\)
relative to frame \(\{a\}\). The graph keeps
\({}^{b}\mathbf v=[2\;1]^{\mathsf T}\) fixed, so the geometric vector
rotates with frame \(\{b\}\) and its coordinates
\({}^{a}\mathbf v(\theta)\) update continuously.

The PDF shows the default slider state \(\theta=\pi/2\,\mathrm{rad}\).
Use the HTML book to vary \(\theta\) interactively.

\begin{figure}[H]

\centering{

\includegraphics[width=0.85\linewidth,height=\textheight,keepaspectratio]{index_files/mediabag/figures/planar_frame_rotation_static.pdf}

}

\caption{\label{fig-planar-frame-rotation-static}Static PDF view of the
frame rotation at \(\theta=\pi/2\) with
\({}^{b}\mathbf v=[2\;1]^{\mathsf T}\) and
\({}^{a}\mathbf v=[-1\;2]^{\mathsf T}\).}

\end{figure}%

\subsubsection{Equivalent NumPy
calculation}\label{equivalent-numpy-calculation}

The following NumPy snippet is the direct software translation of
\({}^{a}\mathbf v=\mathbf R_{ab}(\theta){}^{b}\mathbf v\). The scalar
\texttt{theta} stores \(\theta\) in radians; \texttt{R\_ab} is a NumPy
array with shape \texttt{(2,\ 2)} representing
\(\mathbf R_{ab}(\theta)\in\mathbb R^{2\times2}\); \texttt{v\_b} is a
NumPy array with shape \texttt{(2,)} representing the fixed coordinate
column \({}^{b}\mathbf v=[2\;1]^{\mathsf T}\); and \texttt{v\_a} is a
NumPy array with shape \texttt{(2,)} representing
\({}^{a}\mathbf v\in\mathbb R^2\). The operator \texttt{@} performs
matrix-vector multiplication.

\begin{Shaded}
\begin{Highlighting}[]
\ImportTok{import}\NormalTok{ numpy }\ImportTok{as}\NormalTok{ np}

\NormalTok{theta }\OperatorTok{=}\NormalTok{ np.pi }\OperatorTok{/} \DecValTok{2}  \CommentTok{\# radians; replace this value with the selected angle}

\NormalTok{R\_ab }\OperatorTok{=}\NormalTok{ np.array([}
\NormalTok{    [np.cos(theta), }\OperatorTok{{-}}\NormalTok{np.sin(theta)],}
\NormalTok{    [np.sin(theta),  np.cos(theta)],}
\NormalTok{])}

\NormalTok{v\_b }\OperatorTok{=}\NormalTok{ np.array([}\FloatTok{2.0}\NormalTok{, }\FloatTok{1.0}\NormalTok{])}
\NormalTok{v\_a }\OperatorTok{=}\NormalTok{ R\_ab }\OperatorTok{@}\NormalTok{ v\_b}
\end{Highlighting}
\end{Shaded}

For \(\theta=\pi/2\,\mathrm{rad}\), the calculation gives \texttt{v\_a}
approximately equal to \texttt{{[}-1.0,\ 2.0{]}}. Changing
\texttt{theta} recomputes the same values displayed by the interactive
graph while \texttt{v\_b} remains fixed.

\section{Differential-drive
interpretation}\label{differential-drive-interpretation}

For a planar mobile robot, let frame \(\{w\}\) be fixed to the world and
frame \(\{b\}\) be attached to the robot. Let \(\theta\in\mathbb R\) be
the robot's heading angle, measured counter-clockwise in radians from
the positive \(x_w\)-axis to the positive \(x_b\)-axis. Let
\(v\in\mathbb R\) be the signed forward speed, measured in
\(\mathrm{m\,s^{-1}}\); positive \(v\) means motion along the positive
\(x_b\)-axis.

The body-frame velocity \({}^{b}\mathbf v\in\mathbb R^2\) is

\[
{}^{b}\mathbf{v}
=
\begin{bmatrix}
v\\
0
\end{bmatrix}
\mathrm{m\,s^{-1}}.
\]

Define \(\mathbf R_{wb}(\theta)\in\mathbb R^{2\times2}\) as the rotation
matrix that maps planar coordinates from the robot frame \(\{b\}\) to
the world frame \(\{w\}\). The world-frame velocity
\({}^{w}\mathbf v\in\mathbb R^2\) is therefore

\[
{}^{w}\mathbf{v}
=
\mathbf{R}_{wb}(\theta)\,{}^{b}\mathbf{v}
=
\begin{bmatrix}
\cos\theta & -\sin\theta\\
\sin\theta & \cos\theta
\end{bmatrix}
\begin{bmatrix}
v\\
0
\end{bmatrix}
=
\begin{bmatrix}
v\cos\theta\\
v\sin\theta
\end{bmatrix}
\mathrm{m\,s^{-1}}.
\]

This equation is the coordinate-frame foundation of the
differential-drive kinematic model. It shows why a controller must
distinguish a forward command in the robot frame from motion observed in
the world frame.

\section{Failure modes and engineering
checks}\label{failure-modes-and-engineering-checks}

\subsection{Missing frame labels}\label{missing-frame-labels}

Two identical columns can describe different physical directions when
they are expressed in different frames. Interfaces should name the frame
explicitly.

\subsection{Reversed transformation}\label{reversed-transformation}

Using \(\mathbf{R}_{ba}\) where \(\mathbf{R}_{ab}\) is required rotates
coordinates in the opposite direction. Check the subscript rule before
multiplying.

\subsection{Mixed units}\label{mixed-units}

Degrees passed to a function expecting radians produce valid-looking but
wrong numbers. Convert at the system boundary and use radians
internally.

\subsection{Point-vector confusion}\label{point-vector-confusion}

A point changes under translation; a free vector does not. Applying a
translation to velocity or direction data is a modelling error.

\subsection{Invalid numerical
rotation}\label{invalid-numerical-rotation}

Let \(\mathbf R\in\mathbb R^{n\times n}\) be a numerically stored matrix
intended to represent a rotation, and let
\(\mathbf I_n\in\mathbb R^{n\times n}\) be the identity matrix.
Rounding, integration drift, or corrupted data can make \(\mathbf R\)
non-orthonormal.

For a matrix \(\mathbf M\in\mathbb R^{n\times n}\) with entries
\(m_{ij}\), define the Frobenius norm by

\[
\lVert\mathbf M\rVert_F
\mathrel{:=}
\sqrt{\sum_{i=1}^{n}\sum_{j=1}^{n}m_{ij}^2}.
\]

Let \(\varepsilon_{\mathrm{orth}}>0\) and \(\varepsilon_{\det}>0\) be
declared numerical tolerances. A robust interface checks

\[
\left\lVert\mathbf{R}^\mathsf{T}\mathbf{R}-\mathbf{I}_n\right\rVert_F
\leq
\varepsilon_{\mathrm{orth}}
\]

and

\[
\left|\det(\mathbf R)-1\right|
\leq
\varepsilon_{\det}.
\]

The first condition tests approximate orthonormality; the second tests
approximate preservation of right-handed orientation.

\section*{Quick test B: coordinate frames and
rotations}\label{quick-test-b-coordinate-frames-and-rotations}
\addcontentsline{toc}{section}{Quick test B: coordinate frames and
rotations}

\markright{Quick test B: coordinate frames and rotations}

\begin{enumerate}
\def\labelenumi{\arabic{enumi}.}
\tightlist
\item
  \textbf{Recall:} What do the columns of \(\mathbf{R}_{ab}\) represent?
\item
  \textbf{Explanation:} Explain in words why
  \(\mathbf{R}_{ab}^{-1}=\mathbf{R}_{ba}\).
\item
  \textbf{Derivation:} Starting with
  \({}^{a}\mathbf{v}=\mathbf{R}_{ab}{}^{b}\mathbf{v}\), solve for
  \({}^{b}\mathbf{v}\).
\item
  \textbf{Application:} A robot has heading \(\theta=\pi/2\) and forward
  speed \(v=0.5\,\mathrm{m\,s^{-1}}\). Find its world-frame velocity.
\item
  \textbf{Limitations:} Give two ways a numerically stored matrix
  intended to be a rotation matrix can become invalid.
\item
  \textbf{Derivation:} Let \(\{a\}\) and \(\{b\}\) be right-handed
  orthonormal planar frames. Frame \(\{b\}\) is rotated
  counter-clockwise by \(\theta\,\mathrm{rad}\) relative to frame
  \(\{a\}\), and \(\mathbf R_{ab}(\theta)\in\mathbb R^{2\times2}\) maps
  coordinate columns from frame \(\{b\}\) to frame \(\{a\}\). First
  express the \(x_b\)-axis and \(y_b\)-axis as coordinate columns
  \({}^{a}\mathbf e_{x,b},{}^{a}\mathbf e_{y,b}\in\mathbb R^2\); then
  place those columns side by side to derive \(\mathbf R_{ab}(\theta)\).
\end{enumerate}

\subsection*{Answers and hints}\label{answers-and-hints-5}
\addcontentsline{toc}{subsection}{Answers and hints}

\begin{enumerate}
\def\labelenumi{\arabic{enumi}.}
\item
  The axes of frame \(\{b\}\) expressed in coordinates of frame
  \(\{a\}\).
\item
  Reversing the frame change undoes the original change; for an
  orthonormal rotation matrix, the inverse is its transpose.
\item
  Premultiply by \(\mathbf{R}_{ab}^{-1}=\mathbf{R}_{ba}\):
  \({}^{b}\mathbf{v}=\mathbf{R}_{ba}{}^{a}\mathbf{v}\).
\item
  \([0\;\;0.5]^\mathsf{T}\,\mathrm{m\,s^{-1}}\), up to floating-point
  rounding.
\item
  Examples include accumulated integration error, rounded entries,
  malformed input, unit mistakes in angle construction, and data
  corruption.
\item
  Since the \(x_b\)-axis makes angle \(\theta\) with the \(x_a\)-axis,
  its frame-\(\{a\}\) coordinate column is
  \({}^{a}\mathbf e_{x,b}=[\cos\theta\;\sin\theta]^{\mathsf T}\).
  Right-handed orthonormality places the \(y_b\)-axis at angle
  \(\theta+\pi/2\), so
  \({}^{a}\mathbf e_{y,b}=[\cos(\theta+\pi/2)\;\sin(\theta+\pi/2)]^{\mathsf T}=[-\sin\theta\;\cos\theta]^{\mathsf T}\).
  Because these two coordinate columns are the columns of
  \(\mathbf R_{ab}(\theta)\),

  \[
  \mathbf R_{ab}(\theta)
  =
  \begin{bmatrix}
  \cos\theta&-\sin\theta\\
  \sin\theta&\cos\theta
  \end{bmatrix}.
  \]
\end{enumerate}

\section*{References}\label{references-1}
\addcontentsline{toc}{section}{References}

\markright{References}

\begin{itemize}
\tightlist
\item
  Kevin M. Lynch and Frank C. Park, \emph{Modern Robotics}, Chapter 3,
  especially Sections 3.1 and 3.2.1 on free vectors, frames, rotation
  matrices, and rigid-body motions.
\item
  Steven M. LaValle, \emph{Planning Algorithms}, Chapters 3 and 4 on
  geometric representations, transformations, and configuration spaces.
\end{itemize}

\chapter{Differential-Drive Kinematics and Numerical Integration: Model
and Derivation}\label{sec-differential-drive-model}

\section*{Why this chapter matters}\label{why-this-chapter-matters}
\addcontentsline{toc}{section}{Why this chapter matters}

\markright{Why this chapter matters}

Kinematics describes motion without asking which forces or torques
produce it. For a differential-drive robot, kinematics connects wheel
angular speeds to the time evolution of position and heading. Numerical
integration converts this continuous-time model into sampled updates
used by simulation, odometry, planning, and control.

This chapter develops that chain in four steps: define planar pose and
its derivative, derive the no-sideways-motion model, derive body
velocity from wheel speeds, and compare exact constant-input motion with
Forward Euler integration.

\section*{Prerequisites and
assumptions}\label{prerequisites-and-assumptions-2}
\addcontentsline{toc}{section}{Prerequisites and assumptions}

\markright{Prerequisites and assumptions}

The reader should understand vectors and matrices from
Chapter~\ref{sec-linear-algebra-foundations}, coordinate frames and
rotations from Chapter~\ref{sec-geometry-coordinate-frames}, and
elementary derivatives, integrals, sine, and cosine.

The system considered here is a differential-drive ground robot: a rigid
chassis propelled by left and right drive wheels that are mounted on a
common axle and can be driven independently. The robot operates on a
stationary, level, rigid floor, and the model describes its motion as
viewed from above. Any additional support hardware does not affect the
ideal kinematic model.

Frame \(\{w\}\) is fixed to the floor and is called the world frame.
Frame \(\{b\}\) is rigidly attached to the robot and is called the body
frame. Its origin \(O_b\) lies midway between the two drive-wheel
ground-contact centre lines. The positive \(x_b\)-axis points forward,
the positive \(y_b\)-axis points left, and the positive \(z_b\)-axis
points upward, so frame \(\{b\}\) is right-handed. As the robot moves,
frame \(\{b\}\) translates and rotates relative to the stationary frame
\(\{w\}\).

Both drive wheels have the same effective rolling radius \(r>0\),
measured in metres. The track width \(L>0\), measured in metres, is the
lateral distance between the left and right wheel ground-contact centre
lines.

The kinematic model uses the following assumptions:

\begin{itemize}
\tightlist
\item
  The chassis remains rigid, and its pose changes only through
  horizontal position and rotation about the upward \(z_b\)-axis; its
  height, roll, and pitch remain unchanged.
\item
  Each wheel maintains contact with the floor and rolls without
  longitudinal slip, so the distance travelled in the wheel's rolling
  direction equals the arc length swept out at its effective radius.
\item
  The scalar body-frame lateral component \({}^{b}v_y\), also written
  \(v_y\) when the frame is clear, is measured in \(\mathrm{m\,s^{-1}}\)
  and lies along the \(+y_b\)-direction. The wheel constraints impose
  \({}^{b}v_y=0\) at \(O_b\). The robot can still turn because its
  forward direction changes continuously.
\item
  The effective wheel radius \(r\) and track width \(L\) remain
  constant.
\item
  Wheel deformation, tyre slip, backlash, encoder error, actuator
  dynamics, collisions, and uneven or moving terrain are ignored.
\item
  Length is measured in metres, time in seconds, and angle in radians;
  corresponding velocities are measured in \(\mathrm{m\,s^{-1}}\) and
  \(\mathrm{rad\,s^{-1}}\).
\end{itemize}

Together, these assumptions define an ideal planar kinematic model
rather than a complete physical model. A real robot violates them to
some degree, especially during rapid acceleration, braking, collision,
wheel slip, or motion over compliant or uneven terrain.

Figure Figure~\ref{fig-differential-drive-geometry} summarizes the
frames, wheel dimensions, rolling condition, and lateral-velocity
constraint used by the model.

\begin{figure}

\centering{

\includegraphics[width=0.95\linewidth,height=\textheight,keepaspectratio]{index_files/mediabag/figures/differential_drive_geometry.pdf}

}

\caption{\label{fig-differential-drive-geometry}Differential-drive robot
geometry and ideal rolling constraints.}

\end{figure}%

As a preview of the pure-rolling relation, \(v_i=r\dot\phi_i\) for
\(i\in\{L,R\}\), where \(v_i\) is the signed longitudinal ground speed
of wheel \(i\) and \(\dot\phi_i\) is its signed angular speed.
Section~\ref{sec-wheel-speeds-body-motion} derives how the individual
wheel speeds determine the body-forward speed \(v\) and yaw rate
\(\omega\).

\section{Planar pose and its
derivative}\label{planar-pose-and-its-derivative}

Let \(t\in\mathbb R_{\geq0}\) denote continuous time, measured in
seconds. Define the planar robot pose by

\[
\mathbf q(t)
\mathrel{:=}
\begin{bmatrix}
x(t)\\
y(t)\\
\theta(t)
\end{bmatrix}.
\]

Here, \(x(t),y(t)\in\mathbb R\) are the world-frame coordinates of
body-frame origin \(O_b\), measured in metres, and
\(\theta(t)\in\mathbb R\) is the counter-clockwise heading of frame
\(\{b\}\) relative to frame \(\{w\}\), measured in radians. Angles that
differ by an integer multiple of \(2\pi\) represent the same physical
heading.

The pose derivative is

\[
\dot{\mathbf q}(t)
\mathrel{:=}
\frac{d\mathbf q(t)}{dt}
=
\begin{bmatrix}
\dot x(t)\\
\dot y(t)\\
\dot\theta(t)
\end{bmatrix}.
\]

The components \(\dot x\) and \(\dot y\) are measured in
\(\mathrm{m\,s^{-1}}\), while \(\dot\theta\) is measured in
\(\mathrm{rad\,s^{-1}}\). Although they appear in one column, their
units are not identical.

Define the body-velocity input by

\[
\mathbf u(t)
\mathrel{:=}
\begin{bmatrix}
v(t)\\
\omega(t)
\end{bmatrix},
\]

where \(v(t)\in\mathbb R\) is signed forward speed along \(+x_b\),
measured in \(\mathrm{m\,s^{-1}}\), and \(\omega(t)\in\mathbb R\) is
signed yaw rate about \(+z_b\), measured in \(\mathrm{rad\,s^{-1}}\).
Positive \(\omega\) produces counter-clockwise rotation when viewed from
above.

\subsection{Deriving the planar model}\label{deriving-the-planar-model}

Because the ideal robot has no body-frame lateral velocity, its
translational velocity expressed in frame \(\{b\}\) is

\[
{}^{b}\mathbf v
=
\begin{bmatrix}
v\\
0
\end{bmatrix}
\in\mathbb R^2.
\]

The rotation matrix \(\mathbf R_{wb}(\theta)\in\mathbb R^{2\times2}\)
maps coordinates from frame \(\{b\}\) to frame \(\{w\}\). Therefore,

\[
\begin{bmatrix}
\dot x\\
\dot y
\end{bmatrix}
=
\underbrace{
\begin{bmatrix}
\cos\theta&-\sin\theta\\
\sin\theta&\cos\theta
\end{bmatrix}
}_{\mathbf R_{wb}(\theta)}
\begin{bmatrix}
v\\
0
\end{bmatrix}
=
\begin{bmatrix}
v\cos\theta\\
v\sin\theta
\end{bmatrix}.
\]

Together with \(\dot\theta=\omega\), this gives

\[
\boxed{
\dot{\mathbf q}
=
\mathbf f(\mathbf q,\mathbf u)
=
\begin{bmatrix}
v\cos\theta\\
v\sin\theta\\
\omega
\end{bmatrix}
=
\begin{bmatrix}
\cos\theta&0\\
\sin\theta&0\\
0&1
\end{bmatrix}
\begin{bmatrix}
v\\
\omega
\end{bmatrix}
}.
\]

Here, \(\mathbf f\) is the kinematic state-derivative function: it maps
current pose \(\mathbf q\) and input \(\mathbf u\) to instantaneous pose
derivative \(\dot{\mathbf q}\). The position rates depend on \(\theta\)
because the body forward direction changes relative to the world frame.

\subsection{The no-sideways-motion
constraint}\label{the-no-sideways-motion-constraint}

Define the axle-midpoint linear velocity expressed in the world frame by

\[
{}^{w}\mathbf v
\mathrel{:=}
\begin{bmatrix}
\dot x\\
\dot y
\end{bmatrix}
\in\mathbb R^2,
\]

where \(\dot x\) and \(\dot y\) are its components along \(+x_w\) and
\(+y_w\), respectively, and are measured in \(\mathrm{m\,s^{-1}}\).

The geometric unit vector along the body-frame \(+y_b\)-axis, expressed
using world-frame coordinates, is the second column of
\(\mathbf R_{wb}(\theta)\):

\[
{}^{w}\mathbf e_{y,b}
=
\begin{bmatrix}
-\sin\theta\\
\cos\theta
\end{bmatrix}
\in\mathbb R^2.
\]

Both \({}^{w}\mathbf e_{y,b}\) and \({}^{w}\mathbf v\) are expressed in
frame \(\{w\}\), and \({}^{w}\mathbf e_{y,b}\) has unit length. Their
dot product therefore gives the signed component of the velocity along
the robot's lateral \(+y_b\)-direction:

\[
\begin{aligned}
{}^{b}v_y
&=
({}^{w}\mathbf e_{y,b})^{\mathsf T}{}^{w}\mathbf v\\
&=
\begin{bmatrix}
-\sin\theta & \cos\theta
\end{bmatrix}
\begin{bmatrix}
\dot x\\
\dot y
\end{bmatrix}\\
&=
-\dot x\sin\theta+
\dot y\cos\theta.
\end{aligned}
\]

The left superscript in \({}^{b}v_y\) makes explicit that this scalar is
the \(y\)-component of velocity measured along a body-frame axis, even
though the projection was calculated using world-frame coordinates.

The matrix \(\mathbf R_{wb}(\theta)\) maps coordinate columns from frame
\(\{b\}\) to frame \(\{w\}\). Its columns are orthonormal, so its
inverse equals its transpose. Define the reverse coordinate map by

\[
\begin{aligned}
\mathbf R_{bw}(\theta)
&\mathrel{:=}
\mathbf R_{wb}(\theta)^{-1}\\
&=
\mathbf R_{wb}(\theta)^{\mathsf T}\\
&=
\begin{bmatrix}
\cos\theta & \sin\theta\\
-\sin\theta & \cos\theta
\end{bmatrix}.
\end{aligned}
\]

Therefore, \(\mathbf R_{bw}(\theta)\) maps coordinate columns from frame
\(\{w\}\) to frame \(\{b\}\). The notation uses the same heading
\(\theta\) because both matrices depend on the robot's heading, but the
reverse map applies the opposite rotation:
\(\mathbf R_{bw}(\theta)=\mathbf R_{wb}(-\theta)\).

The complete change from world-frame velocity components to body-frame
velocity components is

\[
{}^{b}\mathbf v
=
\begin{bmatrix}
{}^{b}v_x\\
{}^{b}v_y
\end{bmatrix}
=
\mathbf R_{bw}(\theta){}^{w}\mathbf v
=
\mathbf R_{wb}(\theta)^{\mathsf T}{}^{w}\mathbf v
=
\begin{bmatrix}
\dot x\cos\theta+\dot y\sin\theta\\
-\dot x\sin\theta+\dot y\cos\theta
\end{bmatrix}.
\]

Here, \({}^{b}v_x\) and \({}^{b}v_y\) are the signed forward and lateral
velocity components along \(+x_b\) and \(+y_b\), respectively, and both
are measured in \(\mathrm{m\,s^{-1}}\). Thus, \(\dot x\) and \(\dot y\)
are world-frame components, whereas \({}^{b}v_x\) and \({}^{b}v_y\) are
body-frame components.

The ideal no-sideways-motion assumption is therefore written most
clearly as

\[
\boxed{
{}^{b}v_y
=
-\dot x\sin\theta+
\dot y\cos\theta
=
0
}.
\]

The signed forward speed \(v\) defined earlier equals \({}^{b}v_x\), so
the constrained body-frame velocity and its world-frame representation
are

\[
{}^{b}\mathbf v
=
\begin{bmatrix}
v\\
0
\end{bmatrix},
\qquad
{}^{w}\mathbf v
=
\mathbf R_{wb}(\theta)
\begin{bmatrix}
v\\
0
\end{bmatrix}
=
\begin{bmatrix}
v\cos\theta\\
v\sin\theta
\end{bmatrix}.
\]

Consequently, zero body-frame lateral velocity does not imply
\(\dot y=0\). For example, when \(\theta=\pi/4\), forward motion gives
\(\dot x=\dot y=v/\sqrt{2}\). The angle \(\theta\) specifies
orientation; it does not itself create translation. If \(v=0\) and
\(\omega\ne0\), the robot rotates in place about \(O_b\), so
\(\dot x=\dot y=0\) while \(\theta\) changes.

This equation is a nonholonomic velocity constraint. It restricts the
directions in which the robot can move at the current pose, but it does
not require the pose variables \(x\), \(y\), and \(\theta\) to remain on
a fixed lower-dimensional surface. In contrast, a holonomic condition
such as \(y=0\) would confine every allowable pose to the world
\(x_w\)-axis.

The distinction concerns instantaneous motion versus accumulated motion.
At every instant, the axle midpoint translates only along the robot's
current \(x_b\)-axis, but turning changes the direction of \(x_b\)
relative to the world. For example, a robot initially aligned with
\(+x_w\) can rotate by \(\pi/2\), drive forward a distance \(d>0\), and
rotate back by \(-\pi/2\). Its final heading equals its initial heading,
yet it has acquired a world-frame \(+y_w\) displacement of \(d\). It
never moved along its instantaneous \(y_b\)-axis; the net sideways
displacement resulted from a sequence of allowed rotations and forward
motions.

\subsection*{Quick test A: planar
kinematics}\label{quick-test-a-planar-kinematics}
\addcontentsline{toc}{subsection}{Quick test A: planar kinematics}

\begin{enumerate}
\def\labelenumi{\arabic{enumi}.}
\tightlist
\item
  Define every component of \(\mathbf q=[x\;y\;\theta]^\mathsf T\),
  including its frame and units.
\item
  Why do \(\dot x\) and \(\dot y\) depend on \(\theta\), while \(v\)
  does not?
\item
  Starting with \({}^{b}\mathbf v=[v\;0]^\mathsf T\), derive
  \(\dot y=v\sin\theta\).
\item
  If \(v=1\,\mathrm{m\,s^{-1}}\) and \(\theta=\pi/2\,\mathrm{rad}\),
  find \(\dot x\) and \(\dot y\).
\item
  Does the no-sideways-motion constraint prevent a net sideways
  displacement? Explain.
\end{enumerate}

\subsubsection*{Answers and hints}\label{answers-and-hints-6}
\addcontentsline{toc}{subsubsection}{Answers and hints}

\begin{enumerate}
\def\labelenumi{\arabic{enumi}.}
\tightlist
\item
  \(x\) and \(y\) are the world-frame position coordinates of \(O_b\) in
  metres; \(\theta\) is counter-clockwise body heading relative to the
  world frame in radians.
\item
  \(v\) is defined along the moving \(x_b\)-axis. Its world-frame
  components depend on the orientation of that axis.
\item
  The second row of \(\mathbf R_{wb}(\theta)[v\;0]^\mathsf T\) is
  \(v\sin\theta\).
\item
  \(\dot x=0\,\mathrm{m\,s^{-1}}\) and \(\dot y=1\,\mathrm{m\,s^{-1}}\).
\item
  No.~The constraint forbids instantaneous sideways velocity, but
  combined forward and turning motions can produce a sideways
  displacement.
\end{enumerate}

\section{Differential-drive wheel
kinematics}\label{differential-drive-wheel-kinematics}

Let \(\phi_L(t),\phi_R(t)\in\mathbb R\) be the rotational positions of
the left and right wheels about their axles at time \(t\), measured in
radians. The subscripts \(L\) and \(R\) identify the left and right
wheels. These are rolling angles, not steering angles: the wheel planes
remain fixed relative to the chassis.

Choose \(\phi_L=\phi_R=0\) at an arbitrary reference position. A change
of \(2\pi\) radians is one complete wheel revolution. Although the
marked orientation of a wheel repeats after each revolution, \(\phi_L\)
and \(\phi_R\) are treated as unwrapped accumulated angles in
\(\mathbb R\), so they may increase beyond \(2\pi\) or decrease below
zero.

Their time derivatives \(\dot\phi_L(t)\) and \(\dot\phi_R(t)\) are the
left and right wheel angular speeds, measured in
\(\mathrm{rad\,s^{-1}}\). Positive angular speed is defined to roll the
corresponding wheel and the robot forward.

For a wheel of effective radius \(r\), let \(\Delta\phi\) be a signed
angular change and let \(\Delta s\) be the corresponding signed arc
displacement swept out at the wheel rim. Both \(r\) and \(\Delta s\) are
measured in metres. By the definition of radian measure,

\[
\Delta\phi
=
\frac{\Delta s}{r},
\qquad
\frac{\mathrm m}{\mathrm m}=1.
\]

The metre units cancel, so a radian is dimensionless in the
International System of Units (SI). Nevertheless, angles are labelled in
\(\mathrm{rad}\) and angular speeds in \(\mathrm{rad\,s^{-1}}\) so their
physical meanings remain clear. Rearranging the radian definition gives
the swept arc displacement \(\Delta s=r\Delta\phi\). Under rolling
without slip, this swept arc displacement equals the signed longitudinal
displacement along the ground.

For wheel \(i\in\{L,R\}\), let \(s_i(t)\) be its signed accumulated
longitudinal ground displacement. Over a time interval \(\Delta t>0\),
pure rolling gives \(\Delta s_i=r\Delta\phi_i\). Because \(r\) is
constant, dividing by \(\Delta t\) and taking the limit \(\Delta t\to0\)
gives

\[
v_i
\mathrel{:=}
\dot s_i
=
\lim_{\Delta t\to0}\frac{\Delta s_i}{\Delta t}
=
r\lim_{\Delta t\to0}\frac{\Delta\phi_i}{\Delta t}
=
r\dot\phi_i.
\]

Applying this relation to the left and right wheels gives

\[
v_L=r\dot\phi_L,
\qquad
v_R=r\dot\phi_R,
\]

where \(v_L,v_R\in\mathbb R\) are measured in \(\mathrm{m\,s^{-1}}\).
The unit calculation for either wheel is

\[
\begin{aligned}
\text{metres}\times\frac{\text{radians}}{\text{second}}
&=
\mathrm m\cdot\frac{\mathrm{rad}}{\mathrm s}\\
&=
\mathrm m\cdot\frac{1}{\mathrm s}\\
&=
\mathrm{m\,s^{-1}}.
\end{aligned}
\]

The symbol \(\mathrm{rad}\) is retained to communicate that \(\phi\) is
an angle, even though it contributes no independent physical dimension
in the unit calculation.

\subsection{Relating wheel speeds to body
motion}\label{sec-wheel-speeds-body-motion}

The rigid-body point-velocity relation derived in the
\hyperref[sec-geometry-coordinate-frames]{geometry and coordinate-frames
chapter} is

\[
{}^{b}\mathbf v_P
=
{}^{b}\mathbf v_O
+
{}^{b}\boldsymbol{\omega}\times{}^{b}\mathbf r_{OP},
\]

provided that every vector is expressed in the same orthonormal frame,
here the body frame \(\{b\}\). The vector \({}^{b}\mathbf v_P\) is the
velocity of a rigid-body point \(P\), \({}^{b}\mathbf v_O\) is the
velocity of a reference point \(O\), \({}^{b}\mathbf r_{OP}\) is the
position of \(P\) relative to \(O\), and \({}^{b}\boldsymbol{\omega}\)
is the rigid body angular-velocity vector.

For the differential-drive chassis, choose the reference point as the
axle midpoint \(O=O_b\). Let \(C_L\) and \(C_R\) denote the left and
right wheel centres on the axle. Both points are fixed relative to the
chassis. Under ideal rolling without longitudinal slip, the \(+x_b\)
components of their translational velocities equal the longitudinal
wheel speeds \(v_L\) and \(v_R\) defined above.

Because \(+y_b\) points left, the wheel-centre position vectors are

\[
{}^{b}\mathbf r_{O_bC_L}
=
\begin{bmatrix}0\\L/2\\0\end{bmatrix},
\qquad
{}^{b}\mathbf r_{O_bC_R}
=
\begin{bmatrix}0\\-L/2\\0\end{bmatrix},
\qquad
{}^{b}\mathbf r_{O_bC_L},{}^{b}\mathbf r_{O_bC_R}\in\mathbb R^3.
\]

Their components have units of metres. The axle midpoint moves forward
with scalar speed \(v\), and the chassis yaws about \(+z_b\) with scalar
rate \(\omega\), so

\[
{}^{b}\mathbf v_{O_b}
=
\begin{bmatrix}v\\0\\0\end{bmatrix}
\in\mathbb R^3,
\qquad
{}^{b}\boldsymbol{\omega}
=
\begin{bmatrix}0\\0\\\omega\end{bmatrix}
\in\mathbb R^3.
\]

The components of \({}^{b}\mathbf v_{O_b}\) have units of
\(\mathrm{m\,s^{-1}}\), and the components of
\({}^{b}\boldsymbol{\omega}\) have units of \(\mathrm{rad\,s^{-1}}\).

Applying the rigid-body relation to each wheel-centre point gives

\[
{}^{b}\mathbf v_{C_L}
=
{}^{b}\mathbf v_{O_b}
+
{}^{b}\boldsymbol{\omega}\times{}^{b}\mathbf r_{O_bC_L},
\qquad
{}^{b}\mathbf v_{C_R}
=
{}^{b}\mathbf v_{O_b}
+
{}^{b}\boldsymbol{\omega}\times{}^{b}\mathbf r_{O_bC_R}.
\]

For the left wheel, the rotational contribution is

\[
\begin{aligned}
{}^{b}\boldsymbol{\omega}\times{}^{b}\mathbf r_{O_bC_L}
&=
\begin{bmatrix}0\\0\\\omega\end{bmatrix}
\times
\begin{bmatrix}0\\L/2\\0\end{bmatrix}
=
\begin{bmatrix}-L\omega/2\\0\\0\end{bmatrix}.
\end{aligned}
\]

For the right wheel, it is

\[
\begin{aligned}
{}^{b}\boldsymbol{\omega}\times{}^{b}\mathbf r_{O_bC_R}
&=
\begin{bmatrix}0\\0\\\omega\end{bmatrix}
\times
\begin{bmatrix}0\\-L/2\\0\end{bmatrix}
=
\begin{bmatrix}L\omega/2\\0\\0\end{bmatrix}.
\end{aligned}
\]

Therefore, the wheel-centre velocity vectors are

\[
{}^{b}\mathbf v_{C_L}
=
\begin{bmatrix}v-L\omega/2\\0\\0\end{bmatrix},
\qquad
{}^{b}\mathbf v_{C_R}
=
\begin{bmatrix}v+L\omega/2\\0\\0\end{bmatrix}.
\]

The longitudinal wheel speeds \(v_L\) and \(v_R\) are the \(+x_b\)
components of these vectors. Consequently,

\[
\boxed{
v_L=v-\frac{L}{2}\omega,
\qquad
v_R=v+\frac{L}{2}\omega
}.
\]

The signs follow from the right-hand-rule identities
\(\mathbf e_{z,b}\times\mathbf e_{y,b}=-\mathbf e_{x,b}\) and
\(\mathbf e_{z,b}\times(-\mathbf e_{y,b})=\mathbf e_{x,b}\). For a
positive yaw rate, the robot turns left: rotation reduces the left-wheel
forward speed and increases the right-wheel forward speed. The product
\(L\omega\) has units \(\mathrm{m\,s^{-1}}\) because radians are
dimensionless in dimensional calculations.

Adding the two equations isolates forward speed, while subtracting the
left equation from the right equation isolates yaw rate:

\[
v=\frac{v_R+v_L}{2},
\qquad
\omega=\frac{v_R-v_L}{L}.
\]

Substituting the pure-rolling relations produces

\[
\boxed{
v
=
\frac{r}{2}\left(\dot\phi_R+\dot\phi_L\right),
\qquad
\omega
=
\frac{r}{L}\left(\dot\phi_R-\dot\phi_L\right)
}.
\]

In matrix form,

\[
\boxed{
\begin{bmatrix}
v\\
\omega
\end{bmatrix}
=
\underbrace{
\begin{bmatrix}
r/2&r/2\\
-r/L&r/L
\end{bmatrix}
}_{\mathbf G\in\mathbb R^{2\times2}}
\begin{bmatrix}
\dot\phi_L\\
\dot\phi_R
\end{bmatrix}
}.
\]

The matrix \(\mathbf G\) maps the wheel-speed column
\([\dot\phi_L\;\dot\phi_R]^\mathsf T\) to the body-velocity column
\([v\;\omega]^\mathsf T\). In this chapter, this input-output direction
is called the \textbf{wheel-to-body velocity mapping}: its inputs are
the two wheel angular speeds, and its outputs are the resulting
body-forward speed and yaw rate. It is also called the \textbf{forward
differential-drive velocity kinematics}. Here, \emph{velocity
kinematics} means an algebraic relationship among instantaneous
velocities without integrating the pose or modelling the forces that
produce the motion. The word \emph{forward} describes the calculation
from wheel motion to body motion; it does not mean that the robot must
travel in the positive forward direction, because the signed wheel
speeds and \(v\) may be negative.

The two rows of \(\mathbf G\) have different physical units because the
two outputs have different units.

Because \(r>0\) and \(L>0\), the mapping can be inverted. Solving for
the wheel speeds gives

\[
\boxed{
\dot\phi_L=\frac{v}{r}-\frac{L\omega}{2r},
\qquad
\dot\phi_R=\frac{v}{r}+\frac{L\omega}{2r}
}.
\]

This opposite input-output direction is called the \textbf{body-to-wheel
velocity mapping}, or the \textbf{inverse differential-drive velocity
kinematics}. Its inputs are a desired body-forward speed \(v\) and yaw
rate \(\omega\), and its outputs are the ideal left and right wheel
angular speeds \(\dot\phi_L\) and \(\dot\phi_R\) required by the no-slip
model. The word \emph{inverse} means that this calculation reverses the
wheel-to-body mapping; it does not mean that the robot is commanded to
drive backwards. In a controller, the body-to-wheel mapping converts a
desired chassis motion into wheel-speed commands, whereas the
wheel-to-body mapping can convert measured wheel speeds into an estimate
of chassis motion.

\subsection{Characteristic motions}\label{characteristic-motions}

\begin{itemize}
\tightlist
\item
  If \(\dot\phi_L=\dot\phi_R\), then \(\omega=0\) and the robot moves
  straight.
\item
  If \(\dot\phi_L=-\dot\phi_R\), then \(v=0\) and the robot rotates
  about \(O_b\).
\item
  If one wheel is stationary, the ideal robot follows a circle about
  that wheel contact point.
\item
  If the wheel speeds are unequal and not opposite, the robot follows an
  arc about an instantaneous centre; the signed turning radius is
  derived next.
\end{itemize}

\subsubsection{Instantaneous centre of curvature and signed turning
radius}\label{instantaneous-centre-of-curvature-and-signed-turning-radius}

Assume \(\omega\neq0\) at time \(t\). At that instant, the planar
rigid-body motion can be regarded as rotation about a point \(C\) in the
plane whose instantaneous linear velocity is zero. This point is called
the instantaneous centre of curvature (ICC). The body origin \(O_b\)
therefore moves instantaneously along a circular arc centred at \(C\).

The no-sideways-motion constraint makes the velocity of \(O_b\) tangent
to the arc and parallel to the body \(x_b\)-axis. A circle radius is
perpendicular to its tangent, so the line joining \(O_b\) and \(C\) lies
along the lateral \(y_b\)-axis.

Let \({}^{w}\mathbf p_{O_b},{}^{w}\mathbf p_C\in\mathbb R^2\) be the
world-frame position columns of \(O_b\) and \(C\), respectively. Define
the ordinary nonnegative radius \(\rho\), measured in metres, by

\[
\rho
\mathrel{:=}
\operatorname{distance}(O_b,C)
=
\left\|{}^{w}\mathbf p_C-{}^{w}\mathbf p_{O_b}\right\|_2
\geq0.
\]

The magnitude \(\rho\) records only the distance between the two points.
Define the signed turning radius \(R\in\mathbb R\) by using its sign to
record which side of the robot contains the ICC:

\[
R
\mathrel{:=}
\begin{cases}
+\rho, & C\text{ is to the left of }O_b,\\[4pt]
-\rho, & C\text{ is to the right of }O_b.
\end{cases}
\]

Consequently, \(|R|=\rho\). The relative-position vector from \(O_b\) to
\(C\), expressed in world-frame coordinates, is

\[
{}^{w}\mathbf r_{O_bC}
\mathrel{:=}
{}^{w}\mathbf p_C-{}^{w}\mathbf p_{O_b}
=
R\,{}^{w}\mathbf e_{y,b},
\qquad
{}^{w}\mathbf e_{y,b}
=
\begin{bmatrix}-\sin\theta\\\cos\theta\end{bmatrix}
\in\mathbb R^2.
\]

The unit vector \({}^{w}\mathbf e_{y,b}\) points along the robot lateral
\(+y_b\)-axis and is expressed in world coordinates. Multiplying it by
\(R>0\) places \(C\) to the left of \(O_b\), while multiplying it by
\(R<0\) reverses the direction and places \(C\) to the right. Figure
Figure~\ref{fig-signed-turning-radius} shows both cases. The forward
velocity is tangent to the circular path and therefore perpendicular to
the radius direction.

\begin{figure}

\centering{

\includegraphics[width=0.95\linewidth,height=\textheight,keepaspectratio]{index_files/mediabag/figures/signed_turning_radius.pdf}

}

\caption{\label{fig-signed-turning-radius}Relationship among the body
origin, body axes, instantaneous centre of curvature, tangent velocity,
and signed turning radius.}

\end{figure}%

To derive the relationship among \(R\), \(v\), and \(\omega\), consider
a short interval \(\Delta t>0\) during which the instantaneous values
are treated as constant. Define the signed heading change by

\[
\Delta\theta
\mathrel{:=}
\theta(t+\Delta t)-\theta(t),
\]

where \(\Delta\theta\) is measured in radians. Let
\(\Delta\ell\in\mathbb R\) be the signed distance travelled by \(O_b\)
along the arc, measured in metres; \(\Delta\ell>0\) denotes travel in
the \(+x_b\)-direction and \(\Delta\ell<0\) denotes travel in the
\(-x_b\)-direction.

The radian definition gives arc length equal to radius times angular
change. With the signed quantities just defined, it becomes

\[
\boxed{\Delta\ell=R\,\Delta\theta}.
\]

For a forward left turn, both \(R\) and \(\Delta\theta\) are positive.
For a forward right turn, both are negative. Their product is therefore
the positive forward distance \(\Delta\ell\) in either case; the same
signed equation also remains valid for backward motion.

Divide the signed arc equation by \(\Delta t\):

\[
\frac{\Delta\ell}{\Delta t}
=
R\frac{\Delta\theta}{\Delta t}.
\]

Taking the limit as \(\Delta t\to0\) gives

\[
\frac{d\ell}{dt}
=
R\frac{d\theta}{dt}.
\]

By definition, the signed forward speed is \(v=d\ell/dt\), measured in
\(\mathrm{m\,s^{-1}}\), and the signed yaw rate is
\(\omega=d\theta/dt\), measured in \(\mathrm{rad\,s^{-1}}\). Therefore,

\[
\boxed{v=R\omega}.
\]

When \(\omega\neq0\), division by \(\omega\) yields

\[
\boxed{
R=\frac{v}{\omega}
},
\qquad
\omega\neq0.
\]

The quotient has units of metres because radians are dimensionless in
dimensional calculations. If \(\omega=0\) and \(v\neq0\), the robot
moves straight and has no finite ICC; its turning radius is informally
described as infinite. If \(v=0\) and \(\omega\neq0\), then \(R=0\) and
the robot rotates in place about \(O_b\).

The symbol \(R\) in this subsection is a scalar signed turning radius,
not a rotation matrix.

\subsection*{Quick test B: wheel-to-body
mapping}\label{quick-test-b-wheel-to-body-mapping}
\addcontentsline{toc}{subsection}{Quick test B: wheel-to-body mapping}

\begin{enumerate}
\def\labelenumi{\arabic{enumi}.}
\tightlist
\item
  Define \(\phi_L\), \(\phi_R\), \(\dot\phi_L\), \(\dot\phi_R\), \(r\),
  and \(L\), including their units, and distinguish wheel angle from
  angular speed.
\item
  Why is a radian dimensionless in an SI unit calculation, and why does
  \(r\dot\phi\) have units of metres per second?
\item
  Starting from
  \({}^{b}\mathbf v_P={}^{b}\mathbf v_O+{}^{b}\boldsymbol{\omega}\times{}^{b}\mathbf r_{OP}\),
  use the two wheel-centre position vectors to derive
  \(v_L=v-L\omega/2\) and \(v_R=v+L\omega/2\).
\item
  Why does equal wheel speed produce \(\omega=0\)?
\item
  Derive \(\omega=r(\dot\phi_R-\dot\phi_L)/L\).
\item
  Let \(r=0.10\,\mathrm m\), \(L=0.50\,\mathrm m\),
  \(\dot\phi_L=2\,\mathrm{rad\,s^{-1}}\), and
  \(\dot\phi_R=4\,\mathrm{rad\,s^{-1}}\). Find \(v\) and \(\omega\).
\item
  Define the ICC \(C\), ordinary radius \(\rho\), signed radius \(R\),
  signed distance \(\Delta\ell\), and heading change \(\Delta\theta\).
  Starting from \(\Delta\ell=R\Delta\theta\), derive \(R=v/\omega\) and
  explain the signs for left and right turns.
\item
  Name two physical effects that can make predicted and real motion
  disagree.
\end{enumerate}

\subsubsection*{Answers and hints}\label{answers-and-hints-7}
\addcontentsline{toc}{subsubsection}{Answers and hints}

\begin{enumerate}
\def\labelenumi{\arabic{enumi}.}
\tightlist
\item
  \(\phi_L\) and \(\phi_R\) are unwrapped left and right wheel
  rotational positions in radians; \(\dot\phi_L\) and \(\dot\phi_R\) are
  their angular speeds in \(\mathrm{rad\,s^{-1}}\); \(r\) is effective
  wheel radius in metres; and \(L\) is track width in metres.
\item
  Radian measure is an arc-length-to-radius ratio, so its metre units
  cancel. Therefore, \(r\dot\phi\) has units
  \(\mathrm m\cdot\mathrm{rad\,s^{-1}}=\mathrm{m\,s^{-1}}\).
\item
  In frame \(\{b\}\), use
  \({}^{b}\mathbf r_{O_bC_L}=[0,L/2,0]^\mathsf T\),
  \({}^{b}\mathbf r_{O_bC_R}=[0,-L/2,0]^\mathsf T\),
  \({}^{b}\mathbf v_{O_b}=[v,0,0]^\mathsf T\), and
  \({}^{b}\boldsymbol{\omega}=[0,0,\omega]^\mathsf T\). The cross
  products contribute \([-L\omega/2,0,0]^\mathsf T\) on the left and
  \([L\omega/2,0,0]^\mathsf T\) on the right. Adding
  \({}^{b}\mathbf v_{O_b}\) and selecting the \(+x_b\) components gives
  the two required equations.
\item
  Equal wheel speeds have zero difference, so the yaw-rate equation
  gives \(\omega=0\).
\item
  Subtract \(v_L=v-L\omega/2\) from \(v_R=v+L\omega/2\), then substitute
  \(v_R=r\dot\phi_R\) and \(v_L=r\dot\phi_L\).
\item
  \(v=0.30\,\mathrm{m\,s^{-1}}\) and
  \(\omega=0.40\,\mathrm{rad\,s^{-1}}\).
\item
  The ICC \(C\) is the point about which the robot is instantaneously
  rotating. The radius \(\rho=\operatorname{distance}(O_b,C)\geq0\) is
  an ordinary distance, while \(R=+\rho\) for an ICC on the left and
  \(R=-\rho\) for one on the right. The quantity \(\Delta\ell\) is
  signed forward distance and \(\Delta\theta\) is signed heading change.
  Dividing \(\Delta\ell=R\Delta\theta\) by \(\Delta t\) and taking the
  limit gives \(v=R\omega\), hence \(R=v/\omega\) for \(\omega\neq0\). A
  forward left turn has \(R>0\) and \(\omega>0\); a forward right turn
  has \(R<0\) and \(\omega<0\).
\item
  Examples include wheel slip, unequal wheel radii, tyre deformation,
  backlash, encoder error, an incorrect track width, and uneven ground.
\end{enumerate}

\section{From velocity to a finite pose
change}\label{from-velocity-to-a-finite-pose-change}

The differential equation
\(\dot{\mathbf q}=\mathbf f(\mathbf q,\mathbf u)\) specifies
instantaneous motion. A simulator or odometry system needs the pose
after a finite interval.

Let \(t_k\) denote sample time \(k\), where \(k\in\{0,1,2,\ldots\}\),
and define

\[
\Delta t
\mathrel{:=}
t_{k+1}-t_k>0.
\]

The sampling interval \(\Delta t\) is measured in seconds. Define
\(\mathbf q_k\mathrel{:=}\mathbf q(t_k)\) and assume that \(v\) and
\(\omega\) remain constant on \([t_k,t_{k+1}]\). This is a
zero-order-hold input assumption.

\subsection{Exact straight motion}\label{exact-straight-motion}

If \(\omega=0\), the heading remains constant. Recall that

\[
\begin{bmatrix}
\dot x\\
\dot y \\
\dot \theta \\
\end{bmatrix}
=
\begin{bmatrix}
v\cos\theta\\
v\sin\theta \\
\omega
\end{bmatrix}.
\]

Integrating the three rate equations gives

\[
\boxed{
\begin{aligned}
x_{k+1}&=x_k+v\Delta t\cos\theta_k,\\
y_{k+1}&=y_k+v\Delta t\sin\theta_k,\\
\theta_{k+1}&=\theta_k.
\end{aligned}
}
\]

\subsection{Exact constant-curvature
motion}\label{exact-constant-curvature-motion}

If \(\omega\neq0\), let \(\tau\in[0,\Delta t]\) denote elapsed time
within the interval. Then

\[
\theta(t_k+\tau)=\theta_k+\omega\tau.
\]

Substituting this heading into the position-rate equations and
integrating gives

\[
\begin{aligned}
x_{k+1}-x_k
&=\int_0^{\Delta t}v\cos(\theta_k+\omega\tau)\,d\tau\\
&=\frac{v}{\omega}\left[\sin(\theta_k+\omega\Delta t)-\sin\theta_k\right],\\[4pt]
y_{k+1}-y_k
&=\int_0^{\Delta t}v\sin(\theta_k+\omega\tau)\,d\tau\\
&=\frac{v}{\omega}\left[\cos\theta_k-\cos(\theta_k+\omega\Delta t)\right].
\end{aligned}
\]

Therefore,

\[
\boxed{
\begin{aligned}
x_{k+1}&=x_k+\frac{v}{\omega}\left[\sin(\theta_k+\omega\Delta t)-\sin\theta_k\right],\\
y_{k+1}&=y_k+\frac{v}{\omega}\left[\cos\theta_k-\cos(\theta_k+\omega\Delta t)\right],\\
\theta_{k+1}&=\theta_k+\omega\Delta t.
\end{aligned}
}
\]

The signed turning radius is \(R=v/\omega\), measured in metres, and the
signed angular change is \(\Delta\theta=\omega\Delta t\), measured in
radians. The robot therefore follows a circular arc.

The formulas contain \(v/\omega\) and therefore cannot be evaluated
directly when \(\omega=0\). To express their continuous limit, define
the unnormalised sinc function
\(\operatorname{sinc}:\mathbb R\to\mathbb R\) by

\[
\operatorname{sinc}(z)
\mathrel{:=}
\begin{cases}
\dfrac{\sin z}{z},&z\neq0,\\[6pt]
1,&z=0.
\end{cases}
\]

Here, \(z\in\mathbb R\) is a dimensionless scalar, and the value at zero
follows from \(\lim_{z\to0}\sin z/z=1\). This convention must be stated
because some software libraries instead define a normalised sinc
function containing \(\pi z\).

Using the trigonometric identities

\[
\sin(\theta_k+\Delta\theta)-\sin\theta_k
=
2\cos\left(\theta_k+\frac{\Delta\theta}{2}\right)
\sin\left(\frac{\Delta\theta}{2}\right)
\]

and

\[
\cos\theta_k-\cos(\theta_k+\Delta\theta)
=
2\sin\left(\theta_k+\frac{\Delta\theta}{2}\right)
\sin\left(\frac{\Delta\theta}{2}\right),
\]

the common scale factor in both position increments can be rewritten,
for \(\omega\neq0\), as

\[
\begin{aligned}
\frac{v}{\omega}\,2\sin\left(\frac{\Delta\theta}{2}\right)
&=
v\Delta t\,
\frac{\sin(\Delta\theta/2)}{\Delta\theta/2}\\
&=
v\Delta t\,
\operatorname{sinc}\left(\frac{\Delta\theta}{2}\right),
\end{aligned}
\]

where \(\Delta\theta=\omega\Delta t\) was used in the first equality.
The exact update is consequently equivalent to

\[
\boxed{
\begin{aligned}
x_{k+1}
&=
x_k+
v\Delta t\,
\operatorname{sinc}\left(\frac{\Delta\theta}{2}\right)
\cos\left(\theta_k+\frac{\Delta\theta}{2}\right),\\
y_{k+1}
&=
y_k+
v\Delta t\,
\operatorname{sinc}\left(\frac{\Delta\theta}{2}\right)
\sin\left(\theta_k+\frac{\Delta\theta}{2}\right),\\
\theta_{k+1}
&=
\theta_k+\Delta\theta.
\end{aligned}
}
\]

As \(\omega\to0\), the angular change \(\Delta\theta\to0\),
\(\operatorname{sinc}(\Delta\theta/2)\to1\), and the midpoint heading
\(\theta_k+\Delta\theta/2\to\theta_k\). The position equations therefore
reduce smoothly to the straight-motion update

\[
x_{k+1}=x_k+v\Delta t\cos\theta_k,
\qquad
y_{k+1}=y_k+v\Delta t\sin\theta_k.
\]

Software may use an explicit straight-motion case when \(|\omega|\) is
below a declared tolerance or a sinc implementation designed to remain
numerically stable near zero. Directly computing \(\sin z/z\) at \(z=0\)
would still divide by zero despite the function's well-defined limiting
value.

\subsection*{Quick test C: exact finite
motion}\label{quick-test-c-exact-finite-motion}
\addcontentsline{toc}{subsection}{Quick test C: exact finite motion}

\begin{enumerate}
\def\labelenumi{\arabic{enumi}.}
\tightlist
\item
  What does the zero-order-hold assumption say about \(v\) and
  \(\omega\)?
\item
  Why is \(\theta(t_k+\tau)=\theta_k+\omega\tau\) when \(\omega\) is
  constant?
\item
  Where does the factor \(1/\omega\) in the position integral come from?
\item
  Starting from \(\mathbf q_k=[0\;0\;0]^\mathsf T\), let
  \(v=1\,\mathrm{m\,s^{-1}}\), \(\omega=1\,\mathrm{rad\,s^{-1}}\), and
  \(\Delta t=1\,\mathrm s\). Find the exact next pose.
\item
  Define the unnormalised sinc function and explain how it exposes the
  straight-motion limit.
\item
  Why should software not directly evaluate either \(v/\omega\) at
  \(\omega=0\) or \(\sin z/z\) at \(z=0\)?
\end{enumerate}

\subsubsection*{Answers and hints}\label{answers-and-hints-8}
\addcontentsline{toc}{subsubsection}{Answers and hints}

\begin{enumerate}
\def\labelenumi{\arabic{enumi}.}
\tightlist
\item
  Both inputs remain constant throughout the sample interval.
\item
  Integrating \(\dot\theta=\omega\) for elapsed time \(\tau\) gives
  heading change \(\omega\tau\).
\item
  With \(z=\theta_k+\omega\tau\), \(dz=\omega\,d\tau\), so
  \(d\tau=dz/\omega\).
\item
  \(\mathbf q_{k+1}=[\sin1\;\;1-\cos1\;\;1]^\mathsf T\approx[0.8415\;\;0.4597\;\;1.0000]^\mathsf T\),
  with position in metres and heading in radians.
\item
  The unnormalised function is \(\operatorname{sinc}(z)=\sin z/z\) for
  \(z\neq0\) and \(\operatorname{sinc}(0)=1\). Rewriting the update with
  \(\operatorname{sinc}(\Delta\theta/2)\) makes the limit visible: as
  \(\omega\to0\), \(\Delta\theta\to0\), the sinc factor tends to \(1\),
  and the midpoint heading tends to \(\theta_k\).
\item
  Both direct expressions would divide by zero even though their
  mathematical limits are well defined. Software must select the
  straight case or evaluate sinc with a numerically stable near-zero
  implementation.
\end{enumerate}

\section{Numerical integration with Forward
Euler}\label{numerical-integration-with-forward-euler}

For a general time-varying input, the exact solution may not have a
convenient closed form. The integral form of the state equation over one
interval is

\[
\mathbf q(t_{k+1})
=
\mathbf q(t_k)
+
\int_{t_k}^{t_{k+1}}
\mathbf f(\mathbf q(t),\mathbf u(t))\,dt.
\]

Numerical integration approximates this integral using a finite number
of evaluations of \(\mathbf f\).

\subsection{Derivation from a Taylor
expansion}\label{derivation-from-a-taylor-expansion}

Assume \(\mathbf q(t)\) is sufficiently differentiable near \(t_k\). Its
Taylor expansion is

\[
\mathbf q(t_k+\Delta t)
=
\mathbf q(t_k)
+\Delta t\,\dot{\mathbf q}(t_k)
+\frac{\Delta t^2}{2}\,\ddot{\mathbf q}(\xi_k),
\]

for some \(\xi_k\in(t_k,t_{k+1})\). Substituting
\(\dot{\mathbf q}(t_k)=\mathbf f(\mathbf q_k,\mathbf u_k)\) and
discarding the term proportional to \(\Delta t^2\) gives the Forward
Euler update

\[
\boxed{
\mathbf q_{k+1}^{\mathrm E}
=
\mathbf q_k^{\mathrm E}
+\Delta t\,\mathbf f(\mathbf q_k^{\mathrm E},\mathbf u_k)
}.
\]

The superscript \(\mathrm E\) identifies the Euler approximation, and
\(\mathbf u_k\mathrel{:=}\mathbf u(t_k)\) is the sampled input. Forward
Euler uses the slope at the beginning of the interval and treats it as
constant across the whole interval.

For the planar kinematic model,

\[
\boxed{
\begin{aligned}
x_{k+1}^{\mathrm E}&=x_k^{\mathrm E}+\Delta t\,v_k\cos\theta_k^{\mathrm E},\\
y_{k+1}^{\mathrm E}&=y_k^{\mathrm E}+\Delta t\,v_k\sin\theta_k^{\mathrm E},\\
\theta_{k+1}^{\mathrm E}&=\theta_k^{\mathrm E}+\Delta t\,\omega_k.
\end{aligned}
}
\]

Here, \(v_k\) and \(\omega_k\) are the body-forward speed and yaw rate
used for sample \(k\).

\subsection{Local and accumulated
error}\label{local-and-accumulated-error}

Let \(E(\Delta t)\geq0\) denote a chosen error magnitude produced using
step size \(\Delta t>0\). For an exponent \(p>0\), the statement

\[
E(\Delta t)=O(\Delta t^p)
\qquad\text{as}\qquad
\Delta t\to0
\]

means that there are constants \(C>0\) and \(\Delta t_0>0\), independent
of \(\Delta t\), such that

\[
E(\Delta t)\leq C\Delta t^p
\qquad
\text{for every}\qquad
0<\Delta t\leq\Delta t_0.
\]

The exponent \(p\) is the convergence order. The big-\(O\) statement is
a bound on how the error scales for sufficiently small steps, not an
equation claiming that the error equals \(\Delta t^p\). When the leading
error term is nonzero and \(\Delta t\) is sufficiently small, one
commonly observes the approximate scaling
\(E(\Delta t)\approx K\Delta t^p\) for a constant \(K>0\) that does not
depend on \(\Delta t\).

The local truncation error is the error introduced in one step when that
step begins from the exact state. Forward Euler has local truncation
error \(O(\Delta t^2)\). Over a fixed total time interval, reducing
\(\Delta t\) creates more steps, and their errors accumulate. The
resulting global error of Forward Euler is generally \(O(\Delta t)\). In
the sufficiently-small-step regime where the leading first-order term
dominates, halving \(\Delta t\) therefore approximately halves the
global error; it does not usually quarter it.

\subsection{Worked comparison with the exact
update}\label{worked-comparison-with-the-exact-update}

Use the initial pose and constant input

\[
\mathbf q_0=
\begin{bmatrix}
0\\
0\\
0
\end{bmatrix},
\qquad
v=1\,\mathrm{m\,s^{-1}},
\qquad
\omega=1\,\mathrm{rad\,s^{-1}},
\qquad
\Delta t=1\,\mathrm s.
\]

The exact update is

\[
\mathbf q_1
\approx
\begin{bmatrix}
0.8415\,\mathrm m\\
0.4597\,\mathrm m\\
1.0000\,\mathrm{rad}
\end{bmatrix}.
\]

Forward Euler evaluates the direction only at initial heading
\(\theta_0=0\):

\[
\mathbf q_1^{\mathrm E}
=
\begin{bmatrix}
0+1\cos0\\
0+1\sin0\\
0+1
\end{bmatrix}
=
\begin{bmatrix}
1\,\mathrm m\\
0\,\mathrm m\\
1\,\mathrm{rad}
\end{bmatrix}.
\]

Both methods produce the same heading because \(\omega\) is constant and
the heading equation is linear. The Euler position is wrong because the
forward direction rotates substantially during the one-second interval.
Define one-step position error by

\[
e_p
\mathrel{:=}
\sqrt{(x_1^{\mathrm E}-x_1)^2+(y_1^{\mathrm E}-y_1)^2}.
\]

For this example, \(e_p\approx0.4863\,\mathrm m\). Smaller steps allow
Euler to update the direction more often. When constant-input
differential-drive motion is assumed, however, the exact arc update is
available and more accurate.

\subsection{Choosing and checking an integration
method}\label{choosing-and-checking-an-integration-method}

\begin{itemize}
\tightlist
\item
  Use the exact straight-or-arc update when \(v\) and \(\omega\) are
  held constant and the ideal planar model is appropriate.
\item
  Use a numerical method when the input or model varies within the
  interval and no convenient exact solution is available.
\item
  Reduce \(\Delta t\) when the state changes too much during one step,
  while recognising that smaller steps increase computation.
\item
  Compare a numerical trajectory with a known exact case before trusting
  it on a harder case.
\item
  Compute \(\Delta t\) from the intended time source; do not silently
  assume perfectly periodic execution.
\end{itemize}

\subsection*{Quick test D: Forward
Euler}\label{quick-test-d-forward-euler}
\addcontentsline{toc}{subsection}{Quick test D: Forward Euler}

\begin{enumerate}
\def\labelenumi{\arabic{enumi}.}
\tightlist
\item
  State the Forward Euler update and define every symbol in it.
\item
  In plain language, what approximation does Forward Euler make during
  one step?
\item
  Suppose one Forward Euler step begins from the exact pose and \(v\)
  and \(\omega\) remain constant during the step. Why does Forward Euler
  compute the exact position when \(\omega=0\), but generally not when
  the robot turns with \(\omega\neq0\)?
\item
  What does the global-error statement \(E(\Delta t)=O(\Delta t)\) mean,
  and what approximate change in global error should occur when a
  sufficiently small \(\Delta t\) is halved over the same fixed total
  time interval?
\item
  Give one exact verification case and one physical limitation for an
  odometry integrator.
\end{enumerate}

\subsubsection*{Answers and hints}\label{answers-and-hints-9}
\addcontentsline{toc}{subsubsection}{Answers and hints}

\begin{enumerate}
\def\labelenumi{\arabic{enumi}.}
\tightlist
\item
  \(\mathbf q_{k+1}^{\mathrm E}=\mathbf q_k^{\mathrm E}+\Delta t\,\mathbf f(\mathbf q_k^{\mathrm E},\mathbf u_k)\).
  The terms are the next approximate state, current approximate state,
  step duration, state-derivative function, and sampled input.
\item
  It uses the derivative at the start as if that derivative remained
  constant across the interval.
\item
  When \(\omega=0\), the heading \(\theta\) remains constant. Constant
  \(v\) then gives the constant world-frame velocity
  \([v\cos\theta_k\;\;v\sin\theta_k]^\mathsf T\), so multiplying the
  initial velocity by \(\Delta t\) evaluates the position integral
  exactly. When \(\omega\neq0\) and \(v\neq0\), the heading and
  therefore the world-frame velocity direction change continuously
  during the step. Forward Euler holds the initial direction constant
  and produces a straight tangent displacement instead of the exact arc
  displacement. The word ``generally'' allows special cases such as pure
  rotation with \(v=0\), for which the position remains unchanged and is
  therefore exact.
\item
  The statement \(E(\Delta t)=O(\Delta t)\) means that constants \(C>0\)
  and \(\Delta t_0>0\) exist such that \(E(\Delta t)\leq C\Delta t\)
  whenever \(0<\Delta t\leq\Delta t_0\). It describes first-order
  scaling as \(\Delta t\to0\), not exact equality. In the regime where
  the leading first-order error dominates,
  \(E(\Delta t)\approx K\Delta t\), so replacing \(\Delta t\) by
  \(\Delta t/2\) approximately halves the global error over the same
  fixed total time interval.
\item
  Exact cases include constant straight motion, pure rotation, and
  constant-curvature motion. Physical limitations include slip,
  parameter error, encoder quantisation, and timing error.
\end{enumerate}

\section{Engineering implications and failure
modes}\label{engineering-implications-and-failure-modes}

\begin{enumerate}
\def\labelenumi{\arabic{enumi}.}
\item
  \textbf{Frame inconsistency:} The scalar \(v\) is a body-frame forward
  speed, but \(\dot x\) and \(\dot y\) are world-frame components.
  Assigning \(\dot x=v\) is valid only when \(\theta=0\) modulo
  \(2\pi\).
\item
  \textbf{Sign-convention mismatch:} Swapping left and right wheels or
  using a clockwise-positive yaw convention reverses the predicted turn.
  Under this chapter convention, \(\dot\phi_L<\dot\phi_R\) must produce
  positive counter-clockwise yaw.
\item
  \textbf{Unit mismatch:} Wheel speed may arrive in revolutions per
  minute while the model expects radians per second. Heading may arrive
  in degrees while trigonometric functions expect radians. Convert units
  at explicit interface boundaries.
\item
  \textbf{Incorrect or variable sampling time:} Using an assumed
  \(\Delta t\) when measurements arrive irregularly causes distance and
  angle errors. An update should validate timestamp differences and
  define behaviour for zero, negative, missing, or unreasonably large
  intervals.
\item
  \textbf{Angle representation:} The equations allow \(\theta\) to grow
  without bound, but an interface may require \(\theta\in(-\pi,\pi]\).
  Wrapping changes the numerical representative, not the physical
  orientation. Heading differences require a wrap-aware calculation.
\item
  \textbf{Model error versus integration error:} Reducing \(\Delta t\)
  reduces integration error but cannot correct an inaccurate physical
  model. If a wheel slips, even exact integration of the ideal no-slip
  equations gives the wrong physical trajectory.
\item
  \textbf{Encoder integration and drift:} Wheel encoders measure
  incremental rotation rather than global position. Integrating these
  increments estimates pose relative to an initial condition. Small
  systematic and random errors accumulate, so wheel odometry alone
  generally drifts and requires external sensing for long-term
  correction.
\end{enumerate}

\section*{Cumulative chapter test}\label{cumulative-chapter-test-1}
\addcontentsline{toc}{section}{Cumulative chapter test}

\markright{Cumulative chapter test}

\begin{enumerate}
\def\labelenumi{\arabic{enumi}.}
\tightlist
\item
  Define \(\mathbf q\), \(\mathbf u\), \(\mathbf f\), \(r\), \(L\),
  \(\dot\phi_L\), \(\dot\phi_R\), \(t_k\), and \(\Delta t\), including
  dimensions or units.
\item
  Starting with \({}^{b}\mathbf v=[v\;0]^\mathsf T\), derive
  \(\dot x=v\cos\theta\) and \(\dot y=v\sin\theta\).
\item
  Derive the no-sideways-motion constraint from the world-frame
  coordinate column of the \(y_b\)-axis.
\item
  Derive the wheel-to-body equations for \(v\) and \(\omega\) and the
  body-to-wheel equations for \(\dot\phi_L\) and \(\dot\phi_R\).
\item
  A robot has \(r=0.08\,\mathrm m\), \(L=0.40\,\mathrm m\),
  \(\dot\phi_L=5\,\mathrm{rad\,s^{-1}}\), and
  \(\dot\phi_R=7\,\mathrm{rad\,s^{-1}}\). Find \(v\) and \(\omega\) and
  state the turn direction.
\item
  Derive the exact constant-curvature update by integrating the position
  rates over one interval.
\item
  Derive Forward Euler from the first two Taylor terms and explain the
  discarded term.
\item
  Explain why exact mathematical integration can still produce
  inaccurate physical odometry.
\item
  Design three verification cases that separately reveal a wheel-order
  error, a yaw-sign error, and excessive integration error.
\end{enumerate}

\subsection*{Cumulative answers and
hints}\label{cumulative-answers-and-hints-1}
\addcontentsline{toc}{subsection}{Cumulative answers and hints}

\begin{enumerate}
\def\labelenumi{\arabic{enumi}.}
\tightlist
\item
  \(\mathbf q=[x\;y\;\theta]^\mathsf T\) is planar pose with units
  \(\mathrm m,\mathrm m,\mathrm{rad}\);
  \(\mathbf u=[v\;\omega]^\mathsf T\) contains forward speed and yaw
  rate in \(\mathrm{m\,s^{-1}}\) and \(\mathrm{rad\,s^{-1}}\);
  \(\mathbf f\) maps pose and input to pose derivative; \(r\) is wheel
  radius in metres; \(L\) is track width in metres; \(\dot\phi_L\) and
  \(\dot\phi_R\) are wheel angular speeds in radians per second; \(t_k\)
  is sample time in seconds; and \(\Delta t=t_{k+1}-t_k\) is the
  positive integration interval in seconds.
\item
  Premultiply \({}^{b}\mathbf v\) by \(\mathbf R_{wb}(\theta)\) and read
  the two components.
\item
  Use \({}^{w}\mathbf e_{y,b}=[-\sin\theta\;\cos\theta]^\mathsf T\) and
  set its dot product with \([\dot x\;\dot y]^\mathsf T\) equal to zero.
\item
  Use \(r\dot\phi_L=v-L\omega/2\) and \(r\dot\phi_R=v+L\omega/2\).
  Addition isolates \(v\) and subtraction isolates \(\omega\); solving
  the pair for wheel speeds gives the inverse equations.
\item
  \(v=(0.08/2)(5+7)=0.48\,\mathrm{m\,s^{-1}}\) and
  \(\omega=(0.08/0.40)(7-5)=0.40\,\mathrm{rad\,s^{-1}}\). The positive
  yaw rate is a counter-clockwise or left turn.
\item
  Write \(\theta(t_k+\tau)=\theta_k+\omega\tau\), integrate over
  \(\tau\in[0,\Delta t]\), and use antiderivatives
  \(\sin(\theta_k+\omega\tau)/\omega\) and
  \(-\cos(\theta_k+\omega\tau)/\omega\).
\item
  Taylor expansion gives
  \(\mathbf q(t_k+\Delta t)=\mathbf q(t_k)+\Delta t\dot{\mathbf q}(t_k)+O(\Delta t^2)\).
  Replacing \(\dot{\mathbf q}\) with \(\mathbf f\) and omitting the
  second-order remainder gives Forward Euler. The discarded term
  represents within-step change of the derivative.
\item
  Exact integration eliminates approximation for the stated model and
  input assumption; it cannot eliminate slip, parameter error, encoder
  error, timing error, or other model violations.
\item
  For wheel order, command unequal known wheel speeds and check the turn
  direction. For yaw sign, command pure positive rotation and check that
  heading increases counter-clockwise. For integration error, compare
  with the exact constant-curvature solution at decreasing values of
  \(\Delta t\) and verify convergence.
\end{enumerate}

\section*{References}\label{references-2}
\addcontentsline{toc}{section}{References}

\markright{References}

\begin{itemize}
\tightlist
\item
  Kevin M. Lynch and Frank C. Park, \emph{Modern Robotics: Mechanics,
  Planning, and Control}, Chapter 13, especially the
  wheeled-mobile-robot kinematic models and odometry discussion.
\item
  Steven M. LaValle, \emph{Planning Algorithms}, Chapter 13, especially
  continuous-time differential models, velocity constraints, and
  wheeled-vehicle models.
\end{itemize}

\chapter{Differential-Drive Kinematics and Numerical Integration:
Implementation}\label{sec-kinematics-integration-implementation}

\section*{Purpose and scope}\label{purpose-and-scope}
\addcontentsline{toc}{section}{Purpose and scope}

\markright{Purpose and scope}

This companion translates the differential-drive mathematics derived in
\hyperref[sec-differential-drive-model]{Differential-Drive Kinematics
and Numerical Integration: Model and Derivation} into reusable software
contracts, invariants, validation rules, and tests. It records how those
equations map into dependable software without repeating the
derivations.

\section{Wheel-to-body velocity
mapping}\label{wheel-to-body-velocity-mapping}

Let \(r>0\) be the common effective wheel radius in metres and let
\(L>0\) be the track width in metres. Let \(\dot\phi_L\) and
\(\dot\phi_R\) be the signed left and right wheel angular speeds in
\(\mathrm{rad\,s^{-1}}\), where positive angular speed rolls the
corresponding wheel forward.

The software implements

\[
v
=
\frac{r}{2}
\left(\dot\phi_R+\dot\phi_L\right),
\qquad
\omega
=
\frac{r}{L}
\left(\dot\phi_R-\dot\phi_L\right),
\]

where \(v\) is signed body-forward speed in \(\mathrm{m\,s^{-1}}\) and
\(\omega\) is signed counter-clockwise yaw rate in
\(\mathrm{rad\,s^{-1}}\).

\subsection{Software contract}\label{software-contract}

The \texttt{wheel\_to\_body} operation accepts four finite scalar
inputs: left wheel angular speed, right wheel angular speed, wheel
radius, and track width. The wheel radius and track width must be
strictly positive. Invalid inputs raise \texttt{ValueError}; the
operation does not clamp values or substitute default geometry.

The operation returns an immutable \texttt{BodyVelocity} value
containing \texttt{forward\_speed\_m\_s} and \texttt{yaw\_rate\_rad\_s}.
It is deterministic, does not mutate its inputs, and does not depend on
ROS state, timestamps, configuration files, or hidden global state.

\subsection{Simplified reference
snippet}\label{simplified-reference-snippet}

The following equation-focused snippet shows how the mathematical
factors map to Python expressions. It assumes that all inputs are finite
and that \texttt{wheel\_radius\_m} and \texttt{track\_width\_m} are
strictly positive; the production operation enforces those preconditions
before performing the calculation.

\begin{Shaded}
\begin{Highlighting}[]
\AttributeTok{@dataclass}\NormalTok{(frozen}\OperatorTok{=}\VariableTok{True}\NormalTok{, slots}\OperatorTok{=}\VariableTok{True}\NormalTok{)}
\KeywordTok{class}\NormalTok{ BodyVelocity:}
    \CommentTok{"""Reduced planar body velocity in SI units."""}

\NormalTok{    forward\_speed\_m\_s: }\BuiltInTok{float}
\NormalTok{    yaw\_rate\_rad\_s: }\BuiltInTok{float}


\KeywordTok{def}\NormalTok{ wheel\_to\_body(}
    \OperatorTok{*}\NormalTok{,}
\NormalTok{    left\_angular\_speed\_rad\_s: }\BuiltInTok{float}\NormalTok{,}
\NormalTok{    right\_angular\_speed\_rad\_s: }\BuiltInTok{float}\NormalTok{,}
\NormalTok{    wheel\_radius\_m: }\BuiltInTok{float}\NormalTok{,}
\NormalTok{    track\_width\_m: }\BuiltInTok{float}\NormalTok{,}
\NormalTok{) }\OperatorTok{{-}\textgreater{}}\NormalTok{ BodyVelocity:}
    \CommentTok{"""Apply the wheel{-}to{-}body equations to validated inputs."""}
\NormalTok{    forward\_speed\_m\_s }\OperatorTok{=}\NormalTok{ (}
\NormalTok{        wheel\_radius\_m }\OperatorTok{/} \DecValTok{2}
        \OperatorTok{*}\NormalTok{ (right\_angular\_speed\_rad\_s }\OperatorTok{+}\NormalTok{ left\_angular\_speed\_rad\_s)}
\NormalTok{    )}
\NormalTok{    yaw\_rate\_rad\_s }\OperatorTok{=}\NormalTok{ (}
\NormalTok{        wheel\_radius\_m }\OperatorTok{/}\NormalTok{ track\_width\_m}
        \OperatorTok{*}\NormalTok{ (right\_angular\_speed\_rad\_s }\OperatorTok{{-}}\NormalTok{ left\_angular\_speed\_rad\_s)}
\NormalTok{    )}

    \ControlFlowTok{return}\NormalTok{ BodyVelocity(}
\NormalTok{        forward\_speed\_m\_s}\OperatorTok{=}\NormalTok{forward\_speed\_m\_s,}
\NormalTok{        yaw\_rate\_rad\_s}\OperatorTok{=}\NormalTok{yaw\_rate\_rad\_s,}
\NormalTok{    )}
\end{Highlighting}
\end{Shaded}

\subsection{Invariants}\label{invariants}

\begin{itemize}
\tightlist
\item
  Equal wheel speeds produce zero yaw rate.
\item
  Equal and opposite wheel speeds produce zero forward speed.
\item
  Swapping the left and right wheel speeds preserves forward speed and
  reverses yaw rate.
\item
  Two zero wheel speeds produce zero forward speed and zero yaw rate.
\item
  Positive right-wheel speed relative to left-wheel speed produces
  positive counter-clockwise yaw.
\end{itemize}

\section{Body-to-wheel velocity
mapping}\label{body-to-wheel-velocity-mapping}

The \texttt{body\_to\_wheel} operation implements the inverse of the
wheel-to-body mapping. Its inputs are signed body-forward speed \(v\) in
\(\mathrm{m\,s^{-1}}\), signed counter-clockwise yaw rate \(\omega\) in
\(\mathrm{rad\,s^{-1}}\), wheel radius \(r>0\) in metres, and track
width \(L>0\) in metres.

The required wheel angular speeds are

\[
\dot\phi_L
=
\frac{v}{r}
-
\frac{L\omega}{2r},
\qquad
\dot\phi_R
=
\frac{v}{r}
+
\frac{L\omega}{2r},
\]

where \(\dot\phi_L\) and \(\dot\phi_R\) are the signed left and right
wheel angular speeds in \(\mathrm{rad\,s^{-1}}\).

\subsection{Software contract}\label{software-contract-1}

The operation accepts four finite scalar inputs. The wheel radius and
track width must be strictly positive. Invalid inputs raise
\texttt{ValueError}; the operation does not clamp values or substitute
geometry.

The result is an immutable \texttt{WheelAngularSpeeds} value containing
\texttt{left\_angular\_speed\_rad\_s} and
\texttt{right\_angular\_speed\_rad\_s}.

\subsection{Implementation mapping}\label{implementation-mapping}

The calculation separates forward and turning contributions:

\begin{Shaded}
\begin{Highlighting}[]
\AttributeTok{@dataclass}\NormalTok{(frozen}\OperatorTok{=}\VariableTok{True}\NormalTok{, slots}\OperatorTok{=}\VariableTok{True}\NormalTok{)}
\KeywordTok{class}\NormalTok{ WheelAngularSpeeds:}
    \CommentTok{"""Signed left and right wheel angular speeds in radians per second."""}

\NormalTok{    left\_angular\_speed\_rad\_s: }\BuiltInTok{float}
\NormalTok{    right\_angular\_speed\_rad\_s: }\BuiltInTok{float}


\KeywordTok{def}\NormalTok{ body\_to\_wheel(}
    \OperatorTok{*}\NormalTok{,}
\NormalTok{    forward\_speed\_m\_s: }\BuiltInTok{float}\NormalTok{,}
\NormalTok{    yaw\_rate\_rad\_s: }\BuiltInTok{float}\NormalTok{,}
\NormalTok{    wheel\_radius\_m: }\BuiltInTok{float}\NormalTok{,}
\NormalTok{    track\_width\_m: }\BuiltInTok{float}\NormalTok{,}
\NormalTok{) }\OperatorTok{{-}\textgreater{}}\NormalTok{ WheelAngularSpeeds:}
    \CommentTok{"""Apply the body{-}to{-}wheel equations to validated inputs."""}
\NormalTok{    forward\_component\_rad\_s }\OperatorTok{=}\NormalTok{ forward\_speed\_m\_s }\OperatorTok{/}\NormalTok{ wheel\_radius\_m}
\NormalTok{    turning\_component\_rad\_s }\OperatorTok{=}\NormalTok{ (}
\NormalTok{        track\_width\_m }\OperatorTok{*}\NormalTok{ yaw\_rate\_rad\_s }\OperatorTok{/}\NormalTok{ (}\DecValTok{2} \OperatorTok{*}\NormalTok{ wheel\_radius\_m)}
\NormalTok{    )}
\NormalTok{    left\_angular\_speed\_rad\_s }\OperatorTok{=}\NormalTok{ (}
\NormalTok{        forward\_component\_rad\_s }\OperatorTok{{-}}\NormalTok{ turning\_component\_rad\_s}
\NormalTok{    )}
\NormalTok{    right\_angular\_speed\_rad\_s }\OperatorTok{=}\NormalTok{ (}
\NormalTok{        forward\_component\_rad\_s }\OperatorTok{+}\NormalTok{ turning\_component\_rad\_s}
\NormalTok{    )}

    \ControlFlowTok{return}\NormalTok{ WheelAngularSpeeds(}
\NormalTok{        left\_angular\_speed\_rad\_s}\OperatorTok{=}\NormalTok{left\_angular\_speed\_rad\_s,}
\NormalTok{        right\_angular\_speed\_rad\_s}\OperatorTok{=}\NormalTok{right\_angular\_speed\_rad\_s,}
\NormalTok{    )}
\end{Highlighting}
\end{Shaded}

The forward component contributes equally to both wheels. The turning
component is subtracted from the left wheel and added to the right
wheel. Therefore, a positive yaw rate makes the right wheel turn faster
than the left wheel.

\subsection{Invariants and required
tests}\label{invariants-and-required-tests}

\begin{itemize}
\tightlist
\item
  Zero yaw rate produces equal wheel speeds of \(v/r\).
\item
  Zero forward speed produces equal-and-opposite wheel speeds.
\item
  Reversing the yaw-rate sign swaps the left and right wheel speeds.
\item
  Applying \texttt{body\_to\_wheel} and then \texttt{wheel\_to\_body}
  recovers the original \(v\) and \(\omega\) within the declared
  floating-point tolerance.
\item
  Non-finite inputs and non-positive geometry raise \texttt{ValueError}.
\end{itemize}

\section{Exact constant-input pose
integration}\label{exact-constant-input-pose-integration}

Let the initial planar pose be \((x_k,y_k,\theta_k)\), where \(x_k\) and
\(y_k\) are the world-frame coordinates of the axle midpoint in metres
and \(\theta_k\) is the unwrapped counter-clockwise heading in radians.

Assume that forward speed \(v\) in \(\mathrm{m\,s^{-1}}\) and yaw rate
\(\omega\) in \(\mathrm{rad\,s^{-1}}\) remain constant over a time step
\(\Delta t>0\) in seconds. Define the heading change

\[
\Delta\theta
=
\omega\Delta t.
\]

The exact update is

\[
\begin{aligned}
x_{k+1}
&=
x_k+
v\Delta t\,
\operatorname{sinc}\left(\frac{\Delta\theta}{2}\right)
\cos\left(\theta_k+\frac{\Delta\theta}{2}\right),\\
y_{k+1}
&=
y_k+
v\Delta t\,
\operatorname{sinc}\left(\frac{\Delta\theta}{2}\right)
\sin\left(\theta_k+\frac{\Delta\theta}{2}\right),\\
\theta_{k+1}
&=
\theta_k+\Delta\theta,
\end{aligned}
\]

where the unnormalised sinc function is

\[
\operatorname{sinc}(z)
=
\begin{cases}
\dfrac{\sin z}{z}, & z\neq0,\\[4pt]
1, & z=0.
\end{cases}
\]

The update is exact for the ideal continuous kinematic model only when
\(v\) and \(\omega\) remain constant during the time step. It does not
account for wheel slip, changing commands, actuator dynamics, or model
error.

\subsection{Software contract}\label{software-contract-2}

The \texttt{integrate\_exact} operation accepts an immutable
\texttt{PlanarPose}, finite forward speed, finite yaw rate, and a finite
positive time step. A non-finite input or non-positive time step raises
\texttt{ValueError}.

The operation returns a new \texttt{PlanarPose}. It does not mutate the
initial pose or wrap the resulting heading into a fixed angular
interval.

\subsection{Numerically stable implementation
mapping}\label{numerically-stable-implementation-mapping}

The implementation avoids evaluating \(v/\omega\), which would be
undefined when \(\omega=0\). Instead, it uses the sinc form:

\begin{Shaded}
\begin{Highlighting}[]
\AttributeTok{@dataclass}\NormalTok{(frozen}\OperatorTok{=}\VariableTok{True}\NormalTok{, slots}\OperatorTok{=}\VariableTok{True}\NormalTok{)}
\KeywordTok{class}\NormalTok{ PlanarPose:}
    \CommentTok{"""Pose of the robot\textquotesingle{}s axle midpoint in the world frame.}
\CommentTok{    }
\CommentTok{    Attributes:}
\CommentTok{        position\_x\_m: World{-}frame x{-}coordinate in metres.}
\CommentTok{        position\_y\_m: World{-}frame y{-}coordinate in metres.}
\CommentTok{        heading\_rad: Counter{-}clockwise heading from the world{-}frame}
\CommentTok{            positive x{-}axis in radians.}
\CommentTok{    """}

\NormalTok{    position\_x\_m: }\BuiltInTok{float}
\NormalTok{    position\_y\_m: }\BuiltInTok{float}
\NormalTok{    heading\_rad: }\BuiltInTok{float}
    

\KeywordTok{def}\NormalTok{ \_sinc(value: }\BuiltInTok{float}\NormalTok{) }\OperatorTok{{-}\textgreater{}} \BuiltInTok{float}\NormalTok{:}
    \CommentTok{"""Return the unnormalised sinc function, including its value at zero."""}
    \ControlFlowTok{if}\NormalTok{ value }\OperatorTok{==} \FloatTok{0.0}\NormalTok{:}
        \ControlFlowTok{return} \FloatTok{1.0}

    \ControlFlowTok{return}\NormalTok{ sin(value) }\OperatorTok{/}\NormalTok{ value}


\NormalTok{heading\_change\_rad }\OperatorTok{=}\NormalTok{ yaw\_rate\_rad\_s }\OperatorTok{*}\NormalTok{ time\_step\_s}
\NormalTok{half\_heading\_change\_rad }\OperatorTok{=}\NormalTok{ heading\_change\_rad }\OperatorTok{/} \DecValTok{2}
\NormalTok{midpoint\_heading\_rad }\OperatorTok{=}\NormalTok{ pose.heading\_rad }\OperatorTok{+}\NormalTok{ half\_heading\_change\_rad}

\NormalTok{signed\_displacement\_m }\OperatorTok{=}\NormalTok{ (}
\NormalTok{    forward\_speed\_m\_s}
    \OperatorTok{*}\NormalTok{ time\_step\_s}
    \OperatorTok{*}\NormalTok{ \_sinc(half\_heading\_change\_rad)}
\NormalTok{)}

\NormalTok{next\_pose }\OperatorTok{=}\NormalTok{ PlanarPose(}
\NormalTok{    position\_x\_m}\OperatorTok{=}\NormalTok{(}
\NormalTok{        pose.position\_x\_m}
        \OperatorTok{+}\NormalTok{ signed\_displacement\_m }\OperatorTok{*}\NormalTok{ cos(midpoint\_heading\_rad)}
\NormalTok{    ),}
\NormalTok{    position\_y\_m}\OperatorTok{=}\NormalTok{(}
\NormalTok{        pose.position\_y\_m}
        \OperatorTok{+}\NormalTok{ signed\_displacement\_m }\OperatorTok{*}\NormalTok{ sin(midpoint\_heading\_rad)}
\NormalTok{    ),}
\NormalTok{    heading\_rad}\OperatorTok{=}\NormalTok{pose.heading\_rad }\OperatorTok{+}\NormalTok{ heading\_change\_rad,}
\NormalTok{)}
\end{Highlighting}
\end{Shaded}

When \(\omega=0\), we have \(\Delta\theta=0\) and
\(\operatorname{sinc}(0)=1\). The same implementation therefore reduces
smoothly to straight motion without a separate division by zero.

\subsection{Invariants and required
tests}\label{invariants-and-required-tests-1}

\begin{itemize}
\tightlist
\item
  Zero yaw rate produces the exact straight-motion update.
\item
  Zero forward speed changes only the heading.
\item
  Constant nonzero forward speed and yaw rate produce constant-curvature
  motion.
\item
  The input pose remains unchanged.
\item
  The output heading remains unwrapped.
\item
  Non-finite inputs raise \texttt{ValueError}.
\item
  A zero or negative time step raises \texttt{ValueError}.
\end{itemize}

\section{Forward Euler pose
integration}\label{forward-euler-pose-integration}

Forward Euler approximates motion over a time step \(\Delta t>0\) by
evaluating the pose derivative at the beginning of the step. For initial
pose \((x_k,y_k,\theta_k)\), constant forward speed \(v\), and constant
yaw rate \(\omega\), the update is

\[
\begin{aligned}
x_{k+1}^{\mathrm E}
&=
x_k+\Delta t\,v\cos\theta_k,\\
y_{k+1}^{\mathrm E}
&=
y_k+\Delta t\,v\sin\theta_k,\\
\theta_{k+1}^{\mathrm E}
&=
\theta_k+\Delta t\,\omega.
\end{aligned}
\]

The superscript \(\mathrm E\) identifies a Forward Euler approximation.
Position is calculated entirely from the beginning-of-step heading
\(\theta_k\); the changing heading within the step does not affect that
step's position increment.

\subsection{Software contract}\label{software-contract-3}

The \texttt{integrate\_forward\_euler} operation accepts the same inputs
and validation rules as \texttt{integrate\_exact}: an immutable
\texttt{PlanarPose}, finite forward speed, finite yaw rate, and a finite
positive time step.

The operation returns a new \texttt{PlanarPose}, does not mutate the
input pose, and leaves the heading unwrapped. Invalid inputs raise
\texttt{ValueError}.

\subsection{Implementation mapping}\label{implementation-mapping-1}

\begin{Shaded}
\begin{Highlighting}[]
\NormalTok{next\_pose }\OperatorTok{=}\NormalTok{ PlanarPose(}
\NormalTok{    position\_x\_m}\OperatorTok{=}\NormalTok{(}
\NormalTok{        pose.position\_x\_m}
        \OperatorTok{+}\NormalTok{ time\_step\_s }\OperatorTok{*}\NormalTok{ forward\_speed\_m\_s }\OperatorTok{*}\NormalTok{ cos(pose.heading\_rad)}
\NormalTok{    ),}
\NormalTok{    position\_y\_m}\OperatorTok{=}\NormalTok{(}
\NormalTok{        pose.position\_y\_m}
        \OperatorTok{+}\NormalTok{ time\_step\_s }\OperatorTok{*}\NormalTok{ forward\_speed\_m\_s }\OperatorTok{*}\NormalTok{ sin(pose.heading\_rad)}
\NormalTok{    ),}
\NormalTok{    heading\_rad}\OperatorTok{=}\NormalTok{(}
\NormalTok{        pose.heading\_rad }\OperatorTok{+}\NormalTok{ time\_step\_s }\OperatorTok{*}\NormalTok{ yaw\_rate\_rad\_s}
\NormalTok{    ),}
\NormalTok{)}
\end{Highlighting}
\end{Shaded}

\subsection{Exact and approximate
cases}\label{exact-and-approximate-cases}

When \(\omega=0\), the heading remains constant. The world-frame
velocity is therefore constant, so one Forward Euler step produces the
exact straight-motion position.

When \(v=0\), position remains unchanged and the heading update is exact
because \(\dot\theta=\omega\) is constant.

When both \(v\neq0\) and \(\omega\neq0\), the true path curves during
the time step. Forward Euler instead follows the tangent defined by
\(\theta_k\), so its one-step position is generally approximate.

For repeated integration over a fixed duration, Forward Euler has global
error \(O(\Delta t)\). In this regime, halving \(\Delta t\) should
approximately halve the accumulated error.

\subsection{Invariants and required
tests}\label{invariants-and-required-tests-2}

\begin{itemize}
\tightlist
\item
  Straight constant-speed motion agrees with the exact update.
\item
  Pure rotation changes only the heading.
\item
  A turning step uses the beginning-of-step heading.
\item
  Smaller time steps reduce accumulated turning-position error.
\item
  Halving the time step approximately halves global error in the
  first-order regime.
\item
  Non-finite inputs and non-positive time steps raise
  \texttt{ValueError}.
\end{itemize}

\chapter{ROS 2 Package Development
Workflow}\label{sec-ros2-package-development-workflow}

\section*{Goal}\label{goal}
\addcontentsline{toc}{section}{Goal}

\markright{Goal}

This workflow creates, builds, runs, and tests a small ROS 2 Python
package from scratch. It uses \texttt{my\_robot\_package} as an example
package name and \texttt{my\_node} as an example executable name;
replace them in a real project.

A \textbf{workspace} holds packages developed and built together. A
\textbf{ROS 2 package} combines source code with the metadata needed for
building, installing, and discovery. An \textbf{executable} is a program
installed by a package and started with \texttt{ros2\ run}.

The generated \texttt{my\_node} executable only prints a message. It
demonstrates the package workflow but does not yet communicate with
other ROS 2 nodes.

\section{Source the ROS 2 underlay}\label{source-the-ros-2-underlay}

For Zsh, run:

\begin{Shaded}
\begin{Highlighting}[]
\BuiltInTok{source}\NormalTok{ /opt/ros/jazzy/setup.zsh}
\end{Highlighting}
\end{Shaded}

For Bash, run:

\begin{Shaded}
\begin{Highlighting}[]
\BuiltInTok{source}\NormalTok{ /opt/ros/jazzy/setup.bash}
\end{Highlighting}
\end{Shaded}

Use the script matching the current shell.

\textbf{What happens:} \texttt{/opt/ros/jazzy} is the ROS 2 Jazzy
\textbf{underlay}, meaning the existing installation on which the new
workspace depends. Its setup script adds ROS installation prefixes,
Python packages, and executables to the current shell's search
environment.

The command uses \texttt{source} because the environment of the current
shell must change. If the script ran in a separate \textbf{child
process}, meaning a process started by the shell, its environment
changes would disappear when that child exited.

Verify the environment:

\begin{Shaded}
\begin{Highlighting}[]
\BuiltInTok{printf} \StringTok{\textquotesingle{}\%s\textbackslash{}n\textquotesingle{}} \StringTok{"}\VariableTok{$ROS\_DISTRO}\StringTok{"}
\BuiltInTok{command} \AttributeTok{{-}v}\NormalTok{ ros2}
\BuiltInTok{command} \AttributeTok{{-}v}\NormalTok{ colcon}
\end{Highlighting}
\end{Shaded}

The expected observations are \texttt{jazzy} and executable paths for
\texttt{ros2} and \texttt{colcon}.

\section{Create a workspace}\label{create-a-workspace}

For a new standalone workspace, run:

\begin{Shaded}
\begin{Highlighting}[]
\FunctionTok{mkdir} \AttributeTok{{-}p}\NormalTok{ \textasciitilde{}/ros\_ws/src}
\BuiltInTok{cd}\NormalTok{ \textasciitilde{}/ros\_ws}
\end{Highlighting}
\end{Shaded}

The workspace develops into:

\begin{Shaded}
\begin{Highlighting}[]
\NormalTok{ros\_ws/}
\NormalTok{├── src/       \# Editable source}
\NormalTok{├── build/     \# Generated build workspace}
\NormalTok{├── install/   \# Generated runnable workspace image}
\NormalTok{└── log/       \# Generated build and test records}
\end{Highlighting}
\end{Shaded}

Only \texttt{src/} exists initially. \texttt{colcon} creates the other
directories. Edit \texttt{src/}, not the generated files in
\texttt{build/}, \texttt{install/}, or \texttt{log/}.

\section{Create a Python package}\label{create-a-python-package}

From the workspace root, run:

\begin{Shaded}
\begin{Highlighting}[]
\BuiltInTok{cd}\NormalTok{ src}
\ExtensionTok{ros2}\NormalTok{ pkg create }\DataTypeTok{\textbackslash{}}
  \AttributeTok{{-}{-}build{-}type}\NormalTok{ ament\_python }\DataTypeTok{\textbackslash{}}
  \AttributeTok{{-}{-}license}\NormalTok{ Apache{-}2.0 }\DataTypeTok{\textbackslash{}}
  \AttributeTok{{-}{-}node{-}name}\NormalTok{ my\_node }\DataTypeTok{\textbackslash{}}
\NormalTok{  my\_robot\_package}
\BuiltInTok{cd}\NormalTok{ ..}
\end{Highlighting}
\end{Shaded}

Run this command only when the package does not already exist.

\textbf{What happens:} \texttt{ros2\ pkg\ create} generates
\texttt{src/my\_robot\_package}.'

\begin{itemize}
\tightlist
\item
  \texttt{ament\_python} selects Python packaging through
  \texttt{setuptools}.
\item
  \texttt{Apache-2.0} is recorded in \texttt{package.xml}, and a
  matching \texttt{LICENSE} file is created.
\item
  \texttt{-\/-node-name\ my\_node} creates a minimal Python executable
  and its \texttt{setup.py} entry point.
\end{itemize}

Omit \texttt{-\/-node-name} when creating a library-only package.

For comparison, create a C++ package with:

\begin{Shaded}
\begin{Highlighting}[]
\ExtensionTok{ros2}\NormalTok{ pkg create }\DataTypeTok{\textbackslash{}}
  \AttributeTok{{-}{-}build{-}type}\NormalTok{ ament\_cmake }\DataTypeTok{\textbackslash{}}
  \AttributeTok{{-}{-}license}\NormalTok{ Apache{-}2.0 }\DataTypeTok{\textbackslash{}}
\NormalTok{  robot\_controller}
\end{Highlighting}
\end{Shaded}

The common build types are \texttt{ament\_python} for Python with
\texttt{setuptools} and \texttt{ament\_cmake} for C++ with CMake.

\section{Inspect the generated
package}\label{inspect-the-generated-package}

\begin{Shaded}
\begin{Highlighting}[]
\NormalTok{src/my\_robot\_package/}
\NormalTok{├── LICENSE}
\NormalTok{├── package.xml}
\NormalTok{├── setup.py}
\NormalTok{├── setup.cfg}
\NormalTok{├── resource/}
\NormalTok{│   └── my\_robot\_package}
\NormalTok{├── my\_robot\_package/}
\NormalTok{│   ├── \_\_init\_\_.py}
\NormalTok{│   └── my\_node.py}
\NormalTok{└── test/}
\NormalTok{    ├── test\_copyright.py}
\NormalTok{    ├── test\_flake8.py}
\NormalTok{    └── test\_pep257.py}
\end{Highlighting}
\end{Shaded}

The outer \texttt{my\_robot\_package/} is the ROS package root. The
inner directory is the Python import package used by
\texttt{import\ my\_robot\_package}; the repeated name does not indicate
nested repositories.

{\def\LTcaptype{none} % do not increment counter
\begin{longtable}[]{@{}
  >{\raggedright\arraybackslash}p{(\linewidth - 2\tabcolsep) * \real{0.5000}}
  >{\raggedright\arraybackslash}p{(\linewidth - 2\tabcolsep) * \real{0.5000}}@{}}
\toprule\noalign{}
\begin{minipage}[b]{\linewidth}\raggedright
File or directory
\end{minipage} & \begin{minipage}[b]{\linewidth}\raggedright
Purpose
\end{minipage} \\
\midrule\noalign{}
\endhead
\bottomrule\noalign{}
\endlastfoot
\texttt{package.xml} & Declares package identity, licence, dependencies,
and build type \\
\texttt{setup.py} & Selects Python packages, data files, and executables
to install \\
\texttt{setup.cfg} & Selects the ROS-compatible destination for Python
executables \\
\texttt{resource/my\_robot\_package} & Registers the installed package
in the ament resource index \\
\texttt{my\_robot\_package/} & Contains importable Python source \\
\texttt{test/} & Contains checks run by \texttt{pytest} and
\texttt{colcon\ test} \\
\end{longtable}
}

The resource marker is intentionally empty. Its filename identifies the
package.

\section{Complete the package
metadata}\label{complete-the-package-metadata}

Replace placeholder descriptions and maintainer details in
\texttt{package.xml} and \texttt{setup.py}. Keep their package name,
version, and licence consistent. The \texttt{package.xml} export tells
\texttt{colcon} which build integration to use:

\begin{Shaded}
\begin{Highlighting}[]
\NormalTok{\textless{}}\KeywordTok{export}\NormalTok{\textgreater{}}
\NormalTok{  \textless{}}\KeywordTok{build\_type}\NormalTok{\textgreater{}ament\_python\textless{}/}\KeywordTok{build\_type}\NormalTok{\textgreater{}}
\NormalTok{\textless{}/}\KeywordTok{export}\NormalTok{\textgreater{}}
\end{Highlighting}
\end{Shaded}

A concise \texttt{setup.py} for the example is:

\begin{Shaded}
\begin{Highlighting}[]
\ImportTok{from}\NormalTok{ setuptools }\ImportTok{import}\NormalTok{ find\_packages, setup}

\NormalTok{package\_name }\OperatorTok{=} \StringTok{"my\_robot\_package"}

\NormalTok{setup(}
\NormalTok{    name}\OperatorTok{=}\NormalTok{package\_name,}
\NormalTok{    version}\OperatorTok{=}\StringTok{"0.0.0"}\NormalTok{,}
\NormalTok{    packages}\OperatorTok{=}\NormalTok{find\_packages(exclude}\OperatorTok{=}\NormalTok{[}\StringTok{"test"}\NormalTok{]),}
\NormalTok{    data\_files}\OperatorTok{=}\NormalTok{[}
\NormalTok{        (}
            \StringTok{"share/ament\_index/resource\_index/packages"}\NormalTok{,}
\NormalTok{            [}\SpecialStringTok{f"resource/}\SpecialCharTok{\{}\NormalTok{package\_name}\SpecialCharTok{\}}\SpecialStringTok{"}\NormalTok{],}
\NormalTok{        ),}
\NormalTok{        (}\SpecialStringTok{f"share/}\SpecialCharTok{\{}\NormalTok{package\_name}\SpecialCharTok{\}}\SpecialStringTok{"}\NormalTok{, [}\StringTok{"package.xml"}\NormalTok{]),}
\NormalTok{    ],}
\NormalTok{    install\_requires}\OperatorTok{=}\NormalTok{[}\StringTok{"setuptools"}\NormalTok{],}
\NormalTok{    zip\_safe}\OperatorTok{=}\VariableTok{True}\NormalTok{,}
\NormalTok{    maintainer}\OperatorTok{=}\StringTok{"Your Name"}\NormalTok{,}
\NormalTok{    maintainer\_email}\OperatorTok{=}\StringTok{"you@example.com"}\NormalTok{,}
\NormalTok{    description}\OperatorTok{=}\StringTok{"Example ROS 2 Python package."}\NormalTok{,}
\NormalTok{    license}\OperatorTok{=}\StringTok{"Apache{-}2.0"}\NormalTok{,}
\NormalTok{    entry\_points}\OperatorTok{=}\NormalTok{\{}
        \StringTok{"console\_scripts"}\NormalTok{: [}
            \StringTok{"my\_node = my\_robot\_package.my\_node:main"}\NormalTok{,}
\NormalTok{        ],}
\NormalTok{    \},}
\NormalTok{)}
\end{Highlighting}
\end{Shaded}

The important connections are:

\begin{itemize}
\tightlist
\item
  \texttt{packages} installs importable Python packages.
\item
  \texttt{data\_files} installs the resource marker and
  \texttt{package.xml} where ROS expects them.
\item
  \texttt{console\_scripts} creates an executable that calls a Python
  function.
\end{itemize}

The entry point \texttt{my\_node\ =\ my\_robot\_package.my\_node:main}
maps the executable name \texttt{my\_node} to the function
\texttt{main()} in the Python module
\texttt{my\_robot\_package.my\_node}.

The matching \texttt{setup.cfg} is:

\begin{Shaded}
\begin{Highlighting}[]
\KeywordTok{[develop]}
\DataTypeTok{script\_dir}\OtherTok{=}\StringTok{$base/lib/my\_robot\_package}

\KeywordTok{[install]}
\DataTypeTok{install\_scripts}\OtherTok{=}\StringTok{$base/lib/my\_robot\_package}
\end{Highlighting}
\end{Shaded}

Here, \texttt{\$base} is the package installation prefix.
\texttt{setup.cfg} specifies that executable scripts belong under
\texttt{lib/my\_robot\_package}, where \texttt{ros2\ run} looks for
them. It does not create an executable by itself; the
\texttt{console\_scripts} entry does that.

\section{Resolve dependencies and run fast
tests}\label{resolve-dependencies-and-run-fast-tests}

Resolve dependencies declared by packages under \texttt{src/}:

\begin{Shaded}
\begin{Highlighting}[]
\ExtensionTok{rosdep}\NormalTok{ install }\DataTypeTok{\textbackslash{}}
  \AttributeTok{{-}{-}from{-}paths}\NormalTok{ src }\DataTypeTok{\textbackslash{}}
  \AttributeTok{{-}{-}ignore{-}src} \DataTypeTok{\textbackslash{}}
  \AttributeTok{{-}{-}rosdistro}\NormalTok{ jazzy }\DataTypeTok{\textbackslash{}}
  \AttributeTok{{-}r} \DataTypeTok{\textbackslash{}}
  \AttributeTok{{-}y}
\end{Highlighting}
\end{Shaded}

\texttt{rosdep} ignores dependencies supplied by source packages in the
workspace, maps the remaining dependency keys to operating-system
packages, and installs missing packages.

Run fast tests directly against the source tree:

\begin{Shaded}
\begin{Highlighting}[]
\ExtensionTok{python3} \AttributeTok{{-}m}\NormalTok{ pytest }\DataTypeTok{\textbackslash{}}
\NormalTok{  src/my\_robot\_package/test }\AttributeTok{{-}q}
\end{Highlighting}
\end{Shaded}

Source-level tests provide fast feedback but do not verify ROS metadata
or the installed layout. Continue only when they pass; the count and
duration depend on the package.

\section{Build and install the
package}\label{build-and-install-the-package}

From the workspace root, run:

\begin{Shaded}
\begin{Highlighting}[]
\ExtensionTok{colcon}\NormalTok{ build }\DataTypeTok{\textbackslash{}}
  \AttributeTok{{-}{-}packages{-}select}\NormalTok{ my\_robot\_package }\DataTypeTok{\textbackslash{}}
  \AttributeTok{{-}{-}symlink{-}install}
\end{Highlighting}
\end{Shaded}

\textbf{What happens:} \texttt{colcon} discovers packages under
\texttt{src/}, reads their manifests, orders them by declared
dependencies, and selects the appropriate build integration.
\texttt{ament\_python} then invokes \texttt{setuptools} using
\texttt{setup.py} and \texttt{setup.cfg}.

The logical installed layout is:

\begin{Shaded}
\begin{Highlighting}[]
\NormalTok{install/my\_robot\_package/}
\NormalTok{├── lib/}
\NormalTok{│   ├── python3.12/}
\NormalTok{│   │   └── site{-}packages/}
\NormalTok{│   │       └── my\_robot\_package/}
\NormalTok{│   │           ├── \_\_init\_\_.py}
\NormalTok{│   │           └── my\_node.py}
\NormalTok{│   └── my\_robot\_package/}
\NormalTok{│       └── my\_node}
\NormalTok{└── share/}
\NormalTok{    ├── ament\_index/}
\NormalTok{    │   └── resource\_index/}
\NormalTok{    │       └── packages/}
\NormalTok{    │           └── my\_robot\_package}
\NormalTok{    └── my\_robot\_package/}
\NormalTok{        └── package.xml}
\end{Highlighting}
\end{Shaded}

The Python version in the path depends on the system.

\begin{itemize}
\tightlist
\item
  \texttt{site-packages/my\_robot\_package} supplies importable Python
  code.
\item
  \texttt{lib/my\_robot\_package/my\_node} supplies the executable used
  by \texttt{ros2\ run}.
\item
  The resource marker registers the package with ROS.
\item
  The installed \texttt{package.xml} supplies package metadata.
\end{itemize}

ROS 2 needs this standard runtime layout to locate packages,
executables, metadata, launch files, configuration, Python imports, and
dependency environment hooks.

\subsection{\texorpdfstring{What \texttt{-\/-symlink-install}
changes}{What -\/-symlink-install changes}}\label{what---symlink-install-changes}

Supported files are linked back to editable source rather than copied:

\begin{Shaded}
\begin{Highlighting}[]
\NormalTok{install/.../my\_node.py}
\NormalTok{        ↓ symbolic link}
\NormalTok{src/my\_robot\_package/my\_robot\_package/my\_node.py}
\end{Highlighting}
\end{Shaded}

Python packages may use an \texttt{.egg-link} instead of the copied
directory shown in the logical tree. Consequently, edits to existing
linked Python files may appear without rebuilding, but \texttt{install/}
is still needed for metadata, executables, and environment scripts.

Rebuild after changing \texttt{setup.py}, \texttt{setup.cfg},
\texttt{console\_scripts}, installed data, dependencies, CMake
configuration, or compiled C++ code. Rebuild after structural changes
and before formal verification to avoid stale state.

\section{Source the overlay and run the
package}\label{source-the-overlay-and-run-the-package}

After a successful build, source the script matching the current shell.

\begin{Shaded}
\begin{Highlighting}[]
\BuiltInTok{source}\NormalTok{ install/setup.zsh}
\end{Highlighting}
\end{Shaded}

\begin{Shaded}
\begin{Highlighting}[]
\BuiltInTok{source}\NormalTok{ install/setup.bash}
\end{Highlighting}
\end{Shaded}

This command does not install anything. It adds the workspace's
\texttt{install/} tree as an \textbf{overlay} on the ROS underlay by
updating variables such as \texttt{AMENT\_PREFIX\_PATH},
\texttt{PYTHONPATH}, and \texttt{PATH}.

Sourcing affects only the current shell. A new terminal must source the
ROS underlay and then the workspace overlay again.

Verify package discovery, Python import, and executable discovery:

\begin{Shaded}
\begin{Highlighting}[]
\ExtensionTok{ros2}\NormalTok{ pkg prefix my\_robot\_package}
\ExtensionTok{python3} \AttributeTok{{-}c} \StringTok{"import my\_robot\_package"}
\ExtensionTok{ros2}\NormalTok{ run my\_robot\_package my\_node}
\end{Highlighting}
\end{Shaded}

The final command should print:

\begin{Shaded}
\begin{Highlighting}[]
\NormalTok{Hi from my\_robot\_package.}
\end{Highlighting}
\end{Shaded}

\texttt{ros2\ pkg\ prefix} queries the ament resource index.
\texttt{ros2\ run} uses the first argument as the package name and the
second as the executable name, then starts
\texttt{lib/my\_robot\_package/my\_node}.

\section{\texorpdfstring{Run tests through
\texttt{colcon}}{Run tests through colcon}}\label{run-tests-through-colcon}

\begin{Shaded}
\begin{Highlighting}[]
\ExtensionTok{colcon}\NormalTok{ test }\AttributeTok{{-}{-}packages{-}select}\NormalTok{ my\_robot\_package}
\ExtensionTok{colcon}\NormalTok{ test{-}result }\AttributeTok{{-}{-}verbose}
\end{Highlighting}
\end{Shaded}

\texttt{colcon\ test} runs the package's configured tests and stores
structured results under the build tree. \texttt{colcon\ test-result}
aggregates them. A completed command is not necessarily successful;
verification passes only when the summary reports no errors or failures.

\section{Repeat the daily development
loop}\label{repeat-the-daily-development-loop}

From the workspace root, repeat:

\begin{Shaded}
\begin{Highlighting}[]
\CommentTok{\# 1. Edit source.}
\ExtensionTok{vim}\NormalTok{ src/my\_robot\_package/my\_robot\_package/my\_node.py}

\CommentTok{\# 2. Run fast tests.}
\ExtensionTok{python3} \AttributeTok{{-}m}\NormalTok{ pytest src/my\_robot\_package/test }\AttributeTok{{-}q}

\CommentTok{\# 3. Build and install.}
\ExtensionTok{colcon}\NormalTok{ build }\DataTypeTok{\textbackslash{}}
  \AttributeTok{{-}{-}packages{-}select}\NormalTok{ my\_robot\_package }\DataTypeTok{\textbackslash{}}
  \AttributeTok{{-}{-}symlink{-}install}

\CommentTok{\# 4. Reload the overlay.}
\BuiltInTok{source}\NormalTok{ install/setup.zsh}

\CommentTok{\# 5. Run the executable.}
\ExtensionTok{ros2}\NormalTok{ run my\_robot\_package my\_node}

\CommentTok{\# 6. Run package{-}level verification.}
\ExtensionTok{colcon}\NormalTok{ test }\AttributeTok{{-}{-}packages{-}select}\NormalTok{ my\_robot\_package}
\ExtensionTok{colcon}\NormalTok{ test{-}result }\AttributeTok{{-}{-}verbose}
\end{Highlighting}
\end{Shaded}

Use \texttt{source\ install/setup.bash} in Bash.

\section*{Common failure checks}\label{common-failure-checks}
\addcontentsline{toc}{section}{Common failure checks}

\markright{Common failure checks}

{\def\LTcaptype{none} % do not increment counter
\begin{longtable}[]{@{}
  >{\raggedright\arraybackslash}p{(\linewidth - 2\tabcolsep) * \real{0.5000}}
  >{\raggedright\arraybackslash}p{(\linewidth - 2\tabcolsep) * \real{0.5000}}@{}}
\toprule\noalign{}
\begin{minipage}[b]{\linewidth}\raggedright
Symptom
\end{minipage} & \begin{minipage}[b]{\linewidth}\raggedright
Check
\end{minipage} \\
\midrule\noalign{}
\endhead
\bottomrule\noalign{}
\endlastfoot
\texttt{ros2:\ command\ not\ found} & Source the ROS underlay in the
current shell \\
Package absent from \texttt{colcon\ list} & Check its location under
\texttt{src/} and validate \texttt{package.xml} \\
Python package is not installed & Check \texttt{\_\_init\_\_.py} and the
\texttt{packages} argument in \texttt{setup.py} \\
\texttt{ros2\ pkg\ prefix} cannot find the package & Build and source
the overlay \\
\texttt{ros2\ run} cannot find the executable & Check
\texttt{console\_scripts} and \texttt{setup.cfg}, then rebuild and
source \\
Test status is unclear & Run
\texttt{colcon\ test-result\ -\/-verbose} \\
\end{longtable}
}

\section*{References}\label{references-3}
\addcontentsline{toc}{section}{References}

\markright{References}

\begin{itemize}
\tightlist
\item
  Open Robotics,
  \href{https://docs.ros.org/en/jazzy/How-To-Guides/Developing-a-ROS-2-Package.html}{``Developing
  a ROS 2 package,''} \emph{ROS 2 Jazzy Documentation}.
\item
  Open Robotics,
  \href{https://docs.ros.org/en/jazzy/How-To-Guides/Using-Python-Packages.html}{``Using
  Python Packages with ROS 2,''} \emph{ROS 2 Jazzy Documentation}.
\item
  colcon,
  \href{https://colcon.readthedocs.io/en/released/user/what-is-a-workspace.html}{``What
  is a Workspace?''} \emph{colcon Documentation}.
\end{itemize}

\chapter{\texorpdfstring{Planar Rigid-Body Motion: \(SE(2)\) and
Twists}{Planar Rigid-Body Motion: SE(2) and Twists}}\label{sec-planar-rigid-body-motion}

\section*{Why this chapter matters}\label{why-this-chapter-matters-1}
\addcontentsline{toc}{section}{Why this chapter matters}

\markright{Why this chapter matters}

The differential-drive pose \(\mathbf q=[x\;y\;\theta]^\mathsf T\) is
convenient for storing three coordinates, but it does not by itself
expose how rigid motions compose, how the same motion is expressed in
different frames, or why constant body velocity produces the exact
curved-motion update derived previously. This chapter develops the
planar rigid-motion structure needed to answer those questions. It
begins by distinguishing a pose from its coordinates and representing a
planar pose by a homogeneous transformation in the special Euclidean
group \(SE(2)\).

\section*{Prerequisites and
assumptions}\label{prerequisites-and-assumptions-3}
\addcontentsline{toc}{section}{Prerequisites and assumptions}

\markright{Prerequisites and assumptions}

The reader should understand matrices and matrix multiplication from
Chapter~\ref{sec-linear-algebra-foundations}, coordinate frames and
rotation matrices from Chapter~\ref{sec-geometry-coordinate-frames}, and
planar pose notation from Chapter~\ref{sec-differential-drive-model}.

All frames in this learning block are right-handed planar frames
embedded in three-dimensional physical space with a common vertical
axis. Frame \(\{w\}\) is the fixed world frame and frame \(\{b\}\) is
attached rigidly to the robot body. Distances are measured in metres and
angles in radians.

\section{Pose coordinates and rigid
transformations}\label{pose-coordinates-and-rigid-transformations}

\subsection{A pose is a relative geometric
relationship}\label{a-pose-is-a-relative-geometric-relationship}

A planar rigid body's \textbf{pose} specifies both its position and its
orientation relative to a reference frame. The word \emph{relative} is
essential: saying that a robot is at position \((2,1)\) with heading
\(\pi/2\) is incomplete until the reference frame and sign convention
are known.

Let \(O_w\) and \(O_b\) denote the origins of frames \(\{w\}\) and
\(\{b\}\). The position of \(O_b\) relative to \(O_w\), expressed in
world-frame coordinates, is

\[
{}^{w}\mathbf p_b
=
\begin{bmatrix}
x\\
y
\end{bmatrix}
\in\mathbb R^2,
\]

where \(x\) and \(y\) are measured in metres. The orientation of frame
\(\{b\}\) relative to frame \(\{w\}\) is represented by the rotation
matrix \(\mathbf R_{wb}(\theta)\in\mathbb R^{2\times2}\):

\[
\mathbf R_{wb}(\theta)
=
\begin{bmatrix}
\cos\theta&-\sin\theta\\
\sin\theta&\cos\theta
\end{bmatrix}.
\]

The subscript order \(wb\) means that \(\mathbf R_{wb}\) maps vector
coordinates from frame \(\{b\}\) to frame \(\{w\}\). The pose of
\(\{b\}\) relative to \(\{w\}\) can therefore be described by the pair

\[
\left(\mathbf R_{wb},{}^{w}\mathbf p_b\right).
\]

This pair separates orientation from position. A homogeneous
transformation will combine them into one matrix without confusing their
different roles.

\subsection{\texorpdfstring{Planar rotations and
\(SO(2)\)}{Planar rotations and SO(2)}}\label{sec-se2-planar-rotations}

The set of valid planar rotation matrices is called the \textbf{special
orthogonal group in two dimensions}, written \(SO(2)\) and defined by

\[
SO(2)
\mathrel{:=}
\left\{
\mathbf R\in\mathbb R^{2\times2}
\;\middle|\;
\mathbf R^\mathsf T\mathbf R=\mathbf I_2,
\quad
\det(\mathbf R)=1
\right\}.
\]

The word \emph{orthogonal} refers to
\(\mathbf R^\mathsf T\mathbf R=\mathbf I_2\), which means that the
columns form an orthonormal basis and that
\(\mathbf R^{-1}=\mathbf R^\mathsf T\). The word \emph{special} adds the
condition \(\det(\mathbf R)=1\), excluding reflections whose determinant
is \(-1\).

The word \emph{group} means that the set is equipped with an operation
satisfying four properties: closure, associativity, an identity element,
and an inverse for every element. For \(SO(2)\), the operation is matrix
multiplication.

\textbf{Closure} means that multiplying any two elements of \(SO(2)\)
produces another element of \(SO(2)\). If \(\mathbf R(\alpha)\) and
\(\mathbf R(\beta)\) represent planar rotations through angles
\(\alpha\) and \(\beta\), then

\[
\mathbf R(\alpha)\mathbf R(\beta)
=
\mathbf R(\alpha+\beta)
\in SO(2).
\]

Geometrically, performing one planar rotation and then another still
produces a valid planar rotation; it does not introduce scaling, shear,
or reflection.

\textbf{Associativity} means that the placement of parentheses does not
change the result of composing three rotations:

\[
\left(\mathbf R_1\mathbf R_2\right)\mathbf R_3
=
\mathbf R_1\left(\mathbf R_2\mathbf R_3\right).
\]

Associativity changes only the grouping, not the order, of the matrices.
It should not be confused with commutativity, which would permit the
order to be exchanged. The identity element is \(\mathbf I_2\), and
every \(\mathbf R\in SO(2)\) has the inverse
\(\mathbf R^\mathsf T\in SO(2)\).

Every planar rotation matrix can be written as \(\mathbf R(\theta)\) for
some \(\theta\in\mathbb R\). The angle is not unique because

\[
\mathbf R(\theta+2\pi n)=\mathbf R(\theta),
\qquad n\in\mathbb Z,
\]

where \(\mathbb Z\) is the set of integers. Thus, \(\theta\) is a
\textbf{coordinate} used to describe an orientation, while
\(\mathbf R(\theta)\) represents the \textbf{orientation} itself. Angles
that differ by an integer multiple of \(2\pi\) describe the same element
of \(SO(2)\).

\subsection{Transforming a point by rotation and
translation}\label{sec-se2-transforming-point}

Let \(P\) be a geometric point fixed to the robot. Its position relative
to \(O_b\), expressed in body-frame coordinates, is
\({}^{b}\mathbf p_P\in\mathbb R^2\) and is measured in metres. Its
world-frame position is obtained by first rotating its body-frame
displacement and then adding the world-frame position of the body
origin:

\[
\boxed{
{}^{w}\mathbf p_P
=
\mathbf R_{wb}\,{}^{b}\mathbf p_P
+
{}^{w}\mathbf p_b
}.
\]

This equation follows by tracing the geometric displacement from \(O_w\)
to \(P\). First travel from \(O_w\) to \(O_b\) by \({}^{w}\mathbf p_b\).
Then travel from \(O_b\) to \(P\) by \({}^{b}\mathbf p_P\), but rotate
that second displacement with \(\mathbf R_{wb}\) so that both terms are
expressed in world-frame coordinates before adding them.

A free vector, such as a displacement or velocity vector, has direction
and magnitude but no fixed origin. It therefore changes coordinates
using rotation alone:

\[
{}^{w}\mathbf v
=
\mathbf R_{wb}\,{}^{b}\mathbf v.
\]

Translation appears in the point equation because a point is located
relative to a frame origin. It does not appear in the free-vector
equation because a free vector is not attached to an origin.

\subsection{\texorpdfstring{Homogeneous transformations and
\(SE(2)\)}{Homogeneous transformations and SE(2)}}\label{homogeneous-transformations-and-se2}

The point mapping in Section~\ref{sec-se2-transforming-point} contains
both a matrix multiplication and a vector addition. Homogeneous
coordinates package these two operations into one matrix multiplication.
Define the homogeneous coordinate column of point \(P\) by appending a
\(1\):

\[
{}^{b}\overline{\mathbf p}_P
\mathrel{:=}
\begin{bmatrix}
{}^{b}\mathbf p_P\\
1
\end{bmatrix}
\in\mathbb R^3.
\]

The overline identifies a homogeneous coordinate representation, not a
different geometric point. Define the \textbf{homogeneous
transformation} from frame \(\{b\}\) to frame \(\{w\}\) by

\[
\boxed{
\mathbf T_{wb}
\mathrel{:=}
\begin{bmatrix}
\mathbf R_{wb}&{}^{w}\mathbf p_b\\
\mathbf 0_2^\mathsf T&1
\end{bmatrix}
}
\in\mathbb R^{3\times3},
\]

where \(\mathbf 0_2=[0\;0]^\mathsf T\in\mathbb R^2\). The subscript
order means that \(\mathbf T_{wb}\) maps homogeneous point coordinates
from frame \(\{b\}\) to frame \(\{w\}\):

\[
\begin{aligned}
{}^{w}\overline{\mathbf p}_P
&=
\mathbf T_{wb}\,{}^{b}\overline{\mathbf p}_P\\
&=
\begin{bmatrix}
\mathbf R_{wb}&{}^{w}\mathbf p_b\\
\mathbf 0_2^\mathsf T&1
\end{bmatrix}
\begin{bmatrix}
{}^{b}\mathbf p_P\\
1
\end{bmatrix}\\
&=
\begin{bmatrix}
\mathbf R_{wb}\,{}^{b}\mathbf p_P+{}^{w}\mathbf p_b\\
1
\end{bmatrix}.
\end{aligned}
\]

The upper two entries reproduce the point-transformation equation.
Appending \(1\) causes the translation column to contribute to the
result. Thus, \(\mathbf T_{wb}\) rotates the point's body-frame
displacement into the world frame and then adds the world position of
the body origin.

Using \({}^{w}\mathbf p_b=[x\;y]^\mathsf T\) and
\(\mathbf R_{wb}=\mathbf R_{wb}(\theta)\) gives the component form

\[
\mathbf T_{wb}
=
\begin{bmatrix}
\cos\theta&-\sin\theta&x\\
\sin\theta&\cos\theta&y\\
0&0&1
\end{bmatrix}.
\]

The rotation block is dimensionless, while the first two entries of the
translation column have units of metres. The final row is an algebraic
device that permits rotation and translation to be applied by one matrix
multiplication; it does not represent another physical spatial
dimension.

To preserve the rotation-only mapping for a free vector, define its
homogeneous representation by appending a \(0\):

\[
{}^{b}\overline{\mathbf v}
\mathrel{:=}
\begin{bmatrix}
{}^{b}\mathbf v\\
0
\end{bmatrix}.
\]

Then

\[
\mathbf T_{wb}
\begin{bmatrix}
{}^{b}\mathbf v\\
0
\end{bmatrix}
=
\begin{bmatrix}
\mathbf R_{wb}\,{}^{b}\mathbf v\\
0
\end{bmatrix}.
\]

Appending \(0\) prevents the translation column from contributing. This
algebraic distinction encodes the geometric fact that changing a point's
reference frame requires rotation and translation, whereas changing the
coordinates of a free vector requires only rotation.

The set of all planar homogeneous transformations is the \textbf{special
Euclidean group in two dimensions}:

\[
SE(2)
\mathrel{:=}
\left\{
\begin{bmatrix}
\mathbf R&\mathbf p\\
\mathbf 0_2^\mathsf T&1
\end{bmatrix}
\;\middle|\;
\mathbf R\in SO(2),
\quad
\mathbf p\in\mathbb R^2
\right\}.
\]

Here, \emph{Euclidean} indicates rigid motions that preserve Euclidean
distances and angles. Matrix multiplication is the group operation, the
\(3\times3\) identity \(\mathbf I_3\) is the identity motion, and every
element has an inverse that reverses the coordinate mapping.
Section~\ref{sec-se2-inverting-transformation} derives that inverse
explicitly.

\subsection{Composing transformations}\label{composing-transformations}

Introduce a third frame \(\{c\}\). Let \(\mathbf T_{ab}\in SE(2)\) map
coordinates from frame \(\{b\}\) to frame \(\{a\}\), and let
\(\mathbf T_{bc}\in SE(2)\) map coordinates from frame \(\{c\}\) to
frame \(\{b\}\). Applying the two mappings to a point gives

\[
{}^{a}\overline{\mathbf p}_P
=
\mathbf T_{ab}
\left(
\mathbf T_{bc}\,{}^{c}\overline{\mathbf p}_P
\right)
=
\left(
\mathbf T_{ab}\mathbf T_{bc}
\right)
{}^{c}\overline{\mathbf p}_P.
\]

Therefore, the transform that maps directly from frame \(\{c\}\) to
frame \(\{a\}\) is

\[
\boxed{
\mathbf T_{ac}
=
\mathbf T_{ab}\mathbf T_{bc}
}.
\]

The adjacent frame label \(b\) identifies the intermediate frame.
Expanding the block multiplication gives

\[
\begin{aligned}
\mathbf T_{ab}\mathbf T_{bc}
&=
\begin{bmatrix}
\mathbf R_{ab}&{}^{a}\mathbf p_b\\
\mathbf 0_2^\mathsf T&1
\end{bmatrix}
\begin{bmatrix}
\mathbf R_{bc}&{}^{b}\mathbf p_c\\
\mathbf 0_2^\mathsf T&1
\end{bmatrix}\\
&=
\begin{bmatrix}
\mathbf R_{ab}\mathbf R_{bc}&
\mathbf R_{ab}\,{}^{b}\mathbf p_c+{}^{a}\mathbf p_b\\
\mathbf 0_2^\mathsf T&1
\end{bmatrix}.
\end{aligned}
\]

Consequently,

\[
\boxed{
\mathbf R_{ac}=\mathbf R_{ab}\mathbf R_{bc},
\qquad
{}^{a}\mathbf p_c
=
\mathbf R_{ab}\,{}^{b}\mathbf p_c
+
{}^{a}\mathbf p_b
}.
\]

The translation vectors cannot be added directly until they are
expressed in the same frame. The product first rotates
\({}^{b}\mathbf p_c\) into frame \(\{a\}\) and only then adds
\({}^{a}\mathbf p_b\).

\subsection{Inverting a
transformation}\label{sec-se2-inverting-transformation}

The inverse transform \(\mathbf T_{ba}=\mathbf T_{ab}^{-1}\) must undo
both the translation and the rotation. Starting from the point mapping

\[
{}^{a}\mathbf p_P
=
\mathbf R_{ab}\,{}^{b}\mathbf p_P
+
{}^{a}\mathbf p_b,
\]

subtract \({}^{a}\mathbf p_b\) and premultiply by
\(\mathbf R_{ab}^{-1}=\mathbf R_{ab}^\mathsf T\):

\[
{}^{b}\mathbf p_P
=
\mathbf R_{ab}^\mathsf T
\left(
{}^{a}\mathbf p_P-{}^{a}\mathbf p_b
\right).
\]

Distributing the rotation gives

\[
{}^{b}\mathbf p_P
=
\mathbf R_{ab}^\mathsf T{}^{a}\mathbf p_P
-
\mathbf R_{ab}^\mathsf T{}^{a}\mathbf p_b.
\]

Therefore,

\[
\boxed{
\mathbf T_{ab}^{-1}
=
\mathbf T_{ba}
=
\begin{bmatrix}
\mathbf R_{ab}^\mathsf T&
-\mathbf R_{ab}^\mathsf T{}^{a}\mathbf p_b\\
\mathbf 0_2^\mathsf T&1
\end{bmatrix}
}.
\]

The inverse translation is not generally \(-{}^{a}\mathbf p_b\) because
that column is expressed in frame \(\{a\}\). It must first be negated
and then expressed in frame \(\{b\}\) by \(\mathbf R_{ab}^\mathsf T\).

\subsection{Pose coordinates versus the pose
itself}\label{pose-coordinates-versus-the-pose-itself}

The familiar column

\[
\mathbf q
=
\begin{bmatrix}
x\\y\\\theta
\end{bmatrix}
\]

has the following meaning:

\begin{itemize}
\tightlist
\item
  \(x\in\mathbb R\) is the \(x_w\)-coordinate of the body origin
  \(O_b\), expressed in the world frame and measured in metres;
\item
  \(y\in\mathbb R\) is the \(y_w\)-coordinate of \(O_b\), expressed in
  the world frame and measured in metres; and
\item
  \(\theta\in\mathbb R\) is the counter-clockwise orientation of body
  frame \(\{b\}\) relative to world frame \(\{w\}\), measured in
  radians.
\end{itemize}

Therefore, the position of the body origin relative to the world origin,
expressed in world-frame coordinates, is

\[
{}^w\mathbf p_b
=
\begin{bmatrix}
x\\y
\end{bmatrix}.
\]

The notation \({}^w\mathbf p_b\) means ``the position of the body origin
\(O_b\), expressed in world-frame coordinates.'' Thus, \(\mathbf q\) is
a coordinate description of the pose of frame \(\{b\}\) relative to
frame \(\{w\}\). The same pose can be represented by
\(\mathbf T_{wb}\in SE(2)\):

\[
\boxed{
\mathbf q
=
\begin{bmatrix}
x\\y\\\theta
\end{bmatrix}
\quad\longmapsto\quad
\mathbf T_{wb}(\mathbf q)
=
\begin{bmatrix}
\cos\theta&-\sin\theta&x\\
\sin\theta&\cos\theta&y\\
0&0&1
\end{bmatrix}
}.
\]

The coordinate column is useful for storage, plotting, and differential
equations, but it is not an ordinary geometric vector. Its components
represent different physical quantities, \(\theta\) and
\(\theta+2\pi n\) represent the same orientation, and componentwise
addition does not in general reproduce rigid-motion composition.
Transformations should be composed by matrix multiplication with
explicit frame labels.

\subsection{Worked example}\label{worked-example}

Suppose the body origin is at

\[
{}^{w}\mathbf p_b
=
\begin{bmatrix}
2\\1
\end{bmatrix}
\mathrm m
\]

and the body heading is \(\theta=\pi/2\,\mathrm{rad}\). A point \(P\)
fixed one metre along the body \(+x_b\)-axis has body-frame coordinates

\[
{}^{b}\mathbf p_P
=
\begin{bmatrix}
1\\0
\end{bmatrix}
\mathrm m.
\]

The rotation and homogeneous transformation are

\[
\mathbf R_{wb}\left(\frac{\pi}{2}\right)
=
\begin{bmatrix}
0&-1\\
1&0
\end{bmatrix},
\qquad
\mathbf T_{wb}
=
\begin{bmatrix}
0&-1&2\\
1&0&1\\
0&0&1
\end{bmatrix}.
\]

Applying the transform gives

\[
\begin{aligned}
{}^{w}\overline{\mathbf p}_P
&=
\mathbf T_{wb}{}^{b}\overline{\mathbf p}_P\\
&=
\begin{bmatrix}
0&-1&2\\
1&0&1\\
0&0&1
\end{bmatrix}
\begin{bmatrix}
1\\0\\1
\end{bmatrix}\\
&=
\begin{bmatrix}
2\\2\\1
\end{bmatrix}.
\end{aligned}
\]

Thus, \({}^{w}\mathbf p_P=[2\;2]^\mathsf T\,\mathrm m\). The one-metre
body-frame displacement is rotated into the world \(+y_w\)-direction and
then added to the body origin at \([2\;1]^\mathsf T\,\mathrm m\).

\section*{\texorpdfstring{Quick test A: pose and
\(SE(2)\)}{Quick test A: pose and SE(2)}}\label{quick-test-a-pose-and-se2}
\addcontentsline{toc}{section}{Quick test A: pose and \(SE(2)\)}

\markright{Quick test A: pose and \(SE(2)\)}

\begin{enumerate}
\def\labelenumi{\arabic{enumi}.}
\tightlist
\item
  What geometric information is contained in the pair
  \((\mathbf R_{wb},{}^{w}\mathbf p_b)\), and what does the subscript
  order \(wb\) mean?
\item
  Define \(SO(2)\). What do its orthogonality and determinant conditions
  exclude or guarantee?
\item
  Why do \(\theta\) and \(\theta+2\pi\) describe the same orientation?
\item
  Starting from the block definitions of \(\mathbf T_{ab}\) and
  \(\mathbf T_{bc}\), derive the translation column of
  \(\mathbf T_{ac}=\mathbf T_{ab}\mathbf T_{bc}\).
\item
  Why is the inverse translation
  \(-\mathbf R_{ab}^\mathsf T{}^{a}\mathbf p_b\) rather than simply
  \(-{}^{a}\mathbf p_b\)?
\item
  A body frame has world position \([3\;2]^\mathsf T\,\mathrm m\) and
  heading \(-\pi/2\,\mathrm{rad}\). Find the world coordinates of a
  point whose body coordinates are \([0.5\;0]^\mathsf T\,\mathrm m\).
\item
  Explain why \(\mathbf q_1+\mathbf q_2\) is not generally the
  composition of the two poses represented by \(\mathbf q_1\) and
  \(\mathbf q_2\).
\item
  Name two transformations that cannot be represented by an element of
  \(SE(2)\) and explain why.
\end{enumerate}

\subsection*{Answers and hints}\label{answers-and-hints-10}
\addcontentsline{toc}{subsection}{Answers and hints}

\begin{enumerate}
\def\labelenumi{\arabic{enumi}.}
\tightlist
\item
  \(\mathbf R_{wb}\) describes the orientation of frame \(\{b\}\)
  relative to frame \(\{w\}\), while \({}^{w}\mathbf p_b\) locates
  \(O_b\) relative to \(O_w\) using world coordinates. The order \(wb\)
  means that the rotation maps coordinate columns from \(\{b\}\) to
  \(\{w\}\).
\item
  \(SO(2)=\{\mathbf R\in\mathbb R^{2\times2}\mid\mathbf R^\mathsf T\mathbf R=\mathbf I_2,\det\mathbf R=1\}\).
  Orthogonality preserves lengths and angles and gives
  \(\mathbf R^{-1}=\mathbf R^\mathsf T\); determinant \(+1\) preserves
  orientation and excludes reflections.
\item
  Sine and cosine are \(2\pi\)-periodic, so every entry of
  \(\mathbf R(\theta+2\pi)\) equals the corresponding entry of
  \(\mathbf R(\theta)\).
\item
  Block multiplication gives
  \({}^{a}\mathbf p_c=\mathbf R_{ab}{}^{b}\mathbf p_c+{}^{a}\mathbf p_b\).
  The first term rotates the \(b\)-frame translation into frame
  \(\{a\}\) before the two translations are added.
\item
  The reversed displacement must be expressed in frame \(\{b\}\),
  whereas \({}^{a}\mathbf p_b\) is expressed in frame \(\{a\}\).
  Premultiplication by \(\mathbf R_{ab}^\mathsf T\) performs the
  necessary coordinate change.
\item
  \(\mathbf R_{wb}(-\pi/2)[0.5\;0]^\mathsf T=[0\;-0.5]^\mathsf T\,\mathrm m\),
  so \({}^{w}\mathbf p_P=[3\;1.5]^\mathsf T\,\mathrm m\).
\item
  Rigid-motion composition must account for the frame in which each
  translation is expressed and for the order of rotations and
  translations. Componentwise addition does neither, and angular
  coordinates are periodic.
\item
  Examples include nonuniform scaling, shear, and elastic deformation.
  They do not preserve every distance and angle, so they are not rigid
  motions and do not belong to \(SE(2)\).
\end{enumerate}

\section{Passive coordinate changes and active rigid
motions}\label{passive-coordinate-changes-and-active-rigid-motions}

The same numerical rotation or homogeneous-transformation matrix can
appear in different geometric questions. Correct interpretation requires
identifying what remains fixed, what changes, and the frame in which
every coordinate column is expressed.

\subsection{Passive coordinate change}\label{passive-coordinate-change}

A \textbf{passive coordinate change} keeps the geometric point, vector,
or rigid body unchanged and changes only the coordinate frame used to
describe it. The word \emph{passive} indicates that nothing physically
moves.

For example, let the geometric point \(P\) remain fixed while its
coordinates are changed from frame \(\{b\}\) to frame \(\{w\}\). The
passive point-coordinate transformation is

\[
{}^{w}\overline{\mathbf p}_P
=
\mathbf T_{wb}\,{}^{b}\overline{\mathbf p}_P.
\]

The two coordinate columns \({}^{b}\overline{\mathbf p}_P\) and
\({}^{w}\overline{\mathbf p}_P\) contain different numbers in general,
but both describe the same geometric point \(P\). The transform
\(\mathbf T_{wb}\) accounts for the position and orientation of frame
\(\{b\}\) relative to frame \(\{w\}\).

For a free vector \(\mathbf v\), translation is irrelevant and only the
rotation block is used:

\[
{}^{w}\mathbf v
=
\mathbf R_{wb}\,{}^{b}\mathbf v.
\]

Suppose frame \(\{b\}\) is rotated counter-clockwise by \(\pi/2\)
relative to frame \(\{w\}\) and the vector lies along \(+x_b\). Its
world-frame coordinates are

\[
\begin{aligned}
{}^{w}\mathbf v
&=
\mathbf R_{wb}{}^{b}\mathbf v\\
&=
\begin{bmatrix}
0&-1\\
1&0
\end{bmatrix}
\begin{bmatrix}
1\\0
\end{bmatrix}\\
&=
\begin{bmatrix}
0\\1
\end{bmatrix}.
\end{aligned}
\]

The vector has not rotated physically. It is the same geometric vector
described once using the axes of \(\{b\}\) and once using the axes of
\(\{w\}\).

\subsection{Active rigid-body motion}\label{active-rigid-body-motion}

An \textbf{active transformation} keeps the reference frame fixed and
moves a geometric point, vector, or rigid body. The initial point \(P\)
and final point \(P'\) are different geometric points, but both
coordinate columns are expressed in the same frame.

Let \(\mathbf R_{\Delta}\in SO(2)\) be an applied planar rotation and
let \(\mathbf d_{\Delta}\in\mathbb R^2\) be an applied translation
expressed in world-frame coordinates and measured in metres. Define the
active rigid-motion operator

\[
\mathbf T_{\Delta}
\mathrel{:=}
\begin{bmatrix}
\mathbf R_{\Delta}&\mathbf d_{\Delta}\\
\mathbf 0_2^\mathsf T&1
\end{bmatrix}
\in SE(2).
\]

Applied to point \(P\), it produces the new point \(P'\) according to

\[
\boxed{
{}^{w}\overline{\mathbf p}_{P'}
=
\mathbf T_{\Delta}\,{}^{w}\overline{\mathbf p}_{P}
},
\]

or, using ordinary two-dimensional coordinates,

\[
{}^{w}\mathbf p_{P'}
=
\mathbf R_{\Delta}{}^{w}\mathbf p_P
+
\mathbf d_{\Delta}.
\]

Both sides are expressed in frame \(\{w\}\). The point changes from
\(P\) to \(P'\), unlike a passive coordinate change, where the point
remains \(P\) and the leading frame superscript changes.

For example, let \(\mathbf R_{\Delta}=\mathbf R(\pi/2)\),
\(\mathbf d_{\Delta}=[2\;1]^\mathsf T\,\mathrm m\), and
\({}^{w}\mathbf p_P=[1\;0]^\mathsf T\,\mathrm m\). Then

\[
{}^{w}\mathbf p_{P'}
=
\begin{bmatrix}
0&-1\\
1&0
\end{bmatrix}
\begin{bmatrix}
1\\0
\end{bmatrix}
+
\begin{bmatrix}
2\\1
\end{bmatrix}
=
\begin{bmatrix}
2\\2
\end{bmatrix}
\mathrm m.
\]

Here, the coordinates changed because the geometric point moved, not
because a different coordinate frame was selected.

\subsection{One matrix form, three
interpretations}\label{one-matrix-form-three-interpretations}

A homogeneous transformation may describe a relative pose, perform a
passive coordinate change, or act as an active rigid-motion operator.
The matrix arithmetic can be identical while the geometric meaning
differs.

{\def\LTcaptype{none} % do not increment counter
\begin{longtable}[]{@{}
  >{\raggedright\arraybackslash}p{(\linewidth - 4\tabcolsep) * \real{0.3333}}
  >{\raggedright\arraybackslash}p{(\linewidth - 4\tabcolsep) * \real{0.3333}}
  >{\raggedright\arraybackslash}p{(\linewidth - 4\tabcolsep) * \real{0.3333}}@{}}
\toprule\noalign{}
\begin{minipage}[b]{\linewidth}\raggedright
Interpretation
\end{minipage} & \begin{minipage}[b]{\linewidth}\raggedright
What changes?
\end{minipage} & \begin{minipage}[b]{\linewidth}\raggedright
Representative statement
\end{minipage} \\
\midrule\noalign{}
\endhead
\bottomrule\noalign{}
\endlastfoot
Relative-pose description & Nothing is being operated on; the matrix
records the pose of one frame relative to another. & \(\mathbf T_{wb}\)
describes frame \(\{b\}\) relative to frame \(\{w\}\). \\
Passive coordinate change & The coordinate column and its frame label
change; the geometric point remains fixed. &
\({}^{w}\overline{\mathbf p}_P=\mathbf T_{wb}{}^{b}\overline{\mathbf p}_P\). \\
Active rigid motion & The geometric point changes; the coordinate frame
remains fixed. &
\({}^{w}\overline{\mathbf p}_{P'}=\mathbf T_{\Delta}{}^{w}\overline{\mathbf p}_P\). \\
\end{longtable}
}

The project notation reserves labelled transforms such as
\(\mathbf T_{wb}\) for relative poses and coordinate mappings. A symbol
such as \(\mathbf T_{\Delta}\) must be explicitly declared when the same
matrix form is used as an active motion operator.

\subsection{Composition order}\label{composition-order}

Transformations act on column vectors from the left, so the rightmost
matrix acts first. For a passive coordinate chain,

\[
{}^{a}\overline{\mathbf p}_P
=
\mathbf T_{ab}\mathbf T_{bc}{}^{c}\overline{\mathbf p}_P,
\]

the rightmost transform \(\mathbf T_{bc}\) first changes coordinates
from frame \(\{c\}\) to frame \(\{b\}\), and \(\mathbf T_{ab}\) then
changes them from frame \(\{b\}\) to frame \(\{a\}\). Thus,

\[
\mathbf T_{ac}
=
\mathbf T_{ab}\mathbf T_{bc}.
\]

For two active motions \(\mathbf T_1\) and \(\mathbf T_2\),

\[
{}^{w}\overline{\mathbf p}_{P''}
=
\mathbf T_2\mathbf T_1{}^{w}\overline{\mathbf p}_{P},
\]

motion \(\mathbf T_1\) occurs first and motion \(\mathbf T_2\) occurs
second. Associativity permits the parentheses in a longer product to be
regrouped, but it does not permit the matrix order to be reversed.

\subsection{\texorpdfstring{Why composition in \(SE(2)\) is not
commutative}{Why composition in SE(2) is not commutative}}\label{sec-se2-noncommutative}

Planar rotations alone commute, but planar rigid transformations
generally do not because rotation changes the direction of a
translation. Consider a translation of one metre along \(+x_w\) and a
counter-clockwise rotation of \(\pi/2\) about \(O_w\):

\[
\mathbf T_D
=
\begin{bmatrix}
1&0&1\\
0&1&0\\
0&0&1
\end{bmatrix},
\qquad
\mathbf T_R
=
\begin{bmatrix}
0&-1&0\\
1&0&0\\
0&0&1
\end{bmatrix}.
\]

Translate first and rotate second:

\[
\mathbf T_R\mathbf T_D
=
\begin{bmatrix}
0&-1&0\\
1&0&1\\
0&0&1
\end{bmatrix}.
\]

Rotate first and translate second:

\[
\mathbf T_D\mathbf T_R
=
\begin{bmatrix}
0&-1&1\\
1&0&0\\
0&0&1
\end{bmatrix}.
\]

The products differ. Applied to the world origin
\({}^{w}\overline{\mathbf p}_O=[0\;0\;1]^\mathsf T\), they give

\[
\mathbf T_R\mathbf T_D{}^{w}\overline{\mathbf p}_O
=
\begin{bmatrix}
0\\1\\1
\end{bmatrix},
\qquad
\mathbf T_D\mathbf T_R{}^{w}\overline{\mathbf p}_O
=
\begin{bmatrix}
1\\0\\1
\end{bmatrix}.
\]

Therefore,

\[
\boxed{
\mathbf T_R\mathbf T_D
\neq
\mathbf T_D\mathbf T_R
}.
\]

This noncommutativity is why a transformation sequence must state both
the individual motions and their order.

\subsection{Interpretation failures and engineering
checks}\label{interpretation-failures-and-engineering-checks}

\begin{itemize}
\tightlist
\item
  \textbf{Missing frame labels:} An equation such as
  \(\mathbf p'=\mathbf T\mathbf p\) does not reveal whether
  \(\mathbf T\) changes coordinates or moves the point. Declare the
  interpretation and frames.
\item
  \textbf{Applying transformations in the wrong order:} For
  column-vector notation, transformations act from right to left. In
  \(\mathbf T_2\mathbf T_1\mathbf p\), \(\mathbf T_1\) acts first and
  \(\mathbf T_2\) acts second. Verify an implementation using a point
  whose expected final location is easy to calculate.
\item
  \textbf{Mixing coordinate frames:} Translation columns may be added
  only after they are expressed in the same frame.
\item
  \textbf{Assuming wrong commutativity:} Reordering a rotation and
  translation generally changes the resulting pose.
\end{itemize}

\section*{Quick test B: interpretation and
composition}\label{quick-test-b-interpretation-and-composition}
\addcontentsline{toc}{section}{Quick test B: interpretation and
composition}

\markright{Quick test B: interpretation and composition}

\begin{enumerate}
\def\labelenumi{\arabic{enumi}.}
\tightlist
\item
  Define a passive coordinate change. What changes, and what remains
  fixed?
\item
  Define an active rigid-body motion. What changes, and what remains
  fixed?
\item
  Explain the three interpretations of a homogeneous transformation used
  in this chapter.
\item
  In \(\mathbf T_{ab}\mathbf T_{bc}{}^{c}\overline{\mathbf p}_P\), which
  matrix acts first, and what frame sequence does the product represent?
\item
  For active motions, what sequence is represented by
  \(\mathbf T_2\mathbf T_1{}^{w}\overline{\mathbf p}_P\)?
\item
  Using \(\mathbf T_D\) and \(\mathbf T_R\) from
  Section~\ref{sec-se2-noncommutative}, explain geometrically why
  applying them to the origin in opposite orders gives different
  results.
\item
  Does the commutativity of planar rotations imply that all elements of
  \(SE(2)\) commute? Explain.
\item
  A program receives an unlabelled matrix \(\mathbf T\) and an
  unlabelled coordinate column \(\mathbf p\). What information is
  missing before the multiplication can be interpreted safely?
\end{enumerate}

\subsection*{Answers and hints}\label{answers-and-hints-11}
\addcontentsline{toc}{subsection}{Answers and hints}

\begin{enumerate}
\def\labelenumi{\arabic{enumi}.}
\tightlist
\item
  A passive coordinate change preserves the geometric object and changes
  only the frame used to express its coordinates. The coordinate values
  and frame label may change.
\item
  An active rigid-body motion preserves the reference frame and moves
  the geometric object. The initial and final coordinates describe
  different points or poses in the same frame.
\item
  A homogeneous transformation can describe the pose of one frame
  relative to another, map coordinates of an unchanged object between
  frames, or actively move an object while its coordinates remain
  expressed in one frame.
\item
  \(\mathbf T_{bc}\) acts first, mapping from \(\{c\}\) to \(\{b\}\);
  \(\mathbf T_{ab}\) then maps from \(\{b\}\) to \(\{a\}\). The complete
  product maps from \(\{c\}\) to \(\{a\}\).
\item
  \(\mathbf T_1\) moves the point first, and \(\mathbf T_2\) moves the
  resulting point second.
\item
  Translating the origin first places it at \([1\;0]^\mathsf T\), which
  the subsequent rotation moves to \([0\;1]^\mathsf T\). Rotating the
  origin first leaves it at the origin, after which translation places
  it at \([1\;0]^\mathsf T\).
\item
  No.~Although \(SO(2)\) rotations commute with each other, an \(SE(2)\)
  transformation also contains translation. Rotation changes the
  direction of a translation, so general rotation-translation products
  do not commute.
\item
  The program needs to know whether \(\mathbf p\) represents a point or
  free vector, the frame in which it is expressed, the source and
  destination frames or active-motion interpretation of \(\mathbf T\),
  the units, and the intended multiplication order.
\end{enumerate}

\section{Time-varying rotations and rigid-body point
velocity}\label{time-varying-rotations-and-rigid-body-point-velocity}

A pose describes where a rigid body is at one instant. A time-varying
pose describes motion. This section differentiates the position and
orientation components of \(\mathbf T_{wb}(t)\) and derives the velocity
of any point fixed to the body.

\subsection{Translational and angular
velocity}\label{translational-and-angular-velocity}

Let the planar body pose vary continuously with time:

\[
\mathbf T_{wb}(t)
=
\begin{bmatrix}
\mathbf R_{wb}(t)&{}^{w}\mathbf p_b(t)\\
\mathbf 0_2^\mathsf T&1
\end{bmatrix}
\in SE(2),
\]

where \(t\in\mathbb R_{\geq0}\) is time in seconds. The world-frame
position of the body origin is

\[
{}^{w}\mathbf p_b(t)
=
\begin{bmatrix}
x(t)\\
y(t)
\end{bmatrix}
\in\mathbb R^2.
\]

Its derivative, taken with respect to the fixed world frame and
expressed in world-frame coordinates, is the translational velocity of
\(O_b\):

\[
\boxed{
{}^{w}\mathbf v_b(t)
\mathrel{:=}
\frac{d}{dt}{}^{w}\mathbf p_b(t)
=
\begin{bmatrix}
\dot x(t)\\
\dot y(t)
\end{bmatrix}
}
\in\mathbb R^2.
\]

The entries of \({}^{w}\mathbf v_b\) are measured in
\(\mathrm{m\,s^{-1}}\). The heading is \(\theta(t)\in\mathbb R\),
measured in radians, and its derivative is the signed yaw rate

\[
\boxed{
\omega(t)
\mathrel{:=}
\dot\theta(t)
}
\in\mathbb R,
\]

measured in \(\mathrm{rad\,s^{-1}}\). Positive \(\omega\) denotes
counter-clockwise rotation about the common \(+z\)-axis of the planar
frames.

\subsection{Deriving the rotation-matrix
derivative}\label{sec-se2-rotation-matrix-derivative}

The time-varying planar rotation matrix is

\[
\mathbf R_{wb}(t)
=
\begin{bmatrix}
\cos\theta(t)&-\sin\theta(t)\\
\sin\theta(t)&\cos\theta(t)
\end{bmatrix}.
\]

Differentiate every entry and apply the chain rule:

\[
\begin{aligned}
\dot{\mathbf R}_{wb}(t)
&=
\dot\theta(t)
\begin{bmatrix}
-\sin\theta(t)&-\cos\theta(t)\\
\cos\theta(t)&-\sin\theta(t)
\end{bmatrix}\\
&=
\omega(t)
\begin{bmatrix}
-\sin\theta(t)&-\cos\theta(t)\\
\cos\theta(t)&-\sin\theta(t)
\end{bmatrix}.
\end{aligned}
\]

A square matrix \(\mathbf A\) is \textbf{skew-symmetric} when its
transpose equals its negative:

\[
\boxed{
\mathbf A^\mathsf T
=
-\mathbf A
}.
\]

Transposition reflects the entries across the main diagonal. The
condition therefore means that entries mirrored across the diagonal have
equal magnitudes and opposite signs; every diagonal entry must be zero.
Define the dimensionless planar skew-symmetric generator

\[
\boxed{
\mathbf S_2
\mathrel{:=}
\begin{bmatrix}
0&-1\\
1&0
\end{bmatrix}
}.
\]

The matrix \(\mathbf S_2\) is skew-symmetric because
\(\mathbf S_2^\mathsf T=-\mathbf S_2\). Premultiplication by
\(\mathbf S_2\) rotates any planar coordinate column counter-clockwise
by \(\pi/2\):

\[
\mathbf S_2
\begin{bmatrix}
a\\b
\end{bmatrix}
=
\begin{bmatrix}
-b\\a
\end{bmatrix}.
\]

Direct multiplication gives

\[
\mathbf S_2\mathbf R_{wb}(\theta)
=
\mathbf R_{wb}(\theta)\mathbf S_2
=
\begin{bmatrix}
-\sin\theta&-\cos\theta\\
\cos\theta&-\sin\theta
\end{bmatrix}.
\]

Therefore,

\[
\boxed{
\dot{\mathbf R}_{wb}
=
\omega\mathbf S_2\mathbf R_{wb}
=
\mathbf R_{wb}\omega\mathbf S_2
}.
\]

The two products are equal in the planar case because every element of
\(SO(2)\) commutes with \(\mathbf S_2\). This equality is a special
simplification of planar rotation; matrix order becomes important for
three-dimensional angular velocity.

\subsection{\texorpdfstring{Why \(\dot{\mathbf R}_{wb}\) is not a
rotation
matrix}{Why \textbackslash dot\{\textbackslash mathbf R\}\_\{wb\} is not a rotation matrix}}\label{sec-se2-rotation-derivative-tangent}

Although \(\mathbf R_{wb}(t)\in SO(2)\) at every time \(t\), its
derivative \(\dot{\mathbf R}_{wb}(t)\) is generally not a rotation
matrix. A rotation matrix is a dimensionless map that represents a
finite change of orientation, whereas \(\dot{\mathbf R}_{wb}\) measures
how rapidly that map is changing and has entries with units of
\(\mathrm{s^{-1}}\). The derivative therefore describes an instantaneous
direction of rotational motion rather than another finite rotation.

\subsubsection{The rotation constraint restricts its
derivative}\label{the-rotation-constraint-restricts-its-derivative}

Every planar rotation matrix satisfies

\[
\mathbf R_{wb}^\mathsf T\mathbf R_{wb}
=
\mathbf I_2.
\]

Differentiating this constraint with respect to time gives

\[
\dot{\mathbf R}_{wb}^\mathsf T\mathbf R_{wb}
+
\mathbf R_{wb}^\mathsf T\dot{\mathbf R}_{wb}
=
\mathbf 0_{2\times2}.
\]

Rearranging shows that

\[
\left(
\mathbf R_{wb}^\mathsf T\dot{\mathbf R}_{wb}
\right)^\mathsf T
=
-\mathbf R_{wb}^\mathsf T\dot{\mathbf R}_{wb}.
\]

Thus, \(\mathbf R_{wb}^\mathsf T\dot{\mathbf R}_{wb}\) is
\textbf{skew-symmetric}: its transpose equals its negative. The factor
\(\mathbf R_{wb}^\mathsf T\) expresses the rotational rate relative to
the current body orientation; it does not turn \(\dot{\mathbf R}_{wb}\)
into a finite rotation.

The set of all \(2\times2\) real skew-symmetric matrices is

\[
\mathfrak{so}(2)
\mathrel{:=}
\left\{
\mathbf A\in\mathbb R^{2\times2}
\;\middle|\;
\mathbf A^\mathsf T=-\mathbf A
\right\}
=
\left\{
\alpha\mathbf S_2
\;\middle|\;
\alpha\in\mathbb R
\right\},
\]

where \(\alpha\) is a real scalar and

\[
\mathbf S_2
=
\begin{bmatrix}
0&-1\\
1&0
\end{bmatrix}.
\]

To justify the second equality, begin with a general matrix

\[
\mathbf A
=
\begin{bmatrix}
a&b\\
c&d
\end{bmatrix},
\qquad
a,b,c,d\in\mathbb R.
\]

The condition \(\mathbf A^\mathsf T=-\mathbf A\) requires corresponding
entries to satisfy

\[
a=-a,
\qquad
d=-d,
\qquad
c=-b.
\]

Therefore, \(a=d=0\). Setting \(\alpha=c=-b\) gives

\[
\mathbf A
=
\begin{bmatrix}
0&-\alpha\\
\alpha&0
\end{bmatrix}
=
\alpha\mathbf S_2.
\]

Thus, every \(2\times2\) skew-symmetric matrix is a scalar multiple of
\(\mathbf S_2\). Conversely, for every \(\alpha\in\mathbb R\),

\[
(\alpha\mathbf S_2)^\mathsf T
=
\alpha\mathbf S_2^\mathsf T
=
-\alpha\mathbf S_2,
\]

so every scalar multiple of \(\mathbf S_2\) is skew-symmetric. The two
sets are therefore equal. The notation \(\mathfrak{so}(2)\) is read
``little s o two.''

\subsubsection{Why these matrices form a tangent
space}\label{why-these-matrices-form-a-tangent-space}

The skew-symmetric matrix \(\mathbf S_2\) is tangent to \(SO(2)\) at the
identity rotation \(\mathbf I_2\). More generally, every scalar multiple
\(\alpha\mathbf S_2\), where \(\alpha\in\mathbb R\), is a possible
tangent direction there. This subsection explains the geometric meaning
of that statement and proves that these are all the possible tangent
directions.

A \textbf{tangent direction} to a curve at a point is the limiting
direction of secant lines drawn from that point to nearby points on the
curve. Equivalently, it is the derivative at that point of a
differentiable path that remains on the curve. This familiar definition
applies to \(SO(2)\) because every rotation matrix \(\mathbf R(\theta)\)
is determined by the pair \((\cos\theta,\sin\theta)\) on the unit
circle. Under this correspondence,

\[
\mathbf R(0)=\mathbf I_2
\quad\longleftrightarrow\quad
(\cos0,\sin0)=(1,0).
\]

For a nonzero angular increment \(h\) measured in radians, the secant
direction of the unit circle from \((1,0)\) to \((\cos h,\sin h)\) is

\[
\frac{1}{h}
\left(
\begin{bmatrix}
\cos h\\
\sin h
\end{bmatrix}
-
\begin{bmatrix}
1\\0
\end{bmatrix}
\right).
\]

As \(h\to0\), this secant direction approaches \([0\;1]^\mathsf T\), the
usual tangent direction to the circle at \((1,0)\). The corresponding
difference quotient in matrix coordinates is

\[
\frac{\mathbf R(h)-\mathbf I_2}{h}
=
\begin{bmatrix}
\dfrac{\cos h-1}{h}&-\dfrac{\sin h}{h}\\[6pt]
\dfrac{\sin h}{h}&\dfrac{\cos h-1}{h}
\end{bmatrix}.
\]

The scalar limits

\[
\lim_{h\to0}\frac{\sin h}{h}=1,
\qquad
\lim_{h\to0}\frac{\cos h-1}{h}=0
\]

give

\[
\boxed{
\left.
\frac{d\mathbf R(h)}{dh}
\right|_{h=0}
=
\lim_{h\to0}
\frac{\mathbf R(h)-\mathbf I_2}{h}
=
\mathbf S_2
}.
\]

The second scalar limit follows from \(\cos h-1=-2\sin^2(h/2)\), the
first limit, and the continuity of sine at zero. Therefore,
\(\mathbf S_2\) is tangent to the rotation curve at \(\mathbf I_2\) in
exactly the same sense that \([0\;1]^\mathsf T\) is tangent to the unit
circle at \((1,0)\).

To collect every possible tangent direction, let
\(\mathbf Q(s)\in SO(2)\) be any differentiable rotation curve defined
for the dimensionless real parameter \(s\) near zero, with
\(\mathbf Q(0)=\mathbf I_2\). The \textbf{tangent space of \(SO(2)\) at
\(\mathbf I_2\)} is

\[
T_{\mathbf I_2}SO(2)
\mathrel{:=}
\left\{
\left.
\frac{d\mathbf Q(s)}{ds}
\right|_{s=0}
\;\middle|\;
\mathbf Q(s)\in SO(2),\;
\mathbf Q(0)=\mathbf I_2
\right\}.
\]

The symbol \(T\) denotes a tangent space, and its subscript
\(\mathbf I_2\) identifies the point where the tangent directions are
attached. To prove that this space equals \(\mathfrak{so}(2)\), both set
inclusions must be established: every tangent direction must be shown to
be skew-symmetric, and every skew-symmetric matrix must be shown to
arise from a valid rotation curve.

First, take any \(\mathbf A\in T_{\mathbf I_2}SO(2)\). By definition,
\(\mathbf A=\mathbf Q'(0)\) for some rotation curve \(\mathbf Q(s)\)
through \(\mathbf I_2\). Differentiating
\(\mathbf Q(s)^\mathsf T\mathbf Q(s)=\mathbf I_2\) gives

\[
\mathbf Q'(s)^\mathsf T\mathbf Q(s)
+
\mathbf Q(s)^\mathsf T\mathbf Q'(s)
=
\mathbf 0_{2\times2}.
\]

Now evaluate at (s=0). Because the curve passes through the identity,

\[
\mathbf Q(0)=\mathbf I_2,
\qquad
\mathbf Q'(0)=\mathbf A.
\]

Substitution gives

\[
\mathbf A^\mathsf T+\mathbf A
=
\mathbf 0_{2\times2}.
\]

Thus, \(\mathbf A\) is skew-symmetric and belongs to
\(\mathfrak{so}(2)\). This proves

\[
T_{\mathbf I_2}SO(2)
\subseteq
\mathfrak{so}(2).
\]

Second, take any \(\mathbf A\in\mathfrak{so}(2)\). It has the form
\(\mathbf A=\alpha\mathbf S_2\) for some \(\alpha\in\mathbb R\). The
curve \(\mathbf Q(s)=\mathbf R(\alpha s)\) remains in \(SO(2)\), passes
through \(\mathbf I_2\) at \(s=0\), and satisfies

\[
\mathbf Q'(0)
=
\left.
\frac{d}{ds}\mathbf R(\alpha s)
\right|_{s=0}
=
\alpha\mathbf S_2
=
\mathbf A.
\]

Hence, every element of \(\mathfrak{so}(2)\) is produced by a valid
curve through the identity, which proves

\[
\mathfrak{so}(2)
\subseteq
T_{\mathbf I_2}SO(2).
\]

The two inclusions establish

\[
\boxed{
T_{\mathbf I_2}SO(2)
=
\mathfrak{so}(2)
=
\left\{
\alpha\mathbf S_2
\;\middle|\;
\alpha\in\mathbb R
\right\}
}.
\]

This proves that the skew-symmetric matrices in \(\mathfrak{so}(2)\) are
exactly the instantaneous directions in which a rotation curve can pass
through \(\mathbf I_2\). They describe rotational directions, not finite
rotations. The next subsection shows how these directions provide a
first-order approximation to nearby rotation matrices.

\subsubsection{Lie-group terminology and the local linear
model}\label{lie-group-terminology-and-the-local-linear-model}

Section~\ref{sec-se2-planar-rotations} established that \(SO(2)\) is a
group under matrix multiplication. Its matrices also vary smoothly with
the continuous angle \(\theta\), and its multiplication and inverse
operations are differentiable. A group with such a compatible smooth
structure is called a \textbf{Lie group}. The tangent space at the
identity of a Lie group, together with an operation called the Lie
bracket, is its \textbf{Lie algebra}. Thus, \(\mathfrak{so}(2)\) is the
Lie algebra of \(SO(2)\).

The Lie algebra is useful because it is linear even though the rotation
group is not. For any \(\alpha,\beta,c\in\mathbb R\),

\[
\alpha\mathbf S_2+\beta\mathbf S_2
=
(\alpha+\beta)\mathbf S_2,
\qquad
c(\alpha\mathbf S_2)
=
(c\alpha)\mathbf S_2,
\]

and both results remain in \(\mathfrak{so}(2)\). Instantaneous
rotational directions can therefore be added and scaled using ordinary
linear algebra, whereas adding or scaling finite rotation matrices does
not generally produce valid rotations.

The tangent space is also a local linear approximation of the group. For
a small angle \(h\),

\[
\boxed{
\mathbf R(h)
=
\mathbf I_2
+h\mathbf S_2
+O(h^2)
}.
\]

This result follows more directly from the first-order Taylor expansion
of the matrix-valued function \(\mathbf R(h)\) about \(h=0\):

\[
\mathbf R(h)
=
\mathbf R(0)
+h
\left.
\frac{d\mathbf R(h)}{dh}
\right|_{h=0}
+O(h^2).
\]

The definition of \(\mathbf R(h)\) and the preceding limit give

\[
\mathbf R(0)=\mathbf I_2,
\qquad
\left.
\frac{d\mathbf R(h)}{dh}
\right|_{h=0}
=
\mathbf S_2.
\]

Substituting these two values immediately gives
\(\mathbf R(h)=\mathbf I_2+h\mathbf S_2+O(h^2)\). The derivative is
written \(d\mathbf R/dh\) rather than \(\dot{\mathbf R}\) because \(h\)
is an angular argument, not time. For this matrix equation, \(O(h^2)\)
means that the norm of the Taylor remainder is bounded by \(C|h|^2\) for
some constant \(C>0\) when \(h\) is sufficiently close to zero; this
bound exists because the sine and cosine entries of \(\mathbf R(h)\)
have continuous second derivatives near \(h=0\).

If the same rotation increment is instead parameterised by time as
\(h(t)=\omega t\) with constant angular velocity \(\omega\), then the
time-domain Taylor expansion over a short interval \(\Delta t\) is

\[
\mathbf R(\omega\Delta t)
=
\mathbf R(0)
+\Delta t\,\dot{\mathbf R}(0)
+O(\Delta t^2)
=
\mathbf I_2
+\omega\Delta t\,\mathbf S_2
+O(\Delta t^2),
\]

because \(\dot{\mathbf R}(0)=\omega\mathbf S_2\). This is the
time-parameterised version of the same first-order approximation.

In the angle-domain approximation, \(\mathbf I_2+h\mathbf S_2\) starts
at \(\mathbf I_2\) and changes by \(h\mathbf S_2\) along the tangent
direction. It is generally not an exact rotation matrix, but its
difference from the true rotation is only of order \(h^2\). The
approximation therefore matches both the value and the first derivative
of the rotation curve at \(h=0\).

\subsubsection{Returning to the physical rotation
trajectory}\label{returning-to-the-physical-rotation-trajectory}

For the physical trajectory \(\mathbf R_{wb}(t)\),
Section~\ref{sec-se2-rotation-matrix-derivative} established

\[
\dot{\mathbf R}_{wb}
=
\omega\mathbf R_{wb}\mathbf S_2.
\]

Thus, \(\dot{\mathbf R}_{wb}\) is tangent to \(SO(2)\) at the current
rotation \(\mathbf R_{wb}\), rather than at the identity.
Premultiplication by \(\mathbf R_{wb}^\mathsf T\) carries this direction
back to the tangent space at the identity:

\[
\boxed{
\boldsymbol\Omega_b
\mathrel{:=}
\mathbf R_{wb}^\mathsf T\dot{\mathbf R}_{wb}
=
\omega\mathbf S_2
\in
\mathfrak{so}(2)
}.
\]

Here, \(\boldsymbol\Omega_b\) is the body-frame angular-velocity matrix,
and its entries have units of \(\mathrm{s^{-1}}\). The time-domain
Taylor expansion above connects this instantaneous generator
\(\omega\mathbf S_2\) to a nearby finite rotation. The derivative is not
itself a rotation matrix; it specifies the tangent direction from which
the finite rotation develops over time.

Finally, the \textbf{Lie bracket} measures whether two infinitesimal
matrix motions commute. For matrices it is the commutator

\[
[\mathbf A,\mathbf B]
\mathrel{:=}
\mathbf A\mathbf B-\mathbf B\mathbf A,
\qquad
\mathbf A,\mathbf B\in\mathfrak{so}(2).
\]

For \(\mathfrak{so}(2)\), write \(\mathbf A=\alpha\mathbf S_2\) and
\(\mathbf B=\beta\mathbf S_2\). Then
\(\mathbf A\mathbf B=\mathbf B\mathbf A\), so
\([\mathbf A,\mathbf B]=\mathbf 0_{2\times2}\). The bracket is trivial
for planar rotations because they commute, but it becomes important for
rigid motions in \(SE(2)\) and rotations in \(SO(3)\), where motion
order generally matters.

\subsection{Velocity of a point fixed to the
body}\label{sec-se2-body-point-velocity}

Let \(P\) be a point rigidly fixed to the body. Its body-frame
coordinates \({}^{b}\mathbf p_P\in\mathbb R^2\) are constant, so

\[
\frac{d}{dt}{}^{b}\mathbf p_P
=
\mathbf 0_2.
\]

Its world-frame position is

\[
{}^{w}\mathbf p_P(t)
=
\mathbf R_{wb}(t){}^{b}\mathbf p_P
+
{}^{w}\mathbf p_b(t).
\]

Differentiate both sides:

\[
\begin{aligned}
{}^{w}\mathbf v_P
\mathrel{:=}
\frac{d}{dt}{}^{w}\mathbf p_P
&=
\dot{\mathbf R}_{wb}{}^{b}\mathbf p_P
+
\mathbf R_{wb}\frac{d}{dt}{}^{b}\mathbf p_P
+
{}^{w}\mathbf v_b\\
&=
\dot{\mathbf R}_{wb}{}^{b}\mathbf p_P
+
{}^{w}\mathbf v_b.
\end{aligned}
\]

Define the relative-position vector from \(O_b\) to \(P\), expressed in
world-frame coordinates, by

\[
{}^{w}\mathbf r_{O_bP}
\mathrel{:=}
{}^{w}\mathbf p_P-{}^{w}\mathbf p_b
=
\mathbf R_{wb}{}^{b}\mathbf p_P
\in\mathbb R^2.
\]

Using \(\dot{\mathbf R}_{wb}=\omega\mathbf S_2\mathbf R_{wb}\) gives

\[
\boxed{
{}^{w}\mathbf v_P
=
{}^{w}\mathbf v_b
+
\omega\mathbf S_2{}^{w}\mathbf r_{O_bP}
}.
\]

The first term is the translational velocity shared with the body
origin. The second term is the additional velocity caused by rotation
about the body origin. Because \(\mathbf S_2\) rotates
\({}^{w}\mathbf r_{O_bP}\) by \(\pi/2\), the rotational contribution is
perpendicular to the relative-position vector.

The two panels in Figure~\ref{fig-planar-rigid-body-point-velocity}
separate the position-vector identity from the velocity-vector addition.
In the left panel, the two solid vectors add head to tail to produce the
dashed world-position vector of \(P\). In the right panel, translated
copies of the two velocity contributions form a parallelogram whose
diagonal is \({}^{w}\mathbf v_P\); the right-angle marker applies
specifically to the rotational contribution and
\({}^{w}\mathbf r_{O_bP}\).

\begin{figure}[H]

\centering{

\includegraphics[width=1\linewidth,height=\textheight,keepaspectratio]{index_files/mediabag/figures/planar_rigid_body_point_velocity.pdf}

}

\caption{\label{fig-planar-rigid-body-point-velocity}Position and
velocity relations for a point fixed to a planar rigid body.}

\end{figure}%

The same result can be written with the three-dimensional cross product
by embedding the planar quantities in \(\mathbb R^3\):

\[
{}^{w}\boldsymbol\omega
=
\begin{bmatrix}
0\\0\\\omega
\end{bmatrix},
\qquad
{}^{w}\mathbf r_{O_bP}^{(3)}
=
\begin{bmatrix}
r_x\\r_y\\0
\end{bmatrix}.
\]

Here, \(r_x\) and \(r_y\), measured in metres, are the two entries of
\({}^{w}\mathbf r_{O_bP}\), the superscript \((3)\) marks its
three-dimensional embedding, and \(\times\) denotes the vector cross
product.

Then

\[
{}^{w}\boldsymbol\omega
\times
{}^{w}\mathbf r_{O_bP}^{(3)}
=
\begin{bmatrix}
-\omega r_y\\
\omega r_x\\
0
\end{bmatrix},
\]

whose first two entries equal
\(\omega\mathbf S_2{}^{w}\mathbf r_{O_bP}\). Thus, the planar formula is
the two-dimensional form of the rigid-body velocity relation

\[
\mathbf v_P
=
\mathbf v_{O_b}
+
\boldsymbol\omega\times\mathbf r_{O_bP}.
\]

\subsection{Worked example}\label{worked-example-1}

At one instant, suppose the body heading is
\(\theta=\pi/2\,\mathrm{rad}\), the body-origin velocity is

\[
{}^{w}\mathbf v_b
=
\begin{bmatrix}
0.5\\0
\end{bmatrix}
\mathrm{m\,s^{-1}},
\]

and the yaw rate is \(\omega=2\,\mathrm{rad\,s^{-1}}\). Point \(P\) is
fixed one metre along the body \(+x_b\)-axis:

\[
{}^{b}\mathbf p_P
=
\begin{bmatrix}
1\\0
\end{bmatrix}
\mathrm m.
\]

At \(\theta=\pi/2\), its relative-position vector in the world frame is

\[
{}^{w}\mathbf r_{O_bP}
=
\mathbf R_{wb}\left(\frac{\pi}{2}\right){}^{b}\mathbf p_P
=
\begin{bmatrix}
0\\1
\end{bmatrix}
\mathrm m.
\]

The rotational contribution is

\[
\omega\mathbf S_2{}^{w}\mathbf r_{O_bP}
=
2
\begin{bmatrix}
0&-1\\
1&0
\end{bmatrix}
\begin{bmatrix}
0\\1
\end{bmatrix}
=
\begin{bmatrix}
-2\\0
\end{bmatrix}
\mathrm{m\,s^{-1}}.
\]

Therefore,

\[
\begin{aligned}
{}^{w}\mathbf v_P
&=
{}^{w}\mathbf v_b
+
\omega\mathbf S_2{}^{w}\mathbf r_{O_bP}\\
&=
\begin{bmatrix}
0.5\\0
\end{bmatrix}
+
\begin{bmatrix}
-2\\0
\end{bmatrix}\\
&=
\begin{bmatrix}
-1.5\\0
\end{bmatrix}
\mathrm{m\,s^{-1}}.
\end{aligned}
\]

The point lies one metre above the body origin in world coordinates.
Counter-clockwise rotation therefore moves it instantaneously toward
\(-x_w\) at \(2\,\mathrm{m\,s^{-1}}\), while translation of the body
origin contributes \(0.5\,\mathrm{m\,s^{-1}}\) toward \(+x_w\).

\subsection{Assumptions, failure modes, and
checks}\label{assumptions-failure-modes-and-checks}

\begin{itemize}
\tightlist
\item
  \textbf{Body-fixed point assumption:} The derivation sets
  \(d({}^{b}\mathbf p_P)/dt=\mathbf 0_2\). If \(P\) moves relative to
  the body, the additional term
  \(\mathbf R_{wb}\,d({}^{b}\mathbf p_P)/dt\) must be retained.
\item
  \textbf{Frame consistency:} \({}^{w}\mathbf v_b\) and
  \(\omega\mathbf S_2{}^{w}\mathbf r_{O_bP}\) must be expressed in the
  same frame before they are added.
\item
  \textbf{Yaw sign:} With the stated convention, \(\omega>0\) produces
  counter-clockwise rotation. A sign error reverses the rotational
  velocity of every off-origin point.
\item
  \textbf{Derivative versus rotation:} \(\dot{\mathbf R}_{wb}\) is not a
  rotation matrix. Code must not validate or normalise it as though it
  belonged to \(SO(2)\).
\item
  \textbf{Numerical differentiation:} Estimating \(\dot{\mathbf R}\) or
  velocity from sampled poses amplifies measurement noise and depends on
  an accurate sampling interval.
\end{itemize}

\section*{Quick test C: rotation derivatives and point
velocity}\label{quick-test-c-rotation-derivatives-and-point-velocity}
\addcontentsline{toc}{section}{Quick test C: rotation derivatives and
point velocity}

\markright{Quick test C: rotation derivatives and point velocity}

\begin{enumerate}
\def\labelenumi{\arabic{enumi}.}
\tightlist
\item
  Define \({}^{w}\mathbf v_b\) and \(\omega\), including their frames
  and units.
\item
  Differentiate \(\mathbf R_{wb}(\theta)\) and show that
  \(\dot{\mathbf R}_{wb}=\omega\mathbf S_2\mathbf R_{wb}\).
\item
  What geometric operation does \(\mathbf S_2\) perform on a planar
  vector?
\item
  Why is \(\dot{\mathbf R}_{wb}\) generally not an element of \(SO(2)\)?
\item
  Differentiate \(\mathbf R_{wb}^\mathsf T\mathbf R_{wb}=\mathbf I_2\)
  and use the result to show that
  \(\mathbf R_{wb}^\mathsf T\dot{\mathbf R}_{wb}\) is skew-symmetric.
\item
  Starting from
  \({}^{w}\mathbf p_P=\mathbf R_{wb}{}^{b}\mathbf p_P+{}^{w}\mathbf p_b\),
  derive the velocity of a point fixed to the body.
\item
  Why is the rotational contribution to point velocity perpendicular to
  \({}^{w}\mathbf r_{O_bP}\)?
\item
  How does the velocity equation change if \(P\) moves relative to frame
  \(\{b\}\)?
\end{enumerate}

\subsection*{Answers and hints}\label{answers-and-hints-12}
\addcontentsline{toc}{subsection}{Answers and hints}

\begin{enumerate}
\def\labelenumi{\arabic{enumi}.}
\tightlist
\item
  \({}^{w}\mathbf v_b=d({}^{w}\mathbf p_b)/dt=[\dot x\;\dot y]^\mathsf T\)
  is the velocity of the body origin relative to the fixed world frame
  and expressed in world coordinates, measured in
  \(\mathrm{m\,s^{-1}}\). The scalar \(\omega=\dot\theta\) is the signed
  yaw rate, measured in \(\mathrm{rad\,s^{-1}}\).
\item
  Entrywise differentiation and the chain rule give the matrix displayed
  in Section~\ref{sec-se2-rotation-matrix-derivative}. Direct
  multiplication shows that its coefficient of \(\omega\) equals
  \(\mathbf S_2\mathbf R_{wb}\), so
  \(\dot{\mathbf R}_{wb}=\omega\mathbf S_2\mathbf R_{wb}\).
\item
  \(\mathbf S_2[a\;b]^\mathsf T=[-b\;a]^\mathsf T\), so it rotates the
  coordinate column counter-clockwise by \(\pi/2\) without changing its
  length.
\item
  Its entries have units of \(\mathrm{s^{-1}}\), and it does not
  generally satisfy
  \(\dot{\mathbf R}_{wb}^\mathsf T\dot{\mathbf R}_{wb}=\mathbf I_2\) or
  have determinant \(1\).
\item
  Differentiation gives
  \(\dot{\mathbf R}^\mathsf T\mathbf R+\mathbf R^\mathsf T\dot{\mathbf R}=\mathbf 0\).
  Therefore,
  \((\mathbf R^\mathsf T\dot{\mathbf R})^\mathsf T=\dot{\mathbf R}^\mathsf T\mathbf R=-\mathbf R^\mathsf T\dot{\mathbf R}\).
\item
  Since \(P\) is body-fixed, \(d({}^{b}\mathbf p_P)/dt=\mathbf 0_2\).
  Therefore,
  \({}^{w}\mathbf v_P={}^{w}\mathbf v_b+\dot{\mathbf R}_{wb}{}^{b}\mathbf p_P={}^{w}\mathbf v_b+\omega\mathbf S_2{}^{w}\mathbf r_{O_bP}\).
\item
  \(\mathbf S_2\) rotates the relative-position vector by \(\pi/2\), so
  their dot product is zero. This is the tangential direction of
  instantaneous circular motion about \(O_b\).
\item
  Retain the relative-motion term:
  \({}^{w}\mathbf v_P={}^{w}\mathbf v_b+\dot{\mathbf R}_{wb}{}^{b}\mathbf p_P+\mathbf R_{wb}\,d({}^{b}\mathbf p_P)/dt\).
\end{enumerate}

\section{\texorpdfstring{Planar twists and
\(\mathfrak{se}(2)\)}{Planar twists and \textbackslash mathfrak\{se\}(2)}}\label{planar-twists-and-mathfrakse2}

The previous section described translational and angular velocity
separately. A \textbf{planar twist} combines them into one
representation of the instantaneous motion of a rigid body. This
representation will connect the pose derivative \(\dot{\mathbf T}_{wb}\)
to the differential-drive inputs \(v\) and \(\omega\).

\subsection{Twist coordinates}\label{twist-coordinates}

In this chapter, a planar twist is represented by the coordinate column

\[
\boxed{
\boldsymbol\xi
\mathrel{:=}
\begin{bmatrix}
\nu_x\\
\nu_y\\
\omega
\end{bmatrix}
\in\mathbb R^3
}.
\]

Here, \(\nu_x\) and \(\nu_y\) use the Greek letter \emph{nu}, \(\nu\),
not the Latin letter \(v\). They are \textbf{generic placeholders} for
the \(x\)- and \(y\)-components of the twist's linear part. Because
neither the coordinate frame nor the point chosen as the velocity
reference has been specified yet, these symbols do not yet identify a
unique physical velocity. Together they form the column
\(\boldsymbol\nu=[\nu_x\;\nu_y]^\mathsf T\in\mathbb R^2\), whose entries
are measured in \(\mathrm{m\,s^{-1}}\). The body and spatial cases are
defined in Section~\ref{sec-se2-body-twist} and
Section~\ref{sec-se2-spatial-twist}.

The third entry \(\omega\in\mathbb R\) is angular velocity, measured in
\(\mathrm{rad\,s^{-1}}\). The symbol \(\boldsymbol\xi\) is the Greek
letter \emph{xi}. The word \emph{twist} refers to an instantaneous
rigid-body velocity, not to a pose or a finite displacement.

This chapter uses the component order
\([\nu_x\;\nu_y\;\omega]^\mathsf T\). Some textbooks and software
libraries place angular velocity first, so the component order must
always be declared. Although \(\boldsymbol\xi\) is mathematically a
three-entry real column, its entries have different physical units. A
norm such as \(\lVert\boldsymbol\xi\rVert_2\) therefore has no direct
physical meaning unless a length scale is introduced to compare
translation with rotation.

\subsection{\texorpdfstring{The hat operator and
\(\mathfrak{se}(2)\)}{The hat operator and \textbackslash mathfrak\{se\}(2)}}\label{sec-se2-hat-operator}

Section~\ref{sec-se2-rotation-derivative-tangent} showed that
\(\mathfrak{so}(2)\) contains the tangent directions of the rotation
group \(SO(2)\) at \(\mathbf I_2\). A planar rigid motion also includes
translation, so its instantaneous direction must combine one rotational
component with two translational components.

\subsubsection{From twist coordinates to a
matrix}\label{from-twist-coordinates-to-a-matrix}

The \textbf{hat operator} maps the twist coordinate column to a
\(3\times3\) matrix:

\[
\boxed{
\widehat{\boldsymbol\xi}
\mathrel{:=}
\begin{bmatrix}
0&-\omega&\nu_x\\
\omega&0&\nu_y\\
0&0&0
\end{bmatrix}
=
\begin{bmatrix}
\omega\mathbf S_2&\boldsymbol\nu\\
\mathbf 0_2^\mathsf T&0
\end{bmatrix}
}.
\]

The wide hat in \(\widehat{\boldsymbol\xi}\) denotes this mapping; it
does not denote an estimate. The inverse operation, which recovers
\(\boldsymbol\xi\) from its matrix representation, is commonly called
the \textbf{vee operator} and is written
\((\widehat{\boldsymbol\xi})^\vee=\boldsymbol\xi\).

The upper-left block \(\omega\mathbf S_2\) represents the instantaneous
rotation direction. The upper-right column \(\boldsymbol\nu\) represents
the two-component linear part, and the bottom row is zero because
differentiating the constant bottom row \([0\;0\;1]\) of a homogeneous
transformation gives \([0\;0\;0]\).

\subsubsection{\texorpdfstring{Why these matrices are tangent to
\(SE(2)\)}{Why these matrices are tangent to SE(2)}}\label{why-these-matrices-are-tangent-to-se2}

Recall that an element of \(SE(2)\) is a finite planar homogeneous
transformation

\[
\mathbf G
=
\begin{bmatrix}
\mathbf R&\mathbf p\\
\mathbf 0_2^\mathsf T&1
\end{bmatrix},
\qquad
\mathbf R\in SO(2),
\quad
\mathbf p\in\mathbb R^2.
\]

It rotates and translates planar point coordinates. The \textbf{identity
transformation} is

\[
\mathbf I_3
=
\begin{bmatrix}
\mathbf I_2&\mathbf 0_2\\
\mathbf 0_2^\mathsf T&1
\end{bmatrix}.
\]

It represents zero rotation angle and zero translation, so it leaves
every homogeneous point coordinate
\(\overline{\mathbf p}_P=[\mathbf p_P^\mathsf T\;1]^\mathsf T\in\mathbb R^3\)
unchanged:

\[
\mathbf I_3\overline{\mathbf p}_P
=
\overline{\mathbf p}_P.
\]

Let \(\mathbf G(s)\in SE(2)\) be a differentiable transformation curve
with dimensionless real parameter \(s\) near zero and
\(\mathbf G(0)=\mathbf I_3\). Its derivative \(\mathbf G'(0)\) is a
tangent direction to \(SE(2)\) at \(\mathbf I_3\). The tangent space is
the set of all such derivatives:

\[
T_{\mathbf I_3}SE(2)
\mathrel{:=}
\left\{
\mathbf G'(0)
\;\middle|\;
\mathbf G(s)\in SE(2),\;
\mathbf G(0)=\mathbf I_3
\right\}.
\]

Now define the candidate set

\[
\mathcal C
\mathrel{:=}
\left\{
\begin{bmatrix}
a\mathbf S_2&\mathbf u\\
\mathbf 0_2^\mathsf T&0
\end{bmatrix}
\;\middle|\;
a\in\mathbb R,\;
\mathbf u\in\mathbb R^2
\right\},
\]

where \(\mathcal C\) is merely a temporary name for the proposed
tangent-matrix structure. The goal is to prove

\[
T_{\mathbf I_3}SE(2)
=
\mathcal C.
\]

Equality requires two set inclusions. The first shows that no other
tangent-matrix form is possible. The second shows that every matrix in
the proposed set is actually attainable.

\textbf{First inclusion: every tangent direction has the proposed form.}
Take any \(\mathbf A\in T_{\mathbf I_3}SE(2)\). By the definition of the
tangent space, there is a differentiable curve

\[
\mathbf G(s)
=
\begin{bmatrix}
\mathbf R(s)&\mathbf p(s)\\
\mathbf 0_2^\mathsf T&1
\end{bmatrix}
\in SE(2)
\]

such that \(\mathbf G(0)=\mathbf I_3\) and \(\mathbf A=\mathbf G'(0)\).
The identity condition implies \(\mathbf R(0)=\mathbf I_2\) and
\(\mathbf p(0)=\mathbf 0_2\). Differentiating the block matrix gives

\[
\mathbf A
=
\mathbf G'(0)
=
\begin{bmatrix}
\mathbf R'(0)&\mathbf p'(0)\\
\mathbf 0_2^\mathsf T&0
\end{bmatrix}.
\]

Because \(\mathbf R(s)\in SO(2)\), it satisfies
\(\mathbf R(s)^\mathsf T\mathbf R(s)=\mathbf I_2\). Differentiating this
constraint with respect to \(s\) and evaluating at \(s=0\) gives

\[
\mathbf R'(0)^\mathsf T\mathbf R(0)
+
\mathbf R(0)^\mathsf T\mathbf R'(0)
=
\mathbf 0_{2\times2}.
\]

Using \(\mathbf R(0)=\mathbf I_2\) reduces this equation to

\[
\mathbf R'(0)^\mathsf T
+
\mathbf R'(0)
=
\mathbf 0_{2\times2}.
\]

Thus, \(\mathbf R'(0)\) is skew-symmetric and belongs to
\(\mathfrak{so}(2)\). Section~\ref{sec-se2-rotation-derivative-tangent}
established that every matrix in \(\mathfrak{so}(2)\) has the form
\(a\mathbf S_2\) for some \(a\in\mathbb R\), so
\(\mathbf R'(0)=a\mathbf S_2\). Define
\(\mathbf u\mathrel{:=}\mathbf p'(0)\in\mathbb R^2\). Substitution gives

\[
\mathbf A
=
\begin{bmatrix}
a\mathbf S_2&\mathbf u\\
\mathbf 0_2^\mathsf T&0
\end{bmatrix}.
\]

Therefore, \(\mathbf A\in\mathcal C\). Because \(\mathbf A\) was an
arbitrary tangent direction, this proves

\[
T_{\mathbf I_3}SE(2)
\subseteq
\mathcal C.
\]

\textbf{Second inclusion: every matrix of the proposed form is
attainable.} Take any \(\mathbf A\in\mathcal C\). By the definition of
\(\mathcal C\), there are values \(a\in\mathbb R\) and
\(\mathbf u\in\mathbb R^2\) such that

\[
\mathbf A
=
\begin{bmatrix}
a\mathbf S_2&\mathbf u\\
\mathbf 0_2^\mathsf T&0
\end{bmatrix}.
\]

Construct the curve

\[
\mathbf G_{a,\mathbf u}(s)
\mathrel{:=}
\begin{bmatrix}
\mathbf R(as)&s\mathbf u\\
\mathbf 0_2^\mathsf T&1
\end{bmatrix}.
\]

For every \(s\), \(\mathbf R(as)\in SO(2)\) and
\(s\mathbf u\in\mathbb R^2\), so
\(\mathbf G_{a,\mathbf u}(s)\in SE(2)\). At \(s=0\),

\[
\mathbf G_{a,\mathbf u}(0)
=
\begin{bmatrix}
\mathbf R(0)&\mathbf 0_2\\
\mathbf 0_2^\mathsf T&1
\end{bmatrix}
=
\mathbf I_3.
\]

The chain rule and
\(d\mathbf R(\theta)/d\theta|_{\theta=0}=\mathbf S_2\) give its
derivative at zero:

\[
\begin{aligned}
\mathbf G_{a,\mathbf u}'(0)
&=
\begin{bmatrix}
\left.\dfrac{d}{ds}\mathbf R(as)\right|_{s=0}&
\left.\dfrac{d}{ds}(s\mathbf u)\right|_{s=0}\\[6pt]
\mathbf 0_2^\mathsf T&0
\end{bmatrix}\\
&=
\begin{bmatrix}
a\mathbf S_2&\mathbf u\\
\mathbf 0_2^\mathsf T&0
\end{bmatrix}
=
\mathbf A.
\end{aligned}
\]

The curve is valid, passes through \(\mathbf I_3\), and has the
prescribed derivative \(\mathbf A\). Hence,
\(\mathbf A\in T_{\mathbf I_3}SE(2)\). Because \(\mathbf A\) was an
arbitrary member of \(\mathcal C\), this proves

\[
\mathcal C
\subseteq
T_{\mathbf I_3}SE(2).
\]

\textbf{Conclusion.} The two inclusions prove
\(T_{\mathbf I_3}SE(2)=\mathcal C\). This tangent space is denoted
\(\mathfrak{se}(2)\):

\[
\boxed{
\mathfrak{se}(2)
\mathrel{:=}
T_{\mathbf I_3}SE(2)
=
\mathcal C
=
\left\{
\begin{bmatrix}
a\mathbf S_2&\mathbf u\\
\mathbf 0_2^\mathsf T&0
\end{bmatrix}
\;\middle|\;
a\in\mathbb R,\;
\mathbf u\in\mathbb R^2
\right\}
}.
\]

The notation \(\mathfrak{se}(2)\) is read ``little s e two.'' Because
\(SE(2)\) is a Lie group, its tangent space at the identity is called
its \textbf{Lie algebra}.

This vector space is three-dimensional because each matrix is determined
by three independent real numbers: one scalar \(a\) for rotation and the
two entries of \(\mathbf u\) for translation. Elements of \(SE(2)\)
represent finite poses or finite rigid transformations. Elements of
\(\mathfrak{se}(2)\) represent directions in which such transformations
can change instantaneously from the identity.

\subsubsection{Physical interpretation}\label{physical-interpretation}

When the curve parameter is physical time, the rotational coefficient
becomes angular velocity \(a=\omega\), and the translation derivative
becomes a linear-velocity column \(\mathbf u=\boldsymbol\nu\). The
corresponding tangent matrix is exactly the hatted twist
\(\widehat{\boldsymbol\xi}\).

When \(\widehat{\boldsymbol\xi}\) represents physical velocity, the
entries of its \(2\times2\) block have units of \(\mathrm{s^{-1}}\) and
the entries of its right column have units of \(\mathrm{m\,s^{-1}}\).
The matrix is therefore a structured representation of a physical rate,
not a finite homogeneous transformation.

For a sufficiently short time interval \(\Delta t>0\), a smooth
transformation curve through the identity with
\(\mathbf G'(0)=\widehat{\boldsymbol\xi}\) has the first-order expansion

\[
\mathbf G(\Delta t)
=
\mathbf I_3
+\Delta t\,\widehat{\boldsymbol\xi}
+O(\Delta t^2).
\]

This equation explains the term \textbf{infinitesimal rigid motion}:
\(\widehat{\boldsymbol\xi}\) is not itself a finite transformation in
\(SE(2)\), but it gives the first-order direction from which a nearby
finite transformation develops.

\subsection{Body twist}\label{sec-se2-body-twist}

\subsubsection{Plain-language meaning of pose and
twist}\label{plain-language-meaning-of-pose-and-twist}

A \textbf{pose} describes where a rigid body is and which direction it
is facing relative to a chosen reference frame. In planar motion, the
pose contains two position coordinates and one orientation angle. The
homogeneous transformation \(\mathbf T_{wb}\) represents the pose of
body frame \(\{b\}\) relative to world frame \(\{w\}\).

A \textbf{twist} describes how a rigid body's pose is changing at one
instant. It combines linear velocity and angular velocity in one
coordinate column. The word \emph{twist} does not mean that the rigid
body bends or deforms; the body remains rigid. In this planar chapter, a
twist has two linear-velocity components and one angular-velocity
component.

A \textbf{body twist} is the motion of the body relative to the world,
expressed using the body's instantaneous axes. The phrase \emph{relative
to the world} identifies the reference used to measure the motion,
whereas \emph{expressed using the body axes} identifies the directions
used to report its numerical components. The body frame moves through
the world with the rigid body, but it remains fixed relative to the
body.

For compactness in this derivation, write

\[
\mathbf T
\mathrel{:=}
\mathbf T_{wb}
=
\begin{bmatrix}
\mathbf R&\mathbf p\\
\mathbf 0_2^\mathsf T&1
\end{bmatrix},
\qquad
\dot{\mathbf T}
=
\begin{bmatrix}
\dot{\mathbf R}&\dot{\mathbf p}\\
\mathbf 0_2^\mathsf T&0
\end{bmatrix},
\]

where \(\mathbf R=\mathbf R_{wb}\in SO(2)\) and
\(\mathbf p={}^w\mathbf p_b\in\mathbb R^2\). The entries of
\(\mathbf p\) are measured in metres, and
\(\dot{\mathbf p}={}^w\mathbf v_b\) is the world-frame velocity of the
body origin.

Premultiply \(\dot{\mathbf T}\) by \(\mathbf T^{-1}\):

\[
\begin{aligned}
\mathbf T^{-1}\dot{\mathbf T}
&=
\begin{bmatrix}
\mathbf R^\mathsf T&-\mathbf R^\mathsf T\mathbf p\\
\mathbf 0_2^\mathsf T&1
\end{bmatrix}
\begin{bmatrix}
\dot{\mathbf R}&\dot{\mathbf p}\\
\mathbf 0_2^\mathsf T&0
\end{bmatrix}\\
&=
\begin{bmatrix}
\mathbf R^\mathsf T\dot{\mathbf R}&
\mathbf R^\mathsf T\dot{\mathbf p}\\
\mathbf 0_2^\mathsf T&0
\end{bmatrix}\\
&=
\begin{bmatrix}
\omega\mathbf S_2&\boldsymbol\nu_b\\
\mathbf 0_2^\mathsf T&0
\end{bmatrix}.
\end{aligned}
\]

Here,

\[
\boldsymbol\nu_b
\mathrel{:=}
\mathbf R_{wb}^\mathsf T{}^w\mathbf v_b
=
\begin{bmatrix}
\nu_{x_b}\\
\nu_{y_b}
\end{bmatrix}
\]

is the velocity of the body origin expressed along the body axes \(x_b\)
and \(y_b\). Its entries are measured in \(\mathrm{m\,s^{-1}}\). The
\textbf{body twist} and its matrix representation are therefore

\[
\boxed{
\boldsymbol\xi_b
\mathrel{:=}
\begin{bmatrix}
\nu_{x_b}\\
\nu_{y_b}\\
\omega
\end{bmatrix},
\qquad
\widehat{\boldsymbol\xi}_b
=
\mathbf T_{wb}^{-1}\dot{\mathbf T}_{wb}
}.
\]

Equivalently,

\[
\boxed{
\dot{\mathbf T}_{wb}
=
\mathbf T_{wb}\widehat{\boldsymbol\xi}_b
}.
\]

The subscript \(b\) states that the twist is resolved in the
instantaneous body frame. In planar motion, \(\omega\) has the same
scalar value in the world and body descriptions because both frames
share the same rotation axis; this simplification does not extend to the
three-component angular-velocity coordinates used in spatial motion.

The coordinate conversion

\[
\boldsymbol\nu_b
=
\mathbf R_{wb}^{\mathsf T}{}^w\mathbf v_b
\]

does not create a different physical velocity. The column
\({}^w\mathbf v_b=[\dot x\;\dot y]^\mathsf T\) describes the body
origin's velocity \emph{relative to the world} using the world axes.
Multiplication by \(\mathbf R_{wb}^{\mathsf T}\) \emph{expresses} that
same velocity using the body's forward and sideways axes:

\[
\boldsymbol\nu_b
=
\begin{bmatrix}
\nu_{x_b}\\
\nu_{y_b}
\end{bmatrix}.
\]

Thus, \(\nu_{x_b}\) is the forward component and \(\nu_{y_b}\) is the
sideways component of the world-relative velocity. For example, a robot
facing north and moving north at \(1\,\mathrm{m\,s^{-1}}\) has
world-frame velocity components \([0\;1]^\mathsf T\,\mathrm{m\,s^{-1}}\)
but body-frame components \([1\;0]^\mathsf T\,\mathrm{m\,s^{-1}}\).
These two columns describe the same physical motion.

The body origin does not move relative to its own attached frame. Its
body-frame position is always \({}^b\mathbf p_b=\mathbf 0_2\), so the
derivative of that relative position is also zero. \emph{This zero
relative velocity is different from \(\boldsymbol\nu_b\), which is the
body's velocity relative to the world merely expressed using body axes.}

Nor does the coordinate conversion imply that the body directly feels or
measures constant linear velocity. A robot estimates its motion relative
to the ground or another external reference using information such as
wheel encoders, cameras, lidar, satellite positioning, or sensor fusion.
The body axes then provide convenient forward and sideways directions in
which to express that estimate.

\subsubsection{Physical meaning and use}\label{physical-meaning-and-use}

The two boxed equations describe the same relationship in opposite
directions. The equation

\[
\widehat{\boldsymbol\xi}_b
=
\mathbf T_{wb}^{-1}\dot{\mathbf T}_{wb}
\]

\textbf{extracts motion from a pose trajectory}: if
\(\mathbf T_{wb}(t)\) and its derivative are known, it returns the
instantaneous angular and linear velocity resolved along the moving body
axes. The inverse rotation inside \(\mathbf T_{wb}^{-1}\) removes the
body's current world orientation from the pose rate.

The equation

\[
\dot{\mathbf T}_{wb}
=
\mathbf T_{wb}\widehat{\boldsymbol\xi}_b
\]

\textbf{predicts how the pose changes}: given the current pose and a
body-frame twist, it produces the corresponding world-frame pose rate.
This is the form commonly used to propagate mobile-robot odometry,
simulate motion, and predict a future pose from a commanded or estimated
body velocity.

For an ideal differential-drive robot, the no-sideways-motion constraint
gives \(\nu_{y_b}=0\). The wheel-to-body mapping from
Section~\ref{sec-wheel-speeds-body-motion} therefore supplies

\[
\boxed{
\boldsymbol\xi_b
=
\begin{bmatrix}
v\\
0\\
\omega
\end{bmatrix}
},
\qquad
v
=
\frac{r}{2}
\left(
\dot\phi_L+\dot\phi_R
\right),
\qquad
\omega
=
\frac{r}{L}
\left(
\dot\phi_R-\dot\phi_L
\right).
\]

Here, \(r>0\) is the wheel radius in metres, \(L>0\) is the track width
in metres, and \(\dot\phi_L\) and \(\dot\phi_R\) are the left and right
wheel angular speeds in \(\mathrm{rad\,s^{-1}}\). The resulting
engineering flow is

\[
(\dot\phi_L,\dot\phi_R)
\longrightarrow
\boldsymbol\xi_b
\longrightarrow
\dot{\mathbf T}_{wb}
\longrightarrow
\mathbf T_{wb}(t+\Delta t).
\]

Wheel encoders measure wheel rotation, not the true body twist.
Inferring \(\boldsymbol\xi_b\) from them relies on the wheel geometry
and ideal rolling assumptions; slip, deformation, unequal effective
radii, and encoder errors make the inferred twist differ from the
physical motion.

\subsection{Spatial twist expressed in the world
frame}\label{sec-se2-spatial-twist}

Postmultiplying \(\dot{\mathbf T}\) by \(\mathbf T^{-1}\) produces a
different element of \(\mathfrak{se}(2)\):

\[
\begin{aligned}
\dot{\mathbf T}\mathbf T^{-1}
&=
\begin{bmatrix}
\dot{\mathbf R}&\dot{\mathbf p}\\
\mathbf 0_2^\mathsf T&0
\end{bmatrix}
\begin{bmatrix}
\mathbf R^\mathsf T&-\mathbf R^\mathsf T\mathbf p\\
\mathbf 0_2^\mathsf T&1
\end{bmatrix}\\
&=
\begin{bmatrix}
\dot{\mathbf R}\mathbf R^\mathsf T&
\dot{\mathbf p}-\dot{\mathbf R}\mathbf R^\mathsf T\mathbf p\\
\mathbf 0_2^\mathsf T&0
\end{bmatrix}\\
&=
\begin{bmatrix}
\omega\mathbf S_2&\boldsymbol\nu_w\\
\mathbf 0_2^\mathsf T&0
\end{bmatrix},
\end{aligned}
\]

where

\[
\boxed{
\boldsymbol\nu_w
\mathrel{:=}
\dot{\mathbf p}
-
\omega\mathbf S_2\mathbf p
}.
\]

Write
\(\boldsymbol\nu_w=[\nu_{x_w}\;\nu_{y_w}]^\mathsf T\in\mathbb R^2\),
where \(\nu_{x_w}\) and \(\nu_{y_w}\) are its components along the world
axes and are measured in \(\mathrm{m\,s^{-1}}\).

The \textbf{spatial twist expressed in the world frame} is

\[
\boxed{
\boldsymbol\xi_w
\mathrel{:=}
\begin{bmatrix}
\nu_{x_w}\\
\nu_{y_w}\\
\omega
\end{bmatrix},
\qquad
\widehat{\boldsymbol\xi}_w
=
\dot{\mathbf T}_{wb}\mathbf T_{wb}^{-1}
}.
\]

Equivalently,

\[
\boxed{
\dot{\mathbf T}_{wb}
=
\widehat{\boldsymbol\xi}_w\mathbf T_{wb}
}.
\]

Here, \emph{spatial} means that the twist is represented relative to the
fixed space or world frame; it does not mean that the motion is
three-dimensional.

A crucial distinction is

\[
\boldsymbol\nu_w
\neq
\dot{\mathbf p}
=
{}^w\mathbf v_b
\]

in general. The column \(\boldsymbol\nu_w\) is the linear part of the
spatial twist, not simply the body-origin velocity expressed in world
coordinates. Its meaning follows from the rigid-body point-velocity
equation derived in Section~\ref{sec-se2-body-point-velocity}:

\[
{}^w\mathbf v_P
=
{}^w\mathbf v_b
+
\omega\mathbf S_2{}^w\mathbf r_{O_bP}.
\]

Here, \({}^w\mathbf v_P\in\mathbb R^2\) is the velocity of a body-fixed
point \(P\) relative to the world and expressed in world coordinates,
\({}^w\mathbf v_b=\dot{\mathbf p}\in\mathbb R^2\) is the corresponding
velocity of the body origin \(O_b\), and
\({}^w\mathbf r_{O_bP}\in\mathbb R^2\) is the relative-position vector
from \(O_b\) to \(P\), also expressed in world coordinates. Because

\[
{}^w\mathbf r_{O_bP}
=
{}^w\mathbf p_P-{}^w\mathbf p_b
=
{}^w\mathbf p_P-\mathbf p,
\]

substitution and regrouping give

\[
\begin{aligned}
{}^w\mathbf v_P
&=
\dot{\mathbf p}
+
\omega\mathbf S_2
\left(
{}^w\mathbf p_P-\mathbf p
\right)\\
&=
\dot{\mathbf p}
+
\omega\mathbf S_2{}^w\mathbf p_P
-
\omega\mathbf S_2\mathbf p\\
&=
\underbrace{
\left(
\dot{\mathbf p}
-
\omega\mathbf S_2\mathbf p
\right)
}_{\boldsymbol\nu_w}
+
\omega\mathbf S_2{}^w\mathbf p_P.
\end{aligned}
\]

Using the definition
\(\boldsymbol\nu_w=\dot{\mathbf p}-\omega\mathbf S_2\mathbf p\)
therefore produces the world-frame velocity field

\[
\boxed{
{}^w\mathbf v_P
=
\boldsymbol\nu_w
+
\omega\mathbf S_2{}^w\mathbf p_P
}.
\]

The symbol \(O_w\) denotes the fixed origin of world frame \(\{w\}\).
Its world-coordinate column is \(\mathbf 0_2\). Evaluating the velocity
field at that location by setting \({}^w\mathbf p_P=\mathbf 0_2\) gives
\({}^w\mathbf v_P=\boldsymbol\nu_w\). Thus, \(\boldsymbol\nu_w\) is the
velocity, expressed in world coordinates, assigned by the instantaneous
rigid-motion field to the location \(O_w\). The world origin need not be
a physical point on the robot; the rigid-motion velocity field is being
extended mathematically to that location. This dependence on the chosen
origin explains why changing the reference origin changes the linear
part of a twist.

\subsubsection{Physical meaning and
use}\label{physical-meaning-and-use-1}

As in the body case, the two boxed equations support two complementary
calculations. The equation

\[
\widehat{\boldsymbol\xi}_w
=
\dot{\mathbf T}_{wb}\mathbf T_{wb}^{-1}
\]

extracts the instantaneous rigid-body motion relative to the fixed world
axes and the fixed world origin. The equation

\[
\dot{\mathbf T}_{wb}
=
\widehat{\boldsymbol\xi}_w\mathbf T_{wb}
\]

uses that world-referenced motion to calculate the pose rate. The
multiplication side records the convention: a body twist acts on
\(\mathbf T_{wb}\) from the right, whereas a spatial twist acts from the
left.

The spatial twist describes the velocity of the entire rigid body
through the world-frame velocity field derived above. Once
\(\boldsymbol\xi_w\) is known, substituting the world position
\({}^w\mathbf p_P\) of any body-fixed point \(P\) gives that point's
world-frame velocity. At the body origin, \({}^w\mathbf p_P=\mathbf p\)
and \(P=O_{b}\), so
\(\boldsymbol\nu_w=\dot{\mathbf p}-\omega\mathbf S_2\mathbf p\) gives

\[
\boxed{
{}^w\mathbf v_b
=
\dot{\mathbf p}
=
\boldsymbol\nu_w
+
\omega\mathbf S_2\mathbf p
}.
\]

Therefore,

\begin{equation}\protect\phantomsection\label{eq-se2-spatial-twist-body-origin-velocity}{
\underbrace{{}^w\mathbf v_b}_{\text{velocity at }O_b}
=
\underbrace{\boldsymbol\nu_w}_{\text{velocity field at }O_w}
+
\underbrace{\omega\mathbf S_2\mathbf p}_{\text{change caused by the different location}}.
}\end{equation}

The linear part \(\boldsymbol\nu_w\) and the body-origin velocity
\({}^w\mathbf v_b\) refer to different locations. They are equal only
when the rotational contribution \(\omega\mathbf S_2\mathbf p\) is zero,
such as when \(\omega=0\) or when \(O_b\) currently coincides with
\(O_w\).

The spatial representation is useful when an algorithm uses a common
fixed world frame: for example, when comparing the motions of several
bodies, evaluating point velocities at world-frame locations, or
formulating world-referenced estimation and control. Differential-drive
wheel kinematics more naturally produces a body twist, so such an
application first maps that body twist to the equivalent spatial twist.
The next section derives this conversion. \emph{Both twists describe the
same physical instantaneous motion; only the reference frame and
reference origin used for their coordinates differ.}

\subsection{Relating body and spatial
twists}\label{relating-body-and-spatial-twists}

The body and spatial twists describe the same instantaneous motion, so
one must be convertible into the other. Start from their matrix
definitions:

\[
\widehat{\boldsymbol\xi}_b
=
\mathbf T^{-1}\dot{\mathbf T},
\qquad
\widehat{\boldsymbol\xi}_w
=
\dot{\mathbf T}\mathbf T^{-1}.
\]

The body-twist definition implies
\(\dot{\mathbf T}=\mathbf T\widehat{\boldsymbol\xi}_b\). Substituting
this result into the spatial-twist definition gives

\[
\begin{aligned}
\widehat{\boldsymbol\xi}_w
&=
\dot{\mathbf T}\mathbf T^{-1}\\
&=
\left(
\mathbf T\widehat{\boldsymbol\xi}_b
\right)
\mathbf T^{-1}\\
&=
\mathbf T\widehat{\boldsymbol\xi}_b\mathbf T^{-1}.
\end{aligned}
\]

The product \(\mathbf T\widehat{\boldsymbol\xi}_b\mathbf T^{-1}\) is
called \textbf{conjugation by \(\mathbf T\)}. To find what it does to
the twist coordinates, write the three matrices in block form:

\[
\mathbf T
=
\begin{bmatrix}
\mathbf R&\mathbf p\\
\mathbf 0_2^\mathsf T&1
\end{bmatrix},
\qquad
\widehat{\boldsymbol\xi}_b
=
\begin{bmatrix}
\omega\mathbf S_2&\boldsymbol\nu_b\\
\mathbf 0_2^\mathsf T&0
\end{bmatrix},
\qquad
\mathbf T^{-1}
=
\begin{bmatrix}
\mathbf R^\mathsf T&-\mathbf R^\mathsf T\mathbf p\\
\mathbf 0_2^\mathsf T&1
\end{bmatrix}.
\]

First multiply the two matrices on the left:

\[
\mathbf T\widehat{\boldsymbol\xi}_b
=
\begin{bmatrix}
\mathbf R\omega\mathbf S_2&\mathbf R\boldsymbol\nu_b\\
\mathbf 0_2^\mathsf T&0
\end{bmatrix}.
\]

Postmultiplying this result by \(\mathbf T^{-1}\) gives

\[
\mathbf T\widehat{\boldsymbol\xi}_b\mathbf T^{-1}
=
\begin{bmatrix}
\mathbf R\omega\mathbf S_2\mathbf R^\mathsf T&
-\mathbf R\omega\mathbf S_2\mathbf R^\mathsf T\mathbf p
+\mathbf R\boldsymbol\nu_b\\
\mathbf 0_2^\mathsf T&0
\end{bmatrix}.
\]

As shown in Section~\ref{sec-se2-rotation-matrix-derivative}, every
planar rotation \(\mathbf R\) commutes with \(\mathbf S_2\). Using
\(\mathbf R\mathbf S_2=\mathbf S_2\mathbf R\) and
\(\mathbf R\mathbf R^\mathsf T=\mathbf I_2\) gives
\(\mathbf R\mathbf S_2\mathbf R^\mathsf T=\mathbf S_2\). Therefore,

\[
\widehat{\boldsymbol\xi}_w
=
\begin{bmatrix}
\omega\mathbf S_2&
\mathbf R\boldsymbol\nu_b-\omega\mathbf S_2\mathbf p\\
\mathbf 0_2^\mathsf T&0
\end{bmatrix}.
\]

By definition, the same spatial-twist matrix is

\[
\widehat{\boldsymbol\xi}_w
=
\begin{bmatrix}
\omega\mathbf S_2&\boldsymbol\nu_w\\
\mathbf 0_2^\mathsf T&0
\end{bmatrix}.
\]

Comparing the upper-right columns gives

\[
\boxed{
\boldsymbol\nu_w
=
\mathbf R\boldsymbol\nu_b
-
\omega\mathbf S_2\mathbf p
}.
\]

This result also follows directly from the definition of
\(\boldsymbol\nu_w\) or
\({}^w\mathbf v_b=\boldsymbol\nu_w+\omega\mathbf S_2\mathbf p\) in
Equation~\ref{eq-se2-spatial-twist-body-origin-velocity}. The
body-origin velocity expressed in world coordinates is
\({}^w\mathbf v_b=\mathbf R\boldsymbol\nu_b\), so that equation becomes

\[
\mathbf R\boldsymbol\nu_b
=
\boldsymbol\nu_w
+
\omega\mathbf S_2\mathbf p.
\]

Rearranging gives
\(\boldsymbol\nu_w=\mathbf R\boldsymbol\nu_b-\omega\mathbf S_2\mathbf p\),
which agrees with the block-matrix derivation.

The first term rotates the body-axis components of the linear velocity
into world-axis components. The second term changes the reference
location from the body origin \(O_b\) to the world origin \(O_w\). A
twist therefore does not transform like a free vector, for which
rotation alone would be sufficient.

Now collect the linear and angular coordinates into

\[
\boldsymbol\xi_b
=
\begin{bmatrix}
\boldsymbol\nu_b\\
\omega
\end{bmatrix},
\qquad
\boldsymbol\xi_w
=
\begin{bmatrix}
\boldsymbol\nu_w\\
\omega
\end{bmatrix}.
\]

The preceding component equation can be written as one linear mapping:

\[
\begin{aligned}
\boldsymbol\xi_w
&=
\begin{bmatrix}
\mathbf R&-\mathbf S_2\mathbf p\\
\mathbf 0_2^\mathsf T&1
\end{bmatrix}
\begin{bmatrix}
\boldsymbol\nu_b\\
\omega
\end{bmatrix}\\
&=
\begin{bmatrix}
\mathbf R\boldsymbol\nu_b-\omega\mathbf S_2\mathbf p\\
\omega
\end{bmatrix}.
\end{aligned}
\]

This \(3\times3\) coordinate mapping is called the \textbf{adjoint
representation} of \(\mathbf T\):

\[
\boxed{
\operatorname{Ad}_{\mathbf T}
=
\begin{bmatrix}
\mathbf R&-\mathbf S_2\mathbf p\\
\mathbf 0_2^\mathsf T&1
\end{bmatrix}
\in\mathbb R^{3\times3}
}.
\]

Equivalently, the adjoint is the linear mapping defined by

\[
\widehat{
\operatorname{Ad}_{\mathbf T}\boldsymbol\xi
}
=
\mathbf T\widehat{\boldsymbol\xi}\mathbf T^{-1}.
\]

For \(\mathbf T=\mathbf T_{wb}\), it converts the body-twist coordinates
into the spatial-twist coordinates:

\[
\boxed{
\boldsymbol\xi_w
=
\operatorname{Ad}_{\mathbf T_{wb}}
\boldsymbol\xi_b
}.
\]

\subsection{Differential-drive body
twist}\label{differential-drive-body-twist}

Under the ideal rolling constraints, a differential-drive robot has no
lateral velocity at the body origin. Its body-frame linear velocity is

\[
\boldsymbol\nu_b
=
\begin{bmatrix}
v\\0
\end{bmatrix},
\]

where \(v\) is the signed forward speed in \(\mathrm{m\,s^{-1}}\). Its
body twist is therefore

\[
\boxed{
\boldsymbol\xi_b
=
\begin{bmatrix}
v\\0\\\omega
\end{bmatrix},
\qquad
\widehat{\boldsymbol\xi}_b
=
\begin{bmatrix}
0&-\omega&v\\
\omega&0&0\\
0&0&0
\end{bmatrix}
}.
\]

Substitute the hatted body twist into
\(\dot{\mathbf T}_{wb}=\mathbf T_{wb}\widehat{\boldsymbol\xi}_b\). In
block form,

\[
\mathbf T_{wb}
=
\begin{bmatrix}
\mathbf R_{wb}&\mathbf p\\
\mathbf 0_2^\mathsf T&1
\end{bmatrix},
\qquad
\widehat{\boldsymbol\xi}_b
=
\begin{bmatrix}
\omega\mathbf S_2&
\begin{bmatrix}
v\\0
\end{bmatrix}\\
\mathbf 0_2^\mathsf T&0
\end{bmatrix}.
\]

Multiplying these matrices gives

\[
\begin{aligned}
\mathbf T_{wb}\widehat{\boldsymbol\xi}_b
&=
\begin{bmatrix}
\mathbf R_{wb}&\mathbf p\\
\mathbf 0_2^\mathsf T&1
\end{bmatrix}
\begin{bmatrix}
\omega\mathbf S_2&
\begin{bmatrix}
v\\0
\end{bmatrix}\\
\mathbf 0_2^\mathsf T&0
\end{bmatrix}\\
&=
\begin{bmatrix}
\mathbf R_{wb}\omega\mathbf S_2&
\mathbf R_{wb}
\begin{bmatrix}
v\\0
\end{bmatrix}\\
\mathbf 0_2^\mathsf T&0
\end{bmatrix}.
\end{aligned}
\]

The terms involving the translation column \(\mathbf p\) disappear
because they multiply the zero blocks in \(\widehat{\boldsymbol\xi}_b\).
On the other hand,

\[
\dot{\mathbf T}_{wb}
=
\begin{bmatrix}
\dot{\mathbf R}_{wb}&\dot{\mathbf p}\\
\mathbf 0_2^\mathsf T&0
\end{bmatrix}.
\]

Because these two block matrices are equal, their corresponding
upper-left and upper-right blocks must be equal:

\[
\dot{\mathbf R}_{wb}
=
\mathbf R_{wb}\omega\mathbf S_2,
\qquad
\dot{\mathbf p}
=
\mathbf R_{wb}
\begin{bmatrix}
v\\0
\end{bmatrix}.
\]

Since \(\dot{\mathbf p}={}^w\mathbf v_b\) and

\[
\mathbf R_{wb}
=
\begin{bmatrix}
\cos\theta&-\sin\theta\\
\sin\theta&\cos\theta
\end{bmatrix},
\]

the translation block becomes

\[
\begin{aligned}
{}^w\mathbf v_b
&=
\mathbf R_{wb}
\begin{bmatrix}
v\\0
\end{bmatrix}
=
\begin{bmatrix}
v\cos\theta\\
v\sin\theta
\end{bmatrix}.
\end{aligned}
\]

The rotation derivative derived in
Section~\ref{sec-se2-rotation-matrix-derivative} is
\(\dot{\mathbf R}_{wb}=\mathbf R_{wb}\dot\theta\mathbf S_2\). Comparing
it with \(\dot{\mathbf R}_{wb}=\mathbf R_{wb}\omega\mathbf S_2\) gives
\(\dot\theta=\omega\).

Therefore, the familiar differential-drive state equation is recovered:

\[
\boxed{
\dot{\mathbf q}
=
\begin{bmatrix}
\dot x\\
\dot y\\
\dot\theta
\end{bmatrix}
=
\begin{bmatrix}
v\cos\theta\\
v\sin\theta\\
\omega
\end{bmatrix}
}.
\]

The twist formulation has not introduced a different motion model. It
has expressed the same model in a form that respects rigid-motion
composition and generalises naturally to three-dimensional robotics.

\subsection{Worked example: the two twist
representations}\label{worked-example-the-two-twist-representations}

Suppose

\[
{}^w\mathbf p_b
=
\begin{bmatrix}
2\\1
\end{bmatrix}
\mathrm m,
\qquad
\theta
=
\frac{\pi}{2}\,\mathrm{rad},
\]

and the differential-drive body twist has \(v=0.5\,\mathrm{m\,s^{-1}}\)
and \(\omega=2\,\mathrm{rad\,s^{-1}}\). The body-twist column is

\[
\boldsymbol\xi_b
=
\begin{bmatrix}
0.5\\0\\2
\end{bmatrix},
\]

where the first two entries have units of \(\mathrm{m\,s^{-1}}\) and the
last has units of \(\mathrm{rad\,s^{-1}}\). At \(\theta=\pi/2\), the
velocity of the body origin expressed in the world frame is

\[
\dot{\mathbf p}
=
\mathbf R_{wb}
\begin{bmatrix}
0.5\\0
\end{bmatrix}
=
\begin{bmatrix}
0\\0.5
\end{bmatrix}
\mathrm{m\,s^{-1}}.
\]

The spatial twist's linear part is

\[
\begin{aligned}
\boldsymbol\nu_w
&=
\dot{\mathbf p}
-
\omega\mathbf S_2\mathbf p\\
&=
\begin{bmatrix}
0\\0.5
\end{bmatrix}
-
2
\begin{bmatrix}
0&-1\\
1&0
\end{bmatrix}
\begin{bmatrix}
2\\1
\end{bmatrix}\\
&=
\begin{bmatrix}
2\\-3.5
\end{bmatrix}
\mathrm{m\,s^{-1}}.
\end{aligned}
\]

This does not contradict
\({}^w\mathbf v_b=[0\;0.5]^\mathsf T\,\mathrm{m\,s^{-1}}\). The two
columns refer to different points in the same instantaneous velocity
field. Substituting the body-origin position into the spatial velocity
field recovers its correct velocity:

\[
\boldsymbol\nu_w
+
\omega\mathbf S_2{}^w\mathbf p_b
=
\begin{bmatrix}
2\\-3.5
\end{bmatrix}
+
\begin{bmatrix}
-2\\4
\end{bmatrix}
=
\begin{bmatrix}
0\\0.5
\end{bmatrix}
\mathrm{m\,s^{-1}}.
\]

\subsection{Assumptions, failure modes, and
checks}\label{assumptions-failure-modes-and-checks-1}

\begin{itemize}
\tightlist
\item
  \textbf{Undeclared component order:} A column containing the same
  three numbers can represent different twists if one convention uses
  linear velocity first and another uses angular velocity first.
\item
  \textbf{Confusing a twist with a pose:} \(\boldsymbol\xi\) and
  \(\widehat{\boldsymbol\xi}\) describe instantaneous motion; neither is
  an element of \(SE(2)\) or a finite pose.
\item
  \textbf{Confusing \(\boldsymbol\nu_w\) with \({}^w\mathbf v_b\):} They
  are equal only when \(\omega=0\), \(\mathbf p=\mathbf 0_2\), or their
  correction term happens to vanish.
\item
  \textbf{Rotating a twist like a free vector:} Changing the reference
  origin introduces the translation-dependent term in
  \(\operatorname{Ad}_{\mathbf T}\).
\item
  \textbf{Assuming zero lateral velocity for every rigid body:} The
  condition \(\nu_{y_b}=0\) comes from the ideal differential-drive
  rolling constraint, not from the definition of a planar twist.
\item
  \textbf{Mixing finite and infinitesimal quantities:} In
  \(\mathbf T_{wb}\), the rotation block \(\mathbf R_{wb}\) is
  dimensionless and the translation column \(\mathbf p\) has units of
  metres. In \(\widehat{\boldsymbol\xi}\), the corresponding blocks have
  units of \(\mathrm{s^{-1}}\) and \(\mathrm{m\,s^{-1}}\). Therefore,
  \(\mathbf T_{wb}\) and \(\widehat{\boldsymbol\xi}\) cannot be added
  directly. A twist must first be used as a rate over a time interval
  through the appropriate pose-evolution equation.
\end{itemize}

\section*{Quick test D: planar twists}\label{quick-test-d-planar-twists}
\addcontentsline{toc}{section}{Quick test D: planar twists}

\markright{Quick test D: planar twists}

\begin{enumerate}
\def\labelenumi{\arabic{enumi}.}
\tightlist
\item
  What physical quantities are stored in
  \(\boldsymbol\xi=[\nu_x\;\nu_y\;\omega]^\mathsf T\), and what are
  their units?
\item
  Define the hat operator for the component order used in this chapter.
\item
  What is \(\mathfrak{se}(2)\), and how does it differ from \(SE(2)\)?
\item
  Starting from \(\mathbf T_{wb}\) and \(\dot{\mathbf T}_{wb}\), derive
  \(\widehat{\boldsymbol\xi}_b=\mathbf T_{wb}^{-1}\dot{\mathbf T}_{wb}\).
\item
  Why is the linear part \(\boldsymbol\nu_w\) of the spatial twist
  generally different from \({}^w\mathbf v_b\)?
\item
  State the two equivalent pose-rate equations involving the body and
  spatial twists.
\item
  Derive the differential-drive state equation from
  \(\boldsymbol\xi_b=[v\;0\;\omega]^\mathsf T\).
\item
  A planar twist is \(\boldsymbol\xi=[\nu_x\;\nu_y\;\omega]^\mathsf T\),
  where \(\nu_x\) and \(\nu_y\) are measured in \(\mathrm{m\,s^{-1}}\)
  and \(\omega\) is measured in \(\mathrm{rad\,s^{-1}}\). Why is the
  ordinary Euclidean expression
  \(\lVert\boldsymbol\xi\rVert_2=\sqrt{\nu_x^2+\nu_y^2+\omega^2}\) not a
  physically meaningful twist magnitude? How can a chosen characteristic
  length \(\ell>0\), measured in metres, make the terms comparable?
\end{enumerate}

\subsection*{Answers and hints}\label{answers-and-hints-13}
\addcontentsline{toc}{subsection}{Answers and hints}

\begin{enumerate}
\def\labelenumi{\arabic{enumi}.}
\tightlist
\item
  \(\nu_x\) and \(\nu_y\) are the two components of a linear velocity,
  measured in \(\mathrm{m\,s^{-1}}\); \(\omega\) is angular velocity,
  measured in \(\mathrm{rad\,s^{-1}}\). Their frame and reference point
  must also be declared.
\item
  The hat matrix has \(\omega\mathbf S_2\) as its upper-left
  \(2\times2\) block, \(\boldsymbol\nu\) as its upper-right \(2\times1\)
  block, and a zero bottom row, as displayed in
  Section~\ref{sec-se2-hat-operator}.
\item
  \(SE(2)\) is the group of finite planar rigid transformations. Its Lie
  algebra \(\mathfrak{se}(2)\) is the vector space of structured
  matrices representing infinitesimal rigid motions and instantaneous
  velocities.
\item
  Block multiplication gives upper-left block
  \(\mathbf R^\mathsf T\dot{\mathbf R}\), upper-right block
  \(\mathbf R^\mathsf T\dot{\mathbf p}\), and a zero bottom row.
  Identify \(\mathbf R^\mathsf T\dot{\mathbf R}=\omega\mathbf S_2\) and
  \(\boldsymbol\nu_b=\mathbf R^\mathsf T\dot{\mathbf p}\).
\item
  A spatial twist uses the fixed world origin as its velocity reference.
  Therefore, its linear part is
  \(\boldsymbol\nu_w=\dot{\mathbf p}-\omega\mathbf S_2\mathbf p\); the
  second term accounts for the offset between the world and body
  origins.
\item
  \(\dot{\mathbf T}_{wb}=\mathbf T_{wb}\widehat{\boldsymbol\xi}_b\) and
  \(\dot{\mathbf T}_{wb}=\widehat{\boldsymbol\xi}_w\mathbf T_{wb}\).
\item
  Matrix multiplication gives
  \({}^w\mathbf v_b=\mathbf R_{wb}[v\;0]^\mathsf T=[v\cos\theta\;v\sin\theta]^\mathsf T\)
  and \(\dot\theta=\omega\).
\item
  The expression adds squared linear speeds to a squared angular speed,
  so its terms represent different physical quantities and its numerical
  value depends arbitrarily on the chosen units. A characteristic length
  converts angular speed into the linear-speed scale \(\ell\omega\),
  measured in \(\mathrm{m\,s^{-1}}\), giving a weighted magnitude such
  as \(\sqrt{\nu_x^2+\nu_y^2+(\ell\omega)^2}\). The choice of \(\ell\)
  is application-dependent---for example, it might be related to the
  robot's size---so this is not a unique intrinsic norm for all tasks.
\end{enumerate}

\section*{Cumulative chapter review}\label{cumulative-chapter-review}
\addcontentsline{toc}{section}{Cumulative chapter review}

\markright{Cumulative chapter review}

\begin{enumerate}
\def\labelenumi{\arabic{enumi}.}
\tightlist
\item
  What does the pose \(\mathbf T_{wb}\in SE(2)\) describe, and how does
  it transform the body-frame coordinates of a point into world-frame
  coordinates?
\item
  Distinguish a passive coordinate change from an active rigid motion.
  For column-vector notation, which transformation acts first in
  \(\mathbf T_2\mathbf T_1\mathbf p\)?
\item
  Starting from \(\mathbf R_{wb}^\mathsf T\mathbf R_{wb}=\mathbf I_2\),
  explain why \(\mathbf R_{wb}^\mathsf T\dot{\mathbf R}_{wb}\) is
  skew-symmetric and identify its planar form.
\item
  Derive the world-frame velocity of a body-fixed point \(P\) from
  \({}^w\mathbf p_P={}^w\mathbf p_b+\mathbf R_{wb}{}^b\mathbf p_P\).
\item
  How do \(SE(2)\) and \(\mathfrak{se}(2)\) differ physically and
  mathematically? What does the hat operator do?
\item
  State the physical meanings of the body twist \(\boldsymbol\xi_b\) and
  spatial twist \(\boldsymbol\xi_w\). Why is \(\boldsymbol\nu_w\)
  generally not the velocity of the body origin?
\item
  A robot has \(\theta=\pi/2\),
  \(\mathbf p=[2\;0]^\mathsf T\,\mathrm m\), and
  \(\boldsymbol\xi_b=[1\;0\;0.5]^\mathsf T\), with linear entries in
  \(\mathrm{m\,s^{-1}}\) and angular entry in \(\mathrm{rad\,s^{-1}}\).
  Find \({}^w\mathbf v_b\) and \(\boldsymbol\nu_w\).
\item
  Starting from the differential-drive body twist
  \(\boldsymbol\xi_b=[v\;0\;\omega]^\mathsf T\), recover the state
  equation for \(\dot{\mathbf q}\).
\end{enumerate}

\subsection*{Answers and hints}\label{answers-and-hints-14}
\addcontentsline{toc}{subsection}{Answers and hints}

\begin{enumerate}
\def\labelenumi{\arabic{enumi}.}
\tightlist
\item
  \(\mathbf T_{wb}\) describes the position and orientation of body
  frame \(\{b\}\) relative to world frame \(\{w\}\). It maps a
  body-frame homogeneous point by
  \({}^w\overline{\mathbf p}_P=\mathbf T_{wb}{}^b\overline{\mathbf p}_P\),
  or equivalently
  \({}^w\mathbf p_P=\mathbf R_{wb}{}^b\mathbf p_P+{}^w\mathbf p_b\).
\item
  A passive change describes the same geometric object in different
  coordinates; an active motion moves the object while retaining the
  reference frame. In \(\mathbf T_2\mathbf T_1\mathbf p\), the rightmost
  transformation \(\mathbf T_1\) acts first.
\item
  Differentiation gives
  \(\dot{\mathbf R}_{wb}^\mathsf T\mathbf R_{wb}+\mathbf R_{wb}^\mathsf T\dot{\mathbf R}_{wb}=\mathbf 0_{2\times2}\).
  Therefore, \(\mathbf R_{wb}^\mathsf T\dot{\mathbf R}_{wb}\) equals the
  negative of its transpose and has the planar form
  \(\omega\mathbf S_2\in\mathfrak{so}(2)\).
\item
  Differentiating gives
  \({}^w\mathbf v_P={}^w\mathbf v_b+\dot{\mathbf R}_{wb}{}^b\mathbf p_P\).
  Using \(\dot{\mathbf R}_{wb}=\omega\mathbf S_2\mathbf R_{wb}\) and
  \({}^w\mathbf r_{O_bP}=\mathbf R_{wb}{}^b\mathbf p_P\) gives
  \({}^w\mathbf v_P={}^w\mathbf v_b+\omega\mathbf S_2{}^w\mathbf r_{O_bP}\).
\item
  \(SE(2)\) is the nonlinear group of finite planar poses and rigid
  transformations. \(\mathfrak{se}(2)\) is its three-dimensional tangent
  vector space at the identity and represents infinitesimal motions or
  rates. The hat operator maps a twist coordinate column
  \([\nu_x\;\nu_y\;\omega]^\mathsf T\) to its structured \(3\times3\)
  matrix in \(\mathfrak{se}(2)\).
\item
  \(\boldsymbol\xi_b\) describes world-relative motion using body axes
  and the body origin as its linear-velocity reference.
  \(\boldsymbol\xi_w\) describes the same motion using world axes and
  the world origin. Its linear part is the velocity field at \(O_w\), so
  the velocity at \(O_b\) is instead
  \({}^w\mathbf v_b=\boldsymbol\nu_w+\omega\mathbf S_2\mathbf p\).
\item
  At \(\theta=\pi/2\),
  \(\mathbf R_{wb}[1\;0]^\mathsf T=[0\;1]^\mathsf T\), so
  \({}^w\mathbf v_b=[0\;1]^\mathsf T\,\mathrm{m\,s^{-1}}\). Because
  \(0.5\mathbf S_2[2\;0]^\mathsf T=[0\;1]^\mathsf T\,\mathrm{m\,s^{-1}}\),
  \(\boldsymbol\nu_w={}^w\mathbf v_b-\omega\mathbf S_2\mathbf p=\mathbf 0_2\,\mathrm{m\,s^{-1}}\).
\item
  Substitution into
  \(\dot{\mathbf T}_{wb}=\mathbf T_{wb}\widehat{\boldsymbol\xi}_b\)
  gives \(\dot x=v\cos\theta\), \(\dot y=v\sin\theta\), and
  \(\dot\theta=\omega\), so
  \(\dot{\mathbf q}=[v\cos\theta\;v\sin\theta\;\omega]^\mathsf T\).
\end{enumerate}

\section*{References}\label{references-4}
\addcontentsline{toc}{section}{References}

\markright{References}

\begin{itemize}
\tightlist
\item
  John J. Craig, \emph{Introduction to Robotics: Mechanics and Control},
  4th ed., Chapters 2 and 5, especially frame mappings, transformation
  interpretations, angular velocity, and rigid-body point velocity.
\item
  Peter Corke, \emph{Robotics, Vision and Control: Fundamental
  Algorithms in Python}, 3rd ed., Chapters 2-3, especially time-varying
  pose, angular velocity, and derivatives of orientation and pose.
\item
  Kevin M. Lynch and Frank C. Park, \emph{Modern Robotics: Mechanics,
  Planning, and Control}, Chapters 2-3, especially configuration spaces,
  angular velocity, skew-symmetric representations, body and spatial
  twists, and rigid-body motions. The \(SE(3)\) twist definitions are
  specialised here to planar motion in \(SE(2)\).
\end{itemize}

\chapter{\texorpdfstring{Degrees of Freedom and Spatial Motion:
\(SO(3)\) and
\(SE(3)\)}{Degrees of Freedom and Spatial Motion: SO(3) and SE(3)}}\label{sec-spatial-rigid-body-motion}

\section*{Why this chapter matters}\label{why-this-chapter-matters-2}
\addcontentsline{toc}{section}{Why this chapter matters}

\markright{Why this chapter matters}

The preceding chapter described a rigid body constrained to a plane by
three pose coordinates and represented its configuration by an element
of \(SE(2)\). A robot operating freely in three-dimensional physical
space can translate in three independent directions and rotate about
three independent directions. This chapter develops the corresponding
six-degree-of-freedom model, represents spatial orientation with
\(SO(3)\), combines orientation and position with \(SE(3)\), and extends
planar angular velocity and twists to three dimensions.

These ideas are required for robot arms, aerial vehicles, underwater
vehicles, cameras, lidar sensors, and any system whose orientation
cannot be represented by one planar heading angle.

\section*{Prerequisites and
assumptions}\label{prerequisites-and-assumptions-4}
\addcontentsline{toc}{section}{Prerequisites and assumptions}

\markright{Prerequisites and assumptions}

The reader should understand vectors, matrices, rank, and constraints
from Chapter~\ref{sec-linear-algebra-foundations}; coordinate frames,
cross products, and three-dimensional point transformations from
Chapter~\ref{sec-geometry-coordinate-frames}; and planar rigid
transformations, tangent spaces, and twists from
Chapter~\ref{sec-planar-rigid-body-motion}.

Physical space is modelled as three-dimensional Euclidean space. Every
body considered here is rigid, so the distance between any two
body-fixed points remains constant. Frames are right-handed and
orthonormal. Position is measured in metres, time in seconds, and angles
in radians.

\section{Degrees of freedom and configuration
space}\label{degrees-of-freedom-and-configuration-space}

\subsection{Configuration is a complete geometric
description}\label{configuration-is-a-complete-geometric-description}

A system's \textbf{configuration} is a complete geometric specification
of the position of every labelled material point of the system at one
instant. A \emph{material point} is a particular point fixed in the
physical material of the body; its identity does not change as the body
moves. Configuration describes where the system is, not how fast it is
moving or what forces act on it.

This definition can be written as a map. Let \(\mathcal B\) denote the
set of labelled material points belonging to one planar rigid body.
After attaching frame \(\{b\}\) to the body, each point
\(P\in\mathcal B\) has a constant body-frame coordinate
\({}^{b}\mathbf p_P\in\mathbb R^2\). A configuration is a
\textbf{placement map} \(\chi\) that assigns a world-frame position to
every material point:

\[
\chi:\mathcal B\rightarrow\mathbb R^2,
\qquad
P\longmapsto\chi(P)={}^{w}\mathbf p_P.
\]

The symbol \(\chi\) is the Greek lowercase letter \emph{chi}. Because
the body is rigid, a valid placement must preserve the distance between
every pair of material points:

\[
\left\|\chi(P)-\chi(Q)\right\|_2
=
\left\|{}^{b}\mathbf p_P-{}^{b}\mathbf p_Q\right\|_2
\qquad
\text{for every }P,Q\in\mathcal B.
\]

The set of all placement maps allowed by the system's physical
constraints is its \textbf{configuration space}, abbreviated
\textbf{C-space} and denoted here by \(\mathcal C\). One configuration
is therefore one element \(\chi\in\mathcal C\). Calling a configuration
a \emph{point in C-space} means one abstract element of this set; it
does not mean one physical point of the rigid body or one point in the
world.

\subsubsection{Why one transformation represents the whole rigid
body}\label{why-one-transformation-represents-the-whole-rigid-body}

Listing \(\chi(P)\) separately for every \(P\in\mathcal B\) would be
highly redundant. All body-frame point coordinates are fixed by
rigidity, so a planar rotation \(\mathbf R_{wb}\in SO(2)\) and the world
position \({}^{w}\mathbf p_b\in\mathbb R^2\) of the body-frame origin
determine every value of the placement map:

\[
\boxed{
\chi(P)
=
{}^{w}\mathbf p_P
=
\mathbf R_{wb}{}^{b}\mathbf p_P
+
{}^{w}\mathbf p_b
\qquad
\text{for every }P\in\mathcal B
}.
\]

The homogeneous transformation

\[
\mathbf T_{wb}
=
\begin{bmatrix}
\mathbf R_{wb}&{}^{w}\mathbf p_b\\
\mathbf 0_2^\mathsf T&1
\end{bmatrix}
\in SE(2)
\]

packages the same rotation and translation. For one selected point
\(P\), the multiplication

\[
{}^{w}\overline{\mathbf p}_P
=
\mathbf T_{wb}{}^{b}\overline{\mathbf p}_P
\]

evaluates the placement map only at that point and returns its
world-frame homogeneous coordinate. Applying the same \(\mathbf T_{wb}\)
to every \(P\in\mathcal B\) produces the world positions of all the
body's material points and therefore specifies the complete
configuration. The transformation is not itself the world coordinate of
one point; it is the single rule that locates the entire rigid body.

Every configuration uses the same fixed world frame \(\{w\}\). Two
configurations \(\chi_1\) and \(\chi_2\) differ when their body-origin
positions \({}^{w}\mathbf p_b^{(1)}\) and \({}^{w}\mathbf p_b^{(2)}\)
differ, their rotations \(\mathbf R_{wb}^{(1)}\) and
\(\mathbf R_{wb}^{(2)}\) differ, or both differ. They are the same
configuration only when both the translation and rotation agree,
equivalently when \(\mathbf T_{wb}^{(1)}=\mathbf T_{wb}^{(2)}\).

Once the body frame has been fixed relative to the body, each
orientation-preserving planar placement determines exactly one
\(\mathbf T_{wb}\in SE(2)\), and each \(\mathbf T_{wb}\in SE(2)\)
determines exactly one such placement. The two sets are therefore in
one-to-one correspondence, written

\[
\boxed{
\mathcal C\cong SE(2)
},
\]

where \(\cong\) means that the two sets carry the same configuration
information through this correspondence. Robotics texts consequently
often identify the configuration space of a free planar rigid body with
\(SE(2)\) and write \(\mathcal C=SE(2)\). Here, an element of \(SE(2)\)
is being used as a representation of a complete body placement, even
though the same matrix can also act as an operator that transforms point
coordinates.

\subsubsection{Configuration coordinates are not the configuration
itself}\label{configuration-coordinates-are-not-the-configuration-itself}

For the planar rigid body, one convenient coordinate column is

\[
\mathbf q
=
\begin{bmatrix}
x\\y\\\theta
\end{bmatrix},
\]

where \(x\) and \(y\) are the entries of
\({}^{w}\mathbf p_b=[x\;y]^\mathsf T\) and \(\theta\) is the orientation
of frame \(\{b\}\) relative to frame \(\{w\}\). These three
\textbf{generalised coordinates} construct the transformation
representing the configuration:

\[
\boxed{
\mathbf q
\longmapsto
\mathbf T_{wb}(\mathbf q)
=
\begin{bmatrix}
\cos\theta&-\sin\theta&x\\
\sin\theta&\cos\theta&y\\
0&0&1
\end{bmatrix}
\in SE(2)
}.
\]

Thus, \(\chi\) denotes the complete physical placement,
\(\mathbf T_{wb}\) represents that placement as a rigid transformation,
and \(\mathbf q\) is a coordinate description used to construct the
transformation. These objects contain equivalent placement information
once the world and body frames, axis directions, units, positive
rotation direction, and transformation direction have been defined, but
they have different mathematical types. In particular, \(\theta\) and
\(\theta+2\pi n\) for any integer \(n\) construct the same rotation, so
an unrestricted real-valued angle does not give a unique global
coordinate for each configuration.

A configuration space is different from a physical \textbf{workspace}. A
workspace is a region of physical space that a robot or end-effector can
reach. A C-space contains complete robot configurations. Different joint
arrangements can place an end-effector at the same workspace point while
representing different points in C-space.

\subsection{Definition of degree of
freedom}\label{definition-of-degree-of-freedom}

A \textbf{degree of freedom}, abbreviated \textbf{DOF}, is one
independent local direction in which a configuration can change while
continuing to satisfy all configuration constraints. The number of
degrees of freedom counts these independent allowable directions.

{\def\LTcaptype{none} % do not increment counter
\begin{longtable}[]{@{}
  >{\raggedright\arraybackslash}p{(\linewidth - 4\tabcolsep) * \real{0.4138}}
  >{\raggedright\arraybackslash}p{(\linewidth - 4\tabcolsep) * \real{0.5517}}
  >{\raggedleft\arraybackslash}p{(\linewidth - 4\tabcolsep) * \real{0.0345}}@{}}
\toprule\noalign{}
\begin{minipage}[b]{\linewidth}\raggedright
System
\end{minipage} & \begin{minipage}[b]{\linewidth}\raggedright
Independently chosen coordinates
\end{minipage} & \begin{minipage}[b]{\linewidth}\raggedleft
DOF
\end{minipage} \\
\midrule\noalign{}
\endhead
\bottomrule\noalign{}
\endlastfoot
Point constrained to a straight rail & Distance \(s\) along the rail &
\(1\) \\
Point free in a plane & Position \((x,y)\) & \(2\) \\
Door attached by an ideal hinge & Hinge angle \(\theta\) & \(1\) \\
Free planar rigid body & Position and heading \((x,y,\theta)\) &
\(3\) \\
Free spatial rigid body & Three position and three orientation
coordinates & \(6\) \\
\end{longtable}
}

The door illustrates independence. Once its hinge angle is chosen, the
positions of all its material points are determined. Recording the world
coordinates of its corners would store many scalar values, but the hinge
and rigidity constraints prevent those values from being varied
independently.

\subsubsection{Mathematical definition from allowable small
changes}\label{mathematical-definition-from-allowable-small-changes}

Suppose a candidate configuration is represented by a variable column

\[
\mathbf z
=
\begin{bmatrix}
z_1&\cdots&z_{n_v}
\end{bmatrix}^{\mathsf T}
\in\mathbb R^{n_v},
\]

where \(n_v\in\mathbb N_0\) is the number of stored scalar variables.
Valid configurations must satisfy \(n_c\in\mathbb N_0\) scalar equality
constraints collected in the mapping

\[
\mathbf g:\mathbb R^{n_v}\rightarrow\mathbb R^{n_c},
\qquad
\mathbf g(\mathbf z)=\mathbf 0_{n_c}.
\]

The \textbf{constraint Jacobian}, or constraint-derivative matrix, is
\(\mathbf J_{\mathbf g}(\mathbf z)\in\mathbb R^{n_c\times n_v}\) with
entry

\[
\left[\mathbf J_{\mathbf g}(\mathbf z)\right]_{ij}
\mathrel{:=}
\frac{\partial g_i}{\partial z_j}.
\]

To describe an allowable local direction at a selected valid
configuration \(\mathbf z_0\in\mathbb R^{n_v}\), let \(\mathbf z(s)\) be
any differentiable curve of valid configurations, where
\(s\in\mathbb R\) is a dimensionless curve parameter,
\(\mathbf z(0)=\mathbf z_0\), and
\(\mathbf g(\mathbf z(s))=\mathbf 0_{n_c}\) for every \(s\) near zero.
Define the curve's direction at \(\mathbf z_0\) by
\(\delta\mathbf z\mathrel{:=}\mathbf z'(0)\in\mathbb R^{n_v}\).
Differentiating the constraint equation with respect to \(s\) and
applying the multivariable chain rule gives

\[
\boxed{
\left.
\frac{d}{ds}\mathbf g(\mathbf z(s))
\right|_{s=0}
=
\mathbf J_{\mathbf g}(\mathbf z_0)\,\delta\mathbf z
=
\mathbf 0_{n_c}
}.
\]

The set of all vectors satisfying this linear equation is the null space
of the Jacobian evaluated at \(\mathbf z_0\):

\[
\boxed{
\mathcal N\!\left(\mathbf J_{\mathbf g}(\mathbf z_0)\right)
=
\left\{
\delta\mathbf z\in\mathbb R^{n_v}
\;\middle|\;
\mathbf J_{\mathbf g}(\mathbf z_0)\delta\mathbf z
=
\mathbf 0_{n_c}
\right\}
}.
\]

This is not the same set as the configuration constraint set
\(\{\mathbf z\in\mathbb R^{n_v}\mid\mathbf g(\mathbf z)=\mathbf0_{n_c}\}\).
The constraint set contains valid configurations and may be curved; the
null space contains allowable infinitesimal changes at the one selected
configuration \(\mathbf z_0\) and is a linear space.

At a \textbf{regular configuration}, meaning that the rank of the
constraint Jacobian remains constant in a neighbourhood of
\(\mathbf z_0\), these null-space vectors are exactly the allowable
infinitesimal directions. Each independent null-space direction is one
independent way to change the configuration locally, so

\[
\boxed{
n_{\mathrm{dof}}(\mathbf z_0)
\mathrel{:=}
\dim\mathcal N\!\left(\mathbf J_{\mathbf g}(\mathbf z_0)\right)
}.
\]

Applying the rank-nullity theorem from
Section~\ref{sec-linear-algebra-rank-nullity} gives

\[
\boxed{
n_{\mathrm{dof}}(\mathbf z_0)
=
n_v-\operatorname{rank}\!\left(\mathbf J_{\mathbf g}(\mathbf z_0)\right)
}.
\]

If the \(n_c\) constraints are independent near \(\mathbf z_0\), then
\(\operatorname{rank}(\mathbf J_{\mathbf g}(\mathbf z_0))=n_c\), and the
familiar counting rule follows:

\[
\boxed{
n_{\mathrm{dof}}
=
n_v-n_c
}.
\]

\subsubsection{Example: a point constrained to a
circle}\label{example-a-point-constrained-to-a-circle}

Consider a point with Cartesian coordinate column
\(\mathbf p=[x\;y]^\mathsf T\in\mathbb R^2\) constrained to a circle of
fixed radius \(r_c>0\) metres centred at the world origin. Its two
stored variables satisfy

\[
g(\mathbf p)
\mathrel{:=}
x^2+y^2-r_c^2
=
0.
\]

Here, \(n_v=2\) and there is one scalar constraint. Its constraint
Jacobian is

\[
\mathbf J_g(\mathbf p)
=
\begin{bmatrix}
2x&2y
\end{bmatrix}
\in\mathbb R^{1\times2}.
\]

An allowable change \(\delta\mathbf p=[\delta x\;\delta y]^\mathsf T\)
must satisfy

\[
\begin{bmatrix}
2x&2y
\end{bmatrix}
\begin{bmatrix}
\delta x\\
\delta y
\end{bmatrix}
=
2x\,\delta x+2y\,\delta y
=
0.
\]

For this local calculation, \(\mathbf p=[x\;y]^\mathsf T\) is one
selected, fixed configuration on the circle. Thus, \(x\) and \(y\) are
given constants, while the allowable change \(\delta\mathbf p\) is the
variable being determined. Because \(r_c>0\), the selected point is not
\(\mathbf0_2\), so the row \(\mathbf J_g(\mathbf p)=[2x\;2y]\) is
nonzero and has rank one. The rank-nullity formula therefore gives

\[
\boxed{
n_{\mathrm{dof}}
=
\dim\mathcal N\!\left(\mathbf J_g(\mathbf p)\right)
=
2-1
=
1
}.
\]

The equation \(x\,\delta x+y\,\delta y=0\) also shows the actual
allowable direction: \(\delta\mathbf p\) must be perpendicular to the
fixed radial vector \(\mathbf p\). Every vector perpendicular to
\([x\;y]^\mathsf T\) in \(\mathbb R^2\) is a scalar multiple of
\([-y\;x]^\mathsf T\), so

\[
\mathcal N\!\left(\mathbf J_g(\mathbf p)\right)
=
\left\{
\alpha
\begin{bmatrix}
-y\\x
\end{bmatrix}
\;\middle|\;
\alpha\in\mathbb R
\right\},
\]

Here, \(x\) and \(y\) remain fixed and only the dimensionless scalar
\(\alpha\in\mathbb R\) varies. Choosing \(\alpha\) determines both
components through \(\delta x=-\alpha y\) and \(\delta y=\alpha x\).
Substitution confirms that every such change satisfies the linearised
constraint:

\[
\mathbf J_g(\mathbf p)\,\delta\mathbf p
=
\begin{bmatrix}
2x&2y
\end{bmatrix}
\alpha
\begin{bmatrix}
-y\\x
\end{bmatrix}
=
\alpha(-2xy+2yx)
=
0.
\]

The single varying coefficient \(\alpha\) explicitly represents the one
local degree of freedom. At a different point on the circle, the fixed
values of \(x\) and \(y\) produce a different tangent direction, but the
null space still has dimension one.

The number of stored values is consequently not necessarily the number
of degrees of freedom. For example, a spatial rotation matrix stores
nine scalar entries, but its constraints leave only three independent
rotational directions. Likewise, the circle uses two Cartesian entries
even though its allowable-change space is one-dimensional.

This counting rule applies directly to geometric, or \textbf{holonomic},
equality constraints that restrict configuration. Constraints imposed
only on velocity require separate analysis and do not necessarily reduce
the dimension of the configuration space.

\subsection{Why a free planar rigid body has three degrees of
freedom}\label{why-a-free-planar-rigid-body-has-three-degrees-of-freedom}

Select three non-collinear points \(A\), \(B\), and \(C\) fixed to a
planar rigid body. \emph{Non-collinear} means that the points do not all
lie on one line, so they form a non-degenerate triangle that visibly
represents a two-dimensional rigid shape.

Three points are convenient but not minimal: two distinct labelled
points already determine planar position and orientation, whereas one
point cannot reveal rotation about itself. The triangle is used here
because its fixed pairwise distances make the rigidity constraints
especially clear.

If the three points could move independently, their planar coordinates
would provide six scalar variables:

\[
(x_A,y_A,x_B,y_B,x_C,y_C)
\in
\mathbb R^6.
\]

Rigidity fixes the three pairwise distances

\[
\left\|\mathbf p_A-\mathbf p_B\right\|_2=d_{AB},
\qquad
\left\|\mathbf p_B-\mathbf p_C\right\|_2=d_{BC},
\qquad
\left\|\mathbf p_A-\mathbf p_C\right\|_2=d_{AC},
\]

where \(\mathbf p_A,\mathbf p_B,\mathbf p_C\in\mathbb R^2\) are the
point-coordinate columns and the positive constants \(d_{AB}\),
\(d_{BC}\), and \(d_{AC}\) are the fixed distances measured in metres.
For a non-degenerate triangle, these are three independent scalar
constraints. Therefore,

\[
n_{\mathrm{dof}}
=
6-3
=
3.
\]

These three freedoms can be represented more conveniently by two
translations and one rotation:

\[
\boxed{
\text{free planar rigid body:}
\qquad
(x,y,\theta)
\qquad
\Longrightarrow
\qquad
3\ \mathrm{DOF}
}.
\]

The fixed distances preserve the body's shape; they do not prevent the
complete triangle from translating or rotating as one rigid object.

\subsection{Why a free spatial rigid body has six degrees of
freedom}\label{why-a-free-spatial-rigid-body-has-six-degrees-of-freedom}

Allow the same three non-collinear body-fixed points to move in
three-dimensional space. Each point has three Cartesian coordinates, so
the three points provide nine scalar variables. Rigidity preserves the
three pairwise distances \(d_{AB}\), \(d_{BC}\), and \(d_{AC}\), which
are three independent constraints for a non-degenerate triangle.
Therefore,

\[
n_{\mathrm{dof}}
=
9-3
=
6.
\]

The six freedoms are interpreted as three translational and three
rotational freedoms:

\[
\boxed{
\text{free spatial rigid body:}
\quad
\underbrace{(x,y,z)}_{3\text{ translations}}
\quad+
\underbrace{\text{spatial orientation}}_{3\text{ rotational DOF}}
\quad
\Longrightarrow
\quad
6\ \mathrm{DOF}
}.
\]

The phrase \emph{three rotational degrees of freedom} does not yet
prescribe three particular angles. Euler angles, rotation vectors,
rotation matrices, and unit quaternions are different representations of
the same three-dimensional orientation freedom, each with different
coordinate properties and limitations. The next learning block develops
the rotation-matrix representation \(SO(3)\) before introducing any
minimal angle convention.

\subsection{Constraints, actuators, state variables, and degrees of
freedom}\label{constraints-actuators-state-variables-and-degrees-of-freedom}

Degrees of freedom are often confused with other counts in a robot
model. The following quantities answer different questions:

\begin{itemize}
\tightlist
\item
  \textbf{Configuration DOF:} How many independent real coordinates are
  needed locally to specify where the system is?
\item
  \textbf{Actuator count:} How many independently commanded physical
  inputs does the system have?
\item
  \textbf{Instantaneous mobility:} How many independent velocity
  directions are allowed at one configuration?
\item
  \textbf{State dimension:} How many variables are required by a chosen
  dynamic model to predict future behaviour?
\item
  \textbf{Representation size:} How many numbers are stored by a
  particular coordinate representation?
\end{itemize}

These counts need not be equal. A free spatial rigid body has six
configuration degrees of freedom. A dynamic state that stores both its
pose and velocity may require twelve local scalar coordinates. A
rotation matrix stores nine numbers even though spatial orientation has
only three degrees of freedom, because its entries satisfy
orthonormality and determinant constraints.

A differential-drive robot provides another important distinction. Its
configuration is represented by \(\mathbf T_{wb}\in SE(2)\), its
configuration space is modelled by \(SE(2)\), and it has three degrees
of freedom described locally by \((x,y,\theta)\). Under ideal rolling,
however, its instantaneous body velocity has only two independently
commanded components, \(v\) and \(\omega\), because sideways velocity is
constrained to zero. The robot cannot move sideways instantaneously, but
it can change all three pose coordinates through a sequence of forward
and turning motions. Therefore, its two instantaneous control directions
do not reduce its configuration space to two dimensions.

\subsection{Worked examples}\label{worked-examples}

\subsubsection{A hinged door}\label{a-hinged-door}

A free spatial door panel would have six degrees of freedom. An ideal
revolute hinge permits only rotation about its hinge axis and prevents
the other five independent motions. The constrained door therefore has

\[
n_{\mathrm{dof}}
=
6-5
=
1,
\]

which can be represented by one hinge angle \(\theta\in\mathbb R\)
measured in radians over the mechanically allowed interval.

\subsubsection{A free aerial vehicle}\label{a-free-aerial-vehicle}

Treat the vehicle as one rigid body with no geometric connection to the
environment. Its configuration has six degrees of freedom: three
position freedoms and three orientation freedoms. This statement does
not imply that the vehicle can independently accelerate in all six
directions at one instant. Its actuator arrangement and dynamics
determine which accelerations can actually be produced.

\subsection{Assumptions, limitations, and
checks}\label{assumptions-limitations-and-checks}

\begin{itemize}
\tightlist
\item
  \textbf{Counting redundant constraints:} Several equations may express
  the same geometric restriction. Check the rank of the constraint
  derivatives rather than counting written equations blindly.
\item
  \textbf{Confusing coordinates with physical freedom:} A rotation
  matrix has nine entries but only three rotational degrees of freedom.
  Identify the independent constraints imposed on a representation.
\item
  \textbf{Confusing actuator count with DOF:} A robot may be
  underactuated or may have redundant actuators. State whether a count
  refers to configuration, velocity, inputs, or state.
\item
  \textbf{Applying \(n_v-n_c\) at a singular configuration:} Constraint
  rank can change at singularities. Verify that the constraints are
  regular in the region being analysed.
\item
  \textbf{Treating a velocity constraint as a configuration constraint:}
  Under ideal rolling, a wheel moves without sliding sideways at its
  ground contact. This condition is \textbf{nonholonomic} when it
  restricts the allowable velocities but cannot be replaced by an
  equivalent equation involving only the configuration. For example, a
  differential-drive robot cannot move sideways instantaneously, but
  turning and driving can still change all three coordinates
  \((x,y,\theta)\). The rolling constraint therefore reduces
  instantaneous mobility without reducing the dimension of its C-space;
  analyse valid configurations and valid velocities separately.
\end{itemize}

\section*{Quick test A: configuration and degrees of
freedom}\label{quick-test-a-configuration-and-degrees-of-freedom}
\addcontentsline{toc}{section}{Quick test A: configuration and degrees
of freedom}

\markright{Quick test A: configuration and degrees of freedom}

\begin{enumerate}
\def\labelenumi{\arabic{enumi}.}
\tightlist
\item
  Define configuration and configuration space using the placement map
  \(\chi\).
\item
  For a planar rigid body, distinguish the placement map \(\chi\), the
  transformation \(\mathbf T_{wb}\), and the generalised-coordinate
  column \(\mathbf q\).
\item
  Why does applying \(\mathbf T_{wb}\) to one body-fixed point return
  only that point's world coordinate, while the same transformation
  still represents the body's complete configuration?
\item
  Express the number of degrees of freedom using the null space and rank
  of a constraint Jacobian. Apply the definition to a point constrained
  to a circle.
\item
  Use three non-collinear body-fixed points to explain why a free planar
  rigid body has three degrees of freedom.
\item
  Why does an ideal hinged door have one degree of freedom rather than
  six?
\item
  Distinguish configuration DOF, instantaneous mobility, actuator count,
  and state dimension.
\item
  Why does a differential-drive robot have three configuration degrees
  of freedom even though its ideal body velocity is specified by only
  \(v\) and \(\omega\)?
\item
  Under what condition is the counting rule \(n_{\mathrm{dof}}=n_v-n_c\)
  locally valid?
\end{enumerate}

\subsection*{Answers and hints}\label{answers-and-hints-15}
\addcontentsline{toc}{subsection}{Answers and hints}

\begin{enumerate}
\def\labelenumi{\arabic{enumi}.}
\tightlist
\item
  A configuration is a placement map \(\chi\) assigning a world position
  \(\chi(P)={}^{w}\mathbf p_P\) to every labelled material point \(P\)
  while satisfying the system's geometric constraints. The configuration
  space \(\mathcal C\) is the set of all allowed placement maps.
\item
  The map \(\chi\) is the complete placement of all body points. The
  matrix \(\mathbf T_{wb}\in SE(2)\) represents that placement through
  one rotation and translation. The column
  \(\mathbf q=[x\;y\;\theta]^\mathsf T\) is a set of generalised
  coordinates used to construct \(\mathbf T_{wb}\).
\item
  Multiplication by one point coordinate evaluates the placement rule at
  that one point. Because every material point has a fixed body-frame
  coordinate, applying the same transformation to all such coordinates
  reconstructs the complete placement map; one transformation therefore
  represents the whole rigid body.
\item
  At a selected valid configuration \(\mathbf z_0\), allowable small
  changes satisfy
  \(\mathbf J_{\mathbf g}(\mathbf z_0)\delta\mathbf z=\mathbf 0_{n_c}\),
  so
  \(n_{\mathrm{dof}}(\mathbf z_0)=\dim\mathcal N(\mathbf J_{\mathbf g}(\mathbf z_0))=n_v-\operatorname{rank}(\mathbf J_{\mathbf g}(\mathbf z_0))\).
  For \(g(x,y)=x^2+y^2-r_c^2\), the Jacobian \([2x\;\;2y]\) has rank one
  on a circle with \(r_c>0\), and its null space is spanned by
  \([-y\;x]^\mathsf T\), so \(n_{\mathrm{dof}}=2-1=1\).
\item
  Three planar points provide six scalar coordinates. Rigidity fixes
  their three pairwise distances, giving three independent constraints
  for a non-degenerate triangle. Thus, \(6-3=3\) freedoms remain: two
  translations and one rotation.
\item
  The hinge imposes five independent constraints on a free spatial body
  and leaves only rotation about the hinge axis, so \(6-5=1\).
\item
  Configuration DOF counts independent pose coordinates; instantaneous
  mobility counts allowed velocity directions at one configuration;
  actuator count counts independent command inputs; and state dimension
  counts the variables retained by a dynamic model.
\item
  Its configuration still requires \((x,y,\theta)\), but the rolling
  constraint removes instantaneous sideways velocity. Sequences of
  allowed forward and turning motions can nevertheless change all three
  configuration coordinates.
\item
  The rule applies locally when the equality constraints are independent
  and regular, meaning that their derivative has constant rank in the
  neighbourhood being considered.
\end{enumerate}

\section{\texorpdfstring{Spatial orientation and
\(SO(3)\)}{Spatial orientation and SO(3)}}\label{sec-so3-spatial-orientation}

\subsection{Constructing a spatial rotation
matrix}\label{sec-so3-rotation-construction}

The \textbf{spatial orientation} of body frame \(\{b\}\) relative to
world frame \(\{w\}\) specifies how the three body axes point when
expressed using world-frame coordinates; the locations of the frame
origins are irrelevant to orientation. As established in
Chapter~\ref{sec-geometry-coordinate-frames}, let \(\mathbf e_{x,b}\),
\(\mathbf e_{y,b}\), and \(\mathbf e_{z,b}\) be the geometric unit
vectors along the body axes. Their world-frame coordinate columns are

\[
{}^{w}\mathbf e_{x,b},
\quad
{}^{w}\mathbf e_{y,b},
\quad
{}^{w}\mathbf e_{z,b}
\in\mathbb R^3.
\]

Each column is dimensionless because it describes a unit direction.
Placing the columns side by side defines

\[
\boxed{
\mathbf R_{wb}
\mathrel{:=}
\begin{bmatrix}
{}^{w}\mathbf e_{x,b}&
{}^{w}\mathbf e_{y,b}&
{}^{w}\mathbf e_{z,b}
\end{bmatrix}
\in\mathbb R^{3\times3}
}.
\]

The first column states where the positive \(x_b\)-axis points in world
coordinates, and the other columns do the same for the positive \(y_b\)-
and \(z_b\)-axes. The subscript order \(wb\) means that
\(\mathbf R_{wb}\) maps coordinate columns from frame \(\{b\}\) to frame
\(\{w\}\).

To derive this mapping, let the geometric vector \(\mathbf v\) have
body-frame coordinate column

\[
{}^{b}\mathbf v
=
\begin{bmatrix}
v_{x_b}\\
v_{y_b}\\
v_{z_b}
\end{bmatrix}
\in\mathbb R^3,
\]

where the entries are its scalar components along the body axes and
carry the physical units of \(\mathbf v\). Expanding the vector in the
body basis and then expressing the result in world coordinates gives

\[
\begin{aligned}
{}^{w}\mathbf v
&=
v_{x_b}\,{}^{w}\mathbf e_{x,b}
+
v_{y_b}\,{}^{w}\mathbf e_{y,b}
+
v_{z_b}\,{}^{w}\mathbf e_{z,b}\\
&=
\begin{bmatrix}
{}^{w}\mathbf e_{x,b}&
{}^{w}\mathbf e_{y,b}&
{}^{w}\mathbf e_{z,b}
\end{bmatrix}
\begin{bmatrix}
v_{x_b}\\
v_{y_b}\\
v_{z_b}
\end{bmatrix}.
\end{aligned}
\]

Therefore,

\[
\boxed{
{}^{w}\mathbf v
=
\mathbf R_{wb}\,{}^{b}\mathbf v
}.
\]

This multiplication changes the coordinate description of \(\mathbf v\);
it does not change the geometric vector itself.

\subsection{\texorpdfstring{Rotation constraints and the definition of
\(SO(3)\)}{Rotation constraints and the definition of SO(3)}}\label{sec-so3-definition}

Write the columns of \(\mathbf R_{wb}\) as
\(\mathbf r_1,\mathbf r_2,\mathbf r_3\in\mathbb R^3\). Because they
represent the axes of an orthonormal frame, each column has unit length
and each distinct pair is perpendicular:

\[
\mathbf r_i^\mathsf T\mathbf r_i=1,
\qquad
\mathbf r_i^\mathsf T\mathbf r_j=0
\quad\text{for }i\ne j,
\qquad
i,j\in\{1,2,3\}.
\]

Matrix multiplication collects these pairwise dot products:

\[
\mathbf R_{wb}^\mathsf T\mathbf R_{wb}
=
\begin{bmatrix}
\mathbf r_1^\mathsf T\mathbf r_1&
\mathbf r_1^\mathsf T\mathbf r_2&
\mathbf r_1^\mathsf T\mathbf r_3\\
\mathbf r_2^\mathsf T\mathbf r_1&
\mathbf r_2^\mathsf T\mathbf r_2&
\mathbf r_2^\mathsf T\mathbf r_3\\
\mathbf r_3^\mathsf T\mathbf r_1&
\mathbf r_3^\mathsf T\mathbf r_2&
\mathbf r_3^\mathsf T\mathbf r_3
\end{bmatrix}
=
\mathbf I_3.
\]

A real square matrix satisfying
\(\mathbf R^\mathsf T\mathbf R=\mathbf I_3\) is \textbf{orthogonal}.
This condition permits both right-handed orientations and reflections.
For a right-handed body frame,
\(\mathbf r_1\times\mathbf r_2=\mathbf r_3\), so the
scalar-triple-product formula gives

\[
\det(\mathbf R_{wb})
=
\mathbf r_3^\mathsf T
\left(
\mathbf r_1\times\mathbf r_2
\right)
=
\mathbf r_3^\mathsf T\mathbf r_3
=
1.
\]

The \textbf{special orthogonal group in three dimensions}, read ``S O
three,'' is

\[
\boxed{
SO(3)
\mathrel{:=}
\left\{
\mathbf R\in\mathbb R^{3\times3}
\;\middle|\;
\mathbf R^\mathsf T\mathbf R=\mathbf I_3,
\quad
\det(\mathbf R)=1
\right\}
}.
\]

The word \emph{orthogonal} refers to the orthonormal columns, while
\emph{special} selects determinant \(+1\) and excludes
determinant-\(-1\) reflections. The group operation is matrix
multiplication. For any \(\mathbf R_1,\mathbf R_2\in SO(3)\),

\[
\begin{aligned}
\left(\mathbf R_1\mathbf R_2\right)^\mathsf T
\left(\mathbf R_1\mathbf R_2\right)
&=
\mathbf R_2^\mathsf T
\mathbf R_1^\mathsf T
\mathbf R_1
\mathbf R_2
=
\mathbf I_3,\\
\det\left(\mathbf R_1\mathbf R_2\right)
&=
\det(\mathbf R_1)\det(\mathbf R_2)
=
1.
\end{aligned}
\]

Thus, \(\mathbf R_1\mathbf R_2\in SO(3)\), which is the \textbf{closure}
property. Matrix multiplication is associative, \(\mathbf I_3\) is the
identity orientation, and every element has an inverse in the set. From
\(\mathbf R^\mathsf T\mathbf R=\mathbf I_3\),

\[
\boxed{
\mathbf R^{-1}
=
\mathbf R^\mathsf T
}.
\]

A rotation matrix stores nine scalar entries, but the symmetric equation
\(\mathbf R^\mathsf T\mathbf R=\mathbf I_3\) supplies six independent
scalar constraints: three unit-length equations and three
pairwise-orthogonality equations. The determinant requirement selects
the right-handed component rather than removing another continuous
direction. Consequently,

\[
\boxed{
\dim SO(3)
=
9-6
=
3
}.
\]

Here, \(\dim SO(3)=3\) means that three independent continuous
coordinates are required near any selected orientation. This is the same
local meaning of dimension used in the earlier DOF definition.

\subsection{\texorpdfstring{Example: a positive quarter-turn about the
world
\(z_w\)-axis}{Example: a positive quarter-turn about the world z\_w-axis}}\label{example-a-positive-quarter-turn-about-the-world-z_w-axis}

Suppose frame \(\{b\}\) starts aligned with frame \(\{w\}\) and is
rotated through \(\pi/2\,\mathrm{rad}\) about the positive \(z_w\)-axis
according to the right-hand rule. Its body axes, expressed in world
coordinates, become

\[
{}^{w}\mathbf e_{x,b}
=
\begin{bmatrix}
0\\1\\0
\end{bmatrix},
\qquad
{}^{w}\mathbf e_{y,b}
=
\begin{bmatrix}
-1\\0\\0
\end{bmatrix},
\qquad
{}^{w}\mathbf e_{z,b}
=
\begin{bmatrix}
0\\0\\1
\end{bmatrix}.
\]

Placing these axes into the columns gives

\[
\boxed{
\mathbf R_{wb}
=
\begin{bmatrix}
0&-1&0\\
1&0&0\\
0&0&1
\end{bmatrix}
\in SO(3)
}.
\]

Direct calculation gives
\(\mathbf R_{wb}^\mathsf T\mathbf R_{wb}=\mathbf I_3\) and
\(\det(\mathbf R_{wb})=1\). A vector pointing along the positive body
\(x_b\)-axis has \({}^{b}\mathbf v=[1\;0\;0]^\mathsf T\) in arbitrary
vector units, so

\[
{}^{w}\mathbf v
=
\mathbf R_{wb}{}^{b}\mathbf v
=
\begin{bmatrix}
0\\1\\0
\end{bmatrix}.
\]

The result agrees with the first column of \(\mathbf R_{wb}\): the
positive body \(x_b\)-axis points along the positive world \(y_w\)-axis.

\section*{\texorpdfstring{Quick test B: spatial orientation and
\(SO(3)\)}{Quick test B: spatial orientation and SO(3)}}\label{quick-test-b-spatial-orientation-and-so3}
\addcontentsline{toc}{section}{Quick test B: spatial orientation and
\(SO(3)\)}

\markright{Quick test B: spatial orientation and \(SO(3)\)}

\begin{enumerate}
\def\labelenumi{\arabic{enumi}.}
\tightlist
\item
  What does column \(j\) of \(\mathbf R_{wb}\) represent?
\item
  Starting from orthonormal body axes, derive
  \(\mathbf R_{wb}^\mathsf T\mathbf R_{wb}=\mathbf I_3\).
\item
  Why is \(\det(\mathbf R)=1\) required in addition to
  \(\mathbf R^\mathsf T\mathbf R=\mathbf I_3\)?
\item
  Why does a rotation matrix store nine entries but represent only three
  rotational degrees of freedom?
\item
  For the quarter-turn matrix above, calculate \(\mathbf R_{bw}\) and
  use it to express \({}^{w}\mathbf v=[0\;1\;0]^\mathsf T\) in body
  coordinates.
\end{enumerate}

\subsection*{Answers and hints}\label{answers-and-hints-16}
\addcontentsline{toc}{subsection}{Answers and hints}

\begin{enumerate}
\def\labelenumi{\arabic{enumi}.}
\tightlist
\item
  Columns \(1\), \(2\), and \(3\) are respectively the world-frame
  coordinate columns of the body \(x_b\)-, \(y_b\)-, and \(z_b\)-axes.
\item
  Entry \((i,j)\) of \(\mathbf R_{wb}^\mathsf T\mathbf R_{wb}\) is
  \(\mathbf r_i^\mathsf T\mathbf r_j\). It equals \(1\) when \(i=j\) and
  \(0\) otherwise, producing \(\mathbf I_3\).
\item
  Orthogonality alone permits determinant \(-1\) reflections. The
  determinant-\(1\) condition selects right-handed rigid-body
  orientations.
\item
  Orthonormality imposes six independent scalar constraints on the nine
  entries: three unit-length constraints and three
  pairwise-orthogonality constraints. Thus, \(9-6=3\) independent
  continuous values remain.
\item
  \(\mathbf R_{bw}=\mathbf R_{wb}^{-1}=\mathbf R_{wb}^\mathsf T\).
  Therefore,
  \({}^{b}\mathbf v=\mathbf R_{bw}{}^{w}\mathbf v=[1\;0\;0]^\mathsf T\).
\end{enumerate}

\section{Coordinate-axis rotations and
composition}\label{sec-so3-axis-rotations-composition}

\subsection{Constructing coordinate-axis rotation
matrices}\label{sec-so3-active-axis-rotations}

Begin with the coordinate-change equation derived in
Section~\ref{sec-so3-rotation-construction}:

\[
\boxed{
{}^{w}\mathbf v
=
\mathbf R_{wb}\,{}^{b}\mathbf v
}.
\]

Here, \(\mathbf v\) is one unchanged geometric vector, while the two
columns are its coordinates in body frame \(\{b\}\) and world frame
\(\{w\}\). The same rotation matrix was defined from the world-frame
coordinate columns of the body axes:

\[
\boxed{
\mathbf R_{wb}
=
\begin{bmatrix}
{}^{w}\mathbf e_{x,b}&
{}^{w}\mathbf e_{y,b}&
{}^{w}\mathbf e_{z,b}
\end{bmatrix}
}.
\]

The task is therefore to calculate these three coordinate columns from
the geometry and place them into the matrix. Multiplication by standard
basis columns is not needed for this construction.

Suppose frames \(\{b\}\) and \(\{w\}\) are initially aligned, and then
frame \(\{b\}\) is rotated through \(\theta\in\mathbb R\) radians about
the positive world \(x_w\)-axis. A positive rotation follows the
right-hand rule. Define

\[
c_\theta\mathrel{:=}\cos\theta,
\qquad
s_\theta\mathrel{:=}\sin\theta.
\]

The \(x_b\)-axis remains aligned with \(x_w\), so its world coordinates
are \([1\;0\;0]^\mathsf T\). The \(y_b\)- and \(z_b\)-axes rotate in the
world \(y_wz_w\)-plane. Their projections onto the world axes give

\[
\begin{aligned}
{}^{w}\mathbf e_{x,b}
=
\begin{bmatrix}
1\\0\\0
\end{bmatrix}, \quad
{}^{w}\mathbf e_{y,b}
=
\begin{bmatrix}
0\\c_\theta\\s_\theta
\end{bmatrix}, \quad
{}^{w}\mathbf e_{z,b}
=
\begin{bmatrix}
0\\-s_\theta\\c_\theta
\end{bmatrix}.
\end{aligned}
\]

By definition, the first, second, and third columns of the orientation
matrix are the world-frame coordinates of the \(x_b\)-, \(y_b\)-, and
\(z_b\)-axes, respectively. Substituting the three columns above gives

\[
\mathbf R_x(\theta)
=
\begin{bmatrix}
{}^{w}\mathbf e_{x,b}&
{}^{w}\mathbf e_{y,b}&
{}^{w}\mathbf e_{z,b}
\end{bmatrix}
=
\left[
\begin{array}{c|c|c}
1&0&0\\
0&c_\theta&-s_\theta\\
0&s_\theta&c_\theta
\end{array}
\right].
\]

Reading the array by columns shows the substitution directly:
\([1\;0\;0]^\mathsf T\) is column one,
\([0\;c_\theta\;s_\theta]^\mathsf T\) is column two, and
\([0\;-s_\theta\;c_\theta]^\mathsf T\) is column three. Therefore,

\[
\boxed{
\mathbf R_x(\theta)
=
\begin{bmatrix}
1&0&0\\
0&c_\theta&-s_\theta\\
0&s_\theta&c_\theta
\end{bmatrix}
}.
\]

The column placement can be verified by multiplying
\(\mathbf R_x(\theta)\) by each standard basis column:

\[
\begin{aligned}
\begin{bmatrix}
1\\0\\0
\end{bmatrix}
&=
\mathbf R_x(\theta)
\begin{bmatrix}
1\\0\\0
\end{bmatrix},\\[4pt]
\begin{bmatrix}
0\\c_\theta\\s_\theta
\end{bmatrix}
&=
\mathbf R_x(\theta)
\begin{bmatrix}
0\\1\\0
\end{bmatrix},\\[4pt]
\begin{bmatrix}
0\\-s_\theta\\c_\theta
\end{bmatrix}
&=
\mathbf R_x(\theta)
\begin{bmatrix}
0\\0\\1
\end{bmatrix}.
\end{aligned}
\]

The three input columns select columns one, two, and three of
\(\mathbf R_x(\theta)\), respectively. These products verify the
assembled matrix; the preceding geometric projection determined the
column values.

For a positive rotation about the world \(y_w\)-axis, the \(y_b\)-axis
remains fixed. The right-hand rule sends the positive \(x_b\)-axis
toward negative \(z_w\) and the positive \(z_b\)-axis toward positive
\(x_w\). Their world-frame coordinate columns are therefore

\[
\begin{aligned}
{}^{w}\mathbf e_{x,b}
=
\begin{bmatrix}
c_\theta\\0\\-s_\theta
\end{bmatrix},\quad
{}^{w}\mathbf e_{y,b}
=
\begin{bmatrix}
0\\1\\0
\end{bmatrix},\quad
{}^{w}\mathbf e_{z,b}
=
\begin{bmatrix}
s_\theta\\0\\c_\theta
\end{bmatrix}.
\end{aligned}
\]

For a positive rotation about the world \(z_w\)-axis, the \(z_b\)-axis
remains fixed, the positive \(x_b\)-axis rotates toward positive
\(y_w\), and the positive \(y_b\)-axis rotates toward negative \(x_w\).
Hence,

\[
\begin{aligned}
{}^{w}\mathbf e_{x,b}
=
\begin{bmatrix}
c_\theta\\s_\theta\\0
\end{bmatrix},\quad
{}^{w}\mathbf e_{y,b}
=
\begin{bmatrix}
-s_\theta\\c_\theta\\0
\end{bmatrix},\quad
{}^{w}\mathbf e_{z,b}
=
\begin{bmatrix}
0\\0\\1
\end{bmatrix}.
\end{aligned}
\]

Placing each set of three world-frame body-axis coordinates into the
corresponding matrix columns gives

\[
\boxed{
\mathbf R_y(\theta)
=
\begin{bmatrix}
c_\theta&0&s_\theta\\
0&1&0\\
-s_\theta&0&c_\theta
\end{bmatrix},
\qquad
\mathbf R_z(\theta)
=
\begin{bmatrix}
c_\theta&-s_\theta&0\\
s_\theta&c_\theta&0\\
0&0&1
\end{bmatrix}
}.
\]

Each matrix has right-handed orthonormal columns and therefore belongs
to \(SO(3)\).

\subsubsection{The active-rotation
interpretation}\label{the-active-rotation-interpretation}

The same numerical matrices can also be used as \textbf{active rotation
operators}. An active rotation changes a geometric vector while the
frame used to describe it remains fixed. Let \(\mathbf v\) denote the
vector before rotation and \(\mathbf v'\) the vector after rotation.
Their world-frame coordinate columns satisfy

\[
\boxed{
{}^{w}\mathbf v'
=
\mathbf R_i(\theta)\,{}^{w}\mathbf v,
\qquad
i\in\{x,y,z\}
}.
\]

The prime labels the vector after rotation and does not denote a
derivative. The coordinate-change equation describes one geometric
vector in two frames; the active-rotation equation describes two
geometric vectors in one fixed frame.

\subsection{Composition order: the rightmost rotation acts
first}\label{sec-so3-composition-order}

Let \(\mathbf Q_1,\mathbf Q_2\in SO(3)\) be two active rotation
operators whose inputs and outputs are expressed in the same fixed
frame. If \(\mathbf Q_1\) acts first and \(\mathbf Q_2\) acts second,
then

\[
\mathbf v'
=
\mathbf Q_1\mathbf v,
\qquad
\mathbf v''
=
\mathbf Q_2\mathbf v'
=
\mathbf Q_2\mathbf Q_1\mathbf v.
\]

Thus, in the product \(\mathbf Q_2\mathbf Q_1\mathbf v\), the rightmost
matrix \(\mathbf Q_1\) acts first. This is a consequence of function
composition, not a special convention invented for rotations.

Spatial rotations generally do not commute. To demonstrate this, define

\[
\mathbf R_x
\mathrel{:=}
\mathbf R_x\!\left(\frac{\pi}{2}\right),
\qquad
\mathbf R_y
\mathrel{:=}
\mathbf R_y\!\left(\frac{\pi}{2}\right),
\qquad
\mathbf e_z
=
\begin{bmatrix}
0\\0\\1
\end{bmatrix}.
\]

Rotating first about \(x\) and then about \(y\) gives

\[
\mathbf R_y\mathbf R_x\mathbf e_z
=
\mathbf R_y
\begin{bmatrix}
0\\-1\\0
\end{bmatrix}
=
\begin{bmatrix}
0\\-1\\0
\end{bmatrix}.
\]

Reversing the order gives

\[
\mathbf R_x\mathbf R_y\mathbf e_z
=
\mathbf R_x
\begin{bmatrix}
1\\0\\0
\end{bmatrix}
=
\begin{bmatrix}
1\\0\\0
\end{bmatrix}.
\]

The outputs differ, so

\[
\boxed{
\mathbf R_y\mathbf R_x
\ne
\mathbf R_x\mathbf R_y
}.
\]

This calculation assumes that both rotations are active rotations about
axes of the same fixed frame. Changing which frame supplies the rotation
axes changes the meaning of the sequence.

\subsection{Fixed-frame and body-frame
increments}\label{sec-so3-fixed-body-increments}

Consider a rigid body moving while the world frame \(\{w\}\) remains
fixed. Before the body turns, its three body-fixed axes point in
particular world-frame directions. The body then turns, carrying all its
body-fixed points and axes with it, and the axes point in new
world-frame directions afterward.

An \textbf{orientation} describes this rotational state at one instant;
it answers how the body axes currently point relative to a reference
frame. A \textbf{turn} is the physical rotational motion that changes
one orientation into another over a time interval. Let frame \(\{b\}\)
denote the body axes before the turn and frame \(\{b'\}\) denote the
same body axes afterward. Then \(\mathbf R_{wb}\in SO(3)\) is the
orientation before the turn, whereas \(\mathbf R_{wb'}\in SO(3)\) is the
orientation afterward.

A rotation matrix \(\mathbf Q\in SO(3)\) that describes only the change
from the old orientation to the new orientation is called an
\textbf{incremental rotation}. The word \emph{incremental} describes the
matrix's role as an orientation change; \(\mathbf Q\) is still an
ordinary rotation matrix, and the turn need not be infinitesimally
small. Conceptually,

\[
\boxed{
\text{new orientation}
=
\text{rotation change}
\mathbin{\circ}
\text{old orientation}
},
\]

where \(\circ\) denotes composition: apply the old orientation first and
then apply the rotation change. Matrix multiplication will represent
this composition, but the side on which \(\mathbf Q\) multiplies depends
on the coordinate frame used to describe the change.

First describe the change using world-frame coordinates and denote it by
\(\mathbf Q_w\in SO(3)\). The subscript \(w\) means that \(\mathbf Q_w\)
acts on world-coordinate columns. For any direction fixed in the body,
if \({}^{w}\mathbf v_{\mathrm{old}}\in\mathbb R^3\) and
\({}^{w}\mathbf v_{\mathrm{new}}\in\mathbb R^3\) are its
world-coordinate columns before and after the turn, respectively, then
the definition of \(\mathbf Q_w\) is

\[
{}^{w}\mathbf v_{\mathrm{new}}
=
\mathbf Q_w{}^{w}\mathbf v_{\mathrm{old}}.
\]

Write the current orientation as

\[
\mathbf R_{wb}
=
\begin{bmatrix}
\mathbf r_1&\mathbf r_2&\mathbf r_3
\end{bmatrix},
\]

where \(\mathbf r_1,\mathbf r_2,\mathbf r_3\in\mathbb R^3\) are the
current world-frame coordinate columns of the body \(x_b\)-, \(y_b\)-,
and \(z_b\)-axes. Applying \(\mathbf Q_w\) to the physical body rotates
each of these three axes, so the updated body-axis columns are

\[
{}^{w}\mathbf e_{x,b'}
=
\mathbf Q_w\mathbf r_1,
\qquad
{}^{w}\mathbf e_{y,b'}
=
\mathbf Q_w\mathbf r_2,
\qquad
{}^{w}\mathbf e_{z,b'}
=
\mathbf Q_w\mathbf r_3.
\]

The updated orientation matrix is constructed by placing these new
body-axis coordinates into its columns:

\[
\begin{aligned}
\mathbf R_{wb'}
&=
\begin{bmatrix}
{}^{w}\mathbf e_{x,b'}&
{}^{w}\mathbf e_{y,b'}&
{}^{w}\mathbf e_{z,b'}
\end{bmatrix}\\
&=
\begin{bmatrix}
\mathbf Q_w\mathbf r_1&
\mathbf Q_w\mathbf r_2&
\mathbf Q_w\mathbf r_3
\end{bmatrix}\\
&=
\mathbf Q_w
\begin{bmatrix}
\mathbf r_1&\mathbf r_2&\mathbf r_3
\end{bmatrix}\\
&=
\boxed{\mathbf Q_w\mathbf R_{wb}}.
\end{aligned}
\]

Thus, an increment described using world-frame axes
\textbf{premultiplies} the current orientation, where \emph{premultiply}
means multiply on the left. This product answers the practical question
posed at the start: given the old orientation and an additional
world-described rotation, what is the new orientation?

As an alternative, suppose the increment is described relative to the
current body axes. Let
\(\mathbf Q_b\mathrel{:=}\mathbf R_{bb'}\in SO(3)\) describe the
orientation of the new body frame \(\{b'\}\) relative to the old body
frame \(\{b\}\). Coordinate-frame composition gives

\[
\boxed{
\mathbf R_{wb'}
=
\mathbf R_{wb}\mathbf R_{bb'}
=
\mathbf R_{wb}\mathbf Q_b
}.
\]

The matrices in this composition have different roles:

\begin{itemize}
\tightlist
\item
  \(\mathbf R_{wb'}\) is the resulting orientation of the body relative
  to the world frame.
\item
  \(\mathbf R_{wb}\) is the body's orientation relative to the world
  frame before the turn.
\item
  \(\mathbf R_{bb'}=\mathbf Q_b\) is the orientation of the new body
  frame \(\{b'\}\) relative to the old body frame \(\{b\}\); it
  represents the incremental rotation expressed along the old body axes.
\end{itemize}

Thus, composing the old world-relative orientation with the
body-relative orientation change produces the new world-relative
orientation.

The body-frame increment therefore \textbf{postmultiplies} the current
orientation, where \emph{postmultiply} means multiply on the right.
These rules can be remembered from their frame labels:

\[
\underbrace{\mathbf R_{wb}}_{b\rightarrow w}
\underbrace{\mathbf R_{bb'}}_{b'\rightarrow b}
=
\underbrace{\mathbf R_{wb'}}_{b'\rightarrow w}.
\]

The two update formulas are not interchangeable. Premultiplication
applies a rotation about an axis fixed in the world description;
postmultiplication applies a rotation about an axis attached to the
current body description.

\section*{Quick test C: coordinate-axis rotations and
composition}\label{quick-test-c-coordinate-axis-rotations-and-composition}
\addcontentsline{toc}{section}{Quick test C: coordinate-axis rotations
and composition}

\markright{Quick test C: coordinate-axis rotations and composition}

\begin{enumerate}
\def\labelenumi{\arabic{enumi}.}
\tightlist
\item
  Starting from \({}^{w}\mathbf v=\mathbf R_{wb}{}^{b}\mathbf v\),
  derive the columns of \(\mathbf R_y(\theta)\) from the body axes
  expressed in world coordinates.
\item
  In \(\mathbf Q_3\mathbf Q_2\mathbf Q_1\mathbf v\), which rotation acts
  first?
\item
  Use \(\mathbf e_z\) to show that
  \(\mathbf R_y(\pi/2)\mathbf R_x(\pi/2)\ne\mathbf R_x(\pi/2)\mathbf R_y(\pi/2)\).
\item
  If an orientation increment is expressed about world-frame axes, on
  which side of \(\mathbf R_{wb}\) does it multiply?
\item
  If \(\mathbf Q_b=\mathbf R_{bb'}\) is an increment relative to the
  current body frame, derive the updated orientation \(\mathbf R_{wb'}\)
  from the frame labels.
\item
  Distinguish an active rotation from a coordinate change.
\end{enumerate}

\subsection*{Answers and hints}\label{answers-and-hints-17}
\addcontentsline{toc}{subsection}{Answers and hints}

\begin{enumerate}
\def\labelenumi{\arabic{enumi}.}
\tightlist
\item
  After a positive rotation about \(y_w\), the world-frame body-axis
  columns are \([c_\theta\;0\;-s_\theta]^\mathsf T\),
  \([0\;1\;0]^\mathsf T\), and \([s_\theta\;0\;c_\theta]^\mathsf T\).
  Placing them into the first, second, and third columns gives
  \(\mathbf R_y(\theta)\).
\item
  \(\mathbf Q_1\) acts first, followed by \(\mathbf Q_2\) and then
  \(\mathbf Q_3\).
\item
  The first product maps \(\mathbf e_z\) to \(-\mathbf e_y\), while the
  reversed product maps it to \(\mathbf e_x\); therefore, the operators
  are unequal.
\item
  It premultiplies: \(\mathbf R_{wb}'=\mathbf Q_w\mathbf R_{wb}\).
\item
  The labels compose as
  \((b\rightarrow w)(b'\rightarrow b)=(b'\rightarrow w)\), so
  \(\mathbf R_{wb'}=\mathbf R_{wb}\mathbf R_{bb'}=\mathbf R_{wb}\mathbf Q_b\).
\item
  An active rotation changes a geometric vector while retaining one
  coordinate frame. A coordinate change keeps the geometric vector fixed
  and changes the frame used to express its coordinate column.
\end{enumerate}

\section{Angular velocity as the rate of orientation
change}\label{sec-so3-angular-velocity}

\subsection{From a finite turn to an instantaneous
rate}\label{sec-so3-angular-velocity-definition}

The preceding section compared the body's orientation before and after a
finite turn. Motion also requires a rate that describes how rapidly that
orientation is changing at one instant. Consider a rigid body whose
orientation relative to the fixed world frame is the differentiable
function \(\mathbf R_{wb}(t)\in SO(3)\), where \(t\) is time measured in
seconds.

Over a time interval from \(t\) to \(t+\Delta t\), where \(\Delta t>0\)
is measured in seconds, the body undergoes a relative rotation from
\(\mathbf R_{wb}(t)\) to \(\mathbf R_{wb}(t+\Delta t)\). \textbf{Euler's
rotation theorem}, used here as a stated geometric result, says that
every relative orientation change in \(SO(3)\) can be represented as one
rotation through an angle about an axis direction. Accordingly, let
\(\mathbf a_w(t,\Delta t)\in\mathbb R^3\) be a unit rotation-axis column
and let \(\Delta\theta(t,\Delta t)\in\mathbb R\) be the signed rotation
angle measured in radians. The subscript \(w\) states that the axis
components are expressed along the world axes, and

\[
\left\|\mathbf a_w(t,\Delta t)\right\|_2=1.
\]

The \textbf{rotation axis} in an orientation description specifies only
a direction \(\mathbf a_w\). It does not identify a particular line in
physical space. To specify a particular line, also choose the world
position \(\mathbf q_w\in\mathbb R^3\) of one point on it, measured in
metres. The line is then

\[
\mathcal L
\mathrel{:=}
\left\{
\mathbf q_w+\lambda\mathbf a_w
\;\middle|\;
\lambda\in\mathbb R
\right\}
\subset\mathbb R^3,
\]

where \(\mathcal L\) denotes the axis line and \(\lambda\) is a scalar
distance measured in metres. Because \(\mathbf R_{wb}\) describes
orientation only, it contains the axis direction associated with the
rotation but not the location of a physical pivot line.

The sign of \(\Delta\theta\) follows the right-hand rule about
\(\mathbf a_w\). Replacing \(\mathbf a_w\) by \(-\mathbf a_w\) and
\(\Delta\theta\) by \(-\Delta\theta\) describes the same turn, so their
product is unchanged.

The vector

\[
\mathbf a_w(t,\Delta t)
\frac{\Delta\theta(t,\Delta t)}{\Delta t}
\]

describes the rotation-change rate over that interval: its direction
gives the rotation axis and its signed magnitude gives the angle change
per unit time. Taking the limit as the interval shrinks defines the
body's \textbf{angular velocity relative to the world frame, expressed
in world coordinates}:

\[
\boxed{
\boldsymbol\omega_w(t)
\mathrel{:=}
\lim_{\Delta t\to0^+}
\mathbf a_w(t,\Delta t)
\frac{\Delta\theta(t,\Delta t)}{\Delta t}
\in\mathbb R^3
}.
\]

The unit axis \(\mathbf a_w\) is dimensionless, while
\(\Delta\theta/\Delta t\) has units of \(\mathrm{rad\,s^{-1}}\).
Therefore, \(\boldsymbol\omega_w\) has units of
\(\mathrm{rad\,s^{-1}}\). The limit is assumed to exist because
\(\mathbf R_{wb}(t)\) is differentiable.

Orientation and angular velocity describe different physical quantities.
The matrix \(\mathbf R_{wb}(t)\) is the rotational state at time \(t\),
whereas \(\boldsymbol\omega_w(t)\) is the instantaneous rate at which
that state is changing. Knowing the orientation at one instant does not
determine the angular velocity: a stationary body and a turning body can
momentarily have the same orientation but different angular velocities.

To connect angular velocity to the orientation matrix, write its three
columns explicitly:

\[
\boxed{
\mathbf R_{wb}(t)
=
\begin{bmatrix}
\mathbf r_1(t)&\mathbf r_2(t)&\mathbf r_3(t)
\end{bmatrix}
}.
\]

For each \(i\in\{1,2,3\}\), the dimensionless column
\(\mathbf r_i(t)\in\mathbb R^3\) is the world-coordinate representation
of one unit body axis fixed in the rigid body. Thus,
\(\|\mathbf r_i(t)\|_2=1\). The geometric construction below selects one
column \(\mathbf r_i\); Figure~\ref{fig-so3-body-axis-angular-velocity}
illustrates the choice \(i=1\) and draws the other two columns as dashed
vectors.

When \(\boldsymbol\omega_w\ne\mathbf 0_3\), define its unit direction by
\(\mathbf a_w\mathrel{:=}\boldsymbol\omega_w/\|\boldsymbol\omega_w\|_2\)
and let \(\varphi_i\in[0,\pi]\) be the angle between \(\mathbf a_w\) and
\(\mathbf r_i\). Decompose the selected column into components parallel
and perpendicular to the angular-velocity direction:

\[
\mathbf r_{i,\parallel}
\mathrel{:=}
(\mathbf a_w^\mathsf T\mathbf r_i)\mathbf a_w,
\qquad
\mathbf r_{i,\perp}
\mathrel{:=}
\mathbf r_i-\mathbf r_{i,\parallel}.
\]

For the drawing, place the tail of these direction vectors at a common
coordinate origin \(O\) and call the tip of \(\mathbf r_i\) point
\(P_i\). The point \(P_i\) is a graphical endpoint, not necessarily a
material point on the body. Let \(C_i\) be the tip of
\(\mathbf r_{i,\parallel}\). The segment from \(C_i\) to \(P_i\)
represents \(\mathbf r_{i,\perp}\) and is perpendicular to the
angular-velocity direction.

\begin{figure}

\centering{

\includegraphics[width=0.88\linewidth,height=\textheight,keepaspectratio]{index_files/mediabag/figures/spatial_angular_velocity_geometry.pdf}

}

\caption{\label{fig-so3-body-axis-angular-velocity}Angular velocity and
the motion of one orientation-matrix column. The solid vector
\(\mathbf r_1\) is the first column of \(\mathbf R_{wb}\); the dashed
vectors \(\mathbf r_2\) and \(\mathbf r_3\) are its other columns.}

\end{figure}%

The figure is a two-dimensional projection of three-dimensional
geometry. The orange radius vector \(\mathbf r_{1,\perp}\), whose length
is \(\rho_1\), is perpendicular to \(\boldsymbol\omega_w\), and
\(\dot{\mathbf r}_1\) is perpendicular to \(\mathbf r_{1,\perp}\) in
three dimensions. The right-angle marker at \(P_1\) therefore appears
skewed in the projection even though it denotes a true three-dimensional
right angle.

At the selected instant, the tip \(P_i\) has the same tangent direction
as a point moving on a circle centred at \(C_i\). The circle radius is
the length of the perpendicular component:

\[
\boxed{
\rho_i
=
\|\mathbf r_{i,\perp}\|_2
=
\|\mathbf r_i\|_2\sin\varphi_i
=
\sin\varphi_i
}.
\]

The last equality uses \(\|\mathbf r_i\|_2=1\). Multiplying the radius
by the angular speed gives the speed of the tip:

\[
\|\dot{\mathbf r}_i\|_2
=
\|\boldsymbol\omega_w\|_2\rho_i
=
\|\boldsymbol\omega_w\|_2
\|\mathbf r_i\|_2
\sin\varphi_i.
\]

This is exactly the magnitude of
\(\boldsymbol\omega_w\times\mathbf r_i\), and the cross product's
right-hand-rule direction is the tangent direction shown in the figure.
When \(\boldsymbol\omega_w=\mathbf 0_3\), both the cross product and the
column's rate are zero. Therefore, for each of the three columns,

\[
\boxed{
\dot{\mathbf r}_i(t)
=
\boldsymbol\omega_w(t)\times\mathbf r_i(t),
\qquad
i\in\{1,2,3\}
}.
\]

These equations describe how the three columns of the orientation matrix
change. To combine them into one equation for \(\dot{\mathbf R}_{wb}\),
the cross product can be represented as a matrix multiplication.

\subsection{Cross-product matrices and the orientation
derivative}\label{sec-so3-cross-product-matrix}

Let

\[
\mathbf x
=
\begin{bmatrix}
x_1\\x_2\\x_3
\end{bmatrix}
\in\mathbb R^3.
\]

The \textbf{cross-product matrix} of \(\mathbf x\) is defined by

\[
\boxed{
[\mathbf x]_\times
\mathrel{:=}
\begin{bmatrix}
0&-x_3&x_2\\
x_3&0&-x_1\\
-x_2&x_1&0
\end{bmatrix}
\in\mathbb R^{3\times3}
}.
\]

The notation \([\,\cdot\,]_\times\) means ``form the matrix that
performs a cross product with the enclosed vector.'' For any
\(\mathbf y=[y_1\;y_2\;y_3]^\mathsf T\in\mathbb R^3\), direct
multiplication gives

\[
\begin{aligned}
\left[\mathbf x\right]_\times \mathbf y
&=
\begin{bmatrix}
-x_3y_2+x_2y_3\\
x_3y_1-x_1y_3\\
-x_2y_1+x_1y_2
\end{bmatrix}\\
&=
\boxed{
\mathbf x\times\mathbf y
}.
\end{aligned}
\]

Thus, \([\mathbf x]_\times\) is a linear mapping from an input vector
\(\mathbf y\in\mathbb R^3\) to the output vector
\(\mathbf x\times\mathbf y\in\mathbb R^3\). Its entries have the same
units as the components of \(\mathbf x\).

Transposing the matrix changes the sign of every off-diagonal entry:

\[
[\mathbf x]_\times^\mathsf T
=
-[\mathbf x]_\times.
\]

A square matrix \(\mathbf A\) satisfying
\(\mathbf A^\mathsf T=-\mathbf A\) is called \textbf{skew-symmetric}.
The set of all real \(3\times3\) skew-symmetric matrices is

\[
\boxed{
\mathfrak{so}(3)
\mathrel{:=}
\left\{
\mathbf A\in\mathbb R^{3\times3}
\;\middle|\;
\mathbf A^\mathsf T=-\mathbf A
\right\}
=
\left\{
[\mathbf x]_\times
\;\middle|\;
\mathbf x\in\mathbb R^3
\right\}
}.
\]

The second equality follows from the structure imposed by skew symmetry.
The diagonal entries must be zero, and each entry below the diagonal
must be the negative of the corresponding entry above it. Hence, every
\(\mathbf A\in\mathfrak{so}(3)\) is determined by three scalars and can
be written uniquely as

\[
\mathbf A
=
\begin{bmatrix}
0&-x_3&x_2\\
x_3&0&-x_1\\
-x_2&x_1&0
\end{bmatrix}
=
[\mathbf x]_\times.
\]

As in the planar case developed in
Section~\ref{sec-se2-rotation-derivative-tangent}, \(\mathfrak{so}(3)\)
is the \textbf{Lie algebra} of \(SO(3)\): it is the vector space of
instantaneous rotation directions at the identity rotation
\(\mathbf I_3\). A finite rotation matrix belongs to \(SO(3)\), whereas
an angular-velocity matrix belongs to \(\mathfrak{so}(3)\) and describes
a rate rather than an orientation.

Apply the cross-product matrix to the three column-rate equations from
Section~\ref{sec-so3-angular-velocity-definition}:

\[
\begin{aligned}
\dot{\mathbf R}_{wb}
&=
\begin{bmatrix}
\dot{\mathbf r}_1&\dot{\mathbf r}_2&\dot{\mathbf r}_3
\end{bmatrix}\\
&=
\begin{bmatrix}
[\boldsymbol\omega_w]_\times\mathbf r_1&
[\boldsymbol\omega_w]_\times\mathbf r_2&
[\boldsymbol\omega_w]_\times\mathbf r_3
\end{bmatrix}\\
&=
[\boldsymbol\omega_w]_\times
\begin{bmatrix}
\mathbf r_1&\mathbf r_2&\mathbf r_3
\end{bmatrix}.
\end{aligned}
\]

Because the final matrix of columns is \(\mathbf R_{wb}\),

\[
\boxed{
\dot{\mathbf R}_{wb}
=
[\boldsymbol\omega_w]_\times\mathbf R_{wb}
}.
\]

Here, \(\mathbf R_{wb}\in SO(3)\) is dimensionless,
\(\boldsymbol\omega_w\in\mathbb R^3\) has units of
\(\mathrm{rad\,s^{-1}}\), and both \([\boldsymbol\omega_w]_\times\) and
\(\dot{\mathbf R}_{wb}\) have units of \(\mathrm{s^{-1}}\) because
radians are dimensionless in SI. This equation says that the
world-coordinate rate of every body-axis column is produced by the same
world-expressed angular velocity.

Right-multiplying by \(\mathbf R_{wb}^\mathsf T=\mathbf R_{wb}^{-1}\)
isolates the angular-velocity matrix:

\[
\begin{aligned}
\dot{\mathbf R}_{wb}\mathbf R_{wb}^\mathsf T
&=
[\boldsymbol\omega_w]_\times
\mathbf R_{wb}\mathbf R_{wb}^\mathsf T\\
&=
[\boldsymbol\omega_w]_\times\mathbf I_3.
\end{aligned}
\]

Therefore,

\[
\boxed{
[\boldsymbol\omega_w]_\times
=
\dot{\mathbf R}_{wb}\mathbf R_{wb}^\mathsf T
}.
\]

This formula recovers the angular velocity expressed along the world
axes from the orientation and its derivative.

The same physical angular velocity can instead be expressed along the
moving body axes. Define

\[
\boxed{
\boldsymbol\omega_b
\mathrel{:=}
\mathbf R_{wb}^\mathsf T\boldsymbol\omega_w
\in\mathbb R^3
}.
\]

The subscript \(b\) does not mean that the body rotates relative to
itself. It means that the body's angular velocity relative to the world
is resolved into components along \(x_b\), \(y_b\), and \(z_b\).

A proper rotation preserves cross products. Insert the identity
\(\mathbf I_3=\mathbf R_{wb}^\mathsf T\mathbf R_{wb}\) before
\(\mathbf y\), and then use
\((\mathbf R^\mathsf T\mathbf a)\times(\mathbf R^\mathsf T\mathbf b)=\mathbf R^\mathsf T(\mathbf a\times\mathbf b)\)
for \(\mathbf R\in SO(3)\). Consequently, for every
\(\mathbf y\in\mathbb R^3\),

\[
\begin{aligned}
\left[\boldsymbol\omega_b\right]_\times\mathbf y
&=
(\mathbf R_{wb}^\mathsf T\boldsymbol\omega_w)\times\mathbf y\\
&=
(\mathbf R_{wb}^\mathsf T\boldsymbol\omega_w)
\times
(\mathbf R_{wb}^\mathsf T\mathbf R_{wb}\mathbf y)\\
&=
\mathbf R_{wb}^\mathsf T
\left(
\boldsymbol\omega_w\times(\mathbf R_{wb}\mathbf y)
\right)\\
&=
\mathbf R_{wb}^\mathsf T
[\boldsymbol\omega_w]_\times
\mathbf R_{wb}\mathbf y.
\end{aligned}
\]

Because the two matrix expressions produce the same output for every
\(\mathbf y\), their matrices are equal:

\[
[\boldsymbol\omega_b]_\times
=
\mathbf R_{wb}^\mathsf T
[\boldsymbol\omega_w]_\times
\mathbf R_{wb}.
\]

Substituting
\([\boldsymbol\omega_w]_\times=\dot{\mathbf R}_{wb}\mathbf R_{wb}^\mathsf T\)
then gives

\[
\begin{aligned}
[\boldsymbol\omega_b]_\times
&=
\mathbf R_{wb}^\mathsf T
\dot{\mathbf R}_{wb}
\mathbf R_{wb}^\mathsf T
\mathbf R_{wb}\\
&=
\mathbf R_{wb}^\mathsf T\dot{\mathbf R}_{wb}.
\end{aligned}
\]

The two coordinate descriptions and the corresponding orientation-rate
equations are therefore

\[
\boxed{
\begin{aligned}
\left[\boldsymbol\omega_w\right]_\times
&=
\dot{\mathbf R}_{wb}\mathbf R_{wb}^\mathsf T,
&
\dot{\mathbf R}_{wb}
&=
[\boldsymbol\omega_w]_\times\mathbf R_{wb},\\
[\boldsymbol\omega_b]_\times
&=
\mathbf R_{wb}^\mathsf T\dot{\mathbf R}_{wb},
&
\dot{\mathbf R}_{wb}
&=
\mathbf R_{wb}[\boldsymbol\omega_b]_\times.
\end{aligned}
}.
\]

\subsubsection{Connection to finite fixed-frame and body-frame
increments}\label{connection-to-finite-fixed-frame-and-body-frame-increments}

The multiplication sides follow directly from the finite composition
rules in Section~\ref{sec-so3-fixed-body-increments}. To see this, let

\[
\mathbf R
\mathrel{:=}
\mathbf R_{wb}(t),
\qquad
\mathbf R^+
\mathrel{:=}
\mathbf R_{wb}(t+\Delta t),
\]

where \(\mathbf R\) is the orientation at time \(t\), \(\mathbf R^+\) is
the orientation after a short interval \(\Delta t>0\), and both matrices
belong to \(SO(3)\).

First express the orientation change using the fixed world axes. The
finite composition rule is

\[
\mathbf R^+
=
\mathbf Q_w\mathbf R,
\]

where \(\mathbf Q_w\in SO(3)\) is the incremental rotation expressed
using world coordinates. At \(\Delta t=0\), there is no orientation
change, so \(\mathbf Q_w(t,0)=\mathbf I_3\). Its derivative with respect
to the interval length at zero is the world-expressed angular-velocity
matrix \([\boldsymbol\omega_w(t)]_\times\). The first-order Taylor
expansion is therefore

\[
\mathbf Q_w
=
\mathbf I_3
+
\Delta t[\boldsymbol\omega_w]_\times
+
O(\Delta t^2).
\]

Substituting this expansion into the finite composition rule gives

\[
\begin{aligned}
\mathbf R^+
&=
\left(
\mathbf I_3
+
\Delta t[\boldsymbol\omega_w]_\times
+
O(\Delta t^2)
\right)\mathbf R\\
&=
\mathbf R
+
\Delta t[\boldsymbol\omega_w]_\times\mathbf R
+
O(\Delta t^2).
\end{aligned}
\]

Subtract \(\mathbf R\), divide by \(\Delta t\), and take the limit
\(\Delta t\to0\):

\[
\boxed{
\dot{\mathbf R}_{wb}
=
[\boldsymbol\omega_w]_\times\mathbf R_{wb}
}.
\]

Thus, a rotation change expressed using the fixed world axes
premultiplies the finite orientation, and its world-expressed
angular-velocity matrix likewise multiplies the orientation on the left.

Now express the same orientation change using the old body axes. The
finite composition rule is

\[
\mathbf R^+
=
\mathbf R\mathbf Q_b,
\]

where \(\mathbf Q_b\in SO(3)\) is the incremental rotation expressed
using body coordinates. Similarly, \(\mathbf Q_b(t,0)=\mathbf I_3\), and
its derivative with respect to the interval length at zero is
\([\boldsymbol\omega_b(t)]_\times\). Its first-order Taylor expansion is

\[
\mathbf Q_b
=
\mathbf I_3
+
\Delta t[\boldsymbol\omega_b]_\times
+
O(\Delta t^2).
\]

Substitution gives

\[
\begin{aligned}
\mathbf R^+
&=
\mathbf R
\left(
\mathbf I_3
+
\Delta t[\boldsymbol\omega_b]_\times
+
O(\Delta t^2)
\right)\\
&=
\mathbf R
+
\Delta t\mathbf R[\boldsymbol\omega_b]_\times
+
O(\Delta t^2).
\end{aligned}
\]

Taking the same difference-quotient limit gives

\[
\boxed{
\dot{\mathbf R}_{wb}
=
\mathbf R_{wb}[\boldsymbol\omega_b]_\times
}.
\]

Thus, a rotation change expressed using the moving body axes
postmultiplies the finite orientation, and its body-expressed
angular-velocity matrix likewise multiplies the orientation on the
right. Both equations describe the same physical turn; only the axes
used to express that turn differ.

{\def\LTcaptype{none} % do not increment counter
\begin{longtable}[]{@{}
  >{\raggedright\arraybackslash}p{(\linewidth - 4\tabcolsep) * \real{0.3333}}
  >{\raggedright\arraybackslash}p{(\linewidth - 4\tabcolsep) * \real{0.3333}}
  >{\raggedright\arraybackslash}p{(\linewidth - 4\tabcolsep) * \real{0.3333}}@{}}
\toprule\noalign{}
\begin{minipage}[b]{\linewidth}\raggedright
Coordinates used for the change
\end{minipage} & \begin{minipage}[b]{\linewidth}\raggedright
Finite composition
\end{minipage} & \begin{minipage}[b]{\linewidth}\raggedright
Instantaneous orientation rate
\end{minipage} \\
\midrule\noalign{}
\endhead
\bottomrule\noalign{}
\endlastfoot
Fixed world axes & \(\mathbf R^+=\mathbf Q_w\mathbf R\) &
\(\dot{\mathbf R}=[\boldsymbol\omega_w]_\times\mathbf R\) \\
Moving body axes & \(\mathbf R^+=\mathbf R\mathbf Q_b\) &
\(\dot{\mathbf R}=\mathbf R[\boldsymbol\omega_b]_\times\) \\
\end{longtable}
}

Finally, these rate matrices are necessarily skew-symmetric.
Differentiating the rotation constraint
\(\mathbf R_{wb}^\mathsf T\mathbf R_{wb}=\mathbf I_3\) gives

\[
\dot{\mathbf R}_{wb}^\mathsf T\mathbf R_{wb}
+
\mathbf R_{wb}^\mathsf T\dot{\mathbf R}_{wb}
=
\mathbf 0_{3\times3}.
\]

Therefore,
\((\mathbf R_{wb}^\mathsf T\dot{\mathbf R}_{wb})^\mathsf T=-\mathbf R_{wb}^\mathsf T\dot{\mathbf R}_{wb}\),
confirming that \([\boldsymbol\omega_b]_\times\in\mathfrak{so}(3)\).
Differentiating \(\mathbf R_{wb}\mathbf R_{wb}^\mathsf T=\mathbf I_3\)
similarly confirms that
\([\boldsymbol\omega_w]_\times=\dot{\mathbf R}_{wb}\mathbf R_{wb}^\mathsf T\in\mathfrak{so}(3)\).
These products remove the current finite orientation and express its
instantaneous change at the identity rotation.

\section*{Quick test D: angular
velocity}\label{quick-test-d-angular-velocity}
\addcontentsline{toc}{section}{Quick test D: angular velocity}

\markright{Quick test D: angular velocity}

\begin{enumerate}
\def\labelenumi{\arabic{enumi}.}
\tightlist
\item
  What physical information is contained in the direction and magnitude
  of \(\boldsymbol\omega_w\)?
\item
  Why does knowing \(\mathbf R_{wb}(t)\) at one instant not determine
  \(\boldsymbol\omega_w(t)\)?
\item
  Explain geometrically why a body-axis column satisfies
  \(\dot{\mathbf r}_i=\boldsymbol\omega_w\times\mathbf r_i\).
\item
  Suppose
  \(\boldsymbol\omega_w=[0\;0\;2]^\mathsf T\,\mathrm{rad\,s^{-1}}\) and
  \(\mathbf r_i=[1\;0\;0]^\mathsf T\). Calculate \(\dot{\mathbf r}_i\)
  and interpret its direction.
\item
  Construct \([\mathbf x]_\times\) for \(\mathbf x=[1\;2\;3]^\mathsf T\)
  and verify that it is skew-symmetric.
\item
  Starting from
  \(\dot{\mathbf R}_{wb}=[\boldsymbol\omega_w]_\times\mathbf R_{wb}\),
  isolate \([\boldsymbol\omega_w]_\times\).
\item
  Do \(\boldsymbol\omega_w\) and \(\boldsymbol\omega_b\) describe
  different physical turns? Explain their relationship.
\item
  Why is \(\mathbf R_{wb}^\mathsf T\dot{\mathbf R}_{wb}\) skew-symmetric
  rather than a rotation matrix?
\end{enumerate}

\subsection*{Answers and hints}\label{answers-and-hints-18}
\addcontentsline{toc}{subsection}{Answers and hints}

\begin{enumerate}
\def\labelenumi{\arabic{enumi}.}
\tightlist
\item
  Its direction gives the instantaneous rotation axis according to the
  right-hand rule, and its magnitude gives the angular speed in
  \(\mathrm{rad\,s^{-1}}\).
\item
  Orientation is a state, whereas angular velocity is its rate of
  change. Different orientation trajectories can pass through the same
  \(\mathbf R_{wb}(t)\) with different derivatives.
\item
  The tip of \(\mathbf r_i\) has a tangent velocity about the
  instantaneous rotation-axis direction. The cross product has that
  tangent direction and magnitude
  \(\|\boldsymbol\omega_w\|_2\|\mathbf r_i\|_2\sin\varphi_i\).
\item
  \(\dot{\mathbf r}_i=\boldsymbol\omega_w\times\mathbf r_i=[0\;2\;0]^\mathsf T\,\mathrm{s^{-1}}\).
  The direction initially points toward positive \(y_w\), as expected
  for a positive turn about \(z_w\).
\item
  \([\mathbf x]_\times=\begin{bmatrix}0&-3&2\\3&0&-1\\-2&1&0\end{bmatrix}\),
  and direct transposition gives
  \([\mathbf x]_\times^\mathsf T=-[\mathbf x]_\times\).
\item
  Right-multiply by \(\mathbf R_{wb}^\mathsf T\) and use
  \(\mathbf R_{wb}\mathbf R_{wb}^\mathsf T=\mathbf I_3\) to obtain
  \([\boldsymbol\omega_w]_\times=\dot{\mathbf R}_{wb}\mathbf R_{wb}^\mathsf T\).
\item
  No.~They describe the same angular velocity relative to the world but
  resolve it along different axes:
  \(\boldsymbol\omega_b=\mathbf R_{wb}^\mathsf T\boldsymbol\omega_w\).
\item
  Differentiating \(\mathbf R_{wb}^\mathsf T\mathbf R_{wb}=\mathbf I_3\)
  shows that
  \((\mathbf R_{wb}^\mathsf T\dot{\mathbf R}_{wb})^\mathsf T=-\mathbf R_{wb}^\mathsf T\dot{\mathbf R}_{wb}\).
  It represents an instantaneous rotation rate in \(\mathfrak{so}(3)\),
  not a finite orientation in \(SO(3)\).
\end{enumerate}

\section{Finite rotations from axis--angle
coordinates}\label{sec-so3-axis-angle-exponential}

\subsection{The problem: converting a rotation rate into a finite
turn}\label{the-problem-converting-a-rotation-rate-into-a-finite-turn}

The angular-velocity equations in Section~\ref{sec-so3-angular-velocity}
describe an instantaneous orientation rate. A simulation, estimator, or
controller also needs a finite orientation change over a nonzero
interval. This section solves the forward problem: given a rotation-axis
direction and a signed rotation angle, construct the corresponding
matrix in \(SO(3)\).

Let \(\mathbf a_w\in\mathbb R^3\) be a unit vector expressed in world
coordinates,

\[
\|\mathbf a_w\|_2=1,
\]

and let \(\alpha\in\mathbb R\) be a signed rotation angle measured in
radians. The pair \((\mathbf a_w,\alpha)\) is an \textbf{axis--angle
representation} of the intended orientation change: \(\mathbf a_w\)
gives the positive axis direction, and the sign of \(\alpha\) selects
the direction of rotation by the right-hand rule. This orientation-only
description specifies an axis direction through the coordinate origin;
locating a physical pivot line would additionally require a point, as
explained in Section~\ref{sec-so3-angular-velocity-definition}.

The three-component column

\[
\boxed{
\boldsymbol\varphi_w
\mathrel{:=}
\alpha\mathbf a_w
\in\mathbb R^3
}
\]

is called the \textbf{rotation vector}, or the \textbf{exponential
coordinates of the rotation}. Its direction gives the rotation axis and
its norm gives the angle magnitude:

\[
\|\boldsymbol\varphi_w\|_2
=
|\alpha|.
\]

The word \emph{exponential} does not refer to the definition
\(\boldsymbol\varphi_w=\alpha\mathbf a_w\) itself. It anticipates the
construction derived below: the skew-symmetric matrix
\([\boldsymbol\varphi_w]_\times\) is supplied to the matrix exponential
to recover the finite rotation
\(\mathbf Q_w=\exp([\boldsymbol\varphi_w]_\times)\).

If the same rotation is produced during a time interval \([t_0,t_1]\)
about the fixed axis \(\mathbf a_w\), let \(\alpha(t_0)=0\) and
\(\alpha(t_1)=\alpha\). The world-expressed angular velocity and
accumulated angle then satisfy

\[
\boxed{
\boldsymbol\omega_w(t)
=
\dot\alpha(t)\mathbf a_w
},
\qquad
\alpha
=
\int_{t_0}^{t_1}\dot\alpha(t)\,dt.
\]

Because \(\mathbf a_w\) is constant over this interval,

\[
\boxed{
\boldsymbol\varphi_w
=
\alpha\mathbf a_w
=
\int_{t_0}^{t_1}\boldsymbol\omega_w(t)\,dt
}.
\]

The complete fixed-axis relationship is therefore

\[
\underbrace{
\boldsymbol\omega_w(t)
}_{\text{angular velocity}}
\xrightarrow{\;\text{integrate over }[t_0,t_1]\;}
\underbrace{
\boldsymbol\varphi_w=\alpha\mathbf a_w
}_{\text{accumulated rotation coordinates}}
\xrightarrow{\;\exp([\,\cdot\,]_\times)\;}
\underbrace{
\mathbf Q_w
}_{\text{finite rotation}}.
\]

The vector \(\boldsymbol\varphi_w\) is dimensionless in SI because
radians are dimensionless, although its entries are labelled in radians
to preserve their physical meaning. It represents accumulated finite
rotation coordinates, not angular velocity; angular velocity has units
of \(\mathrm{rad\,s^{-1}}\).

\subsection{The matrix exponential}\label{the-matrix-exponential}

To construct the finite rotation, first consider a differentiable
world-coordinate vector \(\mathbf p(s)\in\mathbb R^3\) rotating about
\(\mathbf a_w\), where \(s\in\mathbb R\) is the accumulated signed
rotation angle, measured in radians and dimensionless in SI. Section
Section~\ref{sec-so3-angular-velocity-definition} derived
\(\dot{\mathbf r}_i=\boldsymbol\omega_w\times\mathbf r_i\) for each
body-axis column; the same vector-rate equation applies to any vector
rotating with the body:

\[
\frac{d\mathbf p}{dt}
=
\boldsymbol\omega_w\times\mathbf p.
\]

If \(s=s(t)\), then \(\boldsymbol\omega_w=\mathbf a_w\,ds/dt\). Applying
the chain rule gives

\[
\frac{d\mathbf p}{dt}
=
\frac{d\mathbf p}{ds}\frac{ds}{dt}
=
\left(
\mathbf a_w\times\mathbf p
\right)
\frac{ds}{dt}.
\]

Cancelling \(ds/dt\) when it is nonzero, with the stationary case
obtained by continuity, gives the vector-rate equation with angle rather
than time as the independent variable:

\[
\frac{d\mathbf p}{ds}
=
\mathbf a_w\times\mathbf p(s)
=
[\mathbf a_w]_\times\mathbf p(s).
\]

The initial condition is \(\mathbf p(0)=\mathbf p_0\), where
\(\mathbf p_0\in\mathbb R^3\) is the vector before the turn. The
coefficient matrix \([\mathbf a_w]_\times\in\mathfrak{so}(3)\) is
constant because the axis is fixed in world coordinates.

For any square matrix \(\mathbf A\in\mathbb R^{n\times n}\), the
\textbf{matrix exponential} is defined by the convergent power series

\[
\boxed{
\exp(\mathbf A)
\mathrel{:=}
\mathbf I_n
+
\mathbf A
+
\frac{\mathbf A^2}{2!}
+
\frac{\mathbf A^3}{3!}
+
\cdots
}.
\]

This definition extends the ordinary scalar exponential to square
matrices. For a constant matrix \(\mathbf A\), term-by-term
differentiation of the series gives

\[
\frac{d}{ds}\exp(s\mathbf A)
=
\mathbf A\exp(s\mathbf A).
\]

The scalar separation step \(d p/p\) cannot be used here because
division by a vector and the ordinary logarithm of a vector are not
defined. Its matrix counterpart is an \textbf{integrating factor},
meaning a matrix chosen so that the differential equation becomes the
derivative of one product. Set
\(\mathbf A\mathrel{:=}[\mathbf a_w]_\times\in\mathbb R^{3\times3}\) and
left-multiply \(d\mathbf p/ds=\mathbf A\mathbf p\) by
\(\exp(-s\mathbf A)\). The product rule gives

\[
\begin{aligned}
\frac{d}{ds}
\left[
\exp(-s\mathbf A)\mathbf p(s)
\right]
&=
-\mathbf A\exp(-s\mathbf A)\mathbf p(s)
+
\exp(-s\mathbf A)\frac{d\mathbf p}{ds}\\
&=
-\mathbf A\exp(-s\mathbf A)\mathbf p(s)
+
\exp(-s\mathbf A)\mathbf A\mathbf p(s)\\
&=
\mathbf 0_3.
\end{aligned}
\]

The last equality uses
\(\mathbf A\exp(-s\mathbf A)=\exp(-s\mathbf A)\mathbf A\), which holds
because the exponential series contains only powers of \(\mathbf A\).
The product is therefore constant. Evaluating it at \(s=0\) and using
\(\exp(\mathbf 0_{3\times3})=\mathbf I_3\) gives

\[
\exp(-s\mathbf A)\mathbf p(s)
=
\mathbf p_0.
\]

Left-multiplying by \(\exp(s\mathbf A)\) and using
\(\exp(s\mathbf A)^{-1}=\exp(-s\mathbf A)\) yields

\[
\boxed{
\mathbf p(s)
=
\exp\!\left(s[\mathbf a_w]_\times\right)\mathbf p_0
}.
\]

This integrating-factor argument derives the solution. Differentiating
the result provides a separate verification:

\[
\begin{aligned}
\frac{d\mathbf p}{ds}
&=
\frac{d}{ds}
\left[
\exp\!\left(s[\mathbf a_w]_\times\right)
\right]
\mathbf p_0\\
&=
[\mathbf a_w]_\times
\exp\!\left(s[\mathbf a_w]_\times\right)
\mathbf p_0\\
&=
[\mathbf a_w]_\times\mathbf p(s),
\end{aligned}
\]

and
\(\mathbf p(0)=\exp(\mathbf 0_{3\times3})\mathbf p_0=\mathbf I_3\mathbf p_0=\mathbf p_0\).
The cross-product-matrix construction is linear, so
\([\boldsymbol\varphi_w]_\times=[\alpha\mathbf a_w]_\times=\alpha[\mathbf a_w]_\times\).
Setting \(s=\alpha\) therefore defines the finite rotation

\[
\boxed{
\mathbf Q_w(\mathbf a_w,\alpha)
\mathrel{:=}
\exp\!\left(\alpha[\mathbf a_w]_\times\right)
=
\exp\!\left([\boldsymbol\varphi_w]_\times\right)
}.
\]

The mapping

\[
\exp:\mathfrak{so}(3)\to SO(3)
\]

is called the \textbf{exponential map} for rotations. The arrow states
that the function accepts an element of \(\mathfrak{so}(3)\) and returns
an element of \(SO(3)\). For the rotation vector introduced above,

\[
\underbrace{
[\boldsymbol\varphi_w]_\times
}_{\in\mathfrak{so}(3)}
\quad\longmapsto\quad
\underbrace{
\exp\!\left([\boldsymbol\varphi_w]_\times\right)
}_{\in SO(3)}.
\]

Thus, the exponential map converts the dimensionless skew-symmetric
matrix form \([\boldsymbol\varphi_w]_\times\) of the exponential
coordinates into a finite rotation matrix \(\mathbf Q_w\). Although
\([\boldsymbol\varphi_w]_\times\) belongs to the tangent space
\(\mathfrak{so}(3)\), it need not represent an infinitesimal change: its
scale is set by the possibly finite angle \(\alpha\). During
time-dependent motion, the genuinely infinitesimal rotation coordinate
accumulated over an infinitesimal time interval \(dt\) is

\[
\boxed{
[d\boldsymbol\varphi_w]_\times
=
[\boldsymbol\omega_w]_\times\,dt
},
\]

where \(\boldsymbol\omega_w\in\mathbb R^3\) is the world-expressed
angular velocity in \(\mathrm{rad\,s^{-1}}\).

To verify that the result is a rotation matrix, set
\(\boldsymbol\Omega\mathrel{:=}[\mathbf a_w]_\times\). Because
\(\boldsymbol\Omega^\mathsf T=-\boldsymbol\Omega\), transposing the
exponential series gives

\[
\mathbf Q_w^\mathsf T
=
\exp(-\alpha\boldsymbol\Omega).
\]

The inverse property of the matrix exponential gives
\(\exp(-\alpha\boldsymbol\Omega)=\exp(\alpha\boldsymbol\Omega)^{-1}\).
Hence,

\[
\mathbf Q_w^\mathsf T\mathbf Q_w
=
\exp(-\alpha\boldsymbol\Omega)
\exp(\alpha\boldsymbol\Omega)
=
\exp(\alpha\boldsymbol\Omega)^{-1}
\exp(\alpha\boldsymbol\Omega)
=
\mathbf I_3.
\]

Thus, \(\mathbf Q_w\) is orthogonal. Taking determinants of
\(\mathbf Q_w^\mathsf T\mathbf Q_w=\mathbf I_3\) gives

\[
\det(\mathbf Q_w)^2
=
1,
\]

so an orthogonal matrix can have determinant only \(+1\) or \(-1\). At
\(\alpha=0\),

\[
\mathbf Q_w(0)
=
\exp(\mathbf 0_{3\times3})
=
\mathbf I_3,
\qquad
\det(\mathbf Q_w(0))
=
1.
\]

As \(\alpha\) changes continuously, every entry of
\(\mathbf Q_w(\alpha)=\exp(\alpha\boldsymbol\Omega)\) also changes
continuously. The determinant is calculated from those entries using
sums and products, so \(\det(\mathbf Q_w(\alpha))\) is continuous. To
change from \(+1\) to \(-1\), it would have to pass through intermediate
values, including zero. This is impossible while \(\mathbf Q_w\) remains
orthogonal because an orthogonal matrix has determinant only \(+1\) or
\(-1\) and is never singular. Therefore, \(\det(\mathbf Q_w)=1\) for
every \(\alpha\), proving \(\mathbf Q_w\in SO(3)\).

\subsection{Deriving Rodrigues' rotation
formula}\label{deriving-rodrigues-rotation-formula}

Rodrigues' formula is a finite closed form of the matrix exponential
derived above. Setting
\(\boldsymbol\Omega\mathrel{:=}[\mathbf a_w]_\times\) in the
exponential-series definition gives

\[
\mathbf Q_w
=
\exp(\alpha\boldsymbol\Omega)
=
\mathbf I_3
+
\alpha\boldsymbol\Omega
+
\frac{\alpha^2\boldsymbol\Omega^2}{2!}
+
\frac{\alpha^3\boldsymbol\Omega^3}{3!}
+
\cdots.
\]

This expression contains infinitely many powers of
\(\boldsymbol\Omega\). The powers can be reduced to a repeating pattern
involving only \(\mathbf I_3\), \(\boldsymbol\Omega\), and
\(\boldsymbol\Omega^2\). To find that pattern, apply the vector
triple-product identity from
Section~\ref{sec-linear-algebra-cross-product},

\[
\mathbf a\times(\mathbf b\times\mathbf c)
=
\mathbf b(\mathbf a^\mathsf T\mathbf c)
-
\mathbf c(\mathbf a^\mathsf T\mathbf b),
\]

with \(\mathbf a=\mathbf b=\mathbf a_w\) and \(\mathbf c=\mathbf y\),
where \(\mathbf y\in\mathbb R^3\). This substitution gives

\[
\begin{aligned}
\boldsymbol\Omega^2\mathbf y
&=
\mathbf a_w\times(\mathbf a_w\times\mathbf y)\\
&=
\mathbf a_w(\mathbf a_w^\mathsf T\mathbf y)
-
\mathbf y(\mathbf a_w^\mathsf T\mathbf a_w)\\
&=
\left(
\mathbf a_w\mathbf a_w^\mathsf T
-
\mathbf I_3
\right)\mathbf y.
\end{aligned}
\]

The second equality is the triple-product identity after the stated
substitution: \(\mathbf a_w^\mathsf T\mathbf y\) is the scalar
projection factor in the first term, while
\(\mathbf a_w^\mathsf T\mathbf a_w\) is the squared norm of the axis.
The last equality uses \(\mathbf a_w^\mathsf T\mathbf a_w=1\), rewrites
\(\mathbf a_w(\mathbf a_w^\mathsf T\mathbf y)\) as
\((\mathbf a_w\mathbf a_w^\mathsf T)\mathbf y\), and rewrites
\(\mathbf y\) as \(\mathbf I_3\mathbf y\). Because the equation holds
for every \(\mathbf y\),

\[
\boxed{
\boldsymbol\Omega^2
=
\mathbf a_w\mathbf a_w^\mathsf T
-
\mathbf I_3
}.
\]

Furthermore,
\(\boldsymbol\Omega\mathbf a_w=\mathbf a_w\times\mathbf a_w=\mathbf 0_3\),
so

\[
\begin{aligned}
\boldsymbol\Omega^3
&=
\boldsymbol\Omega
\left(
\mathbf a_w\mathbf a_w^\mathsf T-
\mathbf I_3
\right)\\
&=
(\boldsymbol\Omega\mathbf a_w)\mathbf a_w^\mathsf T
-
\boldsymbol\Omega\\
&=
-\boldsymbol\Omega.
\end{aligned}
\]

Consequently, the odd powers alternate between \(\boldsymbol\Omega\) and
\(-\boldsymbol\Omega\), while the positive even powers alternate between
\(\boldsymbol\Omega^2\) and \(-\boldsymbol\Omega^2\). Expand the matrix
exponential and group those two power sequences:

\[
\begin{aligned}
\exp(\alpha\boldsymbol\Omega)
&=
\mathbf I_3
+
\alpha\boldsymbol\Omega
+
\frac{\alpha^2}{2!}\boldsymbol\Omega^2
+
\frac{\alpha^3}{3!}\boldsymbol\Omega^3
+
\cdots\\
&=
\mathbf I_3
+
\left(
\alpha-
\frac{\alpha^3}{3!}
+
\frac{\alpha^5}{5!}
-
\cdots
\right)\boldsymbol\Omega\\
&\quad+
\left(
\frac{\alpha^2}{2!}
-
\frac{\alpha^4}{4!}
+
\frac{\alpha^6}{6!}
-
\cdots
\right)\boldsymbol\Omega^2.
\end{aligned}
\]

The first scalar series is \(\sin\alpha\), and the second is
\(1-\cos\alpha\). Therefore,

\[
\boxed{
\mathbf Q_w(\mathbf a_w,\alpha)
=
\exp\!\left(\alpha[\mathbf a_w]_\times\right)
=
\mathbf I_3
+
\sin\alpha[\mathbf a_w]_\times
+
(1-\cos\alpha)[\mathbf a_w]_\times^2
}.
\]

This closed form is \textbf{Rodrigues' rotation formula}. Its geometric
meaning becomes clear by decomposing any vector
\(\mathbf y\in\mathbb R^3\) into a component parallel to the axis and a
component perpendicular to it:

\[
\mathbf y_{\parallel}
\mathrel{:=}
(\mathbf a_w^\mathsf T\mathbf y)\mathbf a_w,
\qquad
\mathbf y_{\perp}
\mathrel{:=}
\mathbf y-
\mathbf y_{\parallel}.
\]

Because
\(\boldsymbol\Omega\mathbf y=\mathbf a_w\times\mathbf y_{\perp}\) and
\(\boldsymbol\Omega^2\mathbf y=-\mathbf y_{\perp}\), Rodrigues' formula
gives

\[
\boxed{
\mathbf Q_w\mathbf y
=
\mathbf y_{\parallel}
+
\cos\alpha\,\mathbf y_{\perp}
+
\sin\alpha
(\mathbf a_w\times\mathbf y_{\perp})
}.
\]

The parallel component remains unchanged, while the perpendicular
component rotates through angle \(\alpha\) in its plane according to the
right-hand rule.

\subsection{\texorpdfstring{Check: rotation about the world
\(z_w\)-axis}{Check: rotation about the world z\_w-axis}}\label{check-rotation-about-the-world-z_w-axis}

Choose

\[
\mathbf a_w
=
\mathbf e_z
=
\begin{bmatrix}0\\0\\1\end{bmatrix},
\qquad
[\mathbf a_w]_\times
=
\begin{bmatrix}
0&-1&0\\
1&0&0\\
0&0&0
\end{bmatrix}.
\]

Squaring the cross-product matrix gives

\[
[\mathbf a_w]_\times^2
=
\begin{bmatrix}
-1&0&0\\
0&-1&0\\
0&0&0
\end{bmatrix}.
\]

Substitution into Rodrigues' formula produces

\[
\begin{aligned}
\mathbf Q_w(\mathbf e_z,\alpha)
&=
\mathbf I_3
+
\sin\alpha
\begin{bmatrix}
0&-1&0\\
1&0&0\\
0&0&0
\end{bmatrix}
+
(1-\cos\alpha)
\begin{bmatrix}
-1&0&0\\
0&-1&0\\
0&0&0
\end{bmatrix}\\
&=
\boxed{
\begin{bmatrix}
\cos\alpha&-\sin\alpha&0\\
\sin\alpha&\cos\alpha&0\\
0&0&1
\end{bmatrix}
}.
\end{aligned}
\]

This is exactly the coordinate-axis rotation \(\mathbf R_z(\alpha)\)
derived in Section~\ref{sec-so3-active-axis-rotations}.

\subsection{Applying a constant angular velocity over a time
interval}\label{applying-a-constant-angular-velocity-over-a-time-interval}

Suppose the world-expressed angular velocity
\(\boldsymbol\omega_w\in\mathbb R^3\) remains constant during a time
interval \(\Delta t>0\). If \(\boldsymbol\omega_w\ne\mathbf 0_3\),
define

\[
\mathbf a_w
\mathrel{:=}
\frac{\boldsymbol\omega_w}{\|\boldsymbol\omega_w\|_2},
\qquad
\alpha
\mathrel{:=}
\|\boldsymbol\omega_w\|_2\Delta t.
\]

The rotation vector over the interval is then

\[
\boldsymbol\varphi_w
=
\alpha\mathbf a_w
=
\Delta t\,\boldsymbol\omega_w.
\]

Because the angular velocity is expressed using the fixed world axes,
the finite update premultiplies the current orientation:

\[
\boxed{
\mathbf R_{wb}(t+\Delta t)
=
\exp\!\left(
\Delta t[\boldsymbol\omega_w]_\times
\right)
\mathbf R_{wb}(t)
}.
\]

If the constant angular velocity is instead expressed using the body
axes, the update postmultiplies:

\[
\boxed{
\mathbf R_{wb}(t+\Delta t)
=
\mathbf R_{wb}(t)
\exp\!\left(
\Delta t[\boldsymbol\omega_b]_\times
\right)
}.
\]

When the angular velocity is zero, define the finite increment as
\(\mathbf I_3\); the orientation then remains unchanged. These finite
updates are exact only when the corresponding world- or body-expressed
angular velocity remains constant throughout the interval.

Axis--angle coordinates are not globally unique. The pairs
\((\mathbf a_w,\alpha)\) and \((-\mathbf a_w,-\alpha)\) describe the
same turn, adding \(2\pi k\) to the angle for any \(k\in\mathbb Z\)
produces the same rotation, and when \(\alpha=0\) the axis direction is
arbitrary. These facts matter when recovering axis--angle coordinates
from a rotation matrix or comparing two numerical orientation
representations.

\section*{Quick test E: axis--angle coordinates and the exponential
map}\label{quick-test-e-axisangle-coordinates-and-the-exponential-map}
\addcontentsline{toc}{section}{Quick test E: axis--angle coordinates and
the exponential map}

\markright{Quick test E: axis--angle coordinates and the exponential
map}

\begin{enumerate}
\def\labelenumi{\arabic{enumi}.}
\tightlist
\item
  What information is stored in the axis--angle pair
  \((\mathbf a_w,\alpha)\), and how does the rotation vector
  \(\boldsymbol\varphi_w\) combine it?
\item
  State the matrix-exponential definition and explain why
  \(\mathbf p(s)=\exp(s[\mathbf a_w]_\times)\mathbf p_0\) solves
  \(d\mathbf p/ds=[\mathbf a_w]_\times\mathbf p\).
\item
  Use the vector triple-product identity to derive
  \([\mathbf a_w]_\times^2=\mathbf a_w\mathbf a_w^\mathsf T-\mathbf I_3\)
  for a unit axis.
\item
  Starting from the matrix-exponential series and
  \([\mathbf a_w]_\times^3=-[\mathbf a_w]_\times\), derive Rodrigues'
  rotation formula.
\item
  Why does Rodrigues' formula leave the vector component parallel to
  \(\mathbf a_w\) unchanged?
\item
  A constant angular velocity
  \(\boldsymbol\omega_w=[0\;0\;2]^\mathsf T\,\mathrm{rad\,s^{-1}}\) acts
  for \(\Delta t=0.25\,\mathrm s\). Find the axis, angle, rotation
  vector, and finite orientation increment.
\item
  Why does a world-expressed constant angular velocity premultiply
  \(\mathbf R_{wb}\), while a body-expressed constant angular velocity
  postmultiplies it?
\item
  Give two reasons why an axis--angle representation of a rotation is
  not unique.
\end{enumerate}

\subsection*{Answers and hints}\label{answers-and-hints-19}
\addcontentsline{toc}{subsection}{Answers and hints}

\begin{enumerate}
\def\labelenumi{\arabic{enumi}.}
\tightlist
\item
  The unit vector \(\mathbf a_w\) gives the positive axis direction in
  world coordinates, and \(\alpha\) gives the signed right-hand-rule
  angle. Their product \(\boldsymbol\varphi_w=\alpha\mathbf a_w\) stores
  the same information in one three-component column.
\item
  \(\exp(\mathbf A)=\mathbf I+\mathbf A+\mathbf A^2/2!+\cdots\).
  Differentiating \(\exp(s\mathbf A)\) term by term gives
  \(\mathbf A\exp(s\mathbf A)\); setting
  \(\mathbf A=[\mathbf a_w]_\times\) produces the required differential
  equation and the identity exponential satisfies the initial condition.
\item
  For every \(\mathbf y\),
  \(\mathbf a_w\times(\mathbf a_w\times\mathbf y)=\mathbf a_w(\mathbf a_w^\mathsf T\mathbf y)-(\mathbf a_w^\mathsf T\mathbf a_w)\mathbf y=(\mathbf a_w\mathbf a_w^\mathsf T-\mathbf I_3)\mathbf y\).
  Equality for every \(\mathbf y\) proves the matrix identity.
\item
  Replace successive odd powers by alternating \([\mathbf a_w]_\times\)
  and \(-[\mathbf a_w]_\times\), and successive positive even powers by
  alternating \([\mathbf a_w]_\times^2\) and
  \(-[\mathbf a_w]_\times^2\). Their scalar coefficients sum to
  \(\sin\alpha\) and \(1-\cos\alpha\), respectively.
\item
  The cross product of the axis with a parallel vector is zero, so both
  \([\mathbf a_w]_\times\mathbf y_{\parallel}\) and
  \([\mathbf a_w]_\times^2\mathbf y_{\parallel}\) vanish.
\item
  \(\mathbf a_w=[0\;0\;1]^\mathsf T\),
  \(\alpha=2(0.25)=0.5\,\mathrm{rad}\), and
  \(\boldsymbol\varphi_w=[0\;0\;0.5]^\mathsf T\,\mathrm{rad}\). The
  increment is
  \(\mathbf Q_w=\mathbf R_z(0.5)=\begin{bmatrix}\cos0.5&-\sin0.5&0\\\sin0.5&\cos0.5&0\\0&0&1\end{bmatrix}\).
\item
  A finite increment expressed using world axes acts on the
  world-coordinate columns and therefore premultiplies. An increment
  expressed using the old body axes composes after the body-to-world
  mapping and therefore postmultiplies.
\item
  Examples are \((\mathbf a_w,\alpha)\sim(-\mathbf a_w,-\alpha)\),
  angles differing by \(2\pi k\) give the same rotation, and the axis is
  arbitrary for zero rotation.
\end{enumerate}

\section*{References}\label{references-5}
\addcontentsline{toc}{section}{References}

\markright{References}

\begin{itemize}
\tightlist
\item
  Kevin M. Lynch and Frank C. Park, \emph{Modern Robotics: Mechanics,
  Planning, and Control}, Chapter 2 for configuration space and degrees
  of freedom, and Chapter 3, especially Sections 3.2.1-3.2.3 on spatial
  rotation matrices, angular velocity, exponential coordinates, and
  Rodrigues' formula.
\item
  John J. Craig, \emph{Introduction to Robotics: Mechanics and Control},
  4th ed., Chapters 1-2, especially degrees of freedom, body-attached
  frames, proper orthonormal rotation matrices, spatial orientation, and
  the equivalent angle--axis representation.
\item
  Peter Corke, \emph{Robotics, Vision and Control: Fundamental
  Algorithms in Python}, 3rd ed., Chapters 2-3, especially rotation
  matrices, the matrix exponential, Rodrigues' formula, rotational
  velocity, and body- and world-frame expressions of spatial motion.
\end{itemize}




\end{document}
